% !TeX program = xelatex
% ================================================================
%  《收获与播种》(Récoltes et Semailles) 中文译本
%  卷三 · 第二部 —— 葬礼(I): 或中国皇帝的新衣
%  Deuxième partie: L'Enterrement (I) ou la robe de l'Empereur de Chine
%  原书 p.255-412, 第 13-16 章
%  --------------------------------------------------------------
%  五卷本之一。共用导言区 common.tex, 封面 cover-art.tex。
%  构建: xelatex ×4 (目录页码需 4 遍才收敛)
% ================================================================
\PassOptionsToPackage{AutoFallBack}{xeCJK}
\documentclass[
    PaperSize=a4paper,
    HeadingFont=hei,
    Biblatex=false,
    BookTitleCN={收获与播种},
    Translator={数译},
    LogoFile=logo.pdf
]{mathtranslation}

% ================================================================
%  《收获与播种》中文译本 —— 各卷共用导言区
%  由 vol-prelude / vol-part1 / vol-part2 / vol-part3 / vol-part4
%  以及全书合订 main-all 共同 \input。基类 mathtranslation.cls v3.2 不改。
%  --------------------------------------------------------------
%  样式均在此改写, 以贴近原书:
%    · 章标题仿原书 "5. Travail et découverte"(块状左对齐黑体)
%    · 节/小节: 黑体左对齐, 编号后带句点
%    · 脚注: 人工编号照录原书——前奏与第一部分页内重新编号,
%      第二至四部各自连续编号并带(*)角标; 第一部分另有 \fnoteP
%      处理 "(n)" 式连续注号
%    · 边码: 原打字稿页码 \origpage{N}(前奏作 A1、P1、L1、罗马数字等)
%    · ⋄ 记号: 段首补写标记 \pmark 与上标形式 \dmark
% ================================================================

% 题名含空格, 不能走类选项(会被 kvoptions 吃掉), 在此赋值
\renewcommand{\BookTitleEN}{Récoltes et Semailles}
\renewcommand{\OriginalAuthor}{Alexandre Grothendieck}
\renewcommand{\TranslationDate}{\today}

% --------------------------------------------------------------
%  罕用字后备(如"祐"): Fandol 不载时退到 SimSun
% --------------------------------------------------------------
\setCJKfallbackfamilyfont{\CJKrmdefault}{{SimSun}}

% --------------------------------------------------------------
%  hyperref 元信息(各卷在 % ================================================================
%  《收获与播种》中文译本 —— 各卷共用导言区
%  由 vol-prelude / vol-part1 / vol-part2 / vol-part3 / vol-part4
%  以及全书合订 main-all 共同 \input。基类 mathtranslation.cls v3.2 不改。
%  --------------------------------------------------------------
%  样式均在此改写, 以贴近原书:
%    · 章标题仿原书 "5. Travail et découverte"(块状左对齐黑体)
%    · 节/小节: 黑体左对齐, 编号后带句点
%    · 脚注: 人工编号照录原书——前奏与第一部分页内重新编号,
%      第二至四部各自连续编号并带(*)角标; 第一部分另有 \fnoteP
%      处理 "(n)" 式连续注号
%    · 边码: 原打字稿页码 \origpage{N}(前奏作 A1、P1、L1、罗马数字等)
%    · ⋄ 记号: 段首补写标记 \pmark 与上标形式 \dmark
% ================================================================

% 题名含空格, 不能走类选项(会被 kvoptions 吃掉), 在此赋值
\renewcommand{\BookTitleEN}{Récoltes et Semailles}
\renewcommand{\OriginalAuthor}{Alexandre Grothendieck}
\renewcommand{\TranslationDate}{\today}

% --------------------------------------------------------------
%  罕用字后备(如"祐"): Fandol 不载时退到 SimSun
% --------------------------------------------------------------
\setCJKfallbackfamilyfont{\CJKrmdefault}{{SimSun}}

% --------------------------------------------------------------
%  hyperref 元信息(各卷在 % ================================================================
%  《收获与播种》中文译本 —— 各卷共用导言区
%  由 vol-prelude / vol-part1 / vol-part2 / vol-part3 / vol-part4
%  以及全书合订 main-all 共同 \input。基类 mathtranslation.cls v3.2 不改。
%  --------------------------------------------------------------
%  样式均在此改写, 以贴近原书:
%    · 章标题仿原书 "5. Travail et découverte"(块状左对齐黑体)
%    · 节/小节: 黑体左对齐, 编号后带句点
%    · 脚注: 人工编号照录原书——前奏与第一部分页内重新编号,
%      第二至四部各自连续编号并带(*)角标; 第一部分另有 \fnoteP
%      处理 "(n)" 式连续注号
%    · 边码: 原打字稿页码 \origpage{N}(前奏作 A1、P1、L1、罗马数字等)
%    · ⋄ 记号: 段首补写标记 \pmark 与上标形式 \dmark
% ================================================================

% 题名含空格, 不能走类选项(会被 kvoptions 吃掉), 在此赋值
\renewcommand{\BookTitleEN}{Récoltes et Semailles}
\renewcommand{\OriginalAuthor}{Alexandre Grothendieck}
\renewcommand{\TranslationDate}{\today}

% --------------------------------------------------------------
%  罕用字后备(如"祐"): Fandol 不载时退到 SimSun
% --------------------------------------------------------------
\setCJKfallbackfamilyfont{\CJKrmdefault}{{SimSun}}

% --------------------------------------------------------------
%  hyperref 元信息(各卷在 \input{common} 之后可用 \hypersetup 覆写)
% --------------------------------------------------------------
\hypersetup{
  pdftitle={收获与播种},
  pdfauthor={Alexandre Grothendieck / 数译 译},
  pdfsubject={Récoltes et Semailles — Réflexions et témoignage sur un passé de mathématicien}
}

% --------------------------------------------------------------
%  标题样式: 仿原书 —— 黑体左对齐, 编号后带句点
% --------------------------------------------------------------
\titleformat{\section}[block]
  {\normalfont\Large\headfont}{\thesection.}{0.5em}{}
\titlespacing*{\section}{0pt}{3.2ex plus 1.2ex minus .2ex}{1.6ex plus .2ex}
\titleformat{\subsection}[block]
  {\normalfont\large\headfont}{\thesubsection.}{0.5em}{}
\titlespacing*{\subsection}{0pt}{2.6ex plus 1ex minus .2ex}{1.2ex plus .2ex}
\titleformat{\subsubsection}[block]
  {\normalfont\normalsize\headfont}{\thesubsubsection.}{0.5em}{}
\titlespacing*{\subsubsection}{0pt}{2.2ex plus .8ex minus .2ex}{1ex plus .2ex}
\titleformat{name=\subsubsection,numberless}[block]
  {\normalfont\normalsize\headfont}{}{0em}{}
\setcounter{secnumdepth}{3}
\setcounter{tocdepth}{3}
% 为含长外文书名与行内公式的段落提供必要的断行余量
\setlength{\emergencystretch}{3em}

% 章标题样式: 原书章标题与节同级排印("13. A) HERITAGE..."),
% 故章标题取块状左对齐黑体, 字号介于节与部之间
\titleformat{\chapter}[block]
  {\normalfont\LARGE\headfont}{\thechapter.}{0.5em}{}
\titlespacing*{\chapter}{0pt}{3.6ex plus 1.4ex minus .2ex}{2ex plus .2ex}
% 目录中章号照原书用阿拉伯数字("13. A) ..." 而非 "第十三章")
\renewcommand{\CTEXthechapter}{\thechapter\hspace{.3em}}

% --------------------------------------------------------------
%  目录观感: 条目与引导点之间、引导点与页码之间各留 0.5em
%  内核 \@dottedtocline 让 \leaders 从标题末字起排, 第一个点几乎
%  贴着末字, 读起来像多标了一个句点; 页码前也没有余量。两处各补
%  0.5em。
% --------------------------------------------------------------
\makeatletter
\patchcmd{\@dottedtocline}{\nobreak\leaders}{\nobreak\hspace{0.5em}\leaders}%
  {\typeout{MT: TOC 引导点前已留 0.5em}}%
  {\typeout{MT-WARN: TOC 引导点前间距补丁未生效}}
\patchcmd{\@dottedtocline}{\hfill\nobreak\hb@xt@\@pnumwidth}%
  {\hfill\nobreak\hspace{0.5em}\hb@xt@\@pnumwidth}%
  {\typeout{MT: TOC 页码前已留 0.5em}}%
  {\typeout{MT-WARN: TOC 页码前间距补丁未生效}}
\makeatother

% --------------------------------------------------------------
%  正文半角标点后的间距
%  本书体例用半角 , ; : ? ! 与 ASCII 括号, 由 xeCJK 按 HalfRight
%  (右半留白) 处理。放宽到 0.45em, 与全角标点右半留白大致齐平。
% --------------------------------------------------------------
\xeCJKsetup{CJKecglue={\hskip 0.45em plus 0.06em minus 0.09em}}

% --------------------------------------------------------------
%  页眉页脚: 偶数页眉右对齐斜体节题; 页码页脚外侧
% --------------------------------------------------------------
\fancyhf{}
\fancyhead[RE]{\small\itshape\rightmark}
\fancyfoot[LE,RO]{\small\thepage}
\renewcommand{\headrulewidth}{0pt}
\renewcommand{\sectionmark}[1]{\markright{\thesection.\ #1}}
\fancypagestyle{plain}{%
  \fancyhf{}%
  \fancyfoot[LE,RO]{\small \thepage}%
  \renewcommand{\headrulewidth}{0pt}%
}

% --------------------------------------------------------------
%  原书版式记号宏
% --------------------------------------------------------------
% 边码: 原打字稿页码(G. 手稿系统里的自引页码), 排在页边
\usepackage{marginnote}
\newcommand{\origpage}[1]{\marginnote{\normalfont\scriptsize p.~#1}}

% ⋄ 记号: 段首/行内正常大小; 上标形式为作者后补文字的位置标记
\DeclareRobustCommand{\pmark}{\ensuremath{\diamond}\hspace{0.15em}}
\DeclareRobustCommand{\dmark}{\textsuperscript{\ensuremath{\diamond}}}

% 脚注人工编号: \fnote{795}{(*)}{文本} —— 编号与(*)角标均照录原书
\newcommand{\fnote}[3]{%
  \begingroup
  \renewcommand{\thefootnote}{#1#2}%
  \footnote{#3}%
  \endgroup
}

% 第一部分的脚注: 页内序号作上标(#1), 注文行首印 "序号(连续注号)",
% 正文在序号之后另行括注同一连续注号 —— \fnoteP{1}{(3)}{注文}
% 原书有三条注无连续注号而改标日期(ch07 注3 = 30 Septembre,
% ch09 注11/12 = 8 août), 第二参数留空, 日期写进注文之首。
\newcommand{\fnoteP}[3]{%
  \begingroup
  \renewcommand{\thefootnote}{#1}\footnotemark
  \renewcommand{\thefootnote}{#1#2}\footnotetext{#3}%
  \endgroup
}

% 正文呼号与注文行首标号不同(原书照录)
\newcommand{\fnotevariant}[4]{%
  \begingroup
  \renewcommand{\thefootnote}{#1#2}\footnotemark
  \renewcommand{\thefootnote}{#1#3}\footnotetext{#4}%
  \endgroup
}

% 大注块标题: \notehead{171³}{斜体评论} —— 正文随后另起
\newcommand{\notehead}[2]{%
  \par\bigskip
  \noindent{\headfont 注~#1}\ifblank{#2}{}{\hspace{0.5em}{\itshape #2}}%
  \par\nobreak\smallskip\nobreak
}

% ∗ ∗ ∗ 分隔
\newcommand{\asterisms}{%
  \par\vspace{1.2ex}%
  \noindent\hspace*{2em}$\ast$\hspace{1.2em}$\ast$\hspace{1.2em}$\ast$%
  \par\vspace{1.2ex}%
}

% 无编号子目的目录条目
\newcommand{\tocsubsub}[1]{%
  \phantomsection
  \addcontentsline{toc}{subsubsection}{#1}%
}

% 部扉: 仿原书部题页
\newcommand{\partdivider}[3]{%
  \clearpage
  \thispagestyle{empty}%
  \phantomsection
  \addcontentsline{toc}{part}{#1\quad #3}%
  \vspace*{6.5cm}
  \begin{center}
    {\fontsize{16pt}{26pt}\selectfont\headfont #1\par}
    \vspace{0.9cm}
    {\fontsize{22pt}{34pt}\selectfont\headfont #2\par}
  \end{center}
  \clearpage
}

% 献辞页: 仿原书部题后的题献, 独占一页且不入目录
\newcommand{\dedication}[1]{%
  \thispagestyle{empty}%
  \vspace*{4cm}%
  \begin{center}#1\end{center}%
  \clearpage
}

% --------------------------------------------------------------
%  各卷"译者说明"的共用体例段(各卷的范围段自写, 后接本段)
% --------------------------------------------------------------
\newcommand{\CommonNoteBody}{%
为便于与原书及学界通行的打字稿页码对勘, 本册保留:

\begin{itemize}
  \item 原书全部章节编号(章~1--18, 节如~5.2、18.2, 子目如~18.1.1)与原书小节标题;
  \item 原书脚注的编号与角标——
        \begin{itemize}
          \item 前奏(第~1--4 章)与第一部的注在各页内重新编号, 正文中的上标即该
                页序号; 第一部紧随其后的括注(如“(3)”)是该部分的连续注号, 与第~12
                章“第一部分的注”里的“注~3”一一对应, 注文行首照原书印作“1(3)”式。
                其中三条注原书未标连续注号而改标日期(ch7 注~3 = 30 Septembre、
                ch9 注~11/12 = 8 août), 本册照录其日期;
          \item 第二部第~13--16~章各自重新编号, 分别为~1--41、1--86、1--132、1--29;
                第三部第~17~章自~1~起至~12, 第~18~章自~1~重新起算并延续至~794;
                末段自~\textbf{795(*)}~续至~1053(*)。各段与原书一一对应;
        \end{itemize}
  \item 页边码“p.~1”“p.~A3”“p.~1066”式小注——这是作者引用自己打字稿原稿时
        所用的页码, 与本排印本页码不同。前奏四章的边码各有前缀: 第~1 章作
        \texttt{A1…A6}, 第~2 章作 \texttt{P1…P65}, 第~3 章作 \texttt{L1…L56},
        第~4 章作罗马数字 \texttt{I…XXII}, 自第一部起作纯数字 \texttt{1…171};
  \item 作者用以标记段落的~$\diamond$~记号, 以及标示后来补写文字的上标
        ~\textsuperscript{$\diamond$}~记号;
  \item 原书大小注(“注~44'”“注~46$^1$”“注~171³”等)的编号与层级。
\end{itemize}

第~12 章“第一部分的注”把第~5 至第~11 章各页脚注汇为一编, 并给每条注配上节题
(如“12.19. 耶稣与十二使徒”), 因此部分注文与正文脚注重复——这是原书的体例, 非
重出之误。这一部分属随笔式的自省文字, 同一句法语在正文脚注与章末汇编注中
各有措辞, 本册有意保留各自的译笔, 不强求统一。

除上述标记外, 译文力求贴近作者驳杂、跳跃而辛辣的文风; 数学名词的译法见各卷
术语表(如~\emph{topos}~不译, \emph{faisceaux pervers}~作“反常层”)。文中的日期
括注(如“(5~月~30~日)”“(1984年3月补写)”)为作者历次补写时的纪日, 一律照录。文中
另有一些对注号、页码的引用(如“n°~1695”“打字稿第~227~页”)沿袭作者打字稿的编号
系统, 与本排印本不尽一致, 亦无法在本书中落实, 悉照原文录出。

全套五卷(另有全书合订一册)的分工是:

\begin{itemize}
  \item \textbf{卷一·前奏}——主题的呈现: 或四个乐章的前奏(原书 p.23--122);
  \item \textbf{卷二·第一部}——愚妄与更新(原书 p.123--254);
  \item \textbf{卷三·第二部}——葬礼(I): 或中国皇帝的新衣(原书 p.255--412);
  \item \textbf{卷四·第三部}——葬礼(II): 或阴阳之钥(原书 p.413--632);
  \item \textbf{卷五·第四部}——葬礼(III): 或四种运算(原书 p.633--929)。
\end{itemize}
}
 之后可用 \hypersetup 覆写)
% --------------------------------------------------------------
\hypersetup{
  pdftitle={收获与播种},
  pdfauthor={Alexandre Grothendieck / 数译 译},
  pdfsubject={Récoltes et Semailles — Réflexions et témoignage sur un passé de mathématicien}
}

% --------------------------------------------------------------
%  标题样式: 仿原书 —— 黑体左对齐, 编号后带句点
% --------------------------------------------------------------
\titleformat{\section}[block]
  {\normalfont\Large\headfont}{\thesection.}{0.5em}{}
\titlespacing*{\section}{0pt}{3.2ex plus 1.2ex minus .2ex}{1.6ex plus .2ex}
\titleformat{\subsection}[block]
  {\normalfont\large\headfont}{\thesubsection.}{0.5em}{}
\titlespacing*{\subsection}{0pt}{2.6ex plus 1ex minus .2ex}{1.2ex plus .2ex}
\titleformat{\subsubsection}[block]
  {\normalfont\normalsize\headfont}{\thesubsubsection.}{0.5em}{}
\titlespacing*{\subsubsection}{0pt}{2.2ex plus .8ex minus .2ex}{1ex plus .2ex}
\titleformat{name=\subsubsection,numberless}[block]
  {\normalfont\normalsize\headfont}{}{0em}{}
\setcounter{secnumdepth}{3}
\setcounter{tocdepth}{3}
% 为含长外文书名与行内公式的段落提供必要的断行余量
\setlength{\emergencystretch}{3em}

% 章标题样式: 原书章标题与节同级排印("13. A) HERITAGE..."),
% 故章标题取块状左对齐黑体, 字号介于节与部之间
\titleformat{\chapter}[block]
  {\normalfont\LARGE\headfont}{\thechapter.}{0.5em}{}
\titlespacing*{\chapter}{0pt}{3.6ex plus 1.4ex minus .2ex}{2ex plus .2ex}
% 目录中章号照原书用阿拉伯数字("13. A) ..." 而非 "第十三章")
\renewcommand{\CTEXthechapter}{\thechapter\hspace{.3em}}

% --------------------------------------------------------------
%  目录观感: 条目与引导点之间、引导点与页码之间各留 0.5em
%  内核 \@dottedtocline 让 \leaders 从标题末字起排, 第一个点几乎
%  贴着末字, 读起来像多标了一个句点; 页码前也没有余量。两处各补
%  0.5em。
% --------------------------------------------------------------
\makeatletter
\patchcmd{\@dottedtocline}{\nobreak\leaders}{\nobreak\hspace{0.5em}\leaders}%
  {\typeout{MT: TOC 引导点前已留 0.5em}}%
  {\typeout{MT-WARN: TOC 引导点前间距补丁未生效}}
\patchcmd{\@dottedtocline}{\hfill\nobreak\hb@xt@\@pnumwidth}%
  {\hfill\nobreak\hspace{0.5em}\hb@xt@\@pnumwidth}%
  {\typeout{MT: TOC 页码前已留 0.5em}}%
  {\typeout{MT-WARN: TOC 页码前间距补丁未生效}}
\makeatother

% --------------------------------------------------------------
%  正文半角标点后的间距
%  本书体例用半角 , ; : ? ! 与 ASCII 括号, 由 xeCJK 按 HalfRight
%  (右半留白) 处理。放宽到 0.45em, 与全角标点右半留白大致齐平。
% --------------------------------------------------------------
\xeCJKsetup{CJKecglue={\hskip 0.45em plus 0.06em minus 0.09em}}

% --------------------------------------------------------------
%  页眉页脚: 偶数页眉右对齐斜体节题; 页码页脚外侧
% --------------------------------------------------------------
\fancyhf{}
\fancyhead[RE]{\small\itshape\rightmark}
\fancyfoot[LE,RO]{\small\thepage}
\renewcommand{\headrulewidth}{0pt}
\renewcommand{\sectionmark}[1]{\markright{\thesection.\ #1}}
\fancypagestyle{plain}{%
  \fancyhf{}%
  \fancyfoot[LE,RO]{\small \thepage}%
  \renewcommand{\headrulewidth}{0pt}%
}

% --------------------------------------------------------------
%  原书版式记号宏
% --------------------------------------------------------------
% 边码: 原打字稿页码(G. 手稿系统里的自引页码), 排在页边
\usepackage{marginnote}
\newcommand{\origpage}[1]{\marginnote{\normalfont\scriptsize p.~#1}}

% ⋄ 记号: 段首/行内正常大小; 上标形式为作者后补文字的位置标记
\DeclareRobustCommand{\pmark}{\ensuremath{\diamond}\hspace{0.15em}}
\DeclareRobustCommand{\dmark}{\textsuperscript{\ensuremath{\diamond}}}

% 脚注人工编号: \fnote{795}{(*)}{文本} —— 编号与(*)角标均照录原书
\newcommand{\fnote}[3]{%
  \begingroup
  \renewcommand{\thefootnote}{#1#2}%
  \footnote{#3}%
  \endgroup
}

% 第一部分的脚注: 页内序号作上标(#1), 注文行首印 "序号(连续注号)",
% 正文在序号之后另行括注同一连续注号 —— \fnoteP{1}{(3)}{注文}
% 原书有三条注无连续注号而改标日期(ch07 注3 = 30 Septembre,
% ch09 注11/12 = 8 août), 第二参数留空, 日期写进注文之首。
\newcommand{\fnoteP}[3]{%
  \begingroup
  \renewcommand{\thefootnote}{#1}\footnotemark
  \renewcommand{\thefootnote}{#1#2}\footnotetext{#3}%
  \endgroup
}

% 正文呼号与注文行首标号不同(原书照录)
\newcommand{\fnotevariant}[4]{%
  \begingroup
  \renewcommand{\thefootnote}{#1#2}\footnotemark
  \renewcommand{\thefootnote}{#1#3}\footnotetext{#4}%
  \endgroup
}

% 大注块标题: \notehead{171³}{斜体评论} —— 正文随后另起
\newcommand{\notehead}[2]{%
  \par\bigskip
  \noindent{\headfont 注~#1}\ifblank{#2}{}{\hspace{0.5em}{\itshape #2}}%
  \par\nobreak\smallskip\nobreak
}

% ∗ ∗ ∗ 分隔
\newcommand{\asterisms}{%
  \par\vspace{1.2ex}%
  \noindent\hspace*{2em}$\ast$\hspace{1.2em}$\ast$\hspace{1.2em}$\ast$%
  \par\vspace{1.2ex}%
}

% 无编号子目的目录条目
\newcommand{\tocsubsub}[1]{%
  \phantomsection
  \addcontentsline{toc}{subsubsection}{#1}%
}

% 部扉: 仿原书部题页
\newcommand{\partdivider}[3]{%
  \clearpage
  \thispagestyle{empty}%
  \phantomsection
  \addcontentsline{toc}{part}{#1\quad #3}%
  \vspace*{6.5cm}
  \begin{center}
    {\fontsize{16pt}{26pt}\selectfont\headfont #1\par}
    \vspace{0.9cm}
    {\fontsize{22pt}{34pt}\selectfont\headfont #2\par}
  \end{center}
  \clearpage
}

% 献辞页: 仿原书部题后的题献, 独占一页且不入目录
\newcommand{\dedication}[1]{%
  \thispagestyle{empty}%
  \vspace*{4cm}%
  \begin{center}#1\end{center}%
  \clearpage
}

% --------------------------------------------------------------
%  各卷"译者说明"的共用体例段(各卷的范围段自写, 后接本段)
% --------------------------------------------------------------
\newcommand{\CommonNoteBody}{%
为便于与原书及学界通行的打字稿页码对勘, 本册保留:

\begin{itemize}
  \item 原书全部章节编号(章~1--18, 节如~5.2、18.2, 子目如~18.1.1)与原书小节标题;
  \item 原书脚注的编号与角标——
        \begin{itemize}
          \item 前奏(第~1--4 章)与第一部的注在各页内重新编号, 正文中的上标即该
                页序号; 第一部紧随其后的括注(如“(3)”)是该部分的连续注号, 与第~12
                章“第一部分的注”里的“注~3”一一对应, 注文行首照原书印作“1(3)”式。
                其中三条注原书未标连续注号而改标日期(ch7 注~3 = 30 Septembre、
                ch9 注~11/12 = 8 août), 本册照录其日期;
          \item 第二部第~13--16~章各自重新编号, 分别为~1--41、1--86、1--132、1--29;
                第三部第~17~章自~1~起至~12, 第~18~章自~1~重新起算并延续至~794;
                末段自~\textbf{795(*)}~续至~1053(*)。各段与原书一一对应;
        \end{itemize}
  \item 页边码“p.~1”“p.~A3”“p.~1066”式小注——这是作者引用自己打字稿原稿时
        所用的页码, 与本排印本页码不同。前奏四章的边码各有前缀: 第~1 章作
        \texttt{A1…A6}, 第~2 章作 \texttt{P1…P65}, 第~3 章作 \texttt{L1…L56},
        第~4 章作罗马数字 \texttt{I…XXII}, 自第一部起作纯数字 \texttt{1…171};
  \item 作者用以标记段落的~$\diamond$~记号, 以及标示后来补写文字的上标
        ~\textsuperscript{$\diamond$}~记号;
  \item 原书大小注(“注~44'”“注~46$^1$”“注~171³”等)的编号与层级。
\end{itemize}

第~12 章“第一部分的注”把第~5 至第~11 章各页脚注汇为一编, 并给每条注配上节题
(如“12.19. 耶稣与十二使徒”), 因此部分注文与正文脚注重复——这是原书的体例, 非
重出之误。这一部分属随笔式的自省文字, 同一句法语在正文脚注与章末汇编注中
各有措辞, 本册有意保留各自的译笔, 不强求统一。

除上述标记外, 译文力求贴近作者驳杂、跳跃而辛辣的文风; 数学名词的译法见各卷
术语表(如~\emph{topos}~不译, \emph{faisceaux pervers}~作“反常层”)。文中的日期
括注(如“(5~月~30~日)”“(1984年3月补写)”)为作者历次补写时的纪日, 一律照录。文中
另有一些对注号、页码的引用(如“n°~1695”“打字稿第~227~页”)沿袭作者打字稿的编号
系统, 与本排印本不尽一致, 亦无法在本书中落实, 悉照原文录出。

全套五卷(另有全书合订一册)的分工是:

\begin{itemize}
  \item \textbf{卷一·前奏}——主题的呈现: 或四个乐章的前奏(原书 p.23--122);
  \item \textbf{卷二·第一部}——愚妄与更新(原书 p.123--254);
  \item \textbf{卷三·第二部}——葬礼(I): 或中国皇帝的新衣(原书 p.255--412);
  \item \textbf{卷四·第三部}——葬礼(II): 或阴阳之钥(原书 p.413--632);
  \item \textbf{卷五·第四部}——葬礼(III): 或四种运算(原书 p.633--929)。
\end{itemize}
}
 之后可用 \hypersetup 覆写)
% --------------------------------------------------------------
\hypersetup{
  pdftitle={收获与播种},
  pdfauthor={Alexandre Grothendieck / 数译 译},
  pdfsubject={Récoltes et Semailles — Réflexions et témoignage sur un passé de mathématicien}
}

% --------------------------------------------------------------
%  标题样式: 仿原书 —— 黑体左对齐, 编号后带句点
% --------------------------------------------------------------
\titleformat{\section}[block]
  {\normalfont\Large\headfont}{\thesection.}{0.5em}{}
\titlespacing*{\section}{0pt}{3.2ex plus 1.2ex minus .2ex}{1.6ex plus .2ex}
\titleformat{\subsection}[block]
  {\normalfont\large\headfont}{\thesubsection.}{0.5em}{}
\titlespacing*{\subsection}{0pt}{2.6ex plus 1ex minus .2ex}{1.2ex plus .2ex}
\titleformat{\subsubsection}[block]
  {\normalfont\normalsize\headfont}{\thesubsubsection.}{0.5em}{}
\titlespacing*{\subsubsection}{0pt}{2.2ex plus .8ex minus .2ex}{1ex plus .2ex}
\titleformat{name=\subsubsection,numberless}[block]
  {\normalfont\normalsize\headfont}{}{0em}{}
\setcounter{secnumdepth}{3}
\setcounter{tocdepth}{3}
% 为含长外文书名与行内公式的段落提供必要的断行余量
\setlength{\emergencystretch}{3em}

% 章标题样式: 原书章标题与节同级排印("13. A) HERITAGE..."),
% 故章标题取块状左对齐黑体, 字号介于节与部之间
\titleformat{\chapter}[block]
  {\normalfont\LARGE\headfont}{\thechapter.}{0.5em}{}
\titlespacing*{\chapter}{0pt}{3.6ex plus 1.4ex minus .2ex}{2ex plus .2ex}
% 目录中章号照原书用阿拉伯数字("13. A) ..." 而非 "第十三章")
\renewcommand{\CTEXthechapter}{\thechapter\hspace{.3em}}

% --------------------------------------------------------------
%  目录观感: 条目与引导点之间、引导点与页码之间各留 0.5em
%  内核 \@dottedtocline 让 \leaders 从标题末字起排, 第一个点几乎
%  贴着末字, 读起来像多标了一个句点; 页码前也没有余量。两处各补
%  0.5em。
% --------------------------------------------------------------
\makeatletter
\patchcmd{\@dottedtocline}{\nobreak\leaders}{\nobreak\hspace{0.5em}\leaders}%
  {\typeout{MT: TOC 引导点前已留 0.5em}}%
  {\typeout{MT-WARN: TOC 引导点前间距补丁未生效}}
\patchcmd{\@dottedtocline}{\hfill\nobreak\hb@xt@\@pnumwidth}%
  {\hfill\nobreak\hspace{0.5em}\hb@xt@\@pnumwidth}%
  {\typeout{MT: TOC 页码前已留 0.5em}}%
  {\typeout{MT-WARN: TOC 页码前间距补丁未生效}}
\makeatother

% --------------------------------------------------------------
%  正文半角标点后的间距
%  本书体例用半角 , ; : ? ! 与 ASCII 括号, 由 xeCJK 按 HalfRight
%  (右半留白) 处理。放宽到 0.45em, 与全角标点右半留白大致齐平。
% --------------------------------------------------------------
\xeCJKsetup{CJKecglue={\hskip 0.45em plus 0.06em minus 0.09em}}

% --------------------------------------------------------------
%  页眉页脚: 偶数页眉右对齐斜体节题; 页码页脚外侧
% --------------------------------------------------------------
\fancyhf{}
\fancyhead[RE]{\small\itshape\rightmark}
\fancyfoot[LE,RO]{\small\thepage}
\renewcommand{\headrulewidth}{0pt}
\renewcommand{\sectionmark}[1]{\markright{\thesection.\ #1}}
\fancypagestyle{plain}{%
  \fancyhf{}%
  \fancyfoot[LE,RO]{\small \thepage}%
  \renewcommand{\headrulewidth}{0pt}%
}

% --------------------------------------------------------------
%  原书版式记号宏
% --------------------------------------------------------------
% 边码: 原打字稿页码(G. 手稿系统里的自引页码), 排在页边
\usepackage{marginnote}
\newcommand{\origpage}[1]{\marginnote{\normalfont\scriptsize p.~#1}}

% ⋄ 记号: 段首/行内正常大小; 上标形式为作者后补文字的位置标记
\DeclareRobustCommand{\pmark}{\ensuremath{\diamond}\hspace{0.15em}}
\DeclareRobustCommand{\dmark}{\textsuperscript{\ensuremath{\diamond}}}

% 脚注人工编号: \fnote{795}{(*)}{文本} —— 编号与(*)角标均照录原书
\newcommand{\fnote}[3]{%
  \begingroup
  \renewcommand{\thefootnote}{#1#2}%
  \footnote{#3}%
  \endgroup
}

% 第一部分的脚注: 页内序号作上标(#1), 注文行首印 "序号(连续注号)",
% 正文在序号之后另行括注同一连续注号 —— \fnoteP{1}{(3)}{注文}
% 原书有三条注无连续注号而改标日期(ch07 注3 = 30 Septembre,
% ch09 注11/12 = 8 août), 第二参数留空, 日期写进注文之首。
\newcommand{\fnoteP}[3]{%
  \begingroup
  \renewcommand{\thefootnote}{#1}\footnotemark
  \renewcommand{\thefootnote}{#1#2}\footnotetext{#3}%
  \endgroup
}

% 正文呼号与注文行首标号不同(原书照录)
\newcommand{\fnotevariant}[4]{%
  \begingroup
  \renewcommand{\thefootnote}{#1#2}\footnotemark
  \renewcommand{\thefootnote}{#1#3}\footnotetext{#4}%
  \endgroup
}

% 大注块标题: \notehead{171³}{斜体评论} —— 正文随后另起
\newcommand{\notehead}[2]{%
  \par\bigskip
  \noindent{\headfont 注~#1}\ifblank{#2}{}{\hspace{0.5em}{\itshape #2}}%
  \par\nobreak\smallskip\nobreak
}

% ∗ ∗ ∗ 分隔
\newcommand{\asterisms}{%
  \par\vspace{1.2ex}%
  \noindent\hspace*{2em}$\ast$\hspace{1.2em}$\ast$\hspace{1.2em}$\ast$%
  \par\vspace{1.2ex}%
}

% 无编号子目的目录条目
\newcommand{\tocsubsub}[1]{%
  \phantomsection
  \addcontentsline{toc}{subsubsection}{#1}%
}

% 部扉: 仿原书部题页
\newcommand{\partdivider}[3]{%
  \clearpage
  \thispagestyle{empty}%
  \phantomsection
  \addcontentsline{toc}{part}{#1\quad #3}%
  \vspace*{6.5cm}
  \begin{center}
    {\fontsize{16pt}{26pt}\selectfont\headfont #1\par}
    \vspace{0.9cm}
    {\fontsize{22pt}{34pt}\selectfont\headfont #2\par}
  \end{center}
  \clearpage
}

% 献辞页: 仿原书部题后的题献, 独占一页且不入目录
\newcommand{\dedication}[1]{%
  \thispagestyle{empty}%
  \vspace*{4cm}%
  \begin{center}#1\end{center}%
  \clearpage
}

% --------------------------------------------------------------
%  各卷"译者说明"的共用体例段(各卷的范围段自写, 后接本段)
% --------------------------------------------------------------
\newcommand{\CommonNoteBody}{%
为便于与原书及学界通行的打字稿页码对勘, 本册保留:

\begin{itemize}
  \item 原书全部章节编号(章~1--18, 节如~5.2、18.2, 子目如~18.1.1)与原书小节标题;
  \item 原书脚注的编号与角标——
        \begin{itemize}
          \item 前奏(第~1--4 章)与第一部的注在各页内重新编号, 正文中的上标即该
                页序号; 第一部紧随其后的括注(如“(3)”)是该部分的连续注号, 与第~12
                章“第一部分的注”里的“注~3”一一对应, 注文行首照原书印作“1(3)”式。
                其中三条注原书未标连续注号而改标日期(ch7 注~3 = 30 Septembre、
                ch9 注~11/12 = 8 août), 本册照录其日期;
          \item 第二部第~13--16~章各自重新编号, 分别为~1--41、1--86、1--132、1--29;
                第三部第~17~章自~1~起至~12, 第~18~章自~1~重新起算并延续至~794;
                末段自~\textbf{795(*)}~续至~1053(*)。各段与原书一一对应;
        \end{itemize}
  \item 页边码“p.~1”“p.~A3”“p.~1066”式小注——这是作者引用自己打字稿原稿时
        所用的页码, 与本排印本页码不同。前奏四章的边码各有前缀: 第~1 章作
        \texttt{A1…A6}, 第~2 章作 \texttt{P1…P65}, 第~3 章作 \texttt{L1…L56},
        第~4 章作罗马数字 \texttt{I…XXII}, 自第一部起作纯数字 \texttt{1…171};
  \item 作者用以标记段落的~$\diamond$~记号, 以及标示后来补写文字的上标
        ~\textsuperscript{$\diamond$}~记号;
  \item 原书大小注(“注~44'”“注~46$^1$”“注~171³”等)的编号与层级。
\end{itemize}

第~12 章“第一部分的注”把第~5 至第~11 章各页脚注汇为一编, 并给每条注配上节题
(如“12.19. 耶稣与十二使徒”), 因此部分注文与正文脚注重复——这是原书的体例, 非
重出之误。这一部分属随笔式的自省文字, 同一句法语在正文脚注与章末汇编注中
各有措辞, 本册有意保留各自的译笔, 不强求统一。

除上述标记外, 译文力求贴近作者驳杂、跳跃而辛辣的文风; 数学名词的译法见各卷
术语表(如~\emph{topos}~不译, \emph{faisceaux pervers}~作“反常层”)。文中的日期
括注(如“(5~月~30~日)”“(1984年3月补写)”)为作者历次补写时的纪日, 一律照录。文中
另有一些对注号、页码的引用(如“n°~1695”“打字稿第~227~页”)沿袭作者打字稿的编号
系统, 与本排印本不尽一致, 亦无法在本书中落实, 悉照原文录出。

全套五卷(另有全书合订一册)的分工是:

\begin{itemize}
  \item \textbf{卷一·前奏}——主题的呈现: 或四个乐章的前奏(原书 p.23--122);
  \item \textbf{卷二·第一部}——愚妄与更新(原书 p.123--254);
  \item \textbf{卷三·第二部}——葬礼(I): 或中国皇帝的新衣(原书 p.255--412);
  \item \textbf{卷四·第三部}——葬礼(II): 或阴阳之钥(原书 p.413--632);
  \item \textbf{卷五·第四部}——葬礼(III): 或四种运算(原书 p.633--929)。
\end{itemize}
}


% --------------------------------------------------------------
%  封面参数
% --------------------------------------------------------------
% ================================================================
%  《收获与播种》统一封面 —— TikZ 原生矢量绘制(不依赖任何外部图片)
%  母题: 暮色中的麦田 —— 犁沟向地平线的灭点收束, 前景麦穗剪影,
%        一位播种者把一把种子撒向天空, 种子沿弧线落入田里。
%  用法(在 \begin{document} 之前, 依次):
%      % ================================================================
%  《收获与播种》中文译本 —— 各卷共用导言区
%  由 vol-prelude / vol-part1 / vol-part2 / vol-part3 / vol-part4
%  以及全书合订 main-all 共同 \input。基类 mathtranslation.cls v3.2 不改。
%  --------------------------------------------------------------
%  样式均在此改写, 以贴近原书:
%    · 章标题仿原书 "5. Travail et découverte"(块状左对齐黑体)
%    · 节/小节: 黑体左对齐, 编号后带句点
%    · 脚注: 人工编号照录原书——前奏与第一部分页内重新编号,
%      第二至四部各自连续编号并带(*)角标; 第一部分另有 \fnoteP
%      处理 "(n)" 式连续注号
%    · 边码: 原打字稿页码 \origpage{N}(前奏作 A1、P1、L1、罗马数字等)
%    · ⋄ 记号: 段首补写标记 \pmark 与上标形式 \dmark
% ================================================================

% 题名含空格, 不能走类选项(会被 kvoptions 吃掉), 在此赋值
\renewcommand{\BookTitleEN}{Récoltes et Semailles}
\renewcommand{\OriginalAuthor}{Alexandre Grothendieck}
\renewcommand{\TranslationDate}{\today}

% --------------------------------------------------------------
%  罕用字后备(如"祐"): Fandol 不载时退到 SimSun
% --------------------------------------------------------------
\setCJKfallbackfamilyfont{\CJKrmdefault}{{SimSun}}

% --------------------------------------------------------------
%  hyperref 元信息(各卷在 % ================================================================
%  《收获与播种》中文译本 —— 各卷共用导言区
%  由 vol-prelude / vol-part1 / vol-part2 / vol-part3 / vol-part4
%  以及全书合订 main-all 共同 \input。基类 mathtranslation.cls v3.2 不改。
%  --------------------------------------------------------------
%  样式均在此改写, 以贴近原书:
%    · 章标题仿原书 "5. Travail et découverte"(块状左对齐黑体)
%    · 节/小节: 黑体左对齐, 编号后带句点
%    · 脚注: 人工编号照录原书——前奏与第一部分页内重新编号,
%      第二至四部各自连续编号并带(*)角标; 第一部分另有 \fnoteP
%      处理 "(n)" 式连续注号
%    · 边码: 原打字稿页码 \origpage{N}(前奏作 A1、P1、L1、罗马数字等)
%    · ⋄ 记号: 段首补写标记 \pmark 与上标形式 \dmark
% ================================================================

% 题名含空格, 不能走类选项(会被 kvoptions 吃掉), 在此赋值
\renewcommand{\BookTitleEN}{Récoltes et Semailles}
\renewcommand{\OriginalAuthor}{Alexandre Grothendieck}
\renewcommand{\TranslationDate}{\today}

% --------------------------------------------------------------
%  罕用字后备(如"祐"): Fandol 不载时退到 SimSun
% --------------------------------------------------------------
\setCJKfallbackfamilyfont{\CJKrmdefault}{{SimSun}}

% --------------------------------------------------------------
%  hyperref 元信息(各卷在 \input{common} 之后可用 \hypersetup 覆写)
% --------------------------------------------------------------
\hypersetup{
  pdftitle={收获与播种},
  pdfauthor={Alexandre Grothendieck / 数译 译},
  pdfsubject={Récoltes et Semailles — Réflexions et témoignage sur un passé de mathématicien}
}

% --------------------------------------------------------------
%  标题样式: 仿原书 —— 黑体左对齐, 编号后带句点
% --------------------------------------------------------------
\titleformat{\section}[block]
  {\normalfont\Large\headfont}{\thesection.}{0.5em}{}
\titlespacing*{\section}{0pt}{3.2ex plus 1.2ex minus .2ex}{1.6ex plus .2ex}
\titleformat{\subsection}[block]
  {\normalfont\large\headfont}{\thesubsection.}{0.5em}{}
\titlespacing*{\subsection}{0pt}{2.6ex plus 1ex minus .2ex}{1.2ex plus .2ex}
\titleformat{\subsubsection}[block]
  {\normalfont\normalsize\headfont}{\thesubsubsection.}{0.5em}{}
\titlespacing*{\subsubsection}{0pt}{2.2ex plus .8ex minus .2ex}{1ex plus .2ex}
\titleformat{name=\subsubsection,numberless}[block]
  {\normalfont\normalsize\headfont}{}{0em}{}
\setcounter{secnumdepth}{3}
\setcounter{tocdepth}{3}
% 为含长外文书名与行内公式的段落提供必要的断行余量
\setlength{\emergencystretch}{3em}

% 章标题样式: 原书章标题与节同级排印("13. A) HERITAGE..."),
% 故章标题取块状左对齐黑体, 字号介于节与部之间
\titleformat{\chapter}[block]
  {\normalfont\LARGE\headfont}{\thechapter.}{0.5em}{}
\titlespacing*{\chapter}{0pt}{3.6ex plus 1.4ex minus .2ex}{2ex plus .2ex}
% 目录中章号照原书用阿拉伯数字("13. A) ..." 而非 "第十三章")
\renewcommand{\CTEXthechapter}{\thechapter\hspace{.3em}}

% --------------------------------------------------------------
%  目录观感: 条目与引导点之间、引导点与页码之间各留 0.5em
%  内核 \@dottedtocline 让 \leaders 从标题末字起排, 第一个点几乎
%  贴着末字, 读起来像多标了一个句点; 页码前也没有余量。两处各补
%  0.5em。
% --------------------------------------------------------------
\makeatletter
\patchcmd{\@dottedtocline}{\nobreak\leaders}{\nobreak\hspace{0.5em}\leaders}%
  {\typeout{MT: TOC 引导点前已留 0.5em}}%
  {\typeout{MT-WARN: TOC 引导点前间距补丁未生效}}
\patchcmd{\@dottedtocline}{\hfill\nobreak\hb@xt@\@pnumwidth}%
  {\hfill\nobreak\hspace{0.5em}\hb@xt@\@pnumwidth}%
  {\typeout{MT: TOC 页码前已留 0.5em}}%
  {\typeout{MT-WARN: TOC 页码前间距补丁未生效}}
\makeatother

% --------------------------------------------------------------
%  正文半角标点后的间距
%  本书体例用半角 , ; : ? ! 与 ASCII 括号, 由 xeCJK 按 HalfRight
%  (右半留白) 处理。放宽到 0.45em, 与全角标点右半留白大致齐平。
% --------------------------------------------------------------
\xeCJKsetup{CJKecglue={\hskip 0.45em plus 0.06em minus 0.09em}}

% --------------------------------------------------------------
%  页眉页脚: 偶数页眉右对齐斜体节题; 页码页脚外侧
% --------------------------------------------------------------
\fancyhf{}
\fancyhead[RE]{\small\itshape\rightmark}
\fancyfoot[LE,RO]{\small\thepage}
\renewcommand{\headrulewidth}{0pt}
\renewcommand{\sectionmark}[1]{\markright{\thesection.\ #1}}
\fancypagestyle{plain}{%
  \fancyhf{}%
  \fancyfoot[LE,RO]{\small \thepage}%
  \renewcommand{\headrulewidth}{0pt}%
}

% --------------------------------------------------------------
%  原书版式记号宏
% --------------------------------------------------------------
% 边码: 原打字稿页码(G. 手稿系统里的自引页码), 排在页边
\usepackage{marginnote}
\newcommand{\origpage}[1]{\marginnote{\normalfont\scriptsize p.~#1}}

% ⋄ 记号: 段首/行内正常大小; 上标形式为作者后补文字的位置标记
\DeclareRobustCommand{\pmark}{\ensuremath{\diamond}\hspace{0.15em}}
\DeclareRobustCommand{\dmark}{\textsuperscript{\ensuremath{\diamond}}}

% 脚注人工编号: \fnote{795}{(*)}{文本} —— 编号与(*)角标均照录原书
\newcommand{\fnote}[3]{%
  \begingroup
  \renewcommand{\thefootnote}{#1#2}%
  \footnote{#3}%
  \endgroup
}

% 第一部分的脚注: 页内序号作上标(#1), 注文行首印 "序号(连续注号)",
% 正文在序号之后另行括注同一连续注号 —— \fnoteP{1}{(3)}{注文}
% 原书有三条注无连续注号而改标日期(ch07 注3 = 30 Septembre,
% ch09 注11/12 = 8 août), 第二参数留空, 日期写进注文之首。
\newcommand{\fnoteP}[3]{%
  \begingroup
  \renewcommand{\thefootnote}{#1}\footnotemark
  \renewcommand{\thefootnote}{#1#2}\footnotetext{#3}%
  \endgroup
}

% 正文呼号与注文行首标号不同(原书照录)
\newcommand{\fnotevariant}[4]{%
  \begingroup
  \renewcommand{\thefootnote}{#1#2}\footnotemark
  \renewcommand{\thefootnote}{#1#3}\footnotetext{#4}%
  \endgroup
}

% 大注块标题: \notehead{171³}{斜体评论} —— 正文随后另起
\newcommand{\notehead}[2]{%
  \par\bigskip
  \noindent{\headfont 注~#1}\ifblank{#2}{}{\hspace{0.5em}{\itshape #2}}%
  \par\nobreak\smallskip\nobreak
}

% ∗ ∗ ∗ 分隔
\newcommand{\asterisms}{%
  \par\vspace{1.2ex}%
  \noindent\hspace*{2em}$\ast$\hspace{1.2em}$\ast$\hspace{1.2em}$\ast$%
  \par\vspace{1.2ex}%
}

% 无编号子目的目录条目
\newcommand{\tocsubsub}[1]{%
  \phantomsection
  \addcontentsline{toc}{subsubsection}{#1}%
}

% 部扉: 仿原书部题页
\newcommand{\partdivider}[3]{%
  \clearpage
  \thispagestyle{empty}%
  \phantomsection
  \addcontentsline{toc}{part}{#1\quad #3}%
  \vspace*{6.5cm}
  \begin{center}
    {\fontsize{16pt}{26pt}\selectfont\headfont #1\par}
    \vspace{0.9cm}
    {\fontsize{22pt}{34pt}\selectfont\headfont #2\par}
  \end{center}
  \clearpage
}

% 献辞页: 仿原书部题后的题献, 独占一页且不入目录
\newcommand{\dedication}[1]{%
  \thispagestyle{empty}%
  \vspace*{4cm}%
  \begin{center}#1\end{center}%
  \clearpage
}

% --------------------------------------------------------------
%  各卷"译者说明"的共用体例段(各卷的范围段自写, 后接本段)
% --------------------------------------------------------------
\newcommand{\CommonNoteBody}{%
为便于与原书及学界通行的打字稿页码对勘, 本册保留:

\begin{itemize}
  \item 原书全部章节编号(章~1--18, 节如~5.2、18.2, 子目如~18.1.1)与原书小节标题;
  \item 原书脚注的编号与角标——
        \begin{itemize}
          \item 前奏(第~1--4 章)与第一部的注在各页内重新编号, 正文中的上标即该
                页序号; 第一部紧随其后的括注(如“(3)”)是该部分的连续注号, 与第~12
                章“第一部分的注”里的“注~3”一一对应, 注文行首照原书印作“1(3)”式。
                其中三条注原书未标连续注号而改标日期(ch7 注~3 = 30 Septembre、
                ch9 注~11/12 = 8 août), 本册照录其日期;
          \item 第二部第~13--16~章各自重新编号, 分别为~1--41、1--86、1--132、1--29;
                第三部第~17~章自~1~起至~12, 第~18~章自~1~重新起算并延续至~794;
                末段自~\textbf{795(*)}~续至~1053(*)。各段与原书一一对应;
        \end{itemize}
  \item 页边码“p.~1”“p.~A3”“p.~1066”式小注——这是作者引用自己打字稿原稿时
        所用的页码, 与本排印本页码不同。前奏四章的边码各有前缀: 第~1 章作
        \texttt{A1…A6}, 第~2 章作 \texttt{P1…P65}, 第~3 章作 \texttt{L1…L56},
        第~4 章作罗马数字 \texttt{I…XXII}, 自第一部起作纯数字 \texttt{1…171};
  \item 作者用以标记段落的~$\diamond$~记号, 以及标示后来补写文字的上标
        ~\textsuperscript{$\diamond$}~记号;
  \item 原书大小注(“注~44'”“注~46$^1$”“注~171³”等)的编号与层级。
\end{itemize}

第~12 章“第一部分的注”把第~5 至第~11 章各页脚注汇为一编, 并给每条注配上节题
(如“12.19. 耶稣与十二使徒”), 因此部分注文与正文脚注重复——这是原书的体例, 非
重出之误。这一部分属随笔式的自省文字, 同一句法语在正文脚注与章末汇编注中
各有措辞, 本册有意保留各自的译笔, 不强求统一。

除上述标记外, 译文力求贴近作者驳杂、跳跃而辛辣的文风; 数学名词的译法见各卷
术语表(如~\emph{topos}~不译, \emph{faisceaux pervers}~作“反常层”)。文中的日期
括注(如“(5~月~30~日)”“(1984年3月补写)”)为作者历次补写时的纪日, 一律照录。文中
另有一些对注号、页码的引用(如“n°~1695”“打字稿第~227~页”)沿袭作者打字稿的编号
系统, 与本排印本不尽一致, 亦无法在本书中落实, 悉照原文录出。

全套五卷(另有全书合订一册)的分工是:

\begin{itemize}
  \item \textbf{卷一·前奏}——主题的呈现: 或四个乐章的前奏(原书 p.23--122);
  \item \textbf{卷二·第一部}——愚妄与更新(原书 p.123--254);
  \item \textbf{卷三·第二部}——葬礼(I): 或中国皇帝的新衣(原书 p.255--412);
  \item \textbf{卷四·第三部}——葬礼(II): 或阴阳之钥(原书 p.413--632);
  \item \textbf{卷五·第四部}——葬礼(III): 或四种运算(原书 p.633--929)。
\end{itemize}
}
 之后可用 \hypersetup 覆写)
% --------------------------------------------------------------
\hypersetup{
  pdftitle={收获与播种},
  pdfauthor={Alexandre Grothendieck / 数译 译},
  pdfsubject={Récoltes et Semailles — Réflexions et témoignage sur un passé de mathématicien}
}

% --------------------------------------------------------------
%  标题样式: 仿原书 —— 黑体左对齐, 编号后带句点
% --------------------------------------------------------------
\titleformat{\section}[block]
  {\normalfont\Large\headfont}{\thesection.}{0.5em}{}
\titlespacing*{\section}{0pt}{3.2ex plus 1.2ex minus .2ex}{1.6ex plus .2ex}
\titleformat{\subsection}[block]
  {\normalfont\large\headfont}{\thesubsection.}{0.5em}{}
\titlespacing*{\subsection}{0pt}{2.6ex plus 1ex minus .2ex}{1.2ex plus .2ex}
\titleformat{\subsubsection}[block]
  {\normalfont\normalsize\headfont}{\thesubsubsection.}{0.5em}{}
\titlespacing*{\subsubsection}{0pt}{2.2ex plus .8ex minus .2ex}{1ex plus .2ex}
\titleformat{name=\subsubsection,numberless}[block]
  {\normalfont\normalsize\headfont}{}{0em}{}
\setcounter{secnumdepth}{3}
\setcounter{tocdepth}{3}
% 为含长外文书名与行内公式的段落提供必要的断行余量
\setlength{\emergencystretch}{3em}

% 章标题样式: 原书章标题与节同级排印("13. A) HERITAGE..."),
% 故章标题取块状左对齐黑体, 字号介于节与部之间
\titleformat{\chapter}[block]
  {\normalfont\LARGE\headfont}{\thechapter.}{0.5em}{}
\titlespacing*{\chapter}{0pt}{3.6ex plus 1.4ex minus .2ex}{2ex plus .2ex}
% 目录中章号照原书用阿拉伯数字("13. A) ..." 而非 "第十三章")
\renewcommand{\CTEXthechapter}{\thechapter\hspace{.3em}}

% --------------------------------------------------------------
%  目录观感: 条目与引导点之间、引导点与页码之间各留 0.5em
%  内核 \@dottedtocline 让 \leaders 从标题末字起排, 第一个点几乎
%  贴着末字, 读起来像多标了一个句点; 页码前也没有余量。两处各补
%  0.5em。
% --------------------------------------------------------------
\makeatletter
\patchcmd{\@dottedtocline}{\nobreak\leaders}{\nobreak\hspace{0.5em}\leaders}%
  {\typeout{MT: TOC 引导点前已留 0.5em}}%
  {\typeout{MT-WARN: TOC 引导点前间距补丁未生效}}
\patchcmd{\@dottedtocline}{\hfill\nobreak\hb@xt@\@pnumwidth}%
  {\hfill\nobreak\hspace{0.5em}\hb@xt@\@pnumwidth}%
  {\typeout{MT: TOC 页码前已留 0.5em}}%
  {\typeout{MT-WARN: TOC 页码前间距补丁未生效}}
\makeatother

% --------------------------------------------------------------
%  正文半角标点后的间距
%  本书体例用半角 , ; : ? ! 与 ASCII 括号, 由 xeCJK 按 HalfRight
%  (右半留白) 处理。放宽到 0.45em, 与全角标点右半留白大致齐平。
% --------------------------------------------------------------
\xeCJKsetup{CJKecglue={\hskip 0.45em plus 0.06em minus 0.09em}}

% --------------------------------------------------------------
%  页眉页脚: 偶数页眉右对齐斜体节题; 页码页脚外侧
% --------------------------------------------------------------
\fancyhf{}
\fancyhead[RE]{\small\itshape\rightmark}
\fancyfoot[LE,RO]{\small\thepage}
\renewcommand{\headrulewidth}{0pt}
\renewcommand{\sectionmark}[1]{\markright{\thesection.\ #1}}
\fancypagestyle{plain}{%
  \fancyhf{}%
  \fancyfoot[LE,RO]{\small \thepage}%
  \renewcommand{\headrulewidth}{0pt}%
}

% --------------------------------------------------------------
%  原书版式记号宏
% --------------------------------------------------------------
% 边码: 原打字稿页码(G. 手稿系统里的自引页码), 排在页边
\usepackage{marginnote}
\newcommand{\origpage}[1]{\marginnote{\normalfont\scriptsize p.~#1}}

% ⋄ 记号: 段首/行内正常大小; 上标形式为作者后补文字的位置标记
\DeclareRobustCommand{\pmark}{\ensuremath{\diamond}\hspace{0.15em}}
\DeclareRobustCommand{\dmark}{\textsuperscript{\ensuremath{\diamond}}}

% 脚注人工编号: \fnote{795}{(*)}{文本} —— 编号与(*)角标均照录原书
\newcommand{\fnote}[3]{%
  \begingroup
  \renewcommand{\thefootnote}{#1#2}%
  \footnote{#3}%
  \endgroup
}

% 第一部分的脚注: 页内序号作上标(#1), 注文行首印 "序号(连续注号)",
% 正文在序号之后另行括注同一连续注号 —— \fnoteP{1}{(3)}{注文}
% 原书有三条注无连续注号而改标日期(ch07 注3 = 30 Septembre,
% ch09 注11/12 = 8 août), 第二参数留空, 日期写进注文之首。
\newcommand{\fnoteP}[3]{%
  \begingroup
  \renewcommand{\thefootnote}{#1}\footnotemark
  \renewcommand{\thefootnote}{#1#2}\footnotetext{#3}%
  \endgroup
}

% 正文呼号与注文行首标号不同(原书照录)
\newcommand{\fnotevariant}[4]{%
  \begingroup
  \renewcommand{\thefootnote}{#1#2}\footnotemark
  \renewcommand{\thefootnote}{#1#3}\footnotetext{#4}%
  \endgroup
}

% 大注块标题: \notehead{171³}{斜体评论} —— 正文随后另起
\newcommand{\notehead}[2]{%
  \par\bigskip
  \noindent{\headfont 注~#1}\ifblank{#2}{}{\hspace{0.5em}{\itshape #2}}%
  \par\nobreak\smallskip\nobreak
}

% ∗ ∗ ∗ 分隔
\newcommand{\asterisms}{%
  \par\vspace{1.2ex}%
  \noindent\hspace*{2em}$\ast$\hspace{1.2em}$\ast$\hspace{1.2em}$\ast$%
  \par\vspace{1.2ex}%
}

% 无编号子目的目录条目
\newcommand{\tocsubsub}[1]{%
  \phantomsection
  \addcontentsline{toc}{subsubsection}{#1}%
}

% 部扉: 仿原书部题页
\newcommand{\partdivider}[3]{%
  \clearpage
  \thispagestyle{empty}%
  \phantomsection
  \addcontentsline{toc}{part}{#1\quad #3}%
  \vspace*{6.5cm}
  \begin{center}
    {\fontsize{16pt}{26pt}\selectfont\headfont #1\par}
    \vspace{0.9cm}
    {\fontsize{22pt}{34pt}\selectfont\headfont #2\par}
  \end{center}
  \clearpage
}

% 献辞页: 仿原书部题后的题献, 独占一页且不入目录
\newcommand{\dedication}[1]{%
  \thispagestyle{empty}%
  \vspace*{4cm}%
  \begin{center}#1\end{center}%
  \clearpage
}

% --------------------------------------------------------------
%  各卷"译者说明"的共用体例段(各卷的范围段自写, 后接本段)
% --------------------------------------------------------------
\newcommand{\CommonNoteBody}{%
为便于与原书及学界通行的打字稿页码对勘, 本册保留:

\begin{itemize}
  \item 原书全部章节编号(章~1--18, 节如~5.2、18.2, 子目如~18.1.1)与原书小节标题;
  \item 原书脚注的编号与角标——
        \begin{itemize}
          \item 前奏(第~1--4 章)与第一部的注在各页内重新编号, 正文中的上标即该
                页序号; 第一部紧随其后的括注(如“(3)”)是该部分的连续注号, 与第~12
                章“第一部分的注”里的“注~3”一一对应, 注文行首照原书印作“1(3)”式。
                其中三条注原书未标连续注号而改标日期(ch7 注~3 = 30 Septembre、
                ch9 注~11/12 = 8 août), 本册照录其日期;
          \item 第二部第~13--16~章各自重新编号, 分别为~1--41、1--86、1--132、1--29;
                第三部第~17~章自~1~起至~12, 第~18~章自~1~重新起算并延续至~794;
                末段自~\textbf{795(*)}~续至~1053(*)。各段与原书一一对应;
        \end{itemize}
  \item 页边码“p.~1”“p.~A3”“p.~1066”式小注——这是作者引用自己打字稿原稿时
        所用的页码, 与本排印本页码不同。前奏四章的边码各有前缀: 第~1 章作
        \texttt{A1…A6}, 第~2 章作 \texttt{P1…P65}, 第~3 章作 \texttt{L1…L56},
        第~4 章作罗马数字 \texttt{I…XXII}, 自第一部起作纯数字 \texttt{1…171};
  \item 作者用以标记段落的~$\diamond$~记号, 以及标示后来补写文字的上标
        ~\textsuperscript{$\diamond$}~记号;
  \item 原书大小注(“注~44'”“注~46$^1$”“注~171³”等)的编号与层级。
\end{itemize}

第~12 章“第一部分的注”把第~5 至第~11 章各页脚注汇为一编, 并给每条注配上节题
(如“12.19. 耶稣与十二使徒”), 因此部分注文与正文脚注重复——这是原书的体例, 非
重出之误。这一部分属随笔式的自省文字, 同一句法语在正文脚注与章末汇编注中
各有措辞, 本册有意保留各自的译笔, 不强求统一。

除上述标记外, 译文力求贴近作者驳杂、跳跃而辛辣的文风; 数学名词的译法见各卷
术语表(如~\emph{topos}~不译, \emph{faisceaux pervers}~作“反常层”)。文中的日期
括注(如“(5~月~30~日)”“(1984年3月补写)”)为作者历次补写时的纪日, 一律照录。文中
另有一些对注号、页码的引用(如“n°~1695”“打字稿第~227~页”)沿袭作者打字稿的编号
系统, 与本排印本不尽一致, 亦无法在本书中落实, 悉照原文录出。

全套五卷(另有全书合订一册)的分工是:

\begin{itemize}
  \item \textbf{卷一·前奏}——主题的呈现: 或四个乐章的前奏(原书 p.23--122);
  \item \textbf{卷二·第一部}——愚妄与更新(原书 p.123--254);
  \item \textbf{卷三·第二部}——葬礼(I): 或中国皇帝的新衣(原书 p.255--412);
  \item \textbf{卷四·第三部}——葬礼(II): 或阴阳之钥(原书 p.413--632);
  \item \textbf{卷五·第四部}——葬礼(III): 或四种运算(原书 p.633--929)。
\end{itemize}
}
                       % 共用导言区
%      % ================================================================
%  《收获与播种》统一封面 —— TikZ 原生矢量绘制(不依赖任何外部图片)
%  母题: 暮色中的麦田 —— 犁沟向地平线的灭点收束, 前景麦穗剪影,
%        一位播种者把一把种子撒向天空, 种子沿弧线落入田里。
%  用法(在 \begin{document} 之前, 依次):
%      % ================================================================
%  《收获与播种》中文译本 —— 各卷共用导言区
%  由 vol-prelude / vol-part1 / vol-part2 / vol-part3 / vol-part4
%  以及全书合订 main-all 共同 \input。基类 mathtranslation.cls v3.2 不改。
%  --------------------------------------------------------------
%  样式均在此改写, 以贴近原书:
%    · 章标题仿原书 "5. Travail et découverte"(块状左对齐黑体)
%    · 节/小节: 黑体左对齐, 编号后带句点
%    · 脚注: 人工编号照录原书——前奏与第一部分页内重新编号,
%      第二至四部各自连续编号并带(*)角标; 第一部分另有 \fnoteP
%      处理 "(n)" 式连续注号
%    · 边码: 原打字稿页码 \origpage{N}(前奏作 A1、P1、L1、罗马数字等)
%    · ⋄ 记号: 段首补写标记 \pmark 与上标形式 \dmark
% ================================================================

% 题名含空格, 不能走类选项(会被 kvoptions 吃掉), 在此赋值
\renewcommand{\BookTitleEN}{Récoltes et Semailles}
\renewcommand{\OriginalAuthor}{Alexandre Grothendieck}
\renewcommand{\TranslationDate}{\today}

% --------------------------------------------------------------
%  罕用字后备(如"祐"): Fandol 不载时退到 SimSun
% --------------------------------------------------------------
\setCJKfallbackfamilyfont{\CJKrmdefault}{{SimSun}}

% --------------------------------------------------------------
%  hyperref 元信息(各卷在 \input{common} 之后可用 \hypersetup 覆写)
% --------------------------------------------------------------
\hypersetup{
  pdftitle={收获与播种},
  pdfauthor={Alexandre Grothendieck / 数译 译},
  pdfsubject={Récoltes et Semailles — Réflexions et témoignage sur un passé de mathématicien}
}

% --------------------------------------------------------------
%  标题样式: 仿原书 —— 黑体左对齐, 编号后带句点
% --------------------------------------------------------------
\titleformat{\section}[block]
  {\normalfont\Large\headfont}{\thesection.}{0.5em}{}
\titlespacing*{\section}{0pt}{3.2ex plus 1.2ex minus .2ex}{1.6ex plus .2ex}
\titleformat{\subsection}[block]
  {\normalfont\large\headfont}{\thesubsection.}{0.5em}{}
\titlespacing*{\subsection}{0pt}{2.6ex plus 1ex minus .2ex}{1.2ex plus .2ex}
\titleformat{\subsubsection}[block]
  {\normalfont\normalsize\headfont}{\thesubsubsection.}{0.5em}{}
\titlespacing*{\subsubsection}{0pt}{2.2ex plus .8ex minus .2ex}{1ex plus .2ex}
\titleformat{name=\subsubsection,numberless}[block]
  {\normalfont\normalsize\headfont}{}{0em}{}
\setcounter{secnumdepth}{3}
\setcounter{tocdepth}{3}
% 为含长外文书名与行内公式的段落提供必要的断行余量
\setlength{\emergencystretch}{3em}

% 章标题样式: 原书章标题与节同级排印("13. A) HERITAGE..."),
% 故章标题取块状左对齐黑体, 字号介于节与部之间
\titleformat{\chapter}[block]
  {\normalfont\LARGE\headfont}{\thechapter.}{0.5em}{}
\titlespacing*{\chapter}{0pt}{3.6ex plus 1.4ex minus .2ex}{2ex plus .2ex}
% 目录中章号照原书用阿拉伯数字("13. A) ..." 而非 "第十三章")
\renewcommand{\CTEXthechapter}{\thechapter\hspace{.3em}}

% --------------------------------------------------------------
%  目录观感: 条目与引导点之间、引导点与页码之间各留 0.5em
%  内核 \@dottedtocline 让 \leaders 从标题末字起排, 第一个点几乎
%  贴着末字, 读起来像多标了一个句点; 页码前也没有余量。两处各补
%  0.5em。
% --------------------------------------------------------------
\makeatletter
\patchcmd{\@dottedtocline}{\nobreak\leaders}{\nobreak\hspace{0.5em}\leaders}%
  {\typeout{MT: TOC 引导点前已留 0.5em}}%
  {\typeout{MT-WARN: TOC 引导点前间距补丁未生效}}
\patchcmd{\@dottedtocline}{\hfill\nobreak\hb@xt@\@pnumwidth}%
  {\hfill\nobreak\hspace{0.5em}\hb@xt@\@pnumwidth}%
  {\typeout{MT: TOC 页码前已留 0.5em}}%
  {\typeout{MT-WARN: TOC 页码前间距补丁未生效}}
\makeatother

% --------------------------------------------------------------
%  正文半角标点后的间距
%  本书体例用半角 , ; : ? ! 与 ASCII 括号, 由 xeCJK 按 HalfRight
%  (右半留白) 处理。放宽到 0.45em, 与全角标点右半留白大致齐平。
% --------------------------------------------------------------
\xeCJKsetup{CJKecglue={\hskip 0.45em plus 0.06em minus 0.09em}}

% --------------------------------------------------------------
%  页眉页脚: 偶数页眉右对齐斜体节题; 页码页脚外侧
% --------------------------------------------------------------
\fancyhf{}
\fancyhead[RE]{\small\itshape\rightmark}
\fancyfoot[LE,RO]{\small\thepage}
\renewcommand{\headrulewidth}{0pt}
\renewcommand{\sectionmark}[1]{\markright{\thesection.\ #1}}
\fancypagestyle{plain}{%
  \fancyhf{}%
  \fancyfoot[LE,RO]{\small \thepage}%
  \renewcommand{\headrulewidth}{0pt}%
}

% --------------------------------------------------------------
%  原书版式记号宏
% --------------------------------------------------------------
% 边码: 原打字稿页码(G. 手稿系统里的自引页码), 排在页边
\usepackage{marginnote}
\newcommand{\origpage}[1]{\marginnote{\normalfont\scriptsize p.~#1}}

% ⋄ 记号: 段首/行内正常大小; 上标形式为作者后补文字的位置标记
\DeclareRobustCommand{\pmark}{\ensuremath{\diamond}\hspace{0.15em}}
\DeclareRobustCommand{\dmark}{\textsuperscript{\ensuremath{\diamond}}}

% 脚注人工编号: \fnote{795}{(*)}{文本} —— 编号与(*)角标均照录原书
\newcommand{\fnote}[3]{%
  \begingroup
  \renewcommand{\thefootnote}{#1#2}%
  \footnote{#3}%
  \endgroup
}

% 第一部分的脚注: 页内序号作上标(#1), 注文行首印 "序号(连续注号)",
% 正文在序号之后另行括注同一连续注号 —— \fnoteP{1}{(3)}{注文}
% 原书有三条注无连续注号而改标日期(ch07 注3 = 30 Septembre,
% ch09 注11/12 = 8 août), 第二参数留空, 日期写进注文之首。
\newcommand{\fnoteP}[3]{%
  \begingroup
  \renewcommand{\thefootnote}{#1}\footnotemark
  \renewcommand{\thefootnote}{#1#2}\footnotetext{#3}%
  \endgroup
}

% 正文呼号与注文行首标号不同(原书照录)
\newcommand{\fnotevariant}[4]{%
  \begingroup
  \renewcommand{\thefootnote}{#1#2}\footnotemark
  \renewcommand{\thefootnote}{#1#3}\footnotetext{#4}%
  \endgroup
}

% 大注块标题: \notehead{171³}{斜体评论} —— 正文随后另起
\newcommand{\notehead}[2]{%
  \par\bigskip
  \noindent{\headfont 注~#1}\ifblank{#2}{}{\hspace{0.5em}{\itshape #2}}%
  \par\nobreak\smallskip\nobreak
}

% ∗ ∗ ∗ 分隔
\newcommand{\asterisms}{%
  \par\vspace{1.2ex}%
  \noindent\hspace*{2em}$\ast$\hspace{1.2em}$\ast$\hspace{1.2em}$\ast$%
  \par\vspace{1.2ex}%
}

% 无编号子目的目录条目
\newcommand{\tocsubsub}[1]{%
  \phantomsection
  \addcontentsline{toc}{subsubsection}{#1}%
}

% 部扉: 仿原书部题页
\newcommand{\partdivider}[3]{%
  \clearpage
  \thispagestyle{empty}%
  \phantomsection
  \addcontentsline{toc}{part}{#1\quad #3}%
  \vspace*{6.5cm}
  \begin{center}
    {\fontsize{16pt}{26pt}\selectfont\headfont #1\par}
    \vspace{0.9cm}
    {\fontsize{22pt}{34pt}\selectfont\headfont #2\par}
  \end{center}
  \clearpage
}

% 献辞页: 仿原书部题后的题献, 独占一页且不入目录
\newcommand{\dedication}[1]{%
  \thispagestyle{empty}%
  \vspace*{4cm}%
  \begin{center}#1\end{center}%
  \clearpage
}

% --------------------------------------------------------------
%  各卷"译者说明"的共用体例段(各卷的范围段自写, 后接本段)
% --------------------------------------------------------------
\newcommand{\CommonNoteBody}{%
为便于与原书及学界通行的打字稿页码对勘, 本册保留:

\begin{itemize}
  \item 原书全部章节编号(章~1--18, 节如~5.2、18.2, 子目如~18.1.1)与原书小节标题;
  \item 原书脚注的编号与角标——
        \begin{itemize}
          \item 前奏(第~1--4 章)与第一部的注在各页内重新编号, 正文中的上标即该
                页序号; 第一部紧随其后的括注(如“(3)”)是该部分的连续注号, 与第~12
                章“第一部分的注”里的“注~3”一一对应, 注文行首照原书印作“1(3)”式。
                其中三条注原书未标连续注号而改标日期(ch7 注~3 = 30 Septembre、
                ch9 注~11/12 = 8 août), 本册照录其日期;
          \item 第二部第~13--16~章各自重新编号, 分别为~1--41、1--86、1--132、1--29;
                第三部第~17~章自~1~起至~12, 第~18~章自~1~重新起算并延续至~794;
                末段自~\textbf{795(*)}~续至~1053(*)。各段与原书一一对应;
        \end{itemize}
  \item 页边码“p.~1”“p.~A3”“p.~1066”式小注——这是作者引用自己打字稿原稿时
        所用的页码, 与本排印本页码不同。前奏四章的边码各有前缀: 第~1 章作
        \texttt{A1…A6}, 第~2 章作 \texttt{P1…P65}, 第~3 章作 \texttt{L1…L56},
        第~4 章作罗马数字 \texttt{I…XXII}, 自第一部起作纯数字 \texttt{1…171};
  \item 作者用以标记段落的~$\diamond$~记号, 以及标示后来补写文字的上标
        ~\textsuperscript{$\diamond$}~记号;
  \item 原书大小注(“注~44'”“注~46$^1$”“注~171³”等)的编号与层级。
\end{itemize}

第~12 章“第一部分的注”把第~5 至第~11 章各页脚注汇为一编, 并给每条注配上节题
(如“12.19. 耶稣与十二使徒”), 因此部分注文与正文脚注重复——这是原书的体例, 非
重出之误。这一部分属随笔式的自省文字, 同一句法语在正文脚注与章末汇编注中
各有措辞, 本册有意保留各自的译笔, 不强求统一。

除上述标记外, 译文力求贴近作者驳杂、跳跃而辛辣的文风; 数学名词的译法见各卷
术语表(如~\emph{topos}~不译, \emph{faisceaux pervers}~作“反常层”)。文中的日期
括注(如“(5~月~30~日)”“(1984年3月补写)”)为作者历次补写时的纪日, 一律照录。文中
另有一些对注号、页码的引用(如“n°~1695”“打字稿第~227~页”)沿袭作者打字稿的编号
系统, 与本排印本不尽一致, 亦无法在本书中落实, 悉照原文录出。

全套五卷(另有全书合订一册)的分工是:

\begin{itemize}
  \item \textbf{卷一·前奏}——主题的呈现: 或四个乐章的前奏(原书 p.23--122);
  \item \textbf{卷二·第一部}——愚妄与更新(原书 p.123--254);
  \item \textbf{卷三·第二部}——葬礼(I): 或中国皇帝的新衣(原书 p.255--412);
  \item \textbf{卷四·第三部}——葬礼(II): 或阴阳之钥(原书 p.413--632);
  \item \textbf{卷五·第四部}——葬礼(III): 或四种运算(原书 p.633--929)。
\end{itemize}
}
                       % 共用导言区
%      % ================================================================
%  《收获与播种》统一封面 —— TikZ 原生矢量绘制(不依赖任何外部图片)
%  母题: 暮色中的麦田 —— 犁沟向地平线的灭点收束, 前景麦穗剪影,
%        一位播种者把一把种子撒向天空, 种子沿弧线落入田里。
%  用法(在 \begin{document} 之前, 依次):
%      \input{common}                       % 共用导言区
%      \input{cover-art}                    % 本文件
%      \CoverSetup{prelude}{卷一}{前奏}%
%                 {主题的呈现——或四个乐章的前奏}{p.23–122}
%  正文中再 \makecover 即可。配色关键字:
%      prelude / part1 / part2 / part3 / part4 / all
%  五卷共用同一幅画, 只换配色与卷次文字, 以保证"一套书"的观感。
% ================================================================

\RequirePackage{tikz}

% --------------------------------------------------------------
%  配色: 六套(前奏 / 四部 / 全书), 各为同一暮色母题的变调
% --------------------------------------------------------------
\newcommand{\CoverPalette}[1]{%
  \ifstrequal{#1}{prelude}{%
    \definecolor{cvSkyTop}{RGB}{10,26,48}%
    \definecolor{cvSkyMid}{RGB}{46,84,116}%
    \definecolor{cvSkyBot}{RGB}{232,164,92}%
    \definecolor{cvSun}{RGB}{255,224,168}%
    \definecolor{cvFieldTop}{RGB}{150,110,52}%
    \definecolor{cvFieldBot}{RGB}{58,36,18}%
    \definecolor{cvFurrow}{RGB}{40,22,10}%
    \definecolor{cvSil}{RGB}{30,18,10}%
    \definecolor{cvCream}{RGB}{247,240,228}%
  }{%
  \ifstrequal{#1}{part1}{%
    \definecolor{cvSkyTop}{RGB}{22,70,100}%
    \definecolor{cvSkyMid}{RGB}{92,150,168}%
    \definecolor{cvSkyBot}{RGB}{240,200,104}%
    \definecolor{cvSun}{RGB}{255,244,200}%
    \definecolor{cvFieldTop}{RGB}{176,140,58}%
    \definecolor{cvFieldBot}{RGB}{84,54,22}%
    \definecolor{cvFurrow}{RGB}{45,28,12}%
    \definecolor{cvSil}{RGB}{34,22,10}%
    \definecolor{cvCream}{RGB}{250,245,232}%
  }{%
  \ifstrequal{#1}{part2}{%
    \definecolor{cvSkyTop}{RGB}{46,20,30}%
    \definecolor{cvSkyMid}{RGB}{128,52,54}%
    \definecolor{cvSkyBot}{RGB}{226,120,64}%
    \definecolor{cvSun}{RGB}{255,206,150}%
    \definecolor{cvFieldTop}{RGB}{132,84,44}%
    \definecolor{cvFieldBot}{RGB}{48,26,18}%
    \definecolor{cvFurrow}{RGB}{36,16,14}%
    \definecolor{cvSil}{RGB}{26,12,12}%
    \definecolor{cvCream}{RGB}{248,236,224}%
  }{%
  \ifstrequal{#1}{part3}{%
    \definecolor{cvSkyTop}{RGB}{30,20,54}%
    \definecolor{cvSkyMid}{RGB}{92,62,124}%
    \definecolor{cvSkyBot}{RGB}{216,132,140}%
    \definecolor{cvSun}{RGB}{255,214,196}%
    \definecolor{cvFieldTop}{RGB}{128,90,74}%
    \definecolor{cvFieldBot}{RGB}{44,26,30}%
    \definecolor{cvFurrow}{RGB}{34,20,30}%
    \definecolor{cvSil}{RGB}{22,14,22}%
    \definecolor{cvCream}{RGB}{246,238,238}%
  }{%
  \ifstrequal{#1}{part4}{%
    \definecolor{cvSkyTop}{RGB}{12,28,28}%
    \definecolor{cvSkyMid}{RGB}{44,86,72}%
    \definecolor{cvSkyBot}{RGB}{186,168,86}%
    \definecolor{cvSun}{RGB}{240,232,168}%
    \definecolor{cvFieldTop}{RGB}{110,104,50}%
    \definecolor{cvFieldBot}{RGB}{30,34,20}%
    \definecolor{cvFurrow}{RGB}{26,32,20}%
    \definecolor{cvSil}{RGB}{14,20,14}%
    \definecolor{cvCream}{RGB}{240,238,226}%
  }{%
    % 默认: 全书(麦浪金褐)
    \definecolor{cvSkyTop}{RGB}{28,40,32}%
    \definecolor{cvSkyMid}{RGB}{128,110,64}%
    \definecolor{cvSkyBot}{RGB}{232,178,88}%
    \definecolor{cvSun}{RGB}{255,236,180}%
    \definecolor{cvFieldTop}{RGB}{168,124,52}%
    \definecolor{cvFieldBot}{RGB}{62,40,20}%
    \definecolor{cvFurrow}{RGB}{42,26,12}%
    \definecolor{cvSil}{RGB}{28,18,10}%
    \definecolor{cvCream}{RGB}{250,242,226}%
  }}}}}%
}

% --------------------------------------------------------------
%  一键设定:  \CoverSetup{配色}{卷次}{卷名}{副题}{原书页码范围}
% --------------------------------------------------------------
\newcommand{\CoverSetup}[5]{%
  \CoverPalette{#1}%
  \newcommand{\VolNum}{#2}%
  \newcommand{\VolName}{#3}%
  \newcommand{\VolSub}{#4}%
  \newcommand{\VolPages}{#5}%
}

% --------------------------------------------------------------
%  一支麦穗:  \cvWheatEar{x}{底 y}{高}{倾角}
%  茎略弯, 麦粒成对左右斜出, 顶端一芒三针
% --------------------------------------------------------------
\newcommand{\cvWheatEar}[4]{%
  \begin{scope}[shift={(#1,#2)}, rotate=#4]
    \draw[cvSil, line width=0.85mm, line cap=round]
      (0,0) .. controls (0.5,{0.35*#3}) and (0.8,{0.72*#3}) .. ({0.55*#3},#3);
    \foreach \cvK in {0,...,6}{%
      \pgfmathsetmacro{\cvT}{0.50+0.50*\cvK/6}
      \pgfmathsetmacro{\cvGY}{\cvT*#3}
      \pgfmathsetmacro{\cvGX}{0.55*\cvT*\cvT*#3}
      \pgfmathsetmacro{\cvGS}{1.0-0.34*\cvT}
      \begin{scope}[shift={(\cvGX,\cvGY)}, rotate=28]
        \fill[cvSil] (0,{1.9*\cvGS}) ellipse [x radius={0.92*\cvGS}, y radius={2.3*\cvGS}];
      \end{scope}
      \begin{scope}[shift={(\cvGX,\cvGY)}, rotate=-28]
        \fill[cvSil] (0,{-1.9*\cvGS}) ellipse [x radius={0.92*\cvGS}, y radius={2.3*\cvGS}];
      \end{scope}
    }
    \begin{scope}[shift={({0.55*#3},#3)}]
      \fill[cvSil] (0,1.5) ellipse [x radius=0.78, y radius=1.95];
    \end{scope}
    \draw[cvSil, line width=0.26mm, line cap=round]
      ({0.55*#3},{#3+2.3}) -- ({0.55*#3-1.7},{#3+9.5});
    \draw[cvSil, line width=0.26mm, line cap=round]
      ({0.55*#3},{#3+2.3}) -- ({0.55*#3+2.1},{#3+10.5});
    \draw[cvSil, line width=0.26mm, line cap=round]
      ({0.55*#3},{#3+2.3}) -- ({0.55*#3},{#3+11.5});
  \end{scope}
}

% --------------------------------------------------------------
%  封面
% --------------------------------------------------------------
\makeatletter
\renewcommand{\makecover}{%
  \begin{titlepage}
    \thispagestyle{empty}
    \begin{tikzpicture}[remember picture, overlay, x=1mm, y=1mm,
                        shift={(current page.south west)}]
      % ===== 天空 =====
      \shade[top color=cvSkyTop, middle color=cvSkyMid, bottom color=cvSkyBot]
        (0,105) rectangle (210,297);
      % ===== 田野 =====
      \shade[top color=cvFieldTop, bottom color=cvFieldBot]
        (0,0) rectangle (210,105);
      % ===== 落日与光晕 =====
      \begin{scope}
        \clip (0,105) rectangle (210,297);
        \fill[cvSun, opacity=0.14] (152,132) circle (46);
        \fill[cvSun, opacity=0.24] (152,132) circle (30);
      \end{scope}
      \fill[cvSun, opacity=0.95] (152,132) circle (16.5);
      % ===== 犁沟: 向地平线灭点收束 =====
      \begin{scope}
        \clip (0,0) rectangle (210,105);
        \foreach \cvI in {1,...,20}{%
          \pgfmathsetmacro{\cvDX}{3.1*pow(\cvI,1.55)}
          \draw[cvFurrow, opacity=0.30, line width=0.32mm] ({105+\cvDX},0) -- (105,105);
          \draw[cvFurrow, opacity=0.30, line width=0.32mm] ({105-\cvDX},0) -- (105,105);
        }
        \foreach \cvY in {8,18,31,47,68,92}{%
          \draw[cvFurrow, opacity=0.15, line width=0.3mm]
            (0,\cvY) .. controls (60,{\cvY+3.5}) and (150,{\cvY-3}) .. (210,{\cvY+1.5});
        }
      \end{scope}
      % 地平线
      \draw[cvSil, opacity=0.45, line width=0.45mm] (0,105) -- (210,105);
      % ===== 播种者剪影(立于地平线) =====
      \begin{scope}[cvSil]
        \draw[line width=1.9mm, line cap=round, line join=round]
          (55.5,116) -- (50,105)  (55.5,116) -- (60.5,105.5);
        \draw[line width=2.3mm, line cap=round, line join=round]
          (55.5,116) -- (55,123.5);
        \draw[line width=1.6mm, line cap=round, line join=round]
          (55.2,123.2) -- (61.5,126.6) -- (68,129.4);
        \draw[line width=1.5mm, line cap=round, line join=round]
          (55.2,123.2) -- (51.4,119.8) -- (48.6,116.2);
        \fill (55.6,127.3) circle (2.35);
      \end{scope}
      % ===== 种子: 沿贝塞尔弧撒出 =====
      \draw[cvCream, opacity=0.13, line width=0.5mm]
        (68,129.4) .. controls (100,168) and (160,162) .. (196,104);
      % 掌心那一把(尚未撒出)
      \foreach \cvI in {1,...,5}{%
        \pgfmathsetmacro{\cvJX}{1.1*sin(\cvI*53)}
        \pgfmathsetmacro{\cvJY}{1.1*cos(\cvI*83)}
        \fill[cvCream, opacity=0.9] ({68+\cvJX},{129.4+\cvJY}) circle (1.0);
      }
      \foreach \cvI in {1,...,26}{%
        \pgfmathsetmacro{\cvT}{0.035+0.92*\cvI/26}
        \pgfmathsetmacro{\cvU}{1-\cvT}
        \pgfmathsetmacro{\cvSX}{\cvU*\cvU*\cvU*68 + 3*\cvU*\cvU*\cvT*100
                                + 3*\cvU*\cvT*\cvT*160 + \cvT*\cvT*\cvT*196}
        \pgfmathsetmacro{\cvSY}{\cvU*\cvU*\cvU*129.4 + 3*\cvU*\cvU*\cvT*168
                                + 3*\cvU*\cvT*\cvT*162 + \cvT*\cvT*\cvT*104}
        \pgfmathsetmacro{\cvR}{1.55-0.85*\cvT}
        \pgfmathsetmacro{\cvJX}{1.7*sin(\cvI*97)}
        \pgfmathsetmacro{\cvJY}{1.7*cos(\cvI*71)}
        \fill[cvCream, opacity={0.88-0.28*\cvT}] ({\cvSX+\cvJX},{\cvSY+\cvJY}) circle (\cvR);
      }
      % ===== 前景麦穗(两排) =====
      \foreach \cvI in {1,...,11}{%
        \pgfmathsetmacro{\cvWX}{4+19*(\cvI-1)+1.6*sin(\cvI*2.1)}
        \pgfmathsetmacro{\cvWH}{23+7*sin(\cvI*1.7)}
        \pgfmathsetmacro{\cvWL}{3.2*sin(\cvI*2.9)}
        \begin{scope}[opacity=0.45]
          \cvWheatEar{\cvWX}{30}{\cvWH}{\cvWL}
        \end{scope}
      }
      \foreach \cvI in {1,...,10}{%
        \pgfmathsetmacro{\cvWX}{9+20.4*(\cvI-1)+1.5*sin(\cvI*1.3)}
        \pgfmathsetmacro{\cvWH}{41+10*sin(\cvI*2.3)}
        \pgfmathsetmacro{\cvWL}{4.5*sin(\cvI*3.7)}
        \cvWheatEar{\cvWX}{30}{\cvWH}{\cvWL}
      }
      % ===== 压暗一层, 统一色调并保证文字可读 =====
      \fill[black, opacity=0.10] (0,105) rectangle (210,297);
      % ===== 底部信息带 =====
      \fill[cvSil, opacity=0.72] (0,0) rectangle (210,30);
      \draw[cvCream, opacity=0.40, line width=0.4mm] (0,30) -- (210,30);
      % ===== 文字 =====
      \node[anchor=north west, inner sep=0, text=cvCream] at (22,268)
        {\fontsize{38pt}{45pt}\selectfont\headfont \ziju{0.08}\BookTitleCN};
      \node[anchor=north west, inner sep=0, text=cvCream] at (22.6,242)
        {\fontsize{15pt}{20pt}\selectfont\itshape \BookTitleEN};
      \draw[cvCream, opacity=0.55, line width=0.5mm] (22,233) -- (120,233);
      \node[anchor=north west, inner sep=0, text=cvCream] at (22,224)
        {\fontsize{11pt}{15pt}\selectfont 一位数学家对往昔的反思与见证};
      \node[anchor=north west, inner sep=0, text=cvCream, opacity=0.78] at (22.4,208)
        {\fontsize{10pt}{14pt}\selectfont \ziju{0.14}\VolNum};
      \node[anchor=north west, inner sep=0, text=cvCream] at (22,199)
        {\fontsize{30pt}{37pt}\selectfont\headfont \ziju{0.1}\VolName};
      \node[anchor=north west, inner sep=0, text=cvCream, opacity=0.92] at (22,180)
        {\fontsize{13pt}{18pt}\selectfont \VolSub};
      \node[anchor=north west, inner sep=0, text=cvCream] at (22,19)
        {\fontsize{11.5pt}{15pt}\selectfont \OriginalAuthor\quad 著};
      \node[anchor=north west, inner sep=0, text=cvCream, opacity=0.85] at (22,7)
        {\fontsize{9.5pt}{13pt}\selectfont \Translator\ 译\quad ·\quad
         \ifdefempty{\TranslationDate}{\today}{\TranslationDate}};
      \node[anchor=north east, inner sep=0, text=cvCream, opacity=0.85] at (188,13)
        {\fontsize{10pt}{14pt}\selectfont 原书 \VolPages};
    \end{tikzpicture}
  \end{titlepage}%
}
\makeatother
                    % 本文件
%      \CoverSetup{prelude}{卷一}{前奏}%
%                 {主题的呈现——或四个乐章的前奏}{p.23–122}
%  正文中再 \makecover 即可。配色关键字:
%      prelude / part1 / part2 / part3 / part4 / all
%  五卷共用同一幅画, 只换配色与卷次文字, 以保证"一套书"的观感。
% ================================================================

\RequirePackage{tikz}

% --------------------------------------------------------------
%  配色: 六套(前奏 / 四部 / 全书), 各为同一暮色母题的变调
% --------------------------------------------------------------
\newcommand{\CoverPalette}[1]{%
  \ifstrequal{#1}{prelude}{%
    \definecolor{cvSkyTop}{RGB}{10,26,48}%
    \definecolor{cvSkyMid}{RGB}{46,84,116}%
    \definecolor{cvSkyBot}{RGB}{232,164,92}%
    \definecolor{cvSun}{RGB}{255,224,168}%
    \definecolor{cvFieldTop}{RGB}{150,110,52}%
    \definecolor{cvFieldBot}{RGB}{58,36,18}%
    \definecolor{cvFurrow}{RGB}{40,22,10}%
    \definecolor{cvSil}{RGB}{30,18,10}%
    \definecolor{cvCream}{RGB}{247,240,228}%
  }{%
  \ifstrequal{#1}{part1}{%
    \definecolor{cvSkyTop}{RGB}{22,70,100}%
    \definecolor{cvSkyMid}{RGB}{92,150,168}%
    \definecolor{cvSkyBot}{RGB}{240,200,104}%
    \definecolor{cvSun}{RGB}{255,244,200}%
    \definecolor{cvFieldTop}{RGB}{176,140,58}%
    \definecolor{cvFieldBot}{RGB}{84,54,22}%
    \definecolor{cvFurrow}{RGB}{45,28,12}%
    \definecolor{cvSil}{RGB}{34,22,10}%
    \definecolor{cvCream}{RGB}{250,245,232}%
  }{%
  \ifstrequal{#1}{part2}{%
    \definecolor{cvSkyTop}{RGB}{46,20,30}%
    \definecolor{cvSkyMid}{RGB}{128,52,54}%
    \definecolor{cvSkyBot}{RGB}{226,120,64}%
    \definecolor{cvSun}{RGB}{255,206,150}%
    \definecolor{cvFieldTop}{RGB}{132,84,44}%
    \definecolor{cvFieldBot}{RGB}{48,26,18}%
    \definecolor{cvFurrow}{RGB}{36,16,14}%
    \definecolor{cvSil}{RGB}{26,12,12}%
    \definecolor{cvCream}{RGB}{248,236,224}%
  }{%
  \ifstrequal{#1}{part3}{%
    \definecolor{cvSkyTop}{RGB}{30,20,54}%
    \definecolor{cvSkyMid}{RGB}{92,62,124}%
    \definecolor{cvSkyBot}{RGB}{216,132,140}%
    \definecolor{cvSun}{RGB}{255,214,196}%
    \definecolor{cvFieldTop}{RGB}{128,90,74}%
    \definecolor{cvFieldBot}{RGB}{44,26,30}%
    \definecolor{cvFurrow}{RGB}{34,20,30}%
    \definecolor{cvSil}{RGB}{22,14,22}%
    \definecolor{cvCream}{RGB}{246,238,238}%
  }{%
  \ifstrequal{#1}{part4}{%
    \definecolor{cvSkyTop}{RGB}{12,28,28}%
    \definecolor{cvSkyMid}{RGB}{44,86,72}%
    \definecolor{cvSkyBot}{RGB}{186,168,86}%
    \definecolor{cvSun}{RGB}{240,232,168}%
    \definecolor{cvFieldTop}{RGB}{110,104,50}%
    \definecolor{cvFieldBot}{RGB}{30,34,20}%
    \definecolor{cvFurrow}{RGB}{26,32,20}%
    \definecolor{cvSil}{RGB}{14,20,14}%
    \definecolor{cvCream}{RGB}{240,238,226}%
  }{%
    % 默认: 全书(麦浪金褐)
    \definecolor{cvSkyTop}{RGB}{28,40,32}%
    \definecolor{cvSkyMid}{RGB}{128,110,64}%
    \definecolor{cvSkyBot}{RGB}{232,178,88}%
    \definecolor{cvSun}{RGB}{255,236,180}%
    \definecolor{cvFieldTop}{RGB}{168,124,52}%
    \definecolor{cvFieldBot}{RGB}{62,40,20}%
    \definecolor{cvFurrow}{RGB}{42,26,12}%
    \definecolor{cvSil}{RGB}{28,18,10}%
    \definecolor{cvCream}{RGB}{250,242,226}%
  }}}}}%
}

% --------------------------------------------------------------
%  一键设定:  \CoverSetup{配色}{卷次}{卷名}{副题}{原书页码范围}
% --------------------------------------------------------------
\newcommand{\CoverSetup}[5]{%
  \CoverPalette{#1}%
  \newcommand{\VolNum}{#2}%
  \newcommand{\VolName}{#3}%
  \newcommand{\VolSub}{#4}%
  \newcommand{\VolPages}{#5}%
}

% --------------------------------------------------------------
%  一支麦穗:  \cvWheatEar{x}{底 y}{高}{倾角}
%  茎略弯, 麦粒成对左右斜出, 顶端一芒三针
% --------------------------------------------------------------
\newcommand{\cvWheatEar}[4]{%
  \begin{scope}[shift={(#1,#2)}, rotate=#4]
    \draw[cvSil, line width=0.85mm, line cap=round]
      (0,0) .. controls (0.5,{0.35*#3}) and (0.8,{0.72*#3}) .. ({0.55*#3},#3);
    \foreach \cvK in {0,...,6}{%
      \pgfmathsetmacro{\cvT}{0.50+0.50*\cvK/6}
      \pgfmathsetmacro{\cvGY}{\cvT*#3}
      \pgfmathsetmacro{\cvGX}{0.55*\cvT*\cvT*#3}
      \pgfmathsetmacro{\cvGS}{1.0-0.34*\cvT}
      \begin{scope}[shift={(\cvGX,\cvGY)}, rotate=28]
        \fill[cvSil] (0,{1.9*\cvGS}) ellipse [x radius={0.92*\cvGS}, y radius={2.3*\cvGS}];
      \end{scope}
      \begin{scope}[shift={(\cvGX,\cvGY)}, rotate=-28]
        \fill[cvSil] (0,{-1.9*\cvGS}) ellipse [x radius={0.92*\cvGS}, y radius={2.3*\cvGS}];
      \end{scope}
    }
    \begin{scope}[shift={({0.55*#3},#3)}]
      \fill[cvSil] (0,1.5) ellipse [x radius=0.78, y radius=1.95];
    \end{scope}
    \draw[cvSil, line width=0.26mm, line cap=round]
      ({0.55*#3},{#3+2.3}) -- ({0.55*#3-1.7},{#3+9.5});
    \draw[cvSil, line width=0.26mm, line cap=round]
      ({0.55*#3},{#3+2.3}) -- ({0.55*#3+2.1},{#3+10.5});
    \draw[cvSil, line width=0.26mm, line cap=round]
      ({0.55*#3},{#3+2.3}) -- ({0.55*#3},{#3+11.5});
  \end{scope}
}

% --------------------------------------------------------------
%  封面
% --------------------------------------------------------------
\makeatletter
\renewcommand{\makecover}{%
  \begin{titlepage}
    \thispagestyle{empty}
    \begin{tikzpicture}[remember picture, overlay, x=1mm, y=1mm,
                        shift={(current page.south west)}]
      % ===== 天空 =====
      \shade[top color=cvSkyTop, middle color=cvSkyMid, bottom color=cvSkyBot]
        (0,105) rectangle (210,297);
      % ===== 田野 =====
      \shade[top color=cvFieldTop, bottom color=cvFieldBot]
        (0,0) rectangle (210,105);
      % ===== 落日与光晕 =====
      \begin{scope}
        \clip (0,105) rectangle (210,297);
        \fill[cvSun, opacity=0.14] (152,132) circle (46);
        \fill[cvSun, opacity=0.24] (152,132) circle (30);
      \end{scope}
      \fill[cvSun, opacity=0.95] (152,132) circle (16.5);
      % ===== 犁沟: 向地平线灭点收束 =====
      \begin{scope}
        \clip (0,0) rectangle (210,105);
        \foreach \cvI in {1,...,20}{%
          \pgfmathsetmacro{\cvDX}{3.1*pow(\cvI,1.55)}
          \draw[cvFurrow, opacity=0.30, line width=0.32mm] ({105+\cvDX},0) -- (105,105);
          \draw[cvFurrow, opacity=0.30, line width=0.32mm] ({105-\cvDX},0) -- (105,105);
        }
        \foreach \cvY in {8,18,31,47,68,92}{%
          \draw[cvFurrow, opacity=0.15, line width=0.3mm]
            (0,\cvY) .. controls (60,{\cvY+3.5}) and (150,{\cvY-3}) .. (210,{\cvY+1.5});
        }
      \end{scope}
      % 地平线
      \draw[cvSil, opacity=0.45, line width=0.45mm] (0,105) -- (210,105);
      % ===== 播种者剪影(立于地平线) =====
      \begin{scope}[cvSil]
        \draw[line width=1.9mm, line cap=round, line join=round]
          (55.5,116) -- (50,105)  (55.5,116) -- (60.5,105.5);
        \draw[line width=2.3mm, line cap=round, line join=round]
          (55.5,116) -- (55,123.5);
        \draw[line width=1.6mm, line cap=round, line join=round]
          (55.2,123.2) -- (61.5,126.6) -- (68,129.4);
        \draw[line width=1.5mm, line cap=round, line join=round]
          (55.2,123.2) -- (51.4,119.8) -- (48.6,116.2);
        \fill (55.6,127.3) circle (2.35);
      \end{scope}
      % ===== 种子: 沿贝塞尔弧撒出 =====
      \draw[cvCream, opacity=0.13, line width=0.5mm]
        (68,129.4) .. controls (100,168) and (160,162) .. (196,104);
      % 掌心那一把(尚未撒出)
      \foreach \cvI in {1,...,5}{%
        \pgfmathsetmacro{\cvJX}{1.1*sin(\cvI*53)}
        \pgfmathsetmacro{\cvJY}{1.1*cos(\cvI*83)}
        \fill[cvCream, opacity=0.9] ({68+\cvJX},{129.4+\cvJY}) circle (1.0);
      }
      \foreach \cvI in {1,...,26}{%
        \pgfmathsetmacro{\cvT}{0.035+0.92*\cvI/26}
        \pgfmathsetmacro{\cvU}{1-\cvT}
        \pgfmathsetmacro{\cvSX}{\cvU*\cvU*\cvU*68 + 3*\cvU*\cvU*\cvT*100
                                + 3*\cvU*\cvT*\cvT*160 + \cvT*\cvT*\cvT*196}
        \pgfmathsetmacro{\cvSY}{\cvU*\cvU*\cvU*129.4 + 3*\cvU*\cvU*\cvT*168
                                + 3*\cvU*\cvT*\cvT*162 + \cvT*\cvT*\cvT*104}
        \pgfmathsetmacro{\cvR}{1.55-0.85*\cvT}
        \pgfmathsetmacro{\cvJX}{1.7*sin(\cvI*97)}
        \pgfmathsetmacro{\cvJY}{1.7*cos(\cvI*71)}
        \fill[cvCream, opacity={0.88-0.28*\cvT}] ({\cvSX+\cvJX},{\cvSY+\cvJY}) circle (\cvR);
      }
      % ===== 前景麦穗(两排) =====
      \foreach \cvI in {1,...,11}{%
        \pgfmathsetmacro{\cvWX}{4+19*(\cvI-1)+1.6*sin(\cvI*2.1)}
        \pgfmathsetmacro{\cvWH}{23+7*sin(\cvI*1.7)}
        \pgfmathsetmacro{\cvWL}{3.2*sin(\cvI*2.9)}
        \begin{scope}[opacity=0.45]
          \cvWheatEar{\cvWX}{30}{\cvWH}{\cvWL}
        \end{scope}
      }
      \foreach \cvI in {1,...,10}{%
        \pgfmathsetmacro{\cvWX}{9+20.4*(\cvI-1)+1.5*sin(\cvI*1.3)}
        \pgfmathsetmacro{\cvWH}{41+10*sin(\cvI*2.3)}
        \pgfmathsetmacro{\cvWL}{4.5*sin(\cvI*3.7)}
        \cvWheatEar{\cvWX}{30}{\cvWH}{\cvWL}
      }
      % ===== 压暗一层, 统一色调并保证文字可读 =====
      \fill[black, opacity=0.10] (0,105) rectangle (210,297);
      % ===== 底部信息带 =====
      \fill[cvSil, opacity=0.72] (0,0) rectangle (210,30);
      \draw[cvCream, opacity=0.40, line width=0.4mm] (0,30) -- (210,30);
      % ===== 文字 =====
      \node[anchor=north west, inner sep=0, text=cvCream] at (22,268)
        {\fontsize{38pt}{45pt}\selectfont\headfont \ziju{0.08}\BookTitleCN};
      \node[anchor=north west, inner sep=0, text=cvCream] at (22.6,242)
        {\fontsize{15pt}{20pt}\selectfont\itshape \BookTitleEN};
      \draw[cvCream, opacity=0.55, line width=0.5mm] (22,233) -- (120,233);
      \node[anchor=north west, inner sep=0, text=cvCream] at (22,224)
        {\fontsize{11pt}{15pt}\selectfont 一位数学家对往昔的反思与见证};
      \node[anchor=north west, inner sep=0, text=cvCream, opacity=0.78] at (22.4,208)
        {\fontsize{10pt}{14pt}\selectfont \ziju{0.14}\VolNum};
      \node[anchor=north west, inner sep=0, text=cvCream] at (22,199)
        {\fontsize{30pt}{37pt}\selectfont\headfont \ziju{0.1}\VolName};
      \node[anchor=north west, inner sep=0, text=cvCream, opacity=0.92] at (22,180)
        {\fontsize{13pt}{18pt}\selectfont \VolSub};
      \node[anchor=north west, inner sep=0, text=cvCream] at (22,19)
        {\fontsize{11.5pt}{15pt}\selectfont \OriginalAuthor\quad 著};
      \node[anchor=north west, inner sep=0, text=cvCream, opacity=0.85] at (22,7)
        {\fontsize{9.5pt}{13pt}\selectfont \Translator\ 译\quad ·\quad
         \ifdefempty{\TranslationDate}{\today}{\TranslationDate}};
      \node[anchor=north east, inner sep=0, text=cvCream, opacity=0.85] at (188,13)
        {\fontsize{10pt}{14pt}\selectfont 原书 \VolPages};
    \end{tikzpicture}
  \end{titlepage}%
}
\makeatother
                    % 本文件
%      \CoverSetup{prelude}{卷一}{前奏}%
%                 {主题的呈现——或四个乐章的前奏}{p.23–122}
%  正文中再 \makecover 即可。配色关键字:
%      prelude / part1 / part2 / part3 / part4 / all
%  五卷共用同一幅画, 只换配色与卷次文字, 以保证"一套书"的观感。
% ================================================================

\RequirePackage{tikz}

% --------------------------------------------------------------
%  配色: 六套(前奏 / 四部 / 全书), 各为同一暮色母题的变调
% --------------------------------------------------------------
\newcommand{\CoverPalette}[1]{%
  \ifstrequal{#1}{prelude}{%
    \definecolor{cvSkyTop}{RGB}{10,26,48}%
    \definecolor{cvSkyMid}{RGB}{46,84,116}%
    \definecolor{cvSkyBot}{RGB}{232,164,92}%
    \definecolor{cvSun}{RGB}{255,224,168}%
    \definecolor{cvFieldTop}{RGB}{150,110,52}%
    \definecolor{cvFieldBot}{RGB}{58,36,18}%
    \definecolor{cvFurrow}{RGB}{40,22,10}%
    \definecolor{cvSil}{RGB}{30,18,10}%
    \definecolor{cvCream}{RGB}{247,240,228}%
  }{%
  \ifstrequal{#1}{part1}{%
    \definecolor{cvSkyTop}{RGB}{22,70,100}%
    \definecolor{cvSkyMid}{RGB}{92,150,168}%
    \definecolor{cvSkyBot}{RGB}{240,200,104}%
    \definecolor{cvSun}{RGB}{255,244,200}%
    \definecolor{cvFieldTop}{RGB}{176,140,58}%
    \definecolor{cvFieldBot}{RGB}{84,54,22}%
    \definecolor{cvFurrow}{RGB}{45,28,12}%
    \definecolor{cvSil}{RGB}{34,22,10}%
    \definecolor{cvCream}{RGB}{250,245,232}%
  }{%
  \ifstrequal{#1}{part2}{%
    \definecolor{cvSkyTop}{RGB}{46,20,30}%
    \definecolor{cvSkyMid}{RGB}{128,52,54}%
    \definecolor{cvSkyBot}{RGB}{226,120,64}%
    \definecolor{cvSun}{RGB}{255,206,150}%
    \definecolor{cvFieldTop}{RGB}{132,84,44}%
    \definecolor{cvFieldBot}{RGB}{48,26,18}%
    \definecolor{cvFurrow}{RGB}{36,16,14}%
    \definecolor{cvSil}{RGB}{26,12,12}%
    \definecolor{cvCream}{RGB}{248,236,224}%
  }{%
  \ifstrequal{#1}{part3}{%
    \definecolor{cvSkyTop}{RGB}{30,20,54}%
    \definecolor{cvSkyMid}{RGB}{92,62,124}%
    \definecolor{cvSkyBot}{RGB}{216,132,140}%
    \definecolor{cvSun}{RGB}{255,214,196}%
    \definecolor{cvFieldTop}{RGB}{128,90,74}%
    \definecolor{cvFieldBot}{RGB}{44,26,30}%
    \definecolor{cvFurrow}{RGB}{34,20,30}%
    \definecolor{cvSil}{RGB}{22,14,22}%
    \definecolor{cvCream}{RGB}{246,238,238}%
  }{%
  \ifstrequal{#1}{part4}{%
    \definecolor{cvSkyTop}{RGB}{12,28,28}%
    \definecolor{cvSkyMid}{RGB}{44,86,72}%
    \definecolor{cvSkyBot}{RGB}{186,168,86}%
    \definecolor{cvSun}{RGB}{240,232,168}%
    \definecolor{cvFieldTop}{RGB}{110,104,50}%
    \definecolor{cvFieldBot}{RGB}{30,34,20}%
    \definecolor{cvFurrow}{RGB}{26,32,20}%
    \definecolor{cvSil}{RGB}{14,20,14}%
    \definecolor{cvCream}{RGB}{240,238,226}%
  }{%
    % 默认: 全书(麦浪金褐)
    \definecolor{cvSkyTop}{RGB}{28,40,32}%
    \definecolor{cvSkyMid}{RGB}{128,110,64}%
    \definecolor{cvSkyBot}{RGB}{232,178,88}%
    \definecolor{cvSun}{RGB}{255,236,180}%
    \definecolor{cvFieldTop}{RGB}{168,124,52}%
    \definecolor{cvFieldBot}{RGB}{62,40,20}%
    \definecolor{cvFurrow}{RGB}{42,26,12}%
    \definecolor{cvSil}{RGB}{28,18,10}%
    \definecolor{cvCream}{RGB}{250,242,226}%
  }}}}}%
}

% --------------------------------------------------------------
%  一键设定:  \CoverSetup{配色}{卷次}{卷名}{副题}{原书页码范围}
% --------------------------------------------------------------
\newcommand{\CoverSetup}[5]{%
  \CoverPalette{#1}%
  \newcommand{\VolNum}{#2}%
  \newcommand{\VolName}{#3}%
  \newcommand{\VolSub}{#4}%
  \newcommand{\VolPages}{#5}%
}

% --------------------------------------------------------------
%  一支麦穗:  \cvWheatEar{x}{底 y}{高}{倾角}
%  茎略弯, 麦粒成对左右斜出, 顶端一芒三针
% --------------------------------------------------------------
\newcommand{\cvWheatEar}[4]{%
  \begin{scope}[shift={(#1,#2)}, rotate=#4]
    \draw[cvSil, line width=0.85mm, line cap=round]
      (0,0) .. controls (0.5,{0.35*#3}) and (0.8,{0.72*#3}) .. ({0.55*#3},#3);
    \foreach \cvK in {0,...,6}{%
      \pgfmathsetmacro{\cvT}{0.50+0.50*\cvK/6}
      \pgfmathsetmacro{\cvGY}{\cvT*#3}
      \pgfmathsetmacro{\cvGX}{0.55*\cvT*\cvT*#3}
      \pgfmathsetmacro{\cvGS}{1.0-0.34*\cvT}
      \begin{scope}[shift={(\cvGX,\cvGY)}, rotate=28]
        \fill[cvSil] (0,{1.9*\cvGS}) ellipse [x radius={0.92*\cvGS}, y radius={2.3*\cvGS}];
      \end{scope}
      \begin{scope}[shift={(\cvGX,\cvGY)}, rotate=-28]
        \fill[cvSil] (0,{-1.9*\cvGS}) ellipse [x radius={0.92*\cvGS}, y radius={2.3*\cvGS}];
      \end{scope}
    }
    \begin{scope}[shift={({0.55*#3},#3)}]
      \fill[cvSil] (0,1.5) ellipse [x radius=0.78, y radius=1.95];
    \end{scope}
    \draw[cvSil, line width=0.26mm, line cap=round]
      ({0.55*#3},{#3+2.3}) -- ({0.55*#3-1.7},{#3+9.5});
    \draw[cvSil, line width=0.26mm, line cap=round]
      ({0.55*#3},{#3+2.3}) -- ({0.55*#3+2.1},{#3+10.5});
    \draw[cvSil, line width=0.26mm, line cap=round]
      ({0.55*#3},{#3+2.3}) -- ({0.55*#3},{#3+11.5});
  \end{scope}
}

% --------------------------------------------------------------
%  封面
% --------------------------------------------------------------
\makeatletter
\renewcommand{\makecover}{%
  \begin{titlepage}
    \thispagestyle{empty}
    \begin{tikzpicture}[remember picture, overlay, x=1mm, y=1mm,
                        shift={(current page.south west)}]
      % ===== 天空 =====
      \shade[top color=cvSkyTop, middle color=cvSkyMid, bottom color=cvSkyBot]
        (0,105) rectangle (210,297);
      % ===== 田野 =====
      \shade[top color=cvFieldTop, bottom color=cvFieldBot]
        (0,0) rectangle (210,105);
      % ===== 落日与光晕 =====
      \begin{scope}
        \clip (0,105) rectangle (210,297);
        \fill[cvSun, opacity=0.14] (152,132) circle (46);
        \fill[cvSun, opacity=0.24] (152,132) circle (30);
      \end{scope}
      \fill[cvSun, opacity=0.95] (152,132) circle (16.5);
      % ===== 犁沟: 向地平线灭点收束 =====
      \begin{scope}
        \clip (0,0) rectangle (210,105);
        \foreach \cvI in {1,...,20}{%
          \pgfmathsetmacro{\cvDX}{3.1*pow(\cvI,1.55)}
          \draw[cvFurrow, opacity=0.30, line width=0.32mm] ({105+\cvDX},0) -- (105,105);
          \draw[cvFurrow, opacity=0.30, line width=0.32mm] ({105-\cvDX},0) -- (105,105);
        }
        \foreach \cvY in {8,18,31,47,68,92}{%
          \draw[cvFurrow, opacity=0.15, line width=0.3mm]
            (0,\cvY) .. controls (60,{\cvY+3.5}) and (150,{\cvY-3}) .. (210,{\cvY+1.5});
        }
      \end{scope}
      % 地平线
      \draw[cvSil, opacity=0.45, line width=0.45mm] (0,105) -- (210,105);
      % ===== 播种者剪影(立于地平线) =====
      \begin{scope}[cvSil]
        \draw[line width=1.9mm, line cap=round, line join=round]
          (55.5,116) -- (50,105)  (55.5,116) -- (60.5,105.5);
        \draw[line width=2.3mm, line cap=round, line join=round]
          (55.5,116) -- (55,123.5);
        \draw[line width=1.6mm, line cap=round, line join=round]
          (55.2,123.2) -- (61.5,126.6) -- (68,129.4);
        \draw[line width=1.5mm, line cap=round, line join=round]
          (55.2,123.2) -- (51.4,119.8) -- (48.6,116.2);
        \fill (55.6,127.3) circle (2.35);
      \end{scope}
      % ===== 种子: 沿贝塞尔弧撒出 =====
      \draw[cvCream, opacity=0.13, line width=0.5mm]
        (68,129.4) .. controls (100,168) and (160,162) .. (196,104);
      % 掌心那一把(尚未撒出)
      \foreach \cvI in {1,...,5}{%
        \pgfmathsetmacro{\cvJX}{1.1*sin(\cvI*53)}
        \pgfmathsetmacro{\cvJY}{1.1*cos(\cvI*83)}
        \fill[cvCream, opacity=0.9] ({68+\cvJX},{129.4+\cvJY}) circle (1.0);
      }
      \foreach \cvI in {1,...,26}{%
        \pgfmathsetmacro{\cvT}{0.035+0.92*\cvI/26}
        \pgfmathsetmacro{\cvU}{1-\cvT}
        \pgfmathsetmacro{\cvSX}{\cvU*\cvU*\cvU*68 + 3*\cvU*\cvU*\cvT*100
                                + 3*\cvU*\cvT*\cvT*160 + \cvT*\cvT*\cvT*196}
        \pgfmathsetmacro{\cvSY}{\cvU*\cvU*\cvU*129.4 + 3*\cvU*\cvU*\cvT*168
                                + 3*\cvU*\cvT*\cvT*162 + \cvT*\cvT*\cvT*104}
        \pgfmathsetmacro{\cvR}{1.55-0.85*\cvT}
        \pgfmathsetmacro{\cvJX}{1.7*sin(\cvI*97)}
        \pgfmathsetmacro{\cvJY}{1.7*cos(\cvI*71)}
        \fill[cvCream, opacity={0.88-0.28*\cvT}] ({\cvSX+\cvJX},{\cvSY+\cvJY}) circle (\cvR);
      }
      % ===== 前景麦穗(两排) =====
      \foreach \cvI in {1,...,11}{%
        \pgfmathsetmacro{\cvWX}{4+19*(\cvI-1)+1.6*sin(\cvI*2.1)}
        \pgfmathsetmacro{\cvWH}{23+7*sin(\cvI*1.7)}
        \pgfmathsetmacro{\cvWL}{3.2*sin(\cvI*2.9)}
        \begin{scope}[opacity=0.45]
          \cvWheatEar{\cvWX}{30}{\cvWH}{\cvWL}
        \end{scope}
      }
      \foreach \cvI in {1,...,10}{%
        \pgfmathsetmacro{\cvWX}{9+20.4*(\cvI-1)+1.5*sin(\cvI*1.3)}
        \pgfmathsetmacro{\cvWH}{41+10*sin(\cvI*2.3)}
        \pgfmathsetmacro{\cvWL}{4.5*sin(\cvI*3.7)}
        \cvWheatEar{\cvWX}{30}{\cvWH}{\cvWL}
      }
      % ===== 压暗一层, 统一色调并保证文字可读 =====
      \fill[black, opacity=0.10] (0,105) rectangle (210,297);
      % ===== 底部信息带 =====
      \fill[cvSil, opacity=0.72] (0,0) rectangle (210,30);
      \draw[cvCream, opacity=0.40, line width=0.4mm] (0,30) -- (210,30);
      % ===== 文字 =====
      \node[anchor=north west, inner sep=0, text=cvCream] at (22,268)
        {\fontsize{38pt}{45pt}\selectfont\headfont \ziju{0.08}\BookTitleCN};
      \node[anchor=north west, inner sep=0, text=cvCream] at (22.6,242)
        {\fontsize{15pt}{20pt}\selectfont\itshape \BookTitleEN};
      \draw[cvCream, opacity=0.55, line width=0.5mm] (22,233) -- (120,233);
      \node[anchor=north west, inner sep=0, text=cvCream] at (22,224)
        {\fontsize{11pt}{15pt}\selectfont 一位数学家对往昔的反思与见证};
      \node[anchor=north west, inner sep=0, text=cvCream, opacity=0.78] at (22.4,208)
        {\fontsize{10pt}{14pt}\selectfont \ziju{0.14}\VolNum};
      \node[anchor=north west, inner sep=0, text=cvCream] at (22,199)
        {\fontsize{30pt}{37pt}\selectfont\headfont \ziju{0.1}\VolName};
      \node[anchor=north west, inner sep=0, text=cvCream, opacity=0.92] at (22,180)
        {\fontsize{13pt}{18pt}\selectfont \VolSub};
      \node[anchor=north west, inner sep=0, text=cvCream] at (22,19)
        {\fontsize{11.5pt}{15pt}\selectfont \OriginalAuthor\quad 著};
      \node[anchor=north west, inner sep=0, text=cvCream, opacity=0.85] at (22,7)
        {\fontsize{9.5pt}{13pt}\selectfont \Translator\ 译\quad ·\quad
         \ifdefempty{\TranslationDate}{\today}{\TranslationDate}};
      \node[anchor=north east, inner sep=0, text=cvCream, opacity=0.85] at (188,13)
        {\fontsize{10pt}{14pt}\selectfont 原书 \VolPages};
    \end{tikzpicture}
  \end{titlepage}%
}
\makeatother

\CoverSetup{part2}{卷三}{第二部}{葬礼(I)——或中国皇帝的新衣}{p.255–412}

\hypersetup{
  pdftitle={收获与播种·卷三·第二部(葬礼 I)},
  pdfauthor={Alexandre Grothendieck / 数译 译}
}

% --------------------------------------------------------------
\begin{document}

\frontmatter
\makecover

% ---------------------------------------------------------------
%  译者说明
% ---------------------------------------------------------------
\chapter*{译者说明}
\markboth{译者说明}{译者说明}
\addcontentsline{toc}{chapter}{译者说明}

本册是《收获与播种》中文译本的第三卷, 收原书的第二部“葬礼(I):或中国皇帝的
新衣”(Deuxième partie: L'Enterrement~(I) ou la robe de l'Empereur de Chine),
即原书第~255 至第~412 页, 含第~13 至第~16~章。本部把“葬礼”这场经验的原始事实
一桩一桩挖出来, 是全书“调查”部分的前半。

本部的脚注按章各自重新编号, 并照原书带~(*)、(**) 等角标; 第~13--16~章分别自
~1~起编至~41、86、132、29。

\CommonNoteBody

\makecontents

\mainmatter

% ================================================================
%  第二部 (第 13-16 章, 原书 p.255-412)
% ================================================================
\partdivider{第二部}{葬礼(I)\\[0.35em]{\normalsize 或中国皇帝的新衣}}{葬礼(I)或中国皇帝的新衣}

% 原书第257页献辞
\dedication{%
献给曾是我朋友的人,\\[0.5em]
既献给少数仍是朋友的人,\\[0.5em]
也献给众多前来为我的葬礼齐声唱和的人。\\[2em]
纪念那场难忘的讨论会……\\[0.5em]
以及全体教众……%
}

\setcounter{chapter}{12}
\markboth{收获与播种·第二部}{}

% ================================================================
%  第 13 章 A) 遗产与继承人 (原书 p.259-280)
%  13.1 身后的学生 | 13.2 孤儿们 | 13.3 时尚——或名人传
%  源文: _work/src2/source.md
% ================================================================
\providecommand{\fnoteplainbody}[3]{%
  \begingroup\renewcommand{\thefootnote}{#1#2}\footnotemark
  \renewcommand{\thefootnote}{#1}\footnotetext{#3}\endgroup
}

\chapter{A)~遗产与继承人}

\section{I.~身后的学生}

\subsection{1.~教学的失败(2)——或创造与自负}

\notehead{44$'$}{[该注由第(I)部分《愚妄与更新》第八章“孤独的历险”第~50~节所召,原书第~227~页]}

\pmark\origpage{173} 这段文字,在读了“一段过去的重量”这最后一节的那位朋友那里“叮”地一声响了\fnote{1}{(*)}{(5月10日)那位朋友不是别人,正是Zoghman~Mebkhout;他慨然允准我揭开我在本注所引那封信(1984年4月2日)的出处上本来觉得应当维持的匿名。}。他写道:“在你许多昔日学生身上,如你所说,那个咄咄逼人、甚至近乎破坏性的‘老板’的形象,依然强烈。于是就有了你的那种印象。”(我想,这里说的“印象”,就是这一节以及补足它的注~n\textsuperscript{o}s~46、47、50 的某些段落里所表达的东西。)他在信的更前面写道:“首先,我认为你暂时离开数学[\,!],离开得好,因为你和学生之间当时有某种隔膜(当然,Deligne除外)。他们都有点被惊呆了……”。

这是第一次,关于我1970年以前“老板”角色的这种论调传到我耳朵里——超出了客套的恭维!还在同一封信的更前面:“……我明白了,你的昔日学生[按:指那些‘1970年以前的’]不太明白数学的\textbf{创造}是怎么回事,而你或许有一份责任……的确,在他们那个年代,问题都已经摆好了……”\fnote{2}{(**)}{(5月10日)前面这段引文被删削得很厉害,是出于对通信人匿名的尊重。完整的引文以及对其真意的评说,见下一条注;起初因缺乏更详尽的信息,我没有领会这段话的真正含义。}。

我的通信人无疑是想说:是我在提“问题”,连同一并要发展的那些观念,而不让我的学生们自己去找到这些问题与观念;也正是在这一点上,我或许遮蔽了他们,使他们不识数学创造工作中那要紧的部分。这恰与一次谈话留下的印象相合\dmark{}\origpage{174}——那是与两位1970年\textbf{之后}的前学生的谈话,前面一条注(注(23iv))里谈到过。的确,我在向我走来的学生身上,首先寻找的是\textbf{合作者},来发展那些在我之内已然成形的直觉与观念,去给一辆已经在那里的大车“推轮子”;他们因而不必从某种虚无中(如我的通信人当年不得不做的那样)去拉动它。然而——让一种可触的、柔软而致密的东西,从不可触之物的迷雾中取得形体——这恰恰从来就是数学工作于我最迷人的方面,也是工作里我尤其感到“创造”在进行的那一部分:某种比一个单纯的“结果”更精微、更要紧的东西的“诞生”。

有时,当我看到昔日学生中的某一位,以轻蔑的态度对待这件弥足珍贵的东西,在他身上抖开 J.H.C.~Whitehead所说的那种“势利”(就是瞧不起自己“本来也会证明”的东西)\fnote{3}{(*)}{见注“年轻人的势利——或纯洁的捍卫者”(n\textsuperscript{o}~27),原书第247页。},我无疑并非全然无关——这样那样地,我脱不了干系。我教学的失败,在1970年之后的一段时期里触目经心,如今又以一种不同而更隐蔽的形状,显现在我第一时期的教学里——尽管按通常的说法,那一期呈现为完完全全的成功!这一点,近年来我不时有所瞥见,也在给几位前学生的信里提到过,而直到如今,没有从他们任何一位那里得到真正的回音。

不过,倘若说我交给学生们做的、他们同我一起做的工作,是纯技术性的、纯粹例行公事式的,无力调动他们的创造才能,这样说怕也不确切。我交到他们手里的,是一些切实而可靠的出发点,他们在其中享有全部的选择自由,可以由此跃起,正如我自己在他们之前曾经做过的那样。我不记得曾给哪位学生出过一个连我自己也不会乐于去做的题目;他们中任何一位与我同行时走过的艰难路段,也都不比我自己在数学家生涯中独自走过的那些更艰难;只要明白工作必须做、也没有别的路,我就没有因此气馁或尥蹶子。

\pmark\origpage{175} 同时我也觉得,我今天所确认的这场失败,其原因比我所出题目的类型、比这些题目在多大程度上含混或分明,更为精微。我在这场失败中的那一份,似乎更应归咎于我对待数学时自身的自负态度;这些态度,我在这场思考中已有机会检视过。即便不曾浸透与某位学生相伴的那项工作本身,它们也必定或多或少浸染了我本人周围的氛围或空气。自负,哪怕以世界上最“不起眼”的方式表达,总是指向一种封闭,一种对事物精微的本质及其美的麻木——无论那是“数学之物”,还是有生命的、我们本可以接待、鼓励、或者也可以从我们高处的优越俯视的人;对我们身边吹拂的气息及其毁伤性的效力——毁伤他人也毁伤我们自己——麻木不仁。

\subsection{2.~一种不义与无力的感觉}

\notehead{44$''$}{[这条注的出现,不遵守写作的时序]}

(5月10日)借着朋友所给的自由引证权——他的信中凡我认为有用的段落都可以引用——我在这里给出一段完整得多的引文\fnote{4}{(*)}{见前一部分“教学的失败(2)——或创造与自负”的第二条脚注,n\textsuperscript{o}~44$'$。},把那段被删削的引文放回它真正的语境:

\begin{quote}
“的确,75到80年间我非常孤立,除了向Verdier提过寥寥几个问题。但那段时间我并不怨你的昔日
学生们,因为没有人真正懂得这条纽带[按:即离散系数与连续系数之间的纽带]的重要性。一切都在
1980年10月改变了:人们发现了这条纽带对半单群的第一个头等重要的应用,即 Kazhdan--Lusztig
重数公式的证明——其中关键地用到了我们所论的范畴等价。这个等价在得到“Riemann--Hilbert对
应”的名字时,没有一句别的评说——毕竟,这再自然不过了!正是在那里,我明白了你的昔日学生
们不很明白数学的\textbf{创造}是什么,而你或许有一份责任。我至今仍有一种不义与无力的感觉。的确,
在他们那个年代,问题都已摆好。这个定理的应用数目,无论在平展拓扑的框架还是超越的框架里,
都很可观,却始终顶着“Riemann--Hilbert对应”的名字!我有种感觉,我的名字配不上这个结果——对\origpage{176}
许多人,尤其对你的昔日学生而言。但正如你从我的工作的引言里能清楚看到的,正是你的“对偶”
形式主义自然地引出了这个结果。不过同你一样,我并不为“可构造离散系数”与晶体系数(或完整
$\mathcal{D}$-模)之间这条纽带的将来担忧。显然,它无论在空间的上同调还是分析中,都适用于许许多
多的领域。”
\end{quote}

正是朋友的信中这一段,启发了我(除本注之外)后来的那道注“当值无名之辈与好上帝定理”。照这封
信的说法,我全然没有料到(如我在适当的地方所解释的),朋友身上那种“不义与无力的感觉”,所针
对的不仅是那种盲目的轻蔑态度——它\textbf{系统性地贬低}他的贡献(这种态度,在我某些昔日学生身
上,后来我也变得颇为熟悉了)——而且是一场真正的诈骗行动:干脆\textbf{变走}一条关键定理的父亲身
份。这一处境,仅仅八天前才向我揭示——关于这一点,见“不义——或归来的意义”及其后各注
(n\textsuperscript{o}s~75至80),它们集结在“讨论会——或Mebkhout层与反常”的标题之下。

\notehead{45}{}

由于我改变了圈子与生活方式,同旧日朋友们会面或其他接触的机会变得稀少。但这并没有妨碍
“拉开距离”的迹象以种种方式表现出来,彼此强弱不一。而在另一些人那里——如Dieudonné、Cartan、
Schwartz,实际上还有所有在我起步时如此善待我的“长者”——我没有感到丝毫这类东西。除他们之
外,我有种感觉:除极少数人之外,我旧日的朋友或学生在数学世界里,其与我之间的关系(无论有无机会表达),在我
退出那个曾经共同的圈子、共同的世界之后,都变得分裂了,带上了“矛盾的双重态度”。

\section{II.~孤儿们}

\subsection{1.~我的孤儿们}

\notehead{46}{[该注由第(I)部分《愚妄与更新》第八章“孤独的历险”第~50~节所召]}

\pmark\origpage{177} 我想借这个机会,在这里说几句关于数学观念与思想的话——在我带到光天化日之下的所有
这些东西当中,它们在我眼里(而且远远领先于其他)具有最广大的分量(注46$_{1}$)\fnote{5}{(*)}{读者可在注~n\textsuperscript{o}~46$_{1}$至46$_{9}$中找到对本注所巡视诸观念的若干更技术性的评说。此外,在我引入的\textbf{诸}观念之外,读者还可在注 n\textsuperscript{o}~88“尸身”里找到关于我著作中(在我“完全做到底”的那一部分著作之内)被其作者视为“执牛耳者”的那一部分的思考。}。首先
是五个紧密相联的关键观念,我将按其特殊性、丰富性(与深度)递增的次序,快速巡视一遍。

首先是同调代数中\textbf{导出范畴}的观念(见注48,原书第274页),以及把它用于一个“通用”的
形式主义——所谓“\textbf{六种运算的形式主义}”(即诸运算 $\overset{L}{\otimes}$、$Lf^*$、$Rf_!$、$R\mathcal{H}om$、
$Rf_*$、$Lf^!$)(注46$_{2}$)——
用于几何学中迄今引入的最重要的几类“空间”的上同调:“代数”空间(如概形、概形叠等)、解析空间
(复解析的,以及刚性的解析的及同类者)、拓扑空间(当然,还在等待各种“温和空间”的语境,肯定还
有其他许多,比如以小范畴的范畴(Cat)作为同伦模型的那一类……)。这一形式主义所涵盖的,既有
离散性质的系数,也有“连续”系数。

这一对偶形式主义及其无所不在性的逐步发现,是通过一场孤独、执着而严格的思考完成的,它从
1956年延续到1963年。正是在这场思考中,导出范畴的观念逐步成形,以及它在同调代数中所担当
的角色的理解也逐步成形。
这套形式主义所缺的,是我的上同调形式主义视野里对“离散系数”与“连续系数”之间那条可
以想见的纽带的理悟——超越局部系统及其以带可积联络的模或晶体模来解释的那个熟知情形之外
的纽带\origpage{178}。这条深刻的纽带,最初在复解析空间的框架里得到表述,是由Zoghman~Mebkhout
(将近二十年后)发现并建立的:以一边由“可构造”离散系数、另一边由“D-模”或“微分算子复形”
观念所构成的导出范畴来表述(见注46$_{3}$)。

将近十年间,由于得不到那些最有资格鼓励他的昔日学生的鼓励——用他们的关切和他们在与我
接触中积累的经验来支持他——Zoghman~Mebkhout在几乎完全的孤立中,坚持他那出色的研究。但这
并没有妨碍他揭示并证明一门正在普遍的冷漠中蹒跚诞生的新晶体理论的\textbf{两把钥匙定理}
\fnote{6}{(*)}{(6月7日)Mebkhout提醒我,在这两个定理之外,还应添上第三个,同样以导出范畴表述,即他(或许
不太恰当地)所谓的 D-模的“\textbf{双对偶定理}”——三者中最难的一个。关于Mebkhout的观念与结果及其
应用的概观,见 Le~Dung~Trang 与 Zoghman~Mebkhout, Introduction to linear differential Systems, Proc. of Symposia in Pure
Mathematics, vol.~40 (1983), part.~2, p.~31--63。}
——两者(这实在是个不好的兆头!)都以导出范畴表述:其一如方才所指,给出“可构造离散系数”与
晶体系数(满足某些“完整性”与“正则性”条件)之间的范畴等价(注48$'$);其二是“\textbf{那}条晶体
的整体对偶定理”,关于从一个光滑复解析空间(不必紧,这意味着额外的可观的技术困难)到一个点
的常值应用。这些是深刻的定理\fnote{7}{(**)}{(5月30日)第二个定理的证明,遇到超越语境惯有的技术困难,
需要求助于所谓“évétesques”技术(原书用语,此处暂保留法文);我猜它可以归入“困难”证明之列。第一个定理的证明是“显然”的——而且深
刻,动用了Heisuke~Hironaka奇点消解的全部力量。如我在“团结”一注(n\textsuperscript{o}~85)倒数第二段所提示的,定理
一旦提炼出来,“随便哪一个”知情的来者都能证明它。可对照“年轻人的势利——或纯洁的捍卫者”
(n\textsuperscript{o}~27)一注中Whitehead的评语。我写那道注的时候,仿佛在某种隐秘的先知的不声不响的口授之下,
还不曾料到现实会在多大的程度上超过我那些怯生生的、摸索式的提示!},为新照亮解析的与概形的
空间(目前为特征零)\origpage{179}的上同调,并许诺这些空间的上同调理论一场规模宏大的更新。它们最终为它们
的作者——在两次进入法国国家科学研究中心的申请被拒之后——换来一个研究员职位(相当于大
学的助理或助教职位)。

这十年间,没有人想到跟Mebkhout说一声“六种运算的形式主义”——我那些学生人人烂熟
\fnote{8}{(*)}{他们是在 SGA~4 与 SGA~5 的讨论班里第一手学到的,或经由 R.~Hartshorne的《留数与对
偶》(Residues and Duality)转手。}、却在任何地方都找不到整理成文文本的形式主义——而他正苦苦对付着超
越语境所带来的巨大的技术困难。他最终是去年从我口中得知它的存在的(以一份看来只有我
一人掌握的公式总表的形式),其时他正不辞辛劳、好意向我讲解他所做的一切;而我,早已不那么在
行上同调了……也没有人想到向他提示:也许先“转向”特征零的概形语境会更“合算”——在那里,
超越语境固有的困难消失,而理论的基本概念问题反而显得愈发清晰。也没有人想到向他指出(或者
哪怕注意到一件我自引入晶体之时就已知道的事
\fnote{9}{(**)}{(5月30日)不过我来得及把它忘了——直到去年与Mebkhout的第二次会面,才凭那次重逢之
力重新想起。(见“冥世重逢”一注,n\textsuperscript{o}~78。)}):光滑空间(解析的或概形的)上的“D-模”,不多
不少,正是“模的晶体”(只要对双方都不问“凝聚”与否的问题)——而后者是一个万能的观念,对带任
何奇点的“空间”与对光滑空间同样行之有效(注46$_{4}$)。

以Mebkhout所展现的才能(与那份少见的勇气),有一点对我说来再清楚不过:倘若处身于一种友善
的氛围,他建立特征零概形的晶体上同调语境中完整的“六种运算”形式主义,将毫无困难,而且会乐
在其中——这样一项宏大计划所需的全部本质观念(包括他自己的,Sato学派与我自己的),在我
看来本已集结在手。以他的才具,这不过是几年功夫的问题,正如平展上同调的一个通用形式主义的
发展是几年功夫的问题(1962--1965)——只要六种运算这条主线(外加换基的两把钥匙定理)已经在
手\origpage{180}。诚然,那是被同侪与见证者的热情与善意的潮流所托举的年月;而不是逆着那些大权
在握者的傲慢自足、逆流而上的劳作……

我来到我想谈的第二对观念:\textbf{概形},以及与它血缘最近的\textbf{topos}。后者是\textbf{site}观念的
更内蕴的版本——site 是我为了使“局部化”的拓扑直觉形式化而最先引入的。(“site”这个词,后来
是Jean~Giraud引入的;他也为使 site 与 topos 的观念具备全部必要的灵活性做了大量工作。)正是代数几
何的种种触目的需要,使我接二连三地引入了概形与 topos。这对观念,蕴含着对代数几何、算术,以及
拓扑的同样宏大的更新之潜能:通过一种\textbf{综合},把这些太久彼此隔绝的“世界”统一在一个共同
的几何直觉里。

以概形的观点、site(或“下降”)的语言,外加十二年关键的奠基工作(还不算我的学生们和其他
志愿者的劳作)对代数几何与算术的更新,二十年来已是既成之事:概形观念、概形的平展上同调观
念(即便还不是平展 topos 与平展叠的观念),终于进入了习惯,进入了公共遗产。

相反,那场同样会把拓扑囊括进来的宏大综合——尽管二十年來,它所需的本质观念与主要
技术工具在我看来已然集结待命\fnote{10}{(*)}{(5月15日)这些“本质观念与主要技术工具”已在 SGA~4
与 SGA~5 两场讨论班的宏大画卷中集结,时在1963至1965年间。这幅画卷的 SGA~5 部分,其成文与出
版所遭遇的离奇变故——十一年之后(1977年)以一种面目全非、满目疮痍的形状问世——为这幅宏
大视野落入“某种时尚”之手的命运,提供了一幅惊心的图像——或者不如说,落入某些我的学生之手:
正是他们最先确立了这种时尚(见下一条脚注)。这些变故及其意义,在最近四个星期的思考中逐步显
露,延续于“伙伴”、“清场”、“格格不入者”、“信号”、“翻转”、“沉默”、“团结”、“蒙骗”、“死者”、“屠杀”
诸注之中(n\textsuperscript{o}s~63$'$'$'$、67、67$'$、68、68$'$与84--88)。}——却还在等待它的时刻。\dmark{}\origpage{181}十五年来(自我告
别数学舞台以来),topos 这个丰饶的统一观念、这个强大的发现工具,被某种时尚\fnoteplainbody{11}{(*)}{(5月13
日)写完这几行之后的六个星期里,思考的继续显明:这种“时尚”最先是由我的某些学生确立的——恰
恰是那些最有条件把某种视野、一些观念与技术手段据为己有的人;他们选择了占有这些劳作的工
具,同时对催生这一切的视野、以及对这视野在其之内诞生的那个人,一并拒认。}打入所谓严肃
观念的冷宫。直到今天,对自己这门科学这种潜在的可观扩展、以及它所提供的新资源,有丝毫觉察的
拓扑学家仍然寥寥。

在这幅更新的视野里,拓扑学家日常摆弄的拓扑空间、微分流形等,连同概形(他听说过)与拓扑、微分或概形的叠(这些却无人谈论),不过是同一类非凡几何对象——\textbf{赋环 topos}(注46$_{5}$)——的诸多化身:拓扑、代数几何与算术
的种种直觉,汇聚于其中的“空间”,共同构成一幅几何视野。各种“模”叠——只要你睁着眼睛,就
会处处遇到(注46$_{6}$)——为它提供了数不清的醒目例子。对它们的深入研究,是深入这些几何对象(或
别的对象,如果有不是几何的对象……)的本质性质的一条头等主线:这些模叠所描述的,正是那些对
象的变动、退化与泛化的样态。然而这份丰富至今无人理会,因为使我们得以精细描述它的观念,不
在通行认可的范畴之列。

这场被拒认的综合\fnote{12}{(**)}{(5月13日)这场综合之遭“拒认”——无论其精神,还是使之成为可能的那
个关键观念——首先出自的,不是别人,恰恰是那项技术手段(概形语言与平展上同调理论的构造)的主
要使用者和受益者:终其全部著作,他受惠于这些手段。这就是Pierre~Deligne。凭他非凡的(源于其非凡
才能的)影响力,以及他相对于我的著作所占据的十分特殊的位置——他像是我的著作的一位隐含
遗产继承人——对我所引入的主要观念(概形观念与平展上同调除外)所设的那道不动声色的系统的
路障,收效极大,无疑在那场\textbf{埋葬}了这些观念的“时尚”的确立中扮演了头等角色;这些观念,已被
迫过了将近十五年的萎顿生活。他的著作,被这层暧昧深深地打上了烙印;这层暧昧,我在接续本注的思
考中第一次瞥见。(见“拒受遗产——或矛盾的代价”,注 n\textsuperscript{o}~47,原书第271页)我离去之后Deligne
著作中这道持续的障碍,这一鲜活而仍然含混的第一知觉,在整个关于\emph{葬礼}的思考中,以惊人的方
式得到了精确与确证;在那场葬礼中,我的朋友扮演着主祭的角色。}还带来一个意外的一面:\dmark{}\origpage{182}某些最
常见空间(注46$_{7}$)的熟知的同伦不变量(或更确切地说,它们的投射有限紧化),被发现配有无人料到的
算术结构,尤其是某些投射有限Galois群的作用……

然而,将近十五年来,在“大世界”里,谁要是贸然说出“topos”这个词——除非是开玩笑,或者他有
逻辑学家的托词——以居高临下的姿态鄙视他,竟成了合乎体面的做法。(逻辑学家是出了名的“与常人不同”,他们的怪癖人们
自会原谅……)以导出范畴的瑜伽来表达拓扑空间的同调与上同调,同样不曾进入拓扑学家中间:对他
们来说,Künneth 公式(对不是域的系数环)至今仍是一个由两条谱序列(或充其量一大串短正合列)
构成的系统,而不是某个合宜范畴中唯一的一个典范同构;换基定理(对真态射,或对光滑态射)他们同
样始终无视——而这些定理(在其近邻平展上同调的框架里)曾构成这门上同调“强力启动”的关键转
折(见注46$_{8}$,原书第270页)。当曾经为发展这门瑜伽出过力的人自己早就把它忘得一干二净,还对那
个想要使用它的可怜人冷眼相待的时候,我就不必惊讶了\fnote{13}{(*)}{(5月13日)后来的思考显明:处境开
始随1981年6月的吕米尼讨论会而改变:在那里,人们看到那些曾经“忘记”(或者不如说,埋葬了……)
这些观念的人,带着它们招摇过市——却依然对那位没有其人便绝不会有那场光辉讨论会的同一“可
怜人”冷眼相待。(关于这场难忘的讨论会,见注 n\textsuperscript{o}s~75与81。)}

第五个观念——也许比其他任何观念都更贴近我的心——是“\textbf{Motif}”。它与前四者不同之处在
于:“\textbf{那}个”好的 Motif 观念(哪怕只在基域上,更不必说任意基底概形),至今尚未成为某个令人
满意的定义的对象,即便为此承认人们会需要的一切“合理的”猜想也无济于事。或者不如说,分明
\dmark{}\textbf{那}条应当首先做出的“合理猜想”,就是\textbf{一条理论之\textbf{存在}}的猜想——它满足这样这样的
数据与性质,而对一个\origpage{183}“在局中”的人\fnote{14}{(*)}{(5月13日)我终于明白:直到今天,唯一(除我
之外)符合这个相当特殊的“多少在局中”的意义的人,是Pierre~Deligne。他有四年之久,一面听“我在
代数几何里所知道的那一点点”,一面日常地做我的 Motif 思考的心腹之交。的确,这些事我跟周围许多
别的同事也谈过,但看来没有一个足够“在行”,能消化那在我之内发展了数年的总体视野,或把
我的提示当作出发点,自己去发展一套视野与纲领(正如我自己当年从Serre的某些思想所给的两三
份“强印象”出发所做的那样)。也许我错了,但我觉得,关心代数簇上同调的人们,只要Deligne——他在
上同调中一言九鼎,而同时又是唯一被认为了解这些 Motif 底细的人——自己还在对 Motif 沉默不表,
就并不处于“把 Motif 当回事”的心理状态。(6月8日)核实之后发现,我最初的 Motif 思考可追溯到六
十年代初——因此它们持续了将近十年。},把它完全显式地陈述出来,本不是什么难事(而且会非常迷
人!)。其实在“离开数学”前不久,我离做这件事只差一步。

从某些方面看,处境与微分积分学英雄时代的“无穷小”颇为相似——不过有两点不同。首先,今天
的我们拥有建造精微数学理论的经验与一套有效的概念装备,那是我们的先辈所缺的。其次,尽管有这
些手段,而这一分明本质的观念出现已二十多年,却没有人屈尊(或敢于,尽管有人不屑于……)
撸起袖子,提炼出 Motif 理论的大致轮廓,正如我们的先辈直截了当地为无穷小计算所做的那样。然而
如今,对 Motif 就像当年对“无穷小”一样清楚:这些生灵是存在的;只要你关心代数簇以及这种簇的族
的上同调,尤其关心它们的“算术”性质,它们就在代数几何中随时显形。也许比对我谈到的那四种观念
更甚,“Motif”——所有观念中最特殊、最丰富的一个——系于一大堆形形色色的直觉:绝非含混,而往
往可以\dmark{}以完美的方式\origpage{184}表述出来(必要时,偶或承认一些 Motif 式的前提)。这些“Motif 式”
直觉中最迷人的一个,在我这里,是“Motif Galois群”:它在某种意义上,使人得以在(绝对意义的)域
与有限型概形的投射有限Galois群上“配一个 Motif 结构”。(为给这一观念以确切意义所需的技术工
作——以那些为 Motif 观念提供暂时基础的“前提”来表述——是在Neantro~Saavedra论“塔纳卡
范畴”的论文里完成的。)

时下的共识,对 Motif 观念比对它的三个受难兄弟(姐妹)(导出范畴、所谓“六种运算”的对偶形式
主义、topos)要稍许客气些:它没有被干脆斥为“涂鸦”\fnote{15}{(*)}{如我在前面一条脚注里提过的,导出范
畴三年前曾享受一场锣鼓喧天的掘墓礼(只是没有提我的名字)。topos 与六种运算还在等待它们的时
刻,Motif 也是——除了两年前被部分开棺出土、换上另一位父亲的那一小块(见注n\textsuperscript{o}s~51、52、59)。(5月13日)},但实际上这毫无区别:只要没办法“定义”一个 Motif、没法“证明”点什么,严肃的人
们就只能绝口不谈(怀着最大的遗憾——这话是听说了的,可人要么严肃要么不严肃……)。当然,只要
人们宣称谈它本身都不严肃,那就永远别指望建造 Motif 理论、就它“证明”点什么!

可那几个在局中的人(以及制造时尚的人)自己心里非常清楚:借助那些保密的前提,可以证明许多
东西。也就是说,从今天起——实际上,自从这一观念在Weil猜想的余波中出现起(它们倒是被Deligne
证明了,这总算是个加分项!)——\textbf{Motif 的瑜伽}就实实在在地存在。但它有着一门\textbf{秘密科学}的
地位:入门者肯定寥寥\fnote{16}{(**)}{(5月13日)我现在似乎终于明白,截至1982年,所谓“极少数入门者”,其实仅有Deligne一个人。的确,他披露了这门“秘密科学”中透过某些重要结果显现出来的部分;这些结果包含在这套瑜伽之中,他随着自己能够证明它们而逐一公布,借此获取功劳,同时隐瞒仍然保密的灵感来源。可是,十五年来竟无人行动起来,最终投身于一门规模宏大的Motif理论,这分明说明,我们的时代远远没有无穷小计算英雄时代那种大胆的活力!}。它尽管“不严肃”,却仍然\dmark{}\origpage{185}使这些罕见的入门者能够在许
许多多上同调的处境中说出“人们有权从中期待什么”。由此衍生出许许多多的直觉与局部猜想;借助“瑜伽”提供的理解之光,人们有时事后便能用手头现成的手段处理它们。Deligne的若干项工作受到这套瑜伽启发\fnote{17}{(*)}{(5月13日)我终于对文献稍有了解。如今我看到,Deligne的全部著作都扎根于这套瑜伽。我的文献抽样以及其他交叉核对使我推测:在他的全部著作中,唯一提及这一来源的地方,是1970年《霍奇理论I》中的一句简短文字,将我与Serre一并提到。(见注n\textsuperscript{o}s~$78'_{1}$、$78'_{2}$。)},其中一项(如果我没有弄错)是他最早发表的工作:证明代数簇的射影光滑态射的Leray谱序列退化(为证明所需,取特征零)。这个结果是由有关“权”的考虑所
提示的;这类考虑具有算术性质——这些正是典型的“Motif 式”考虑,只不过被表达在 Motif 的“几何”的术语里。Deligne证明
这个命题靠的是Lefschetz--Hodge理论,而对那动机(49)(如果我记得不错)只字不提——没有它,谁也不会想
到去怀疑如此难以置信的东西!

Motif 的瑜伽,恰恰最初正是从我得自Serre的这套“权的瑜伽”之中诞生的\fnote{18}{(**)}{我从Serre那里(六十
年代初?)得到的,是一个观念或直觉:有一个重要的东西要理解!它起了一次初始推动的作用;此后几
年的思考紧随其后,先围绕一个“权的瑜伽”,不久又围绕一个 Motif 瑜伽。}。是他让我领略了“Weil猜想”的全部魅力(后来有了“Deligne定理”)。他向我解释了(在假定所论特征下存在奇点消解的
前提下)人们如何能凭权的瑜伽,给任意域上每个概形(不必光滑也不必真)配上“虚拟 Betti 数”
——这一点当时给我极深的震动(注46$_{9}$)。我相信,正是这个观念,是那场围绕权的瑜伽的思考的起
点;那场思考在随后几年里,一直在我奠基成文的任务的边缘上继续。(这也是我在七十年代重新拾起
的思路,借助任意基底概形上的“虚Motif”观念,旨在至少为虚Motif建立“六种运算”的一个形式主义
\origpage{186}。)如果说这些年里我一直对Deligne(我的得天独厚的对话者)以及凡愿听的人\fnote{19}{(*)}{(4月10日)我觉得,Deligne是唯一“听进去”的人——而且他小心地把听进去的东西的专属特权留给了自己。诚然,另一方面,写下这些收尾文字时,我“落在”了事件后面:两年前,Motif瑜伽已有一场部分的开棺出土,却对我的角色不着一字!对此可参见由一项意外发现所引发的注n\textsuperscript{o}~50、51、59——那项发现,对这场持续了十二年的下葬的意味,投下一束(至少对我而言)意外的光。直到那时,我还只是相当混乱地觉察到某种下葬的存在,却不曾费心凑近去看看……}讲这个
Motif 瑜伽,那当然不是为了叫他们把它当作一门只属于他们自己的秘密科学。($\Rightarrow$
注47,原书第271页)

\notehead{46$_{1}$}{}

除了以当年我对Riemann--Roch定理所给的表述(以及我据此给它的两种证明),还有它的各种变体而
引入的观念与观点之外,我几乎不作例外。如果我记得不错,这类变体是在 SGA~5 讨论班1965/66年
的最后一次报告里展开的;而那次报告连同同一讨论班的其他若干报告,都遗失得一干二净。最有意思的一种,我看是可构造离散系数的一条变形;我不知道它此后在文献里是否得到过显豁\fnote{20}{(**)}{(6月6日)我重新找到了它——以一个近旁的形态,顶着“Deligne--Grothendieck猜想”这个动听的名字——在MacPherson1974年发表的一篇文章里。详情见注n\textsuperscript{o}~87$_{1}$。}。值得注意:这一条同样容许一个“Motif”变形,它实质上等于断言:对互异的素数$\chi$(与剩余特征互素),与$\chi$-adic可构造层相配的“示性类”(在正则概形$Y$的Chow环里),当这些层出自同一个“词”——同一个motif(比方说,都是对一个给定的$f:X\to Y$的$R^if_!(\mathbb{Z}_\chi)$)——时,全都相等。

\notehead{46$_{2}$}{}

这个形式主义,可以看作同调中一种“整体对偶”形式主义的某种精粹;以其最“有效”的形
状——撇清了所论“空间”与应用的任何光滑性前提,或态射的真性——它还需要用一个局部对偶形
式主义来补足:在其中人们区分所谓“\textbf{对偶化}”的对象或“复形”\dmark{}\origpage{187}(在 $Lf^{!}$ 运算下稳定
的观念)——并由此有一个“\textbf{双对偶定理}”(以 $R\mathcal{H}om$ 表述),适用于满足合宜的有限
性条件(关于次数,以及局部层面的凝聚性或“可构造性”)的系数。当我说“六种运算的形式主义”时,
我通过后者所暗指的,是对偶性——其“局部”与“整体”两个方面。

深入理解上同调中对偶性的第一步,是六种运算的形式主义在第一个重要情形——诺特概形及上同调凝聚的模复
形——中的逐步展开。第二步是发现(在概形的平展上同调中)这个形式主义对离散系数同样适用。这
两个极端情形,足以确立这种对偶性之\textbf{无处不在}的信念——对Poincaré型“对偶”的一切几何处
境的信念;这一信念,此后又得到Verdier、拉米、吕盖等人的工作的印证。当那股十五年
来压在这套形式主义上、阻碍其大规模使用的\textbf{阻碍}崩解之后,它必将在任意类型的系数上再
度得到确认。

这种无处不在,在我看来是一个分量极重的\textbf{事实}。它使我们不能不感到Poincaré对偶与Serre对偶之间有一种深刻的统一性——Mebkhout最终以所需的一般性确立了它。这种无处不在,使“六种运算的形式主义”成为
同调代数中理解“全方位”上同调对偶现象的基础结构之一\fnote{21}{(*)}{有兴趣的读者,可在本卷的附录里找到这一形式主义的一个概要。}。至于这套颇为精微的结构竟未被行业内
的人显式陈述出来——连“三角范畴”的“好”观念(Verdier的版本仍是一种十分暂时且不充分的形状)
也不曾有人提出——这一事实毫不改变什么;拓扑学家乃至偶或关心上同调的几何学家,仍然竞相无视
这套对偶形式主义的存在,以及奠基它的导出范畴的语言,也同样无济于事。

\notehead{46$_{3}$}{}

D-模与微分算子复形的观点,是由Sato修引入的,最初由他本人及其学派发展,所循的\dmark{}\origpage{188}观点(据我的理解)与Mebkhout所遵循的相当不同;而后者的进路更贴近我的进路。

“可构造性”的诸种观念——对(复解析的、实解析的、分片线性的)“离散”系数——据我记忆,是由我在五十年代末前后首次提炼出来的(几年后我在平展上同调的框架里又重提了它们)。当时我提过一个关于这种观念之稳定性的问题:它对实或复解析空间之间的一个真态射的高阶直像是否稳定——并且不知道这种稳定性在复解析情形是否已被确立\fnote{22}{(*)}{(5月25日)它已由 J.L.~Verdier确立,见“好的参考文献”一注(n\textsuperscript{o}~82)。}。在实解析情形,我当年设想的观念本来就不是对的,因为没有Hironaka的实亚解析集观念——它带有经由直像的本质先行稳定性。至于 $R\mathcal{H}om$ 一类的局部运算,当时清楚的是:在特征零优秀概形的框架里(利用Hironaka奇点消解)确立可构造系数之稳定性的论证,在复解析情形原封不动地成立,双对偶定理亦然(见 SGA~5~I)。分片线性的框架里,自然的稳定性与双对偶定理是“容易的练习”——我当年很高兴做它们,以此检验对偶形式主义的“无处不在”,那正是在平展上同调起步之时(其一大意外恰恰是发现了这种无处不在)。

回到半解析的情形:这一方向上稳定性定理(经由六种运算的可构造系数)的“好”框架,分明是“温和空间”的框架(见《纲领概要》第5、6节)。

\notehead{46$_{4}$}{}

当然,D-模的观点,连同 $\mathcal{D}$ 是一个凝聚环层的事实,为模的晶体揭示了一种比我所惯用的更隐蔽的“凝聚性”观念;它在不必光滑的空间(解析的或概形的)上仍有意义。称之为“\textbf{M-凝聚性}”($M$ 即Mebkhout)才算公道。于是,对一个多少“在局中”(而且健全的数学直觉丝毫不爽)的人来说,以下应当相当显然:那条推广了光滑情形“微分算子复形”的\dmark{}“好的系数范畴”\origpage{189},不应是别的,正应是模的晶体的“M-凝聚”导出范畴(一个晶体复形称为 M-凝聚的,若其各上同调对象如此)。它在没有光滑性假设下仍有合理的意义,并且应当同时涵盖寻常的“连续”(凝聚)系数理论与“可构造”离散系数理论(对后者引入合宜的完整性与正则性假设)。倘若我的视野是对的,Sato--Mebkhout理论相对于先前已知的晶体语境而言,其两个新的概念配料,就是模的晶体的这种 M-凝聚性观念,以及关于 M-凝聚晶体复形的(更深刻的)完整性与正则性条件。这些观念一旦到手,第一桩本质的工作,就是在晶体语境中发展六种运算的形式主义,以便涵盖二十多年前我已发展的两个特殊情形(寻常凝聚与离散)——而我某些前学生上同调学家,早就为了想必更重要的任务把它们忘得一干二净……

Mebkhout其实早已通过研读我的著作,最终得知“晶体”观念的存在;他也感到他的观点(至少在特征零)应当为这个观念提供一条好的进路——但这点建议落入了聋子的耳朵。从心理上说,很难设想他会投身于那桩势在必行的浩大奠基工程:他所处的,恰是那些以上同调权威自居的人的高傲冷漠的气候之中,而他们本是最有资格鼓励——或者说打消……——的人。

\notehead{46$_{5}$}{}

(5月13日)这里说的,主要是以一个\textbf{交换局部环}赋环的 topos。以拓扑空间上这样一个环层的数据来描述“簇”结构的想法,最先由 H.~Cartan引入,又被Serre在其经典工作 FAC(凝聚代数层)中采用。正是这项工作,为引我走向“概形”观念的思考提供了初始推动。Cartan的进路(经Serre采用)所缺的,为涵盖迄今出现的所有类型的“空间”或“簇”所缺的,正是 topos 的观念(即恰好是这样的“东西”:在其上“集合层”的观念有意义,并且具备那些熟悉的性质)。

\notehead{46$_{6}$}{}

\dmark{}作为 topos 的其他出色的例子——它们不是寻常的空间,\origpage{190}而且以“通行”观念似乎也没有令人满意的替代品——我要指出:一个拓扑空间对一局部等价关系所得的商 topos(比如流形的叶状结构,此时商 topos 甚至是一个“多簇”(multiplicité),即局部是一个簇);对几乎任何一种数学结构(至少那些“以有限投射极限与任意归纳极限来表述的”)的“分类”topos。当人们取一种“簇”结构(拓扑的、微分的、实或复解析的、Nash 的,等等,甚至取一个给定基底上光滑的概形结构)时,每种情形都会出现一个特别诱人的 topos,配得上“普遍簇”(所论那一种的)之名。它的同伦不变量(尤其它的上同调——配得上所论簇类的“分类上同调”之名)早就应当被研究和熟知;但到现在,这条路还根本没有被走过……

\notehead{46$_{7}$}{}

这说的是这样一些空间 $X$:其同伦型被“以自然的方式”描述为一个复代数簇的同伦型。后者于是可定义在复数域的一个子域 $K$ 上,$K$ 是素域 $\mathbb{Q}$ 的有限型扩张。投射有限Galois群 $\mathrm{Gal}(\overline{K}/K)$ 于是自然地作用在 $X$ 的投射有限同伦不变量上。往往(比如 $X$ 是一个奇数维同伦球面时)可取 $K$ 为素域 $\mathbb{Q}$。

\notehead{46$_{8}$}{}

(5月13日)当我在Serre的论文 FAC 里学到代数几何的最初几课(正是它把我“触发”向概形)时,换基这个观念本身在代数几何里还几乎无人知晓,除非在换基域这个特殊情形。随着概形语言的引入,这一运算变成了代数几何中最常用的运算,它在每个时刻都会出现。这一运算在拓扑学里至今仍几乎无人知晓(除一些很特殊的情形),在我看来,是拓扑学相对于来自代数几何的思想与技术的孤立的一个典型征兆(在许多其他征兆之中),也是“几何”拓扑学那套不充分基础的一个顽固遗产。

\notehead{46$_{9}$}{}

\dmark{}(6月5日)Serre的想法是:\origpage{191}对一个域 $K$ 上任意有限型概形 $X$,应能配上整数
$$h^i(X)\quad(i\in\mathbb{N})$$
他称之为“虚拟 Betti 数”,使得:
\begingroup\renewcommand{\labelenumi}{\textbf{\alph{enumi})}}
\begin{enumerate}
\item 对 $Y$ 一个闭子概形、$U$ 其补开集,
\[
h^i(X) = h^i(Y) + h^j(U).
\]
{\small [译按:原书此式右端印作 $h^j(U)$,疑应为 $h^i(U)$。]}
\item 对 $X$ 射影光滑,有
\[
h^i(X) = X\text{ 的第 }i\text{ 个 Betti 数}\text{。}
\]
{\small [译按:原书左端印作 $h(X)$;此处按前后定义保留译稿的 $h^i(X)$。]}
\end{enumerate}\endgroup
(比如经由与 $k$ 的特征互素的 $\chi$-进上同调定义)。如果承认 $k$ 上代数概形的奇点消解,那么立刻可见诸 $h^i(X)$ 由这些性质唯一确定。这样一个映射 $X\to(h^i(X))_{i\in\mathbb{N}}$($k$ 固定)的\textbf{存在},利用具真支撑的上同调形式主义,可本质上归约为基域有限的情形。在 $\mathbb{Q}_\chi$ 上有限维的、$\mathrm{Gal}(\overline{k}/k)$ 连续作用的向量空间的“Grothendieck 群”里工作,取 $X$ 的 $\chi$-adic(具真支撑的)Euler--Poincaré特征在该群中的像,$h^i(X)$ 即指 $EP(X,\mathbb{Q}_\chi)$ 的“权 $i$ 分量”的虚拟秩,其中权的观念由Weil猜想导出,外加奇点消解的一个弱形式。即使没有消解,Serre的想法也因Weil猜想的强形式(Deligne在“Conjectures de Weil II”中建立)而得以实现。

我沿着这条路做过启发式的思考,通向“虚拟相对概形”的六种运算形式主义——基域 $k$ 被换成一个多少任意的基底概形 $S$——以及这类($S$ 上有限表现的)虚拟概形的各种“示性类”观念。于是,我(为简化回到基域的情形)被引向考虑比Serre的更细的整数不变量,记作 $h^{p,q}(X)$,满足与上面 a)、b) 类似的性质,并且由通常的公式
$$h^i(X)=\sum_{p+q=i}h^{p,q}(X)$$
给出Serre的虚拟 Betti 数。

\subsection{2.~拒受遗产——或矛盾的代价}

\notehead{47}{[该注直接续注46,见 13.2.1]}


\pmark 值得注意的是,我方才巡视的五种观念之中,有四种(恰恰是那些被当作“不严肃之事”的观念)关系到上同调,\origpage{192}而且首先是\textbf{概形与代数簇的上同调}。不管怎样,这四种观念全都是代数簇的上同调理论的需要向我提示的——先是连续系数,后是离散系数。这就是说,在1955至1970这十五年间,我诸项工作的一个主要动机与一条不变的主导旋律,正是代数簇的上同调。

引人注目的是,这也正是Deligne至今仍视为其灵感主要来源的主题——如果我相信去年IHES小册子就此所说的话\fnote{23}{(*)}{(5月12日)然而,我刚发现,那本小册子里没有任何东西会让读者疑心:我的著作竟与代数簇的上同调、或与任何别的东西的上同调有什么相干!对此可参见即日所写的“葬礼悼词(1)——或恭维”一注(n\textsuperscript{o}~98)。所论小册子,即“有益的挣脱”一注(n\textsuperscript{o}~42)脚注里提到、并在刚提到的“葬礼悼词”注中得到稍细考察的那一本。}。我得知此事时颇为惊讶。诚然,当Deligne(在他关于Ramanujan猜想的漂亮工作(译按:原书人名拼作Ramanuyam,此处按通常拼写订为Ramanujan.)之后)发展了他对Hodge理论的卓越推广时,我还“在场上”,而且完全“在行”。那首先——姑且从这里开始!——对他和对我都一样,是朝着在复数域上以构造性方式建立Motif观念迈出的第一步。在我1970年“转向”之后的最初几年,我自然也听到了Deligne证明Weil猜想的消息(这同时证明了Ramanujan猜想),还有紧随其后的正特征下的“蛮横的”Lefschetz定理。我对他的期待不比这更少!我甚至确信,他应当同时证明了那些\textbf{标准猜想}——那是我在六十年代末提出的,作为(至少)在一域上奠基“半单”Motif观念的第一步,并把这些Motif的若干预期性质翻译成$\chi$-进上同调性质与代数闭链群性质的语言。Deligne后来告诉我,他对Weil猜想的证明肯定无法用来证明那些(更强的)标准猜想,而且他对如何着手也毫无想法。这应当已有十年左右了。从那时起,我再没有听说,在对代数簇上同调之“Motif”(或“算术”)方面\dmark{}\origpage{193}的理解上,发生过其他真正决定性的进展。了解Deligne的手段,我曾默然断定他的主要兴趣必然已转向别处——故而我读到情况并非如此时,深感惊讶。

有一点在我看来是无可怀疑的:至少二十年以来,若不同时在某种程度上扮演“Grothendieck的继承人”,人们已几乎不可能在代数簇上同调的理解上做出大面积革新的工作。Zoghman~Mebkhout已经切身领教了这一点,而(在某种意义上)Carlos~Contou-Carrère亦然——他很快明白,换题目才是上策(注47$_{1}$)。在那些人们无法回避的事情之中,恰恰首推在尽可能贴近Motif的各种系数语境里(这些Motif目前扮演着一种理想“地平线”的角色)发展那著名的“六种运算形式主义”:特征零的晶体系数(沿着Sato学派与Mebkhout的路线,做成“Grothendieck风味”),或特征$p$的晶体系数(主要由Berthelot、Katz、Messing以及一群显然干劲十足的更年轻的研究者研究);Deligne式的“分层pro-模”(promodules stratifiés)(它们显现为凝聚D-模或D-凝聚晶体这一“ind”观念的对偶化变体,即“pro”变体);最后是“Hodge--Deligne”系数(它们看来与Motif一样好,只除了一点:其定义是超越的,且局限于复数域上有限型的基底概形)……在另一端,则摆着这样一个任务:把Motif观念本身从笼罩着它的迷雾中(且不无原因……)提取出来,并且如果可能,去攻克像“标准猜想”这样精确的问题。(关于后者,我曾设想过,除其他途径外,发展一域上射影光滑簇的“中间Jacobi簇”理论,作为或许能得到迹的正性公式的一个办法——那公式是标准猜想的一个本质配料。)

直到我“退出数学”的那一刻,这些任务与问题还在我手里灼烧——滚烫而多汁的东西,其中任何一件、在任何时刻都从不曾在我面前显形为一堵“墙”、一个终止点\fnoteplainbody{24}{(*)}{(5月25日)然而,那本大名鼎鼎的庆典小册子里,一管我自信认得出的匿名手笔,却正是在友好地暗示这一点。参见“葬礼悼词(2)”一注,它接着前一条脚注所引的“葬礼悼词(1)”。}。它们是一个灵感与实质\dmark{}\origpage{194}永不枯竭的源泉——在那里只需朝涌出的地方一拉(而且它“四处涌出”!),如期而至的与不期而至的都会到来。以我这般有限的手段,但只要在工作中不分心,我清楚地知道,只要肯下功夫,一天、一年或十年,人能做成多少事。我也知道——因为我曾在一个他不分心的年代亲眼见过——Deligne的手段是什么,只要他愿意投入,他一天、一周或一个月能做成什么。可是没有人,连Deligne也不能,长此以往地一面从高处俯视那归根结底需要探查的对象本身,连同其语言与整套由某位先行者为这一目的而发展的工具库(何况还是在后者的协助之下,更不必说众多同样参与了劳作的人……),一面做出丰饶的、深刻革新的工作(注59)。

我也想到$M_{g,\nu}$(于$\mathrm{Spec}\,\mathbb{Z}$上)的“Deligne--Mumford”紧化——它参数化亏格$g$、带$\nu$个标记点的连通光滑代数曲线。引入这些对象\fnote{25}{(*)}{载 Pub.~Math.~36, 1969, 第75--110页。评论见注n\textsuperscript{o}~63$_{1}$。},是为了解决在全特征下证明模空间$M_{g,\nu}$的连通性这一问题,办法是从特征零做特化论证。在我看来,这些对象$M_{g,\nu}$(连同群$\mathrm{Sl}(2)$)是我在数学中遇到过的最美、最迷人的对象(注47$_{2}$)。单单它们的存在,带着如此完美的性质,在我眼中就是一桩奇迹(而且是得到了透彻理解的奇迹),其分量远非那桩本待证明的连通性事实可比。对我而言,它们以精髓的方式蕴含了代数几何中最本质的东西,即(几乎所有)一切代数曲线(在一切可以想象的基底域上)的全体——而这些曲线恰恰是一切其他代数簇的终极建筑材料。但所论这类对象——“$\mathrm{Spec}\,\mathbb{Z}$上真且光滑的多重簇”——仍逸出“通行”范畴之外,也就是说,逸出那些人们(出于从不检视的原因)肯于“承认”的范畴之外。常人谈起它们至多闪烁其词,而且一脸抱歉,仿佛还在做什么“泛泛的无稽之谈”,尽管人们小心地会说“stack”或“叠”,以避讳“topos”或“多重簇”这个禁忌的字眼。毫无疑问,这正是这些独一无二的珍宝自十多年前被引入以来,除了我本人在一些未发表的研讨班笔记\dmark{}里之外,再无人(就我所知)研究或利用的原因。取而代之的是,\origpage{195}人们继续要么用“粗糙的”模簇空间,要么用模多重簇的那些有希望是真概形的有限覆盖——然而无论前者还是后者,都不过是它们所自出的那些完美珍宝的某种苍白而跛脚的影子,而后者实际上仍遭放逐……

Deligne的四项工作——关于Ramanujan猜想、关于混合Hodge结构、关于模多重簇的紧化(与Mumford合作)、关于Weil猜想——每一项都是我们对代数簇之认识的一次革新,因而同时也是一个新的起点。这些奠基性的工作相继出现于几年之间(1968--73)。然而将近十年过去,这些巨大的界标并未成为向着预感与未知发起新一轮冲击的跳板,成为更广泛革新的手段。它们通向的是一种阴郁停滞的局面(注47$_{3}$)。这肯定不是因为十年前各家各人手头的那些“手段”像被施了魔法般消失了;也不是因为我们触手可及的那些美好事物突然化为乌有。而是:世界美是不够的——还得有人肯为之欢喜……

\notehead{47$_{1}$}{}

我在这里想到Contou-Carrère在五六年前一个颇有希望的开端:相对局部Jacobi簇的理论,它们与整体Jacobi簇(所谓“广义Jacobi簇”)的联系——对任意概形上的光滑且不必真的曲线概形——以及与Cartier的交换形式群及典型曲线理论的联系。除去Cartier的一次鼓励性回应之外,最有资格赏识Contou-Carrère第一篇注记的人们,对它的接待如此冷淡,以致作者从此绝口不提他储备着的那第二篇,并赶紧换了题目(尽管如此仍未能躲过其他的噩运)\fnote{26}{(*)}{(6月8日)见“棺材3——或有点太相对的Jacobi簇”一注(n\textsuperscript{o}~95)的子注(951)。}。

是我向他建议了局部与整体Jacobi簇这个题目的,作为迈向一个纲领的第一步——那纲领可上溯至五十年代末,尤其指向一个任意维的“阿代尔”对偶复形理论,由局部Jacobi簇构成(对应局部环\dmark{}\origpage{196}任意维的情形),以类比于诺特概形的剩余复形(由其一切局部环的对偶化模构成)。我这一上同调对偶纲领的这一部分(连同其他部分),由于六十年代涌入的种种显得更为紧迫的任务,多少被搁置进了被遗忘的角落。

\notehead{47$_{2}$}{}

确切说来,是所有这些多重簇的族所嵌入的那座“Teichmüller塔”,以及这座塔以基本群胚表述的离散或投射有限范式,构成我在数学中遇到过的最丰富、最迷人的单一对象。群$\mathrm{Sl}(2,\mathbb{Z})$,连同其投射有限紧化的“算术”结构(即$\mathrm{Gal}(\overline{\mathbb{Q}}/\mathbb{Q})$在其上的作用),可视为这座塔的“投射有限版本”的主要建筑材料。译按:原书此处群名的第二个字形似 $\chi$,这里依前文 $\mathrm{Sl}(2)$ 的记号书写。关于此可参见《纲领概要》中的指示(且待《数学反思》诸卷中专论这一主题的卷帙)。

\notehead{47$_{3}$}{}

这一“阴郁停滞”的观察,并不是一个对近十年间围绕概形与代数簇上同调的主要事件了如指掌的人深思熟虑的意见。这是一个“局外人”的单纯的整体\textbf{印象},它是我1982与1983年间同Illusie、Verdier、Mebkhout的谈话与通信中(除其他来源外)得出的。这一印象当然在许许多多方面有待细化。比如,Deligne1980年发表的“Conjectures de Weil II”,就代表一个实质性的新进展,虽说就其主要结果而言算不上一桩意外。看来特征$p>0$的晶体上同调也有进展,更不必说围绕相交上同调的那阵“抢滩”——它终于迫使某些人(尽管不情愿)回到导出范畴的语言,甚至使他们想起那些久遭否认的父系渊源……

\section{III.~时尚——或名人的生活}

\subsection{1.~本能与时尚——或强者的法则}

\notehead{48}{[该注由注46所召,原书第265页]}

众所周知,导出范畴的理论归于 J.L.~Verdier。在他着手\origpage{197}我提议的那项奠基工作之前,我只是以启发的方式、以这些范畴的一个临时定义(后来证明是对的)同它们打交道,对其内部本质结构也只有一个同样临时的直觉(这一直觉在预期的语境里被证明技术上是错的:“映射锥”并不\textbf{函子地}依赖于导出范畴中定义它的那个箭头——它只在相差一个不唯一同构的意义上定义它)。我在五十年代末发展的凝聚层的对偶理论(即凝聚框架里的“六种运算”形式主义)\fnote{27}{(*)}{其中还缺一个对非真态射的$Rf_!$运算(具真支撑的上同调);它是六七年后由Deligne引入的,有赖于他引入的凝聚pro-模语境——在我看来这是一个重要的新观念(后来被他成功地纳入其分层pro-模理论)。},其全部意义唯有在关于导出范畴观念的一项奠基工作之上才得以显豁,而那项工作是Verdier后来完成的。

Verdier学位论文的文本(迟至1967年才答辩),约二十页,在我看来是迄今所写的对导出范畴语言最好的导引,它把这一语言安置在其本质使用的语境之中(其中好几项使用出自Verdier本人)。它只是一项正在撰写的工作的导言,而那项工作最终也是后来才写成的。我可以自诩——即便不是唯一的一个——也是极少数能够作证曾亲手捧读过这部著作的人之一:仅凭这篇导言,其作者被授予国家博士头衔的正当性就得到了确立!这部著作(或者曾经是——我不知道是否还有一份副本存于某处……)是迄今唯一一部按导出范畴的观点给出同调代数的系统奠基的文本。

\pmark 也许只有我一个人在惋惜:无论是那篇导言还是那奠基的正文本身,都未曾\origpage{198}发表\fnote{28}{(*)}{(5月25日)在这些文字写成之后,我发现Verdier学位论文的最初胚芽——写于1963年,即答辩前四年——最终于1967年发表。对此可参见“同伙”与“赊账的学位论文与全险”两注,n\textsuperscript{o}~63$'''$与81。}——以致使用导出范畴语言所需的本质技术行装,散见于文献中三个不同的地方\fnote{29}{(**)}{这三处是:Hartshorne那场著名的凝聚对偶研讨班,其中有我在五十年代下半叶发展的对偶理论至今唯一发表的部分;SGA~4里Deligne的一两个报告;Illusie那部浩繁的学位论文中的一两章。}。这样一个系统参考文本的阙如——其分量本可比肩Cartan--Eilenberg的经典著作——在我眼中,既是1970年我退出数学舞台之后导出范畴形式主义所遭冷落的一个典型\textbf{原因},也是一个典型的\textbf{征兆}。

诚然,早在1968年就已显露(在SGA~5中发展的一种上同调迹理论的需要之际):原初形态的导出范畴观念,以及相应的三角范畴观念,对某些需要而言是不够的,更深彻的奠基工作尚待去做。Illusie在其学位论文中引入“滤导出范畴”,沿这一方向迈出了有用但仍属谨慎的一步(主要为了迹的问题)。1970年我的离去,似乎成了同调代数基础、乃至与之血肉相连的Motif理论基础的全部反思戛然而止、永不再续的信号(注48$_{1}$)。然而就前者而言,为大面积奠基所需的全部本质观念,似乎在我离去之前的年头里就已到手(注48$_{2}$)。(包括“求导者”(dérivateur)这个关键观念,或曰“导出范畴制造机”——它看来是至今所遇到的诸三角范畴之下共同的那个更丰富的对象;这一观念将近二十年后才在一个非加性框架里得到少许发展,载《追寻栈》第二卷的一章。)此外,尚待完成的奠基工作有相当一部分已由Verdier、Hartshorne、Deligne、Illusie做成,这些工作可以原封不动地用于一项综合,把已获的观念纳入“求导者”这个更开阔的视野。

\pmark 诚然,这十五年来\fnote{30}{(*)}{(5月24日)这“逝去的十五年”应当细化——参见注n\textsuperscript{o}~47$_{3}$,以及内容更详的“赊账的学位论文与全险”一注(n\textsuperscript{o}~81)。}对导出范畴观念本身的冷落——在某些人那里,它近于对一段过往的否认——正顺应了某种时尚:\origpage{199}这种时尚对一切奠基的反思都投以轻蔑的目光,无论其多么紧迫\fnote{31}{(**)}{(5月25日)关于这一时尚之出现与延续背后的诸种力量,可参见“掘墓人——或全体教众”一注,n\textsuperscript{o}~97。}。但从另一方面看,我也再清楚不过:平展上同调的发展——今天“人人”都在不加思索地使用它(哪怕只是通过已故Weil猜想的途径……)——若没有导出范畴、六种运算、site与topos的语言(它们最初正是为这一目的而发展)所构成的概念行装,更不必说SGA~1与SGA~2,是绝无可能的。同样再清楚不过的是:今天代数簇上同调理论中可以观察到的那种停滞,若非我昔日某些学生在那些年月里宁肯追随一场他们自己正是始作俑者的时尚、而不肯听从其作为数学家的健全本能——那场时尚久已借着他们的支持获得了法律的效力——原本既不可能出现,更不可能安营扎寨。

\notehead{48$_{1}$}{}

同样的话(带有某些保留)也可用于我那整个代数几何奠基纲领:它只有很小一部分得到实现,而且随我的离去戛然而止。这一停顿给我的震动,尤在于对偶纲领——我视之为特别多汁的部分。Zoghman~Mebkhout的著作,顶风冒雨地继续着,却仍处在这条纲领的脉络里(并因料想不到的观念的汇入而焕然一新)。1976年Carlos~Contou-Carrère的工作(注47$_{1}$,原书第273页已提及)亦然——那是他谨慎地无限期搁置了的工作。此外还有一项关于曲面fppf上同调对偶的工作(Milne)。我所知者,尽于此矣。

诚然,我从未想到要为1955至1970年间为我自行凸显的那项长期工作纲领写一篇概要,一如我为最近十二年写下《纲领概要》。我想原因很简单:从未出现过一个特别的机缘(如同如今我进入CNRS的申请),足以\dmark{}\origpage{200}促成这样一项陈述性的工作。在收入《模型的历史》(《数学反思》II)第I章附录的致Larry~Breen的信(1975年)里,可以找到一些指示,涉及我1970年前日程表上的某些理论(尤其是对偶理论)——这些理论至今仍在等待有人伸出双臂,把它们接引进公共的遗产之中。

\notehead{48$_{2}$}{}

同样的话对Motif理论也成立,只是有一点不同:它多半注定要在一段时期内保持其猜想形态。

\subsection{2.~当值无名之辈与好上帝定理}

\notehead{48$'$}{[该注由注46所召]}

按惯例,人们以完成了提炼并确立一项理论之关键定理的人的名字来命名这些定理;而Zoghman~Mebkhout的名字,看来却被判定配不上这条基本定理——那是四年顽强而孤独的工作(1975--79)的结晶,逆着当时的时尚与其长辈们的轻蔑而行。当这定理的分量再也无法被忽视的那一天,这些人兴致勃勃地把它称作“Riemann--Hilbert定理”,我完全相信(尽管Riemann与Hilbert肯定未曾要求过这份荣耀……)他们这样做自有极好的理由。毕竟(一旦需要之感——即理解一般离散系数与连续系数之间精确关系的需要——顶着普遍的冷漠显现出来,一旦它经由细致而耐心的工作得到提炼与明确,一旦经过相继的几个阶段那正确的陈述终于被提炼出来,一旦它白纸黑字地写下并被证明,而最终,当这件孤独的果实之定理在最不被人期待之处证明了自身——经过了这一切之后)这定理显得如此显然(对那些“本会证明它”的人来说,简直可以径直说“平凡”……),实在犯不着让一个模糊的当值无名氏的名字来占据记忆!

受了这一先例的鼓舞,我提议从今往后把任何一项理论中真正自然、真正基本的定理称作“亚当与夏娃定理”,或者干脆追溯到更远,把荣耀归于当得荣耀之处,径直称之为“\textbf{好上帝定理}”\fnote{32}{(*)}{在我作为数学家的一生中,我从未有过这份乐趣:在一名学生那里激发、哪怕仅仅是鼓励出一篇含有“好上帝定理”的学位论文——至少不曾有过深度与分量可比的。}。

\pmark 据我所知,除我之外,Deligne是Mebkhout之前唯一一个感到如下问题之价值所在的人:\origpage{201}在一个比分层模更宽广的框架里理解离散系数与连续系数的关系,以便能以“连续”的术语诠释任意的“可构造”系数。沿这一方向的第一番尝试,是Deligne1968或1969年在IHES主持的一场(未发表的)研讨班的内容:他在其中引入“分层pro-模”的观点,并给出一个比较定理(在复数域上),把超越的离散上同调与相伴的De~Rham型上同调联系起来——后者对任意特征零基底域上有限型的概形都有意义。(显然,他那时还不知道他那两位遥远的先行者Riemann与Hilbert的卓越结果……)因此,比起Verdier\fnote{33}{(*)}{Verdier作为Zoghman~Mebkhout学位论文的名义导师(并以此身份“恩准过几次讨论”),看来正是(除Mebkhout本人之外)围绕这条基本定理之父系渊源所发生的隐匿——以及在由Mebkhout所发展的D-模观点为代数簇上同调理论开启的革新中其“学生”应得的信誉——的主要当事人。然而我不知道他是否曾为此比Deligne更多一点不安。}或Berthelot\fnote{34}{(**)}{(5月25日)写这些文字时,我(略有踌躇地)没有把我的朋友Luc~Illusie的名字放进这份“本最有资格”向Zoghman~Mebkhout给予那些本属当然的鼓励的我方学生名单里。我当时未曾留意我内心某种不安——它本可以教会我:我这是在暗暗为一位我心怀感情的人开脱,装作卸下他那与我的其他“上同调学家学生”同样背负的责任。},Deligne都处在一个格外有利的位置,足以赏识Mebkhout1975年所投入的研究方向的全部价值,以及此后Mebkhout诸结果——尤其是“好上帝定理”——的价值;这条定理以连续系数的语言,给出了比他本人所提炼的更精妙、更深刻的离散系数领会。但这并未妨碍Mebkhout不得不在一种令人难堪的精神孤立中继续其工作,也未妨碍属于他的信誉(我更愿意说:尤其属于他的信誉)至今——五年之后\fnote{35}{(***)}{(5月25日)其实,这一隐匿首先是Deligne与Verdier本人的作品。对此可参见“不义——或归来的意义”一注,n\textsuperscript{o}~75。}——仍遭藏掖。

\subsection{3.~罐装的权与十二年的秘密}

\notehead{49}{[该注由注46所召]}

\pmark 核对之后(载 Publications~Mathématiques~35, 1968),我发现:在《Lefschetz定理与谱序列的退化判据》一文临近结尾处\origpage{202},有三行文字影射了某些“权方面的考虑”,正是它们使我(以一个略欠一般性的形式)猜想出了该文的主要结果。我怀疑这句谜语式的影射能否对任何人有用,或者在当时能被Serre和我之外的第二个人读懂——反正我们俩本来就已知情\fnote{36}{(*)}{(4月29日)对这篇文章的更细考察——它从不止一个角度看都富有教益——见“驱逐”一注(n\textsuperscript{o}~63)。}。

借机指出:早在六十年代末,在Weil猜想的余波里,一部非常精确的“权的瑜伽”——包括对$R^if_*$与$R^if_!$这类运算下权的行为的刻画——对我说来(因而也对Deligne)已是熟知的了。这部瑜伽的一部分最终在Deligne的“Conjectures de Weil II”(《Publications Mathématiques》1980)中得到确立(在$\ell$-进系数层的语境里;更自然的Motif框架尚待其确立)。如无差错,在前后相隔的这十二年间\fnote{37}{(**)}{(4月19日)我刚收到并饶有兴味地读了Deligne的论著目录,从中发现,早在1974年,他在温哥华大会上的一次报告中就谈到了“权”——如此,“权周围的秘密”是六年,而非十二年。不过在我看来,这一秘密与围绕Motif的类似秘密(1970--1982的十二年间)密不可分。这一秘密的意味,在今天的反思中、在n\textsuperscript{o}~51--52之后那则长篇双注里,刚刚被一束新的光照亮。},文献中没有留下一处踪迹——哪怕是最简短、最不完全的陈述——述说这部(当时还完全是猜想性的)权的瑜伽;在这整段时间里,它始终是寥寥(两三个?)知情人的专属特权\fnote{38}{(***)}{(5月25日)从反思中陆续出现的全部信息来看,这几个“两三个知情人”似乎到头来只有Deligne一人:看来他煞费苦心,把这部得自我的瑜伽的专属收益留给自己,直到1974年(见上一条脚注)——那时机缘已熟,他尽可以把它当作自己地里长出的观念来呈献,既不提我,也不提Serre(见注n\textsuperscript{o}~$78'_{1}$、$78'_{2}$)。

(1985年4月18日)在这些文字写成之后,我又有机会得知Deligne在尼斯国际数学家大会(1970)上的报告“Théorie de Hodge I”(《会议文集》t.1,第425--430页)。与我据手中残缺信息而有理由相信的相反,这篇文章早在1970年就陈述了权的瑜伽的相当一部分。关于这些观念的来源,他只满足于对Serre的一篇文章(它与本题其实无关)以及“Grothendieck的猜想性Motif理论”做了一句谜语式的、纯礼节性的提及。(可与注n\textsuperscript{o}~$78'_{1}$、$78'_{2}$对读。)至于权的观念在$R^if_!$与$R^if_*$这类运算下的行为这一关键问题,则连提都没有提,而且在1980年前引“Conjectures de Weil II”一文之前也不会提;在该文主要定理的关系中我的名字未被提及,同样,在前一条脚注(一年前的今天)所引的那次报告“代数簇上同调中的权”里,Serre的名字和我的名字也都未被提及。}。然而这部瑜伽构成理解代数簇上同调“算术”性质的第一把本质钥匙:它既是在一给定情形中辨认方向的\textbf{手段},也是做出种种预测的手段——其可靠性从未落空——\origpage{203}而同时,正因如此,它也代表着代数簇上同调理论中摆在我们面前最紧迫、最迷人的\textbf{任务}之一。这部瑜伽直到最终被确立(至少在某些重要方面)之前几乎无人知晓,在我看来,正是“\textbf{信息封锁}”的一个格外触目的例证:扮演这一角色的,往往恰恰是那些身居其位、肩负其职、理应确保信息广为传播的人\fnote{39}{(*)}{对此另可参见第32、33节“数学家的伦理”与“注记——或新伦理(1)”,以及与之相关的两注:“共识性的操守与信息控制”与“年轻人的势利,或纯洁的捍卫者”,n\textsuperscript{o}~25、27。}。

\subsection{4.~进步是挡不住的!}

\notehead{50}{[该注由第(I)部分《愚妄与更新》第八章“孤独的历险”第50节所召]}

我在这方面的最初经验,是我为设法让Yves~Ladegaillerie关于曲面同痕定理的学位论文获得出版而付出的那些徒劳努力所结出的意外果实——那项工作,说实话,丝毫不逊于我曾充当“导师”的那十一篇国家博士论文(“1970年以前”的,当然!)。如果我没记错,这些努力持续了整整一年多,主角中有相当一批我昔日的朋友(当然,还少不了我昔日的一个学生)\fnote{40}{(**)}{对此可参见“棺材2——或被切成一段段的剖视”一注,n\textsuperscript{o}~94。}。其主要情节,今天在我眼前仍如一场场通俗笑剧!

这也是我头一回遭遇到某种新精神与新习俗(它们在我昔日朋友圈里已成风气)——在这次反思中我已在这里或那里有机会影射过它们。正是在那一年(即1976年),我头一次(但不是最后一次)获悉:如今,认认真真去证明那些人人都在用、而先行者们从来满足于默认的精微事项(在这里,即曲面拓扑中不存在野性现象)\fnote{41}{(***)}{对此另可参见“注记——或新伦理”那一情节(第33节)。那则大名鼎鼎的“注记”的过错,恰恰在于把一些向来含糊其辞的观念与陈述挑明了;而我却曾一直含蓄地使用它们,来建立那些以我的名字命名的结果,而且将近二十五年来,人人都在毫不脸红地使用这些结果(这一点,那两位鼎鼎大名的同行其实心知肚明)。(6月8日)详情另见“棺材4——或没有鲜花也没有花冠的topos”一注(n\textsuperscript{o}~96)。所谓“以我的名字命名的结果”,是关于某些整体与局部投射有限基本群的生成与有限表现的结果,它们曾在SGA~1等处得到“证明”,所用的下降技术在没有理论上的、细致的论证的情况下仍属启发性质;那样的论证,在Olivier~Leroy那项(显然“无法发表的”)关于topos基本群的Van~Kampen型定理的工作中完成了。},被视为不够严肃(至少就一个无名之辈而言……)。或者,去证明一个把若干已知的深刻定理\dmark{}\origpage{204}作为特例或推论涵盖于其内的结果(这显然表明,那个所谓的新结果只能是已知结果的一个特例或一个容易的推论)。或者,仅仅是在一个结果的陈述里,或是在用一种语言描述另一种情形时,肯费心把自然的假设仔细表述出来(这是可叹的“无谓折腾”的标志),而不肯把自己局限于投某位高格调意见人士之所好的那种特例。(去年我还见到有人责备Contou-Carrère,说他在其学位论文里不该把一般概形、而只把基底域当作立脚点——同时又赏脸承认,他想必是在其临时导师的坚持下才不得不如此。不过,这样说话的人对门道也足够“在行”,本应知道:即便把自身限定在复数域上,证明的需要也会逼着人引入一般的基底概形……)

当今某种时尚的迷途,竟至于不仅鄙薄细致的证明(乃至干脆鄙薄证明本身),甚至常常鄙薄成形的陈述与定义。照纸张与被填鸭的读者的耐心目前的行情,不久之后,人们大概就不再肯为如此昂贵的奢侈品劳神了!顺着当前的趋势外推,应当能够预言这样一个时刻:届时出版物中将不再有人费神去明确陈述定义或定理,人们从此满足于用一些代码词来指称它们,把填补空白的差事留给自己那不知疲倦、天才横溢的读者,随他按自己的悟性行事。审稿人的差事也将随之大为轻松:他只需到“名人录”(Who is Who)年鉴里查一查,作者是否以可信著称(反正谁也没法反驳一个由空白与省略号组成的精彩论文),抑或相反,是一个不便启齿的无名氏——那他早就(如同今天已然如此、且由来已久)被当然地扫地出门……

% Original notes 27 and 28 have plain callouts but starred body labels.
\providecommand{\fnoteplainmark}[3]{%
  \begingroup\renewcommand{\thefootnote}{#1}\footnotemark
  \renewcommand{\thefootnote}{#1#2}\footnotetext{#3}\endgroup
}
\chapter{B)~皮埃尔与Motif}

\section{IV.~Motif——或一次诞生的下葬}

\subsection{一个梦的回忆——或Motif的诞生……}

\notehead{51}{[该注由注46所召]}

\pmark (4月19日)写下这些文字(它们是“我的孤儿们”一注,即n\textsuperscript{o}~46的收尾)还不到一个月,我已发现它们稍稍落在了事件的后面!我刚收到“Hodge Cycles, Motives and Shimura Varieties”(LN~900),作者为Pierre~Deligne、James~S.~Milne、Arthur~Ogus与Kuang-Yen~Shih——那是Deligne惠寄的,还附了一份他的论著目录。\origpage{205}这部汇集六个文本的文集1982年出版,是1970年以来一个有趣的新事实:Motif观念出现在标题里,并在正文中占有了一席之地——虽然还极其谦逊,主要是经由“MotifGalois群”的观念。当然,距离一部Motif理论的全景图还很远;十五二十年来,这幅全景图一直在等待那位勇敢的数学家,肯把它“描绘出来”——足够宽广,以充当一代或几代算术几何学家的灵感之源、引导线索与地平线,他们将有幸确立其有效性(或无论如何,发现Motif之实在的底细……)(注53)。

好像也是从1982年起\fnote{1}{(*)}{(5月25日)我又落后了——这一回是一年:转折发生在1981年6月的吕米尼会议;参见“不义——或归来的意义”一注,n\textsuperscript{o}~75。},时尚的风向对着导出范畴多少开始转了;Zoghman~Mebkhout(在一阵或许有点欣快的高谈里)已经看到它们即将“席卷数学的所有领域”。它们的用处——任何一个消息灵通的人,单凭数学本能,早在六十年代初就已看得很清楚——如今才开始被承认;这(在我看来)主要得归功于Mebkhout孤身的努力:整整七年,他顶着一场专横的时尚,以只信自己本能者的勇气,扛下了为人开路、承受初试代价这桩吃力不讨好的差事……

\pmark 引人注目的是,读这部第一次让(在我退出数学舞台十二年之后)Motif观念低调重返“通行”数学观念之列的出版物,任何不知情的读者都无从疑心:\origpage{206}我这个人,竟与这个长期遭禁的观念的诞生、与围绕它发展起来的一部丰富而精确的“瑜伽”,有过任何关联——在那里,它(以极其残缺的形态)宛如凭空出现,不提任何父系渊源(注51$_{1}$)。

仅仅三周之前,我花了一两页篇幅大谈Motif瑜伽——就连我的一个“孤儿”,而且是我比任何一个都更挂心的那个,也表明我当时显然完全没有说到点子上!毫无疑问,我是在做梦:当我自觉得记起多年孕育的岁月——一种先是微弱而难以捉摸、随后逐月逐年丰富与明确的视野,在那顽强努力的尝试中,试图把握诸种已知的上同调理论(注54)所各自体现的那个“motif”、那共同的本质——它们每一种都以自己的语言,谈论着作为其直接可触的表现之一的那个“motif”的本性。毫无疑问,我还在做梦:当我记起Serre的某个直觉给我留下的强烈印象——Serre被引向把一个投射有限Galois群(一个本质上属于离散性质的——或至少依定义可归结为一个\textbf{有限}群投射系统的——对象)看作这样一个系统的源头:一个庞大的$\ell$-进\textbf{解析}群、乃至$\mathbb{Q}_\ell$上\textbf{代数}群(经由合宜的代数包络)构成的投射系统;这些群甚至往往是约化群——随之,解析群与代数群的(Lie式)直觉与方法的整套武库也被引入。这一构造对每个素数$\ell$都有意义;而我感到(或者我梦见我感到……)这里有一个待探测的奥秘:对不同素数,这些代数群之间的关系如何;它们必定全都出自同一个代数群投射系统,位于所有这些基底域唯一自然的公共子域之上,即特征零的“绝对”之域$\mathbb{Q}$。既然我喜爱做梦,我就继续梦见:我记起自己经由一项肯定只是梦的工作(因为我什么也没有“证明”)进入了这个隐约瞥见的奥秘;我记起自己终于懂得,Motif观念如何提供了理解这一奥秘的钥匙——如何仅仅凭着一个范畴的在场(这里是给定基底概形\dmark{}\origpage{207}上的“光滑”Motif的范畴,比如给定基底域上的Motif),它带有与域$k$上一个代数pro-群的线性表示范畴上所见相似的内蕴结构(代数pro-群观念的魅力,先前也是Serre向我揭示的),人们就真能重建出这样一个pro-群(一旦有了一个合宜的“纤维函子”),并把“抽象”范畴诠释为它的线性表示范畴。

通往一部“MotifGalois理论”的这一进路,是我多年前为描述一个拓扑空间或概形(乃至任意一个topos——不过在这里我感到,自己要刺痛那些“topos并不能使他们开心”的娇贵的耳朵……)的基本群所找到的那条进路向我提示的:即经由所论“空间”的平展覆盖的范畴,以及其上的纤维函子。而“MotifGalois群”的\textbf{语言}本身(我也满可以叫它们Motif“基本群”——两类直觉,自五十年代末以来,对我是同一回事……),连同“纤维函子”的语言(它们与上文所说的那些“显明的化身”——即适用于一给定Motif范畴的各个“上同调理论”——恰相对应),正是为表达这些群的深层本性而造的语言,并且显然提示了它们与Galois群及寻常基本群的直接联系。

我还记得在纤维函子这场游戏里的乐趣与惊叹,还有在Galois群作用下、使不同纤维函子通过“扭转”彼此过渡的torseur——在一场格外具体而迷人的境况里,重新找回Giraud著作中所发展的非交换上同调的整套观念:由纤维函子构成的gerbe(这里的底是平展topos,或者更好,$\mathbb{Q}$的fpqc topos——非平凡而有趣的topos,若曾有过的话!),以及束住这条gerbe的“带”(lien,以群或代数pro-群为值),连同这条带的诸种化身——经由各样代数群或代数pro-群实现,对应于gerbe的各个“截面”,也就是各个上同调函子。比如,特征零概形的各个复点,经由相应的Hodge函子,生成了gerbe的那么多截面,以及从一个截面到另一个截面的torseur;这些torseur连同其上运作的pro-群,带有卓越的代数-几何结构,表达着Hodge上同调的特有结构——不过在这里\dmark{}\origpage{208}我抢在了Motif之梦的另一卷前面……那是这样的年代:今天引领时尚的人们尚未宣布,topos、gerbe之类不逗他们乐,谈这些就是胡扯(当然,就算他们这么说,也丝毫不会妨碍我在它们所在之处认出topos与gerbe……)。而如今十二年又过去了,同一些人却装出发现并教导的样子说:gerbe(即便还不是topos),确实与代数簇的上同调、甚至与阿贝尔积分的周期,大有关涉……

在这里我本可以唤起另一段记忆的梦(或另一个梦的记忆……),它也生于一记“强印象”(我分明深陷于主观性之中!):Serre就Weil猜想背后某种“哲学”所作的若干评论。把这些猜想翻译成$\ell$-进系数($\ell$变动)的上同调语言,令人隐约感到,相应的上同调带有卓越的结构——“权滤过”的结构\fnote{2}{(*)}{(1985年1月24日)对这段走了样的记忆的更正,见注n\textsuperscript{o}~164(I4)及子注n\textsuperscript{o}~164$_{1}$,其中对“权的瑜伽”的渊源有更准确的说明。}。不同$\ell$-进上同调所共有的那个“motif”,想必正是这一本质算术结构的终极载体;由此,这一结构取得了几何的面貌——即几何对象“motif”上的一种卓越结构。把它称作“工作”(尽管这分明仍旧不折不扣是猜谜的片段)肯定又是我在自欺——当需要“猜出”(唯一的向导,是一个正在成形的视野的内在自洽,借助于此处或彼处已知的或猜出的零星要素……)一个Motif的各个“化身”的上同调所特有的结构:权滤过\dmark{}\fnote{3}{(**)}{(1985年2月28日)这里我的头脑有一点点混乱。实际上,说的是与之紧密相连的按“层”(niveaux)的过滤。}在其中如何体现——从Hodge化身开始(在Hodge--Deligne理论尚未问世的时代,而且不无原因……\fnote{4}{(***)}{那是一个年轻的Deligne恐怕还没有在数学语境里听说过“概形”一词、也没听说过“上同调”一词的时候。(这些观念,是他自1965年起在与我的接触中认识的。)})。这使我在(梦中)得以看到,Tate关于代数闭链的猜想(这已是第三记“强印象”,它在Motif之梦里激励着做梦者!)与Hodge猜想(注55)汇入同一幅宏阔的图景,并从中提炼出两三条同一成色的猜想;我对某些人谈起过它们,他们想必已把它们忘了,因为我从此再没听到有人提起,正如再没人提起“标准猜想”一样\dmark{}。无论如何,\origpage{209}它们只是猜想(而且并未发表……)。其中一条并不涉及某个特定的上同调理论,而是给出:一域上非奇异射影簇的Motif上同调上的权滤过,经由该簇自身按给定余维数的闭子集所作的几何过滤的直接诠释(余维数扮演“重量”的角色)\fnote{5}{(*)}{(1985年2月28日)实际上,这里是按“层”(niveaux)的过滤(参见前一条脚注)。}。还有那项“猜出”权在六种运算下的行为的工作(我本该给“工作”打上引号,可又下不了决心!)(那些工作早已失传……)。在那里,我同样从不曾有发明之感,而始终是发现——不如说,是听事物向我言说,当我肯手握笔杆费心去听的时候。它们所说的,有一种不容错辨的精确。

然后有了第三个“Motif之梦”,它像前两个梦的联姻——这一回,是要诠释Motif范畴所带的各种附加结构(其中最首要的,恰恰就是权滤过):用MotifGalois群上的结构,以及以这些群为结构群、用来“扭转”一个纤维函子以(典范地)得到任一其他纤维函子的torseur\dmark{}\fnote{6}{(**)}{正如某个“空间”$X$在两个“点”$x$与$y$处的基本群$\pi_1(x)$、$\pi_1(y)$,经由“扭转”——即经由从$x$到$y$的道路类的群$\pi_1(x,y)$这个torseur——即可彼此约化……}的语言。我记得,在这里更谈不上什么猜谜,而是规规矩矩的数学翻译。它们是一道道未发表的关于代数群线性表示的“练习题”,我兴味盎然地做了几天、几星期,分明感到自己正一步步逼近一个多年来令我心醉神迷的奥秘!必须以表示的语言去把握与表述的观念里最精微的,或许是Motif的“极化”观念——我由Hodge理论取得启发,力图从中澄清在Motif语境中仍有意义的东西。这一反思,应当就发生在我思考“标准猜想”之表述的前后——两种思考\dmark{}\origpage{210}皆受启于Serre(永远是他!)关于Weil猜想的“kählérien”类似物的想法。在这样的境况里,当事物本身向我们低语:它们隐藏的本性为何,我们经由何种手段能最精妙、最忠实地表达它——尽管许多本质事实似乎还远在证明的射程之外——单纯的直觉告诉我们:把事物执拗地向我们低语的一切,径直白纸黑字写下来;我们越肯听写,写下的就越清晰!无须为完备的证明或构造操心——在这一阶段背负这类要求,无异于自绝于一项大面积发现工作中最精微、最本质的阶段:视野的诞生——它从一派表面的虚无中取得形态与实体。单单\textbf{书写}、\textbf{命名}、\textbf{描述}这一行为——哪怕起初只是描述那些难以捉摸的直觉、或那些迟迟不肯成形的“疑点”——就具有一种\textbf{创造性的力量}。当认识的热情投入理智所能把握的事物之中,这正是它手中的头号工具。在这一类事情的发现之路上,这一工作是所有环节中最具创造性的环节:它永远先于证明,并为我们提供证明的手段——或者不如说,没有它,“证明”某事的问题根本无从提出,因为凡是触及本质的东西,还没有一样得到表述与看见。

单凭表述努力的力量,无形的开始成形,变得可以审察;明摆着是错的与可能的得以澄清;而其中那些与已知的或猜出的事物之整体如此完美地协调的,又转而成为正在诞生的视野的坚实可靠的要素。这一视野在表述的工作中日益丰富、日益明确。十样仅仅被猜到的东西——没有一样(比如Hodge猜想)单独令人信服——却相互映照、相互补足,似乎汇入同一种尚带神秘色彩的和谐,便在这和谐中获得视野的力量。即便这十样到头来全被证明是错的,那通向这一临时视野的工作也没有白做;它使我们隐约瞥见、并容我们稍稍深入的那种和谐,不是幻象,而是实在,召唤着我们去认识它。正是经由这一工作,唯有经由这一工作,我们才得以与这一实在、这一隐藏而完美的和谐亲密接触。当我们懂得:事物之为此物自有其理\dmark{}\origpage{211},我们的天职是认识它们而非支配它们,那么,一个错误曝光的日子,就是欢腾的日子(注56)——正如某项证明使我们毫无疑义地确知:我们曾如此想象的某事,当真是实在本身的忠实而真正的表达——那一天同样是欢腾的日子。

在这两种情形里,这样的发现都是一项工作的报偿,没有那工作,发现无从发生。但它或许要待经年累月的努力之后才姗姗而来,或许我们永远听不到那最后的底细——那是留给后来者的——而工作本身就是它自己的报偿,在每一瞬间都因这一瞬间向我们揭示的东西而丰盈。

\notehead{51$_{1}$}{}

(6月5日)然而Zoghman~Mebkhout刚刚提请我注意:前引书第261页,在Deligne一篇“重写并补充一封致Langlands的信”的文章里,提到了“Grothendieck的Motif”。那里写道:“这里要说的不是Grothendieck的Motif——他那样用代数闭链来定义的Motif——而是\textbf{绝对Hodge Motif},同样用绝对Hodge闭链来定义。”这里点名“Grothendieck的Motif”(不加着重),不是作为灵感之源,而是为了与之划清界限,并强调说的是\textbf{另一回事}(人们特意加了着重)。这一撇清格外引人注目,因为Hodge猜想的成立(我假定Deligne,一如其信函式文章的任何读者——从最初的收信人Langlands起——都知道这条猜想)将蕴含这两个观念\textbf{同一}!!

当然,早在1964年——我发展MotifGalois群观念的那一年——我就清楚地知道,可以在同一个模型上发展“Hodge Motif”观念,以及相应的“Galois-Hodge Motif群”观念;后者后来由Tate独立引入(我说不好在他之前还是之后),并得到了“Hodge-Tate群”(联系于一个Hodge结构)之名。这桩粗劣的欺骗(但出自这样一位显赫人物,它似乎并不叫任何人感到不适)在于:以一个可笑的借口——即通往该观念的技术进路(经由绝对Hodge闭链,而非代数闭链)与我所(极临时地)采取的进路不同(倘若Hodge猜想是错的)——把一个新颖而深刻观念的父系渊源干脆整个地藏掖起来,连同我围绕Motif发展起来的一整片丰富的直觉织体。我发展了近十年之久的这部瑜伽\dmark{}\origpage{212},自1968年肇始,就是Deligne著作中灵感的主要来源。它的多产与作为发现工具的力量,早在我1970年离去之前就昭然若揭;而它的同一性,不依赖于为确立它某一部分有限内容之有效而采取的任何技术进路。Deligne有幸提炼出这样两条互不依赖任何猜想的进路。但他却没有那份诚实点明其灵感之源;相反,从1968年起他就极力对所有人遮掩它,以便把它的收益独占为己有,只等1982年(以默许的方式)认领这份信誉。

\subsection{下葬——或新父亲}

\notehead{52}{}

回到Motif之梦,我记得自己还曾把它大声梦见。诚然,梦的工作按其本性是孤独的劳作——但这项顽强地持续了多年的工作(在一项吞噬了我大部分时间的浩繁奠基写作的边缘上进行)的种种波折,有一位日复一日的见证者,比Serre亲近得多——Serre只满足于远远地看这些事……\fnote{7}{(*)}{(5月25日)不过,我对Motif的反思,其实早于Deligne的出现。我关于MotifGalois理论的手写笔记,标注的日期是1964年。}。关于这位日复一日的心腹之交,我在回顾中写道:他在六十年代中期“多少扮演过学生的角色”,而我“给他讲过我在代数几何里知道的那一点点”。我本可以补上一句:我连按这个词的通常意义我所不曾“知道”的东西也讲给他听——那些数学的“梦”(以Motif为主题,也以别的为主题),它们总能在他那里找到一只专注的耳朵和一颗警醒的心灵,和我一样渴求理解。

诚然,当我写道Pierre~Deligne曾“多少扮演过学生的角色”时,这仍旧全然是一种主观的印象(注57),据我所知,没有任何书面的、至少是印刷的痕迹可资佐证——那种痕迹本可以让随便哪个人疑心,Deligne从我嘴里学到过什么——而在这里,我却乐于记起:我同他谈数学,从来没有一次不从中学到点什么。(甚至在我同他不再谈数学之后,我仍继续经由他学习那些更困难、或许也更重要的事情,包括在我写下这些文字的今天……)。

\pmark 新近,一位第三者(他不知怎么猜到了这事或许会让我感兴趣!)告知我,Deligne和别人有一部文本,其中谈到Motif,或者至少谈到“塔纳卡范畴”;我向Deligne提了一句,他由衷地惊讶,我居然会对这类东西感兴趣。\origpage{213}翻阅他惠寄的那一册,我确信他的惊讶完全有理。显然,我这个人,同书中所论的主题全然无涉。至多在导言里一句话带过:某些“标准猜想”(当年是我提出的,天知道为什么)对一域上Motif范畴的结构会有一个推论……好奇的读者想多知道些就为难了:他在这整部书里,找不到关于这些猜想的任何说明或出处——它们此后就再没被提起;也找不到对唯一一篇我解释如何以标准猜想在域上构造Motif范畴的已刊文本的提及;也找不到1970年以前唯一另一篇谈到Motif的已刊文本——出自Demazure之手(如果我没记错,是在一场Bourbaki研讨班上),它遵循我那条特设的构造原则,只是视角稍有不同……\fnote{8}{(*)}{核实之后,我发现,除了关于标准猜想的几页(Algebraic Geometry, Bombay, 1968, Oxford Univ.~Press (1969), 第193--199页),我没有任何已刊的数学文本谈到Motif。在Demazure的报告(Bourbaki研讨班n\textsuperscript{o}~365,1969/70,继Manin的俄文报告之后)里,提到我1967年在IHES所作的几次报告,它们(我猜想)应当构成Motif 视野的第一幅总草图。关于标准猜想及其与Weil猜想的关系,有一篇比孟买会议通告更详细的报告,出自Kleiman之手(Algebraic Cycles and the Weil conjectures, 载 Dix exposés sur la cohomologie des schémas, Masson-North Holland, 1968, 第359--386页)。在1970年以前,除我之外,我没有见到任何关于标准猜想的反思——尤其朝着对它们的证明。存心无视这些关键猜想(我在孟买的草图里说过,我认为它们连同优秀概形的奇点消解,是代数几何里最重要的未决问题)的做法,在我看来,很大程度上造成了代数簇上同调理论给我的那种停滞印象——就传到我耳中的种种回响而言。}

\pmark 连Neantro~Saavedra——他有幸位列我“1970年以前的学生”——都得到了应有的引用\origpage{214}。他跟我做过一篇学位论文,论的是我记得我称之为“刚性张量范畴”、而他称之为“塔纳卡范畴”的东西。人们至今仍在纳闷:Saavedra是靠了怎样奇迹般的运气,恰好预见了十年后才绽开的DeligneMotif理论的需要!实际上,在他的论文里,他做的恰恰就是\textbf{那}项在技术上构成MotifGalois理论之钥匙的工作,一如 J.L.~Verdier的论文原则上就是\textbf{那}项在技术上构成上同调六运算形式主义之钥匙的工作。有一个(此外还有别的)对Saavedra有利的不同:他费心发表了自己的工作;他确实不曾有Hartshorne、Deligne与Illusie联袂的笔杆,来替他免除这道手续。然而十年之后,Saavedra的论文被从头到尾、几乎一字不易地复制进那部卓越的文集里,这一次执笔的是Deligne与Milne。倘若只是为了更正Saavedra工作中的两个具体之点(注58),这或许并非必要。但万事必有其所以然,我想我辨出了Deligne亲自不辞这份辛劳的缘由\fnote{9}{(*)}{对此可参见“白板”一注中的反思,n\textsuperscript{o}~67。}\fnote{10}{(**)}{(6月8日)而当事情涉及带有我的影响之痕迹的工作时,就更是如此——对此可参见“注记——或新伦理”那一情节,第33节。}:那分明同他自己那条在出版事宜上把严格推到极点的标准背道而驰;而众所周知,一涉及别人,他把这条标准执行得一丝不苟……

至于诸观念与Motif瑜伽本身的父系渊源,对不知情的读者(而知情的读者已开始变得稀少,终将一个个寿终正寝……)来说,这一父系渊源无可置疑——无须在这里去惊动遥远的Hilbert与Riemann,更不必去惊动上帝。倘若那位显赫的作者——他关于阿贝尔簇上绝对Hodge闭链的漂亮结果,俨然是Motif理论的起点,干脆说,是Motif理论的诞生——对其父系渊源只字不提,那是一种令他生辉的谦虚\dmark{}\origpage{215},与行业的惯例与伦理完美合拍:按后者的要求,应当把荣耀归于它分明所当得之处的差事(如有必要的话),留给别人:留给合法的父亲……

\notehead{53}{}

被这个孤儿的坎坷际遇所触动,并疑心不会再有别人去做这项工作——这项工作,至今显然只有我一个人感到其需要与规模——我料想,那位“勇敢的数学家”不是别人,正是我自己,待我走完《追寻栈》之路之后(我预计它还要占去我大约一年)。

\notehead{54}{}

从那以后,代数簇又出现了两种新的上同调理论(除Hodge--Deligne理论之外——后者是Hodge上同调在“Motif”精神里的自然延伸),即Deligne的“分层pro-模”理论,以及尤其是晶体理论——Sato--Mebkhout式的“D-模”版本,配上前文说过的“好上帝定理”(又名Mebkhout定理)所带来的新的照亮。这条通往可构造离散系数的进路,多半注定要取代Deligne先前的版本,因为它无疑更适合表达与De~Rham上同调的关系。况且,这些新理论并不在给定概形上光滑Motif的范畴上提供新的纤维函子,而是(尚需一项比迄今所做更深彻的奠基工作)提供一种方式,以精确把握复数域上有限型概形上一个Motif(不必光滑)的“Hodge”化身,或特征零域上有限型概形上的“De~Rham”化身。此外很有可能:$\mathbb{C}$上有限型概形的Hodge--Deligne系数理论(显然始终未见成文),终将显现为被包含在Sato--Mebkhout式晶体系数理论(同样未见成文;并另附一项关键的滤过数据)之中,或者更确切地说,显现为后者与$\mathbb{Q}$-向量空间的可构造离散系数理论的某种交……至于澄清Mebkhout式晶体理论与Berthelot及其他人在正特征下发展的理论之间的关系,那是Mebkhout早在1978年之前就已感到的任务,处在一种普遍冷漠的气候里;而在我看来,它是近期摆在我们面前、关乎对代数簇“那一个”上同调(唯一而不可分,即Motif的上同调!)之理解的最迷人的任务之一。

\notehead{55}{}

\dmark{}我尽管做梦,但我关于Motif与Hodge结构之关系的梦,却让我在不经意间把手指按在了Hodge“广义”猜想——按Hodge最初表述的样子——的一处不自洽上,并把它代之以一个更正的版本;这一回(我敢打赌),它应当既不会比代数闭链上那个“寻常的”Hodge猜想更错,也不会比它少错。\origpage{216}

\subsection{一场大屠杀的前奏}

\notehead{56}{}

我尤其想到——恰恰在代数簇上同调的语境里——Griffiths对一种诱人的想法之错误性的发现:长期以来人们以为,一个同调等价于零的闭链,会有一个代数等价于零的倍数。对一个全新现象的这一发现,当时给我的触动之深,竟使我花了整整一星期的功夫去力求吃透Griffiths的例子,把他的构造(那是超越的,在$\mathbb{C}$上)搬移成一种“尽可能一般”的构造,并且在任意特征的域上成立。这一推广并非一目了然,靠的是(如果我没记错)Leray谱序列与Lefschetz定理。

(6月16日)这一反思,曾是我有机会在平展语境里发展“Lefschetz束”的上同调理论的机缘。我这方面的笔记,已由 N.~Katz在SGA~7~II研讨班(由 P.~Deligne与 N.~Katz主持)的报告XVII、XVIII、XX里详细阐发(他费心指涉了这些笔记,而且紧随其步伐)。而在 P.~Deligne所写的该卷导言里——其中说,该卷的关键结果是报告XV(平展上同调中的Picard--Lefschetz公式)与XVIII(Lefschetz束理论)——作者却小心地不提:在这套Lefschetz束的“关键理论”里,我也有我的一份。读这篇导言,得到的印象是:对该卷所发挥的主题,我毫无贡献。

SGA~7那场漫长的研讨班,1967--69年间接续了1960至1967年间在我推动下发展的SGA~1至SGA~6诸研讨班,曾由Deligne与我共同主持,是我以消没闭链群的一个系统理论发令起跑的。由于各报告由不同的志愿者执笔、拖沓经年,研讨班的两个分卷(SGA~7~I与SGA~7~II)迟至1973年才出版,经手的是\dmark{}\origpage{217}Deligne。尽管在研讨班当时,大家已商定它将以一场共同主持的研讨班的面目出现;我离去之后,Deligne却向我表示了他那个(在我看来很奇怪的)愿望:把研讨班\textbf{一分为二},一部分I以我主持的名义出现,另一部分以他和Katz的名义。我如今在其中辨出一场“战役”,它预演了“SGA~4½之役”——那场战役的目标(除别的以外)是要让SGA~1至SGA~7整个奠基丛书显出这样一副面孔:仿佛它是一套随手采撷的文本汇编,我这个人只在其中扮演一个偶发的、乃至多余的角色——尽管就其精神与构思而言,这套丛书一如EGA《代数几何原理》丛书,是与我个人密不可分的。这一倾向在SGA~4½分卷里表现得非常清楚,甚至粗暴;尤其表现在对SGA~5研讨班的屠杀里——该分卷与那场屠杀密不可分。对此可参见(除别的以外)“白板”与“屠杀”两注,n\textsuperscript{o}~67与87,尤其是“遗骸……”(n\textsuperscript{o}~88)。

(6月17日)SGA~7研讨班的总体构思(我在其中绝不曾区分什么“I”“II”两部分,至今也仍不区分)出自我这里;另一方面,Deligne作出了重要贡献(在我1969年写的Deligne工作评述报告里有记载,见该报告n\textsuperscript{o}~13、14条),其中就研讨班的需要而言最关键的,是Picard--Lefschetz公式——经由从一个已知的超越情形做特化论证而证明。把研讨班切成两部分,无论在数学上还是就双方的贡献而言,都毫无道理——SGA~7的两个“块块”里,Deligne和我都有实质性的贡献。

当然,倘若Deligne继续我开创的SGA奠基丛书——它远未跑完自己的赛程!——我会喜出望外。这场“SGA~7之役”根本不是什么延续;我感到它是一记粗暴的“锯子”(或电锯……),\textbf{终结}了SGA丛书——用的正是一个摆明要与我个人划清界限的分卷,尽管它维系于我的著作,并且像其余各卷一样带着我的印记。当我的个人痕迹在其中被尽量藏掖的同时,对我著作的口吻,尚未达到“SGA~4½之役”那种几乎不加掩饰的轻蔑——后者是对SGA~4与SGA~5研讨班之统一体更深重的一记锯痕,也是对其中未刊的SGA~5部分大加劫掠的手段与借口;劫下的碎块,在Deligne与Verdier之间均分……

\notehead{57}{}

\dmark{}请允许我赶紧补充:同样的评语,也适用于另一位大手笔的数学家——十年之后,我斗胆说过(在注n\textsuperscript{o}~19里)他也“多少扮演过学生的角色”。 \origpage{218}

\notehead{58}{}

这让我想起:Springer的《Lecture Notes》(它曾发表过六七篇在我指导下完成的“1970年以前”的博士论文)从未肯发表Yves~Ladegaillerie“1970年以后”的那篇论文(理由:他们不发表博士论文!)。人们倒可以说,它把Saavedra的论文发表了第二遍……我还曾向Deligne提起Ladegaillerie那条四处碰壁的漂亮同位结果(心里暗暗指望他会出手相助发表)——却没能引起他的兴趣(理由:他对曲面拓扑不在行……)。

幕落……

\subsection{新伦理(2)——或抢夺的集市}

\notehead{59}{}

(4月20日)写下这些文字——它们记录了一桩矛盾及其代价——已有几个星期;我惊讶地发现,当事人早在两年前就已找到一条最简单的路子来“化解”那桩矛盾——难就难在想得到!人们可以称之为“预先下葬法”(读者可在昨天写成的双注(注50、51)里了解它,那是发现之后情绪未凉之时写的)。我感到遗憾的是:这位被预先下葬的死者,冷不防又重新出现在那有名的“数学舞台”上(它有时候,分明更像一个抢夺的集市……),恐怕会给这项光辉方法的干净利落的应用,带来技术上的麻烦!在前面的注里(“共识性的操守——与信息控制”,n\textsuperscript{o}~6),我(还有些含混地)感到:科学行业里那条最获普遍认同的操守规则,在掌管科学信息控制的人们不尊重“每个科学工作者有权让自己的观念与结果为人所知”这一权利时,“就形同虚设”。正是在反思的这个节点上,我还费心相当详细地描述了一个个案:对这一权利的藐视,在我眼里昭然若揭;而且我分明感到,这藐视已逼近对那条首要规则——获得普遍共识的那一条——的藐视。(见“注记——或新伦理”,第30节。)

\pmark \origpage{219}这已不是我头一回感到这种特殊的痛楚:当我看到这条首要规则的\textbf{精神}遭人藐视,而藐视者却处于“暂停”的豁免状态——无论靠的是其地位(高到不容置疑!)、其才具,还是形式上的漫不经心。我试着在相关于所引第30节的那则注里(“年轻人的势利——或纯洁的捍卫者”)勾勒这种痛楚。当人们可以放胆藐视我在那里谈到的那些“显而易见”的事,并且(如今我可以补上一句)以同样的精神,也藐视那些既未证明、也未以已刊并为众人所知的“猜想”的形式注册备案的事情(或许是深刻的事情),那么人们同样可以(鉴于其“分量之轻”\fnote{11}{(*)}{“好上帝定理”(又名Mebkhout定理)的遭遇,正是如此。(6月8日)而且,一如对Motif瑜伽那样,还费心巧妙地制造出对此物拥有父权的表象,却从不明言!就个案而言,可参见“变戏法的人”一注,n\textsuperscript{o}~75$''$;就这条光辉方法的一般路数或风格而言,可参见“暂停!”一注,n\textsuperscript{o}~77,以及紧随其后的“据为己有与轻蔑”一注,n\textsuperscript{o}~59$'$。})把它们视作公共财产(不言而喻,是平凡的),因而,到时候,以最大的漫不经心与最坦然的良心,把它们当作“自己的”——当然,前提是没有人想到要去据为己有一篇十页或一百页(或哪怕十行)的硬核证明——它确立一个“自己本来不会证明”的结果(注59$'$)。我不曾想到,自己对“形同虚设”这一点,竟感触得这样准、说得这样中,因为我有幸目睹:上面那个案子里那条模糊的“界限”,被人轻松越过——而且可以肯定,越过时同样怀着最坦然的良心,\textbf{鉴于其分量之轻:那是一个梦},而且甚至未经证明(尤其,未经\textbf{发表}……)\fnote{12}{(**)}{人们何必有所顾忌,既然事情本身分明表明:如今普遍的共识认为这完全正常——至少出自一位格调如此之高的人物!所谓“坦然的良心”,不多不少,正是这样一种感觉:觉得自己与自己所处行业里通行的共识相合。}。

幸好我还有点招架之力——必要的时候,我多少能表达出我的感受与想说的话;我(不论对错)挣得了一份信誉,因而,当我有话要说的时候,有一线希望被倾听,或者,若我觉得有必要,把它发表出来。相反,我此刻更痛切地体会到那种“不义与无力之感”——当一个人\dmark{}\origpage{220}受了损害却无处申诉,在“一切尽在掌握”的人们的专断面前,自觉手脚都被捆住——而后者还随心所欲地支配一切。

诚然,在我做数学家的生涯里,我也曾以同样坦然的良心有过可鄙的行为;在反思中我已有机会谈到那些个案——它们是这场反思从遗忘与从未检视过的暧昧的迷雾里重新唤起的。探查它们,我终于懂得:今天(而且由来已久)学生轻松愉快地超过老师,我无须惊讶;凡是与我有感情或情谊相连的人,我也无须否认。但对我、对所有的人来说,把猫就叫作猫,是健康的——无论这只猫出自我家,还是出自别人家。

\subsection{据为己有与轻蔑}

\notehead{59$'$}{}

(6月8日)对这件事,就我的朋友Pierre~Deligne而言,我已完全不这么确信了,因为我有机会确认:他终于滑进了对$\ell$-进上同调工具的“默认父权”游戏——也就是我所说的对平展上同调的“掌控”。从“SGA~4½之役”(其中我的名字还被提起,只是带着一种做作的、漫不经心的轻蔑,轻蔑于我著作中的这一核心部分——他自己的那一部分正源出于此),到“葬礼悼词”(其中凡“上同调”一词本身,凡与我的名字相关处,一概禁用),有一段令人瞩目的演变。(初期阶段见“白板”与“格格不入者”两注,终局阶段见“葬礼悼词(1)、(2)”两注。)

作为这场步步升级中的几个中间阶段:1981年,有那篇“值得纪念的文章”,论的是所谓“反常”层(对此可参见“不义——或归来的意义”与“暂停!”两注,n\textsuperscript{o}~75与77);次年,有LN~900里Motif的开棺出土(而“葬礼悼词”则再下一年,即1983年登场)。在我得以观察到的所有这些情形以及一些规模较小的情形里,那种使Deligne能够以完全坦然的良心据有他人观念之信誉的内在态度与“方法”,就是\textbf{轻蔑}(它部分地保持缄默,却又被巧妙地暗示出来)——对着人们即将据为己有的那个“轻”东西:它“轻”到甚至不值一提,尽管人们转过头去就要用它做出真正有力的东西——Weil猜想,所谓“反常”层的理论……一旦战役完成,据为己有成为既成事实、并为众人所接受,总有的是时间来修正口径,带着已据为己有的东西,谦虚地招摇过市。同一项贡献,只要似乎还沾着一个有待下葬者的名字,就遭到漫不经心的轻蔑;\dmark{}\origpage{221}一旦被他本人据为己有,就被捧上云端($\ell$-进上同调,Motif,Mebkhout的瑜伽尚待其时),或者被某位好哥们据为己有(导出范畴的瑜伽,对偶的瑜伽——由Verdier据有,Deligne积极怂恿)。

\section{V.~我的朋友皮埃尔}

\subsection{孩子}

\notehead{60}{}

\dmark{}(4月21日)再拾起这个关于回忆的梦——它不只是对一个视野之诞生的回忆……\origpage{223}我清楚地记得(尽管我忘掉了那么多事情!)那种每一次都焕然一新的乐趣:同那个很快就不止是心腹之交的人交谈——对于一切令我困惑的东西,对于日复一日在我与数学的恋情中豁亮起来、令我心醉的东西,他比任何“学生”都更是知交。他那份始终警醒的兴趣,他了解任何事情的从容(“仿佛他一直就知道……”),对我而言,是一泓永不枯竭的欢喜之泉。他的倾听是完美的,由那种也在我身上跳动着的、求理解的渴望所推动——一种高度醒着的倾听,是心心相印的标志。他的评点,总是迎着我自己的直觉或保留而来;即便并非如此,它们也往往在我力图穿透重重迷雾去把握的那个实在上,投下一道意想不到的光。如我在别处说过的,我提出的问题,他往往应声就有答案,或者在其后的几天、几星期里把答案发展出来。这就是说,倾听是双向的:当他反过来向我解释他找到的答案时——也就是解释事情的所以然,而那所以然,总是带着那种完美的自然、那种从容显现;同样的从容,曾常令我心醉地见于我的某些长辈,如Schwartz与Serre(同样也见于Cartier)。正是这同一种单纯,这同一种“自明性”,是我在理解数学事物的过程中始终追求的。无须明言:凭着这条进路与这份要求,他和我显然“同属一家”。

从我们相识之初,我就分明感到:他那(人们所谓的)“天资”,有一种极其少见的品质,远在我所拥有的那份寒微的手段之上;而就求理解的热情、就对待数学事物之理解的那份严格而言,我们却在同一个音高上。我还隐约感到——当时我尚无法把它表述出来——我在他身上觉察到的这股“力量”(我在自己身上也感到,只是较弱):那种“看见”人人视而不见的自明之事的力量,乃是童年的力量,是孩童眼里的天真。他身上有某种属于孩子的东西,比我认识的任何别的数学家都更外露;这绝非偶然。他给我讲过:有一天,我记得还在中学,他曾兴致勃勃地用定义去验算乘法表(顺水推舟,连加法表也一并验了)\dmark{}\origpage{224},对1到9这些数。他当然不指望有什么意外——如果说真有意外(愉快的,一如既往……),那就是:这证明居然可以在短短几页里、漂漂亮亮又完完整整地做成,前后不过半个钟头的事。当他笑着给我讲这件事时,我分明感到,那半个钟头用得其所——这件事,我今天比当时懂得更透。这段小故事当时打动了我,甚至震动了我(我想我未曾形于颜色)——我从中感到一种\textbf{内在的自主}、一种相对于既有知识的自由的征兆;这种自由,在我童年与数学的关系中,自最初接触起就已在场(注69)\fnote{13}{(*)}{而且我觉得,这份自由在我做数学家的一生中从未全然隐没,如今又像童年时那样在场了。两三年前,我向我的朋友重提乘法表那桩童年小事。我感到他为这一回忆的唤起而局促——它分明已不再符合他心中的自我形象。这份局促我并不真正惊讶,却仍然难过:又一件我明知道、却仍难于承认的事,再次得到了确认……}。这种彼此互为特权对话者的关系——我想,那会儿我们几乎天天见面\fnote{14}{(**)}{至少在我住在比雷时是如此,他当时住在IHES的一间单间。自1967年(我搬到马西)之后,我想我们仍一周会面一两次,至少直到我不再投身数学为止。}——持续了五年,从1965年(如果我没记错)到1969年(含)。我还记得那一年写一份详尽的评述报告时的乐趣:当时我提议把他增选为我自建校起(1958年)就一直工作的那个机构——我的大部分数学著作就是在那里完成的——的教授。那份报告我没有存本了(注64):我在其中评阅了我想有十几篇我朋友的著作,当时几乎都未发表(有许多后来也没发表),而在我看,其中大部分——如果不是全部——都当得起一篇优秀国家博士论文的主体分量。呈递这份雄辩的报告,我感到比呈递一份关于我自己工作的报告更骄傲、更欢喜(那种事,我一生只做过两次,而且每次都是勉强自己……)。其中许多著作,是对我提出的问题的回答(其中唯一已刊的,是前面提到的那篇关于概形真光滑态射之Leray谱序列退化的工作(注63))。而最重要的两篇\dmark{}\origpage{225}却是对Deligne自己提出的问题的回答;显然,其分量与一篇“优秀的国家博士论文”不可同日而语。那就是他关于Ramanujan猜想的工作(原书拼作Ramanuyam)(发表于Bourbaki研讨班),以及关于混合Hodge结构的工作——亦称“Hodge--Deligne理论”。

有一件事很奇怪,而且是我写这份流光溢彩的报告时远远没有料到的:不到一年之后,我就要离开那个机构——我正准备在那里增选我年轻而令人惊叹的朋友,并且打算在那里终老。还有(如今我把这两桩双生的情节放在一起想)另一件怪事,也绝不像是单纯的“偶然”:一两个月前,正是这位(如今已不年轻!)朋友,向我宣布了他自己要从同一机构离去——而恰好一年之前,我才重新拾起一种规律的数学活动,以一种不期然的“重返”走上数学舞台(虽还不是“上流社会”……)

不止一次,我在《收获与播种》里有机会谈到我的离去——谈那次“有益的挣脱”——以及紧随其后的“觉醒”;正是后者使那一情节成为我生命中的一个关键转折。在其后那些紧张的年月里,数学家的世界,连同我在其中爱过的人们,以及数学本身最令我心醉神迷的那个东西,都变得极其遥远——仿佛淹没在另一个“自我”的回忆的迷雾里,那个自我仿佛死去千年……

但无论在这桩情节之前,还是在这第一个大转折之后的年月里,我都知道:那个曾是(一点点\fnote{15}{(*)}{至于我为何有这份顾忌,不肯把(太!)耀眼的Deligne当作我的一个学生,见“格格不入者”一注(n\textsuperscript{o}~67)。})我的学生、(在很大程度上的)心腹之交与朋友的人,只须顺着体内一个游戏着的、渴求认识的孩子那股自发的冲劲,就能发现并唤起崭新而出人意表的世界,就能探测它们、认识它们内在的本性——并由此把它们向他自己的同类、也向他自己揭示。所以,在我离去之后(绝无回归之念!),倘若我看见“一位勇敢的数学家”、一位富有灵感的人,开始大笔勾勒(权作开端……)我隐约瞥见的那幅宏阔图景——而我尚且只画下了一组局部的、临时的草图——那个人只能是他:他把做成这件事的一切都握在手里!勾勒这幅大面积的第一图景——一位“总设计师”把关于\dmark{}\origpage{226}代数簇上同调的已知与猜出的本质事物,纳入一个共同的视野——对于那个心中视野已然呼之欲出、只待破除“尚未写下”之迷雾的人,是几个月的工作,甚至用不上几年。(不妨在此后经年、甚至经代——倘若需要几代的话——将它重做、深化,直到Motif之实在的底细被充分理解与确立。)而且我毫不怀疑,这项不久前还“在我手里灼烧”的工作,随时都会有人去做,至少在其后两三年内,趁着它还热乎。我离去之后,的确只剩一个人,凭着他对认识的冲劲本身,注定要做这件灼热而迷人的工作。不妨,在“总设计师”写就并经检验、大厦的营造或深或浅地推进之后,把继续这项工程的差事——无论它多么迷人——留给别人,自己投身别的冒险,在那个数学事物的世界里:那里的每一个转弯,都显露着一个崭新而无边世界的许诺,只要我们有睁开的新眼睛去看……

当我的生活还在那座温暖的科学暖房里流转——它把尘世的喧嚣隔绝在外——而当Deligne发展他对Hodge理论的推广时(那应当是1968或1969年),在我们之间,有件事不言而喻:那项工作,是把“Motif图景”的某个\textbf{部分}加以实现、检验与明确的一个头等步骤;而那幅图景,作为一个整体,从未被白纸黑字写下\fnote{16}{(*)}{至于那套Hodge--Deligne理论后来(据我所知)从未超出这第一次尝试的阶段,从未扩展成一套“Hodge--Deligne系数”(及其上的“六种运算”)理论、架在复数域上有限型概形之上——这与另一件怪事密不可分:那幅宏阔的“Motif图景”从未被勾勒出来,而它的存在本身,直到今天还被小心地讳莫如深……}。在我告别暖房之后的年月里,当数学对我已相当遥远时,我得知Weil猜想最终获得证明,毫不意外。(若说有意外,那是“标准猜想”没有在同一步里得到证明——尽管人们提出它们,恰恰就是为了逼近Weil猜想,同时也是作为在一域上至少建立半单Motif理论的一个手段\fnote{17}{(**)}{只是在这最后几年里,我才模糊地(但近来更精确地!)意识到:“标准猜想”,连同它们为之提供了第一条“构造性”进路的Motif观念本身,都被\textbf{埋葬}了——缘由如今在我眼前显得格外清楚。(亦可与前面一条脚注对读。)}。)我很清楚:无论是这第一击\dmark{}\origpage{227}——那初试身手的、Hodge式的一般系数理论——还是对若干关键猜想(在若干其他多少为人所知的猜想之中)的这项证明,都还没有让他使出全副本领——差得还远。而我并不着急,因为我的注意力的绝大部分,已倾注在别处。(→注61)

\subsection{下葬}

\notehead{61}{}

我有幸亲眼见到一股童年的冲劲初次绽放,带着大面积舒展的许诺。在其后的十五年里,我渐渐确认:这一许诺始终被不断推迟。他身上有那样一种精微的东西,我曾懂得去感觉、去辨认(尽管那会儿我对那么多事情都无感!)——一种与脑力全然异质的东西(脑力既碾压人也穿透人……),一件对一切真正的创造性的工作而言头等本质的东西。这种东西,我在别人身上有时也感到过;但在我认识的任何数学家身上,它都未曾以可与之相比的力量显现。而我(想当然地)指望,这东西会继续在他身上绽放、转化,并毫不费力地经由一部独一无二的作品表达出来——在那部作品里,我不过是一个寒微的先行者。然而又一桩怪事(而在这么多“怪事”之间,必定有一条简单而深刻的纽带):我眼看这“精微的东西”、这既非膂力亦非脑力的“力量”,随年月流逝而渐渐消退,仿佛被\textbf{埋葬}在一层层越来越厚的沉积之下——那是一层层\textbf{别的东西}的沉积,我对它们知之甚稔——世上最寻常的东西!它未必与脑力、与老到的经验或在某一学科里磨炼出的嗅觉不能相安;这些东西,凭其作品的积累——也许光灿夺目,肯定自有其力量与美——可以博得一些人的叹服、另一些人的畏惧,或兼而有之。但这并不是我当初说“舒展”或“绽放”时所想的。\textbf{那个}绽放,我所想的那一个,是天真的果实——渴望认识,并且随时准备为这个取之不竭的世界里或大或小的事物的美而欢喜,或为它的某个部分(比如数学事物的那个宏阔世界……)的美而欢喜。唯有它,才有深刻革新的力量:无论是自我的革新,还是对世上事物之认识的革新。在我看来,把这一点实现得圆满的,正是谦逊的Riemann\fnote{18}{(*)}{Riemann(1826--1866)的著作,装满区区一册薄薄的文集,约十篇论文(诚然,他四十来岁就去世了),其中大部分包含一些简单而本质的观念,深刻革新了他那个时代的数学。}。这份真正的绽放,与轻蔑无缘\dmark{}\origpage{228}:对他人(对那些自感远在自己之下的人……)的轻蔑,对那些太“小”、太自明以致不屑一顾之事的轻蔑,对那些被认定低于自己正当期望之事的轻蔑;或者还有对某个\textbf{梦}的轻蔑——那梦执拗地向我们诉说着我们口口声声说爱的事物……它与轻蔑无缘,正如它与滋养轻蔑的自负无缘。

诚然,凭着那令人叹服的“天资”,但更凭着那件不打动任何人、却\textbf{创造}着的精微之物,“学生”注定要远远超过“老师”。我毫不怀疑,在我离开那个地方——我在那里曾见证一次如此美丽的起飞——之后的年月里,Deligne将使一部宏阔而深刻的作品得以全副展开,而我曾是它的先行者之一。这样一部作品的回响,岁月流转,不会不传到我这里;而我自己,在远离数学的其他追寻中跋涉,对他将要发现的那些崭新世界的全部分量与全部美,只能有不太完整的领会。

但学生若在内心深处\textbf{否认}老师,努力在人前人后抹去老师的一切痕迹(无论这份馈赠是为了更好还是为了更坏……),学生就无法真正超过老师——正如儿子若否认父亲,就无法真正超过父亲。这一课,我主要是经由我与我的孩子们的关系学到的,但后来也经由我与昔日某些学生的关系;尤其是与那一个人——在所有人之中,我始终有所顾忌、不肯以“学生”之名相称的那一个;因为从相识那一刻起,我就分明感到,我有东西要向他学习——一如他也有东西要向我学\fnote{19}{(*)}{(6月14日)关于我身上这份执拗的蓄意——贬低我所能给的东西、否认一段师生关系的实在——见“格格不入者”一注,n\textsuperscript{o}~67$'$。明摆着:我的朋友从我的接触中所学到的东西(“仿佛他一直就知道”,当然!),与我经由他所学到的,没有共同的尺度。倘若我继续高强度的数学投资直到今天、我们之间保持着规律的数学接触,情形想必会是两样。}

但只是在相识将近十年之后——在1975年之后,尤其自从我开始沉思我所经历的、我所目睹之事的意味——我才开始感到那个继续为我所珍的人身上的这份\textbf{枷锁}。我也隐约感到:这份对我这个人、对我在其人生关键年月里所扮演角色的秘密否认,更深地,也是对他\textbf{自己}的否认。(想必每当我们否认、并想抹去\dmark{}\origpage{229}某件分明发生过、而其果实有待我们去摘取的事情之时,都是如此。)

然而,由于对“数学界时下在做什么”、对他自己在其中所做的事情\fnote{20}{(*)}{自1970年以来,我收到过Deligne的四份抽印本,都当即匆匆翻过(一如我至今仍不时收到的大多数抽印本)。凭这点东西,难以对一部数学著作——哪怕只就其总轮廓或主要主题——形成概念。}多少不“在局中”,在几个星期前的反思之前,我从不曾掂量:这份枷锁,对他倾其所有投注的事情——他的数学著作——也压得有多重。诚然,八九年来,我不止一次看到,一份蓄意的轻蔑口吻(对我)或鄙夷(对一些他有能力使之气馁的别人)仿佛把单纯的常识、把数学家的健全本能一笔勾销(注66)。而且他并非我旧学生中——带或不带引号——唯一一个让我目睹这类态度对着我所珍视的人(或对着别人)的。但在别人身上,它从不曾这样痛切地触动我。在刚过去的两个月的反思里,我不止一次影射过这份经验——“它是我做数学家的一生中所曾经历的最苦的经验”;我也说过,这场《收获与播种》的反思到头来就它教给了我什么。那份创痛如此剧烈,它就一个始终为我所珍的人教给我分量如此之重的东西(而我继续回避它同时就我自己、就我的过去所教给我的……),以至于这种态度会对他、乃至对那些被弄得气馁或受了羞辱者的数学“创造力”产生多大影响,反而成了全然次要、甚至不足挂齿的问题。

“拒受遗产——或矛盾的代价”那条注,是我头一份写下的反思:盘点经年累月、此处彼处以断片形式回到我这里的东西——既关于“行内现状”,也关于那位我认识得那么好又那么少的人的著作。那也是头一回,我在一眼之间终于看清:那份他携行了大概十五年以上的拒绝,在其数学著作本身中付出的全部“\textbf{代价}”、造成的全部重负。写那条注时,我其实还“落在后面”:因为两年以来(而且没有“人”认为值得通知我),Motif早已走出它们被囚禁十二年的秘密……而今天,在我写下这场对我数学家过去的反思的这一(我相信)最终阶段之时——在我\dmark{}\origpage{230}大致了解那部标志着这次悄然“重返”的值得纪念的文集之后两天——对这份沉重负担的知觉,变得触目惊心:那正是生来该飞翔的人所乐于日复一日、千回百转拖着的重量——一次轻柔轻盈的飞翔,欢快而无畏,迎向未知,为了他的喜悦,也为了托起他的风的喜悦……\fnote{21}{(*)}{我绝无意暗示:“飞翔”、发现世界,是少数特殊者的特权;肯定,我们人人天生受此召唤!只是这份能力,很少有机会哪怕在极有限的方向上(比如数学工作)绽放。但在某个人身上,我曾有幸看到这样一份能力(在“数学”的方向上)格外夺目,如奇迹般被保存,其后却随岁月而倒退。}

倘若他不飞翔,倘若他满足于做一个受人仰慕与畏惧的人,靠不断积累自己优于他人的证据,我无须为他担忧。倘若他拖着他乐于拖的重量,他肯定从中找到满足——一如我自己曾乐于拖重量,并且今天仍在拖那些我在路上尚未学会割舍的重量。我所能给他的,无论好的一面还是坏的一面,他取其所愿。他的选择只属于他自己,我无须为其担忧;甚至无须在此裁决它们是最好的还是最坏的(注62)。对一个人是“最好”的,对另一个人是“最坏”——有时对同一个人也是(只要他肯变,这诚然少见……)。

但我们所做的选择、表达这些选择的行为(即便我们的言语常常否认它们),我们都须自行承担风险与后果。它们常带给我们所期待的报偿(我们把它当作“最好”来领受);而这些报偿本身,有时终于显出其反面(我们把它驳回为“最坏”,甚至当作一桩冒犯)。当人们终于懂得:反面并非冒犯,往往便把它看作一笔该付的价钱,付得老大不情愿。但也会有这样的领悟:某些反面,并不是无情的收银员——为了快活时光,不情愿也罢,总得付账;它们是耐心而执拗的信使,不倦地回来,永远带着同一个讯息;一个不合时宜、不断被拒的讯息——因为比反面本身更甚的,是它那份谦卑的讯息总遭驳回,那份讯息在我们眼前才显为“最坏”:比一千个反面更坏,常常比一千次死亡、比整个宇宙的毁灭更坏——对整个宇宙,我们已经毫不在乎……

\dmark{}\origpage{231}终于有一天,当我们乐于接纳那份讯息,眼睛霍然睁开,看见:曾被畏为“最坏”的,是一场\textbf{解脱},一场宏大的释负——而那份压得我们喘不过气、我们蓦然卸下的重量,恰恰是昨天我们攀住不放、当作“最好”的东西。


\subsection{事件}

\notehead{62}{}

(4月21日)人们会对我说:既然我无须为他担忧,那为什么还要为一个只关乎当事人与我俩的个人关系,铺陈上这么多页!

我之所以感到需要这场对一段关系的某些重要侧面的回顾性反思,是一个具体的、与我切身相关的事件使然(尽管我是带着两年的时差才得知它的)。另一方面,这一事件处在公共领域之中,比那些成名的数学家(如Deligne,或我自己)对待名气较小者或初学者的日常行为举止,还要显豁得多(虽然后者对他人生活的影响,往往与眼下这个个案完全不在一个量级)。所说的事件(即Lecture~Notes的LN~900那部“值得纪念的文集”的出版,别名“下葬文集”),连同围绕着它的一切,在我看来——无论对错——是\textbf{不健全的}。在我看来,给某些来龙去脉作一份周详的见证,直探我今日所感知的事情的根底,对所有人都是有益的,首先要从“当事人”本人做起。经由这一见证与这场反思,我并不想说服任何人相信任何事情(那实在太累人,何况毫无希望!\fnote{22}{(*)}{(5月25日)如果我在这里觉得有必要重申,想说服人是“太累人”而且“毫无希望”的,那恐怕是因为:在我身上某个地方,说服的意图毕竟实实在在地在场,而且也被觉察到了。从4月19日(我得知“值得纪念的文集”LN~900那一天)到4月30日的整个反思,都带着一种内在紧张、乃至分裂的状态——面对一桩全然出乎意料的“事件”的冲击,我尽力消化它带来的讯息。这一紧张,最终随着4月30日的“事物的回归”一注(n\textsuperscript{o}~73)而解开:反思终于回到我自己的身上,随即为我提供了这则讯息的明摆着的钥匙。}),而只是想去理解那些我身陷其中的事件与境况。倘若它们能促使别人在老生常谈之外作一番真正的反思,这份见证就不算白写。

\subsection{驱逐}

\notehead{63}{}

\pmark\origpage{232} (4月22日)这篇文章\fnote{23}{(*)}{即Deligne论谱序列退化与Lefschetz定理的文章(Publications~Mathématiques~35, 1968),“罐装的权与十二年的秘密”一注(n\textsuperscript{o}~49)中已提及。}1968年发表于《Publications Mathématiques》,即我离开数学家世界的前两年。它的出发点曾是一条猜想——我同Deligne谈起过:关于谱序列的一个退化性质,在当时看上去相当不可思议,却经由“算术”的途径、作为Weil猜想的一个推论,变得可信起来。这一动机本身就有极大的趣味,因为它表明:从隐含在Weil猜想里的那部“权的瑜伽”(这部瑜伽最早由Serre在某些重要方面隐约瞥见)里,人们能榨出多少东西。从那时起,我就经常把它应用于各种类似的情形,从“算术的”论证榨取“几何的”(对代数簇的上同调而言)结论。只要Weil猜想尚未确立,这些论证就始终是启发性的,但它们仍然具有巨大的证明力,代表着一种头等的\textbf{发现手段}。Deligne对这条特殊猜想的“几何”证明——借助Lefschetz定理(当时只在特征零确立)——在一个全然不同的方向上另有其趣味,且不提它的头一项功劳:不依赖任何猜想。两条进路所指示的那种联系——一边是Weil猜想(以及当时在我看来其最迷人方面的权的瑜伽),另一边是Lefschetz定理——联系着两件看似互不相干的事情;这联系本身就极富教益。

对我眼下的处境而言,这里有趣的一点——直到今天它才向我显出全部意味——是:这篇文章的读者,几乎无从猜到我参与了其主结果的最初动机,更绝无可能从这篇文章里得知那动机是什么。(另可参见注49开头。)为这样一个结果作陈述,\textbf{自然而然}的路数(我确信,连作者本人最初也是这样)本应是:从那条(的确引人注目的)猜想出发,指出找到的第一条理由——它同样引人注目,而这正是终于把那部著名的权的瑜伽“推销”出去的大好机会,它本身的分量远比该文的主结果更重\dmark{}\origpage{233}\fnote{24}{(*)}{恰恰是这部瑜伽,此后六年里(据我看来)一直保持着秘密!(6月7日)而且(如后来所显现的),它当时被Deligne“记在自己名下”地呈献出来,对Serre与我都不置一词。(见注n\textsuperscript{o}~78\textsubscript{1}、78\textsubscript{2}。)};然后接着谈“Lefschetz定理”的观点\fnote{25}{(**)}{(6月17日)用(“奶牛”)Lefschetz定理来证明谱序列的退化,这一想法出于Blanchard;不过他只在一条苛刻得罕见地得到满足的假设之下——即纤维的有理上同调构成的局部系统是平凡的——才得到退化定理。我熟悉Blanchard的工作,自然向Deligne谈起过;他因此从Blanchard的想法取得启发来完成他的证明,尽管他没有读过那篇文章。Serre对Blanchard证明的记忆比我好,他向Deligne指出:他的证明实际上是Blanchard证明的一个不费力的改写。Deligne在其评注2.10里承认了这一点。但这条引用了Serre的评注,写法却给人这样的印象:仿佛他是事后才得知Blanchard的想法的——事实绝非如此。于是,这篇文章的两条主要\textbf{源头}都被藏掖了:一边是\textbf{算术的动机},它使人得以预见Blanchard结果的一个重大强化;另一边是Blanchard的\textbf{证明想法}——他把后者优雅地加以改写,得到一个Blanchard恐怕不敢指望的结果,也因此从未试图用自己的方法去“得到”它。},它使人得以在略更一般的条件下(基底概形任意,不必真且光滑于一个域)证明最初的猜想,但只限于特征零。而文章实际采取的陈述,却从同调代数的一般的漂亮道理开始(想必十分漂亮,而且以作者惯有的优雅陈述出来)——那些道理,他想必早已同所有人一样忘光了,俨然是Lefschetz定理公理化的路数。主结果(当然,也是唯一人人记得的结果)以系X的面目出现在文章中部;而快到结尾的某处,在“评注2.9”里,读者不知为何,“权”这个词与我的名字才被提起……

文章发表时给我的印象,我已毫无记忆——我在局中,想必只是匆匆浏览了一眼。我肯定感到了一种“保持距离”的意图,但也感到:我的朋友不愿冒被人看作某位“大师”的信徒(或“驹子”)的风险,这原是再自然不过的事\fnote{26}{(***)}{(5月26日)关于我身上的这种态度,可参见紧随本注之后的“高升”一注(n\textsuperscript{o}~63$'$)。(6月8日)把它与我朋友那套据有他人观念的独门手法放在一起看——这里就是第一个典型的例子——我进而明白:我朋友的动机,根本不是要在一个显赫的“大师”面前保持什么“自主”,而恰恰是要藏掖他人的观念在其观念之生成中的作用,以待来日把这些他人的观念也一并据为己有(在第二步里)。(对此可参见“变戏法的人”与“据为己有与轻蔑”两注,n\textsuperscript{o}~75$''$与59$'$。)至于在我的朋友身上,这一癖好得以畅行无阻地发展,我自身应负的那一份责任,可参见“高升”“暧昧”与“格格不入者”三注(n\textsuperscript{o}~63$'$、63$''$、67$'$),其中显出我对这位光彩夺目的年轻人Deligne所曾表现出的某种纵容。}。诚然,倘若他心中有那份对自身力量的安然自信\dmark{}\origpage{234},他原会毫不犹豫地写出一部分量更重、对人人更有用(肯定也对他自己更有用)的著作,而不必担心别人看不出他的本色……(注65)

次年,他第一部大面积的著作——论混合Hodge理论——发表时,情形也多少类似。(我当时认为,这项工作的分量,可与Hodge理论本身相比,把它看作一套“Hodge--Deligne系数”理论的起点——那套理论可惜始终不曾问世……)如我所说,对他和对我都再明显不过:这项工作的“动机”在于我此前数年抵达的Motif瑜伽——它是使这部瑜伽获得有形实现的第一番尝试。在他的著作里点明这层联系,依我看(当时想必也这样觉得),本会立刻使这项工作获得比它凭自身优点已具有的更宽广的分量;同时,这又是一次机会,把读者的注意力引向Motif的实在——它在Hodge结构的实在之后,随处可以感觉到(注63$_{1}$)。

只有带着事后的认识,这些省略才在六年的权的瑜伽沉默\fnoteplainmark{27}{(*)}{(1985年4月19日)关于“六年”与“十二年”的更正,对权而言见第302页脚注(***)(1985年4月18日所写的部分)。}、十二年的Motif沉默(不妨说,是禁令)\fnoteplainmark{28}{(*)}{(1985年4月19日)关于“六年”与“十二年”的更正,对Motif而言见子注“预先出土”(n\textsuperscript{o}~168$_{1}$)。}之后,在这场以LN~900下葬文集为形式的罕见的“重返”,在Hodge--Deligne理论惊艳的开局之后的停滞之中,显出全部意味……但没有人能怀着殡仪人的心肠做出大事!

无论如何,倘若1970年我离开IHES时心智更成熟一些,那么从那一刻起,我就会一清二楚:在那个五年间一直扮演着我最亲近之友的角色的人身上,存在着一种对我而言深刻的暧昧。况且,在一个平静安适的机构里,在同侪间彬彬有礼的情面外观之下,我的离去,最终对人人都是方便的——缘由我自信事后已能辨出,而且各人并不相同。\dmark{}\origpage{235}显然,这一离去对我的年轻朋友再合适不过:他新近就了位,本只需同我站在一边(面对另外三位终身同事迟疑的冷淡),就足以翻转一个未定的局面。当时我不明白这一切意味着什么,是因为我分明不愿去理解那些其实足够清楚、甚至足以自明的事情!一如我生命中常常发生的那样,那时我心中有一种焦虑(我从不曾以这个名字称呼它!),它向我提示:在一个有形而简单的实在,与一幅我不肯放手的实在图景之间,出现了“脱节”:即我在即将告别的机构里所扮演角色的图景,而更甚的,或许是我与朋友之间那层关系的图景。正是这种拒绝面对一项无可置疑的实在,以及作为这桩矛盾之征兆、我却死死攀住的焦虑,使得那次“有益的挣脱”的情节,在当下显得如此痛苦\fnote{29}{(*)}{关于这一情节,见注n\textsuperscript{o}~42。}。

老实说,由于从未曾为这段关系专门写过一篇反思(除了一些散见于给我朋友的书信里的反思的开端,而且没有一封得到回音……),我此前不曾注意到:我的朋友对我之关系中矛盾态度的最初征兆(诚然隐约,却绝不会错)至少可以上溯到1968年,即“大转折”前两年。那是一个这段关系显得完美无缺的时刻:在数学层面上是无云的契合,处身于一段单纯而亲切的友情之中。人们此刻尽可以讥讽那些关于天真、创造的孩子的美好“长篇大论”!

然而我深知,那种契合是一个\textbf{实在},绝非幻象;一如那“精微的东西”是实在——那股创造的力,其后的作品只给了它一个苍白的反照。“天真”与“冲突”是两个可以触摸的实在,略一警醒的知觉就能辨认,绝不是什么概念;而且在我看来,它们按其本性相互排斥,一个排斥另一个。然而毫无疑问,这两个实在,在不同的层面上\fnote{30}{(**)}{在另外两三次机缘里,我得以确认:在同一个人身上、在某个时刻,竟有这等共存——包括我自己的身上,在有些时刻。},共存于我朋友对我的关系之中。在我所说的那个年代,“冲突”似乎并未干扰数学的创造性——至少在\dmark{}\origpage{236}孤独中做的工作里、或在面对面的交谈中进行的工作里没有。诚然,在我方才谈到的那两篇文章里——它们毕竟是这项工作最切实的果实之列——“冲突”的印记已经分明可见。而带着十五年的事后的认识,经由这些天与这些星期的反思,我看清了:这道印记(无论多么隐约),以一种触目惊心的方式,预演了冲突将逐步压倒那股最初的冲劲、并在岁月里剥除其最稀有的本质——那造就大命运的东西——时所采取的特殊形式\fnote{31}{(*)}{这一节如此高贵的抒情,使我有点脱离了尘世的实在。倘若我在这里称这道“印记”为“隐约的”,那是因为我自己还裹在一层厚茧里,还难舍那副我视若珍宝的眼罩!一旦终于摆脱了它,我便看清:所说的“印记”,是一场粗劣的藏掖;我不肯看见它,是出于我身上的某种纵容——6月1日的“暧昧”一注(n\textsuperscript{o}~63$''$)里,我对此有清楚的认识。至于“冲突对我年轻而光彩的朋友那股最初冲劲的压制”,我几乎把它说成一场令人惋惜的宿命,仿佛这位可怜人是全然不自愿的受害者,从而——唉!——也一并取消了“大命运”的收益。然而,他对自己的命运,负有与我对我自己的命运同样的责任。倘若早在我离去之前,他就选择了做他导师的掘墓人(权且从这里开始)的角色,倘若种种情势(其中包括时代精神)有利于这一选择,把“打谁都行的大老板”的角色让他尽情享受,那么,他也同样选择了把威望与权力所能给人的特权——包括(悄悄地)碾压与巧取的特权——品尝到最后一滴。鱼与熊掌不可兼得;按照事情的本性,经由这一选择(在这选择里,他自有一些好同伴),他失去的是那些更精微、更不趋时的东西的收益……(6月初的一则未注日期的脚注)。}

\notehead{63$_{1}$}{}

(5月26日)亦可与注60末尾脚注(*)处的评语对读:那里确认了Hodge--Deligne理论的自然的发育所遭到的“封锁”——由于对我引入的某些观念核(这里是六种运算——Motif与它们密不可分)的拒斥态度;那态度与这里考察的同属一类,因此早在“Hodge理论I与II”发表时就已显形。

同样这种态度——竭尽所能(乃至过之!)抹去我的影响的任何痕迹——也见于与Mumford合作的那项工作(注n\textsuperscript{o}~47已提及):论模多重簇的Mumford--Deligne紧化。(这项工作同样先于我的离去。)它利用一条从$\mathbb{C}$上的拓扑结果(经由超越途径得知)过渡到特征$p>0$的结果的原理——那条原理,是我在五十年代末为基本群理论引入的\dmark{}\origpage{237}。早在六十年代初,我就建议利用这一方法,在全特征下证明模簇的连通性\fnote{32}{(*)}{(1984年9月)核实之后,这一情节确确实实在前引工作的导言里(第75页)得到了记载。}。然而这一想法碰到了一些技术困难,它们曾使Mumford却步,而在他们那项工作里被优雅地克服:办法是引入\textbf{模多重簇},以及一种具有完美性质的“紧化”。模多重簇这个想法本身,至少“在字里行间”,可以在我在Cartan研讨班上的“Teichmüller”报告里找到——那还是site与topos的语言尚不存在的年代。Deligne使用的语言本身——在他本可使用一整套为表达这类情形量身定做的site、topos、多重簇的语言的地方,却用了“algebraic stack”——相当清楚地显示了(带着事后的认识,并且借着后来规模大得多的几场“战役”之光照亮)抹去这项光辉工作中所运用的某些主要观念之来源的意图。肯定正是这种态度(我在“拒受遗产——或矛盾的代价”一注,即n\textsuperscript{o}~47里,头一回预感到它)产生了那记“电锯效应”,斩断了其后对模多重簇的反思——尽管在我看来,模多重簇属于迄今所提炼出的一切“具体的”数学对象之中最美、最基本之列。

顺带指出:我在五十年代末引入的那些论证,使人得以(有赖于Mumford--Deligne紧化)不仅在全特征下证明模多重簇的连通性,而且确定它们的“素于$p$的基本群”:即寻常Teichmüller群的“素于$p$的投射有限紧化”。

\subsection{高升}

\notehead{63$'$}{}

(5月10日)带着再多不到三个星期的事后认识,我如今看清:我当时那种自以为“体谅”的态度——对着那“再自然不过”的保持距离的意图——实际上是一种洞察力的欠缺,是对我年轻而光彩的朋友的一份纵容。倘若那时我信任自己健全的知觉,而不肯被那些含糊的老生常谈弄得眼花缭乱、哄骗自己——它们摆出“体谅”的姿态,甚至摆出“宽宏大量”的姿态(“总不能因为他不把我的名字捧出来,我就说他几句……”)——我那时就会看清我如今(十六年后)看清的东西。我可以称之为:对读者、对我自己、也对他自己,缺乏一份正直。简简单单地看事情,毫不害怕把它们\dmark{}\origpage{238}叫出它们的名字,我本会有能力像现在这样坦然谈论它们;而我的朋友当时本可从中取法——或至少他会懂得:即便凭他那般手段,他的长辈们(或至少其中一位)也期待他在工作中具有与他们自己相同的那份正直。于是我看清:在那次机缘里——它在我退出数学舞台之前,彼时我绝非“局外”,而且对年轻朋友无疑还有某种道德上的分量——我经由当时表现出的这份\textbf{松懈}\fnote{33}{(*)}{(5月28日)这里用“纵容”一词,比那个有点躲闪的“松懈”,更能表达我那份态度的性质。在我与我年轻而光彩的朋友的关系中的这份纵容,在昨天的反思里显现得更清楚;见“格格不入者”一注,n\textsuperscript{o}~67$'$。},没有尽到我对他的责任。这一松懈,在“Hodge理论II”发表时得到了印证:那是Deligne的学位论文工作,其中对Motif、对我,都不置一词。诚然,到那时,数学和我朋友其人,对我都已十分遥远,仿佛隔着一层浓雾!

从我得以观察到的我朋友的演变——精神上的与数学上的(这两个侧面血肉相连)——我如今看清:在我与他相识、并被他的智力手段、他的目光的敏锐、他在数学上理解的敏捷所打动的那个时刻,我丝毫没有在他身上辨出成熟的欠缺;也没有(后来)辨出,他社会地位的扶摇直上——在不过四年的工夫里,从一个无名的学生,一跃而成为数学世界的明星、终身教授,在一个已然声名显赫的机构里被授予可观的特权与权力——将对他产生怎样的影响。我不后悔为他铺平了这条路,还让它走得更快——但我确认:那是出于我自身的眼力与成熟之不足。我为他尽的这份“效劳”,算不得效劳。只要我的朋友自己尚未把那场收成——他在我漫不经心的协助下为自己预备的收成——进行到底,这就不曾是一场“效劳”。

\subsection{暧昧}

\notehead{63$''$}{}

(6月1日)自从在我与我的朋友皮埃尔的关系中确认了那份“松懈”(或按其间出现的更贴切的说法,那份“纵容”)以来的三个星期里,我已有机会在反思中更清楚地认识到我身上的某种不够严格、某种纵容。它们最初显现于我与那个人的关系——在我那里,我比对任何人都更把他当作“格格不入者”来对待——但也显现于我与另一些数学家的关系,在后者眼里我扮演着长辈的角色。迄今我在这个方向上侦察到的\dmark{}\origpage{239},表现为我身上的一种暧昧——恐怕也在那位扮演学生角色的人身上:在那些场合,他把得自我的观念与方法、乃至我交给他的一项完整工作的详细施工图据为己有,而不明确指出来源,有时甚至只字不提。这类情形相当常见,既见于六十年代,也见于我离去之后直到最近这几年的年月。我觉得,在所有这些情形里,在某个层面上,我感到了那份暧昧;它表现为一丝不安的阴影,直到最近这些天才头一回被检视。使我入局、参与某种心照不宣,并使我掠过那丝不安、从不予理会的动机,在于\textbf{顺应}我对自己、对所谓“宽宏大量”应有的样子所抱的某种形象。真正的宽宏,并非生于一种顺应成见的心态,也不是出于想要成为“宽宏的人”、并在自己与他人面前显得宽宏的心思。被压抑的不安,每一次都是明明白白的征兆:这份“宽宏”是做作的,是一种\textbf{姿态},而不是宽宏之真正的、毫无保留的自发赠予。

在这份不安里,我辨出两种来源不同的成分。一种来自“老板”、来自那个始终不得满足的“我”:他没能两头都赢——既分享一项他知道其中有自己(或大或小)一份的工作的信誉,同时又配得上某种品牌形象——其中赫然印着“宽宏大量”这条老生常谈的标签(此外还有许多别的标签)。另一种成分来自“孩子”、来自我身上那个不为姿态与门面所欺的人:他朴素地感到这境况的虚伪\fnote{34}{(*)}{(6月5日)当我在这里说,不安(部分地)来自“孩子”时,这种说法给人的是一幅不真实的图景。并不是对一个虚伪境况的天真知觉造成了什么不安。不安是一个\textbf{抵抗}的征兆——抵抗那份知觉,是(在某些层面上分明被知觉到的)实在(这里是:一个虚伪的境况)与我所攀住的实在\textbf{图景}之间的“脱节”(在这个案子里,图景即:我正在“宽宏大量”,而且不能不这样!)——我为了图景的利益,把那份不合时宜的知觉\textbf{推开}、压抑。在这个案子里,一旦我放弃抵抗、容许知觉进入意识之眼的视野,“不安”便随之止息,连同那个虚伪的境况。我本想说:“这是就涉及我当下的虚伪境况而言,而不是就处在过去的境况而言。”但转念一想,我发现:我方才谈到的那些“过去的”虚伪境况,直到今天仍以本来的面目在身,或者至少直到三天前的反思——仅仅因为它们从未被检视、从而从未被解决。我成了它们的囚徒,以至于一有机会,就机械地重演同样的境况。我对自己的“观”的力量的知识(我在“欲与观”一节,n\textsuperscript{o}~36里谈过它)当时对我毫无用处,因为我缺乏对日复一日身陷其中的境况、对知觉与其“拣择”的不停息的游戏的留意——那孩子与老板的游戏使知觉沉默……}。不但对我自己虚伪,而且\dmark{}\origpage{240}对他人也虚伪。

总之,我的“宽宏”在于入局:入一场游戏,游戏中对方把他得自他人的观念呈献为自己的,也就是呈献一幅他自己与某种实在的图景——而他和我都心知肚明,那是假的。于是我们成了同谋,共谋一桩可称之为“作弊”的事,双方各得其所。至少按“我那个年代”通行的共识,这是一桩“作弊”——而那些共识,依我看,今天人们仍在嘴边念叨。倘若被当作“我庇护的人”的发现来使用的是别人的观念,我肯定不会入这样的局\fnote{35}{(*)}{“我庇护的人”这个说法,曾被我昔日的一个学生用来称呼我当时的、刚在数学里做出漂亮东西的一个学生;它曾使我咬牙切齿。然而,细想之下,我此刻正在考察的这份暧昧的境况,确立的恰恰是一种虚伪的关系:其中一方俨然是另一方的“受庇护者”。}。然而,我给予默许,让诞生于我的观念被呈献为他人之物,并不改变事情的本质——唯一的区别是:这回我们是两个人在作弊,而不是一个人。撇开涉及我个人的这一面(我自己参与一场作弊,一种违背我自己也自认服膺的共识的行为)不谈,有一点再清楚不过:怂恿他人作弊(即便这作弊看似只由我们自掏腰包——而事实绝非如此),或至少怂恿他人对一个他自己也假装服膺的共识持一种暧昧的态度、实则阳奉阴违,这里没有任何宽宏可言。真正的宽宏,本性上造福所有人,首先要造福它于其中显现的那个人、它所施与的那个人。我的这份暧昧态度——在人身上诱发或怂恿暧昧,并使我得以摆出“宽宏”的架势,而按正常的道理,对方这时应当显出几分作弊者的模样(而事实上我们俩都在作弊)——这份态度,对我既非善举,对他亦然。

只需把事情检视一番,明摆着的道理就会显现,甚至无须求助于经验、求助于“事件的教训”。然而,终究是事件把我引到了这番检视,使我终于发现一个明摆着的道理——三十\dmark{}\origpage{241}年前我同样有能力发现它,那时地平线上还没有出现一个学生,要跟我学一门手艺,并在与我的接触中浸染于这门手艺练习中的某种精神。我已有机会谈到工作之中的那种“严格”——我相信自己做出过表率(见“严与苛”一节,n\textsuperscript{o}~26)。但今天我还察觉到,在“工作”本身之外,也有一份严格之不足,表现为暧昧,表现为我所说的纵容。我觉得,我的这份暧昧,不是从我的任何一位长辈那里传来的——他们(我相信)对我都有与自己对自己相当的要求。在这一特定态度的暧昧之外,我还在我这个人本身之中辨出一种暧昧——在《收获与播种》第一部分,我已有不止一次机会谈到它。这份暧昧,随着1976年禅观的发现开始消解;而它的某些征兆——表现为已然成习的态度与行为(尤其在我与学生的关系里)——想必一直延续到今天。

显然,我身上这份暧昧,在我的某些学生那里找到了肥沃的土壤。那种经由默契做成的事,如今似乎已成了“上流”数学界习俗中的一个基调:在那里,浑水摸鱼(有或没有“当事人”的同意)、乃至有组织的劫掠(当放胆为之的人属于那不可撼动的精英时),似乎已成了寻常得无人再表惊讶的惯例——而同时,人人也都绝口不提。我身上的“老板”很想划清界限、揭发、愤慨——可这样一来,我不过是在延续我身上那份同样的暧昧,它的丰饶收成,我今天有目共睹。

\subsection{同伙}

\notehead{63$'''$}{}

(4月24日)\fnote{36}{(*)}{本注出自“本能与时尚——或强者的法则”一注(n\textsuperscript{o}~48)的一条脚注——在那则注里,我曾断言Verdier关于导出范畴的工作从未发表,却不知道其学位论文的一个“第0状态”已于1977年发表。至于Verdier围绕着那本应构成其学位论文工作的理论所玩出的那些古怪的腾挪,其全貌可参见“赊账的学位论文与全险”一注,n\textsuperscript{o}~81。}两天前翻阅刚收到的Mebkhout的一份抽印本,我撞见一处对 J.L.~Verdier一项工作的引证,题为“导出范畴,第0状态”,发表于SGA~4½(Lecture~Notes n\textsuperscript{o}~569,第262--311页)。我到今天才得知这一发表,情有可原:此前我从未有过捧起这一卷的荣幸——无论Verdier,还是其作者Deligne\dmark{}\origpage{242},都不曾费心在出版时或后来寄我一本。我不知道,当年与我共同组成评审委员会、仅凭一篇17页的导言(至今未发表)就把“理学博士”头衔授予 J.L.~Verdier的 C.~Chevalley与 R.~Godement,十年之后,有没有份收到这部不同寻常的“学位论文”的“第0状态”(这一回有50页)!我记得有一天曾亲手捧读过一部几百页的严肃奠基工作,它满可以理直气壮地算一篇优秀的博士论文,并且大体对应于我1960年前后向Verdier提议的那项奠基工作——只除了一点:彼时“三角范畴”框架(他用来表达导出范畴的内蕴结构)业已显出不够用。

几乎无须说:我的名字,在这份学位论文的“第0状态”里无处可见。人们纳闷:它到那里去做什么。众所周知,导出范畴是Verdier引入的,为的是让他得以发展拓扑空间的所谓“Poincaré--Verdier对偶”与解析空间的所谓“Serre--Verdier对偶”;只待某位当值无名之辈\fnote{37}{(*)}{关于这位可疑人物的几点情况,可参见“当值无名之辈与好上帝定理”一注(n\textsuperscript{o}~48$'$)。}出头,为他自己发展一个把二者综合起来的理论,并且理所应当地(无名的学生至少也得做到这一步!)被称作“Poincaré--Serre--Verdier对偶”。经过这一切,我剩下的只有随大流,做几处势在必行的调整,以便在概形的平展上同调或凝聚上同调那个再特殊不过的框架里,发展Poincaré--Verdier对偶与Serre--Verdier对偶……

我刚刚才得知(图书馆到底有用!)SGA~4½的内容\fnote{38}{(**)}{关于这一卷,可参见“白板”一注,n\textsuperscript{o}~67。}——在那里,人们还不吝让我以合著者、或不如说Deligne的“合作者”(原文如此)的身份露面(既不费心告知我,更不费心征求我意见)。显然,这是五年后问世的“下葬文集”的先声——那部文集,几天前我有幸拜读(见由这一事件激发的注n\textsuperscript{o}~50、51及以下)。但我无需亲手捧起那部“预葬”文集,连同那篇不肯自报家门的幽灵论文的铁证,便早在去年就已懂得:这份“学位论文”的下一个状态,除了我,不会有别人来写。于是我就着手这项工程,即\textbf{《追寻栈》}\dmark{}\origpage{243},从我那位显赫的前学生十七年前停下来的地方接着做。

\subsection{授职}

\notehead{64}{}

(4月25日)然而昨天,我在系里的办公室里找回了一份副本。其实是前后相隔一年、写于1968年4月(?)与1969年4月的两份报告。我用十七页的篇幅,评阅了他在IHES三年科学活动中进行的十五项工作。其中有关于Ramanujan猜想的工作(原书拼作Ramanuyam),关于模site的紧化的工作,以及Hodge理论的推广。报告里评阅的这些著作(单是我刚点名的几项就足够了)见证了一种惊人的创造性,以完美的从容、仿佛游戏般地铺展。把Weil猜想的证明——那还是紧随这第一次向未知的冲刺之后——放到一边,我觉得,他后来的著作,只给了那次独一无二的起飞一副苍白的图像:一个手段出众的年轻心灵,享受着使其绽放的种种同样出众的条件。然而,不得不相信,这些“出众的条件”中的某种东西,必定喂养了另一股力量——一股与求知的冲动异质的力量;它最终反客为主、取而代之,使最初的冲劲改道并被吞并。而显然,这个“某种东西”与我这个人有关……\fnote{39}{(*)}{(5月26日)至于我身上的某种纵容——它喂养了这个“某种东西”——可参见比本注晚两周的“高升”一注(n\textsuperscript{o}~63$'$)。}

这篇简短的评述报告(我打算作为附录收入本卷)从不止一个角度看都饶有趣味,包括数学的角度(评阅的有些著作,至今仍未发表)。在报告的几处,我预言:Deligne只勾勒了轮廓、处理了关键之点的某些工作,将由未来的学生们去发展。这些学生从未出现,缘由是他后来对常人世界的关系起了变化\fnote{40}{(**)}{在我还常在IHES与他为伴的年月里(尤其在我的研讨班上),Deligne对其他数学家、尤其对来参加研讨班的青年研究者(往往是初学者)的关系,带着亲切。我在那里看到的对他人思想的开放——哪怕它表达得笨拙甚至混乱——同我们数学的两两对谈里一样。他有一种能力:循着他人的思想,进入他人的图像与语言;这能力我一直欠缺,而且(依我看)使他比我远更适合“老师”的角色——善于激发他人身上天职与创造性的绽放。}。在我评阅的那些想法之中,据我所知\dmark{}\origpage{244}唯一曾被别人发展出来的(那人由此俨然成了Deligne的学生),是上同调下降理论——由圣多纳在SGA~4里发展(因此还处在最初冲劲的时期),此后成了上同调武库里最常用的工具之一。有一处既有趣又有代表性的细节:后来成为Deligne文章的四项工作里有三项\fnote{41}{(*)}{四项工作中唯一不受我直接影响的一项,是从Weil猜想推出Ramanujan猜想的工作;它处在一个研究方向(模形式的方向)上,那是我数学修养里最严重的“空洞”之一。其余三项是:Leray谱序列的退化,Hodge--Deligne理论,模多重簇(与Mumford合作)——见“驱逐”一注(n\textsuperscript{o}~63)与子注n\textsuperscript{o}~63$_{1}$。},我在行文之间,以动人的细心,让读者感觉到这些工作与我引入的想法、我提出的问题的关系——仿佛要抢在前面,预先对付作者日后在他的文章里对这些内容的沉默(这几篇文章,在我写报告的时候,没有一篇已发表,而且我相信,甚至还没有成文)。

\subsection{死结}

\notehead{65}{}

(4月26日)同样清楚的是:把一部大面积的“瑜伽”(权的瑜伽,进而Motif的瑜伽)留在自己怀里——关于它,我对除他之外的人也不止一次谈起过,但唯有他亲密地消化了它、把握了它的全部分量——这本身授予他一种额外的“优势”:作为一件无可比拟的发现工具的独家持有者,用以理解代数簇的上同调。然而我不认为,在我还完完全全在场、活跃于数学世界、毫无迹象预示我无限期离去的那个时刻,这一诱惑起过什么决定性作用。它必定是随着、或在我的离去之后才出现;我的离去,是攫取一份遗产的“天赐良机”——那份遗产本是完全理当归他所有!——只要藏起遗产,连同它的来源。

正是在这里,我又一次看见那个深藏的矛盾的死结,在一个极端而格外触目的个案里,远远超出一切个案:我想谈的是无知、不屑、深深埋藏的怀疑——它们围绕着我们自己身上所怀的那股创造之力:那笔独一无二的遗产,比一个人所能传递的一切都更珍贵。正是这种无知、这种对我们身上最珍贵、最稀有者的阴险的异化\dmark{}\origpage{245},使我们得以嫉妒在他者身上知觉到的力量,使我们得以贪求那力量在他人身上的果实与外部印记——那股我们已在自己身上遗忘的力量。一旦这份嫉妒、这份\textbf{取而代之}的欲望扎下根来、找到滋蔓的机会,把本可用于创造性绽放的能量引向这一欲望,我们身上的异化就变得更深,安营扎寨。我们越接近“取代、排挤、炫目”这个觊觎的“目标”,就越远离、越割断我们身上那股精微的力量,也就剪断我们自己创造冲劲的翅膀。在我们顽强地抬高自己的持久努力中,我们久已忘记飞翔,忘记我们生来就是为飞翔的。

在与我的关系里,从我们相识那天起,我就感到我的朋友完全自在,没有任何征兆使我疑心:他会为我的名声或我这个人感到丝毫的震撼或目眩,或者他心中有任何未曾说出的怀疑——无论是关于他在数学领域的天赋,还是关于任何别的。同样,依我看,他也确实从我和我所处的环境——包括我的家庭——得到了一种友好的、亲切的接待,足以使他自在。但那份吸引我、也吸引别人的单纯而似乎毫无疙瘩的自然,肯定不是等到这次相识才出现并绽放的。他这个人散发的、使他如此可亲的印象,是一种和谐的平衡:他对数学的爱好,丝毫不带吞噬人心的女神的形象。在他旁边,我多少显得像个不知悔改的“书呆子”,甚至“粗鲁迟钝的家伙”——我还记得,他含蓄地惊讶于我与身边的大自然、与四季的节律之间缺乏深入的接触;而我穿行其间,几乎视而不见……

然而,那份深藏的“怀疑”——当时我绝无能力知觉它(甚至今天,处在类似的境况里,恐怕也未必)——必定在我的朋友身上,早在我们相识之前就已存在。带着事后的认识,我毫无歧义地看到它的第一个征兆早在1968年;其后这些年里,又有更清晰的征兆\fnote{42}{(*)}{(5月10日)实际上,还有一个“十分清晰”的征兆,早在1966年;见注n\textsuperscript{o}~82的脚注(*)(原书第329页)。}。不过这些都是“间接的”征兆——从我第一手观察到的东西里,没有一个以怀疑、缺乏自信的面目出现;恰恰相反,而且年复一年越来越如此:一种自满,一种\dmark{}\origpage{246}蓄意的轻蔑口吻,乃至鄙夷。但这样的“反面”恰恰揭示其对面——它们成双,反面是正面的影子。

我还经由第三者得知:对某位他从未有机会亲近相识的、声名显赫(而且以难相处著称)的数学家,他曾处于巨大的紧张期待之中,生怕那位大人物不把他看作与自己的伟大相称。这一证词,与我在年轻朋友身上亲眼所见的截然相反,以至于我很难相信它(那是1973年)。带着事后的认识,它却与我另行得知的那些分裂的征兆相互印证,而且全都指向同一个方向。这份分裂,以及我在其中扮演的角色——作为一桩在相识之前想必还弥散无形的冲突的某种凝固剂——在通常情形下,在与一个曾(以这样或那样的意义)是“老师”、或至少是传递者或托付者的人的关系演变里,恐怕会一直隐而不显。这样,我的离去,就成了一个所有人都未曾察觉、而我或许是唯一知情的冲突的\textbf{显影剂}。

而我今天的“归来”,则是第二个显影剂,无疑更加不合时宜。除了它迄今已就我自己的过去、我的现在、以及我曾爱过而至今仍与之相连的人们所教给我的,我全然无法想象它还将向我揭示什么。也不知道它将向那个人揭示什么——最近一个星期以来,他正处在我反思的这一最终阶段的中心;上个月我曾给这一阶段命名(而且我没想到说得这样准……),叫作“\textbf{一段过去的重量}”。

\subsection{两个转折}

\notehead{66}{}

(4月25日)我的朋友皮埃尔对我的关系中那种蓄意的轻蔑与对抗,只限于数学与职业的层面。个人的关系直到今天仍是一种情谊与敬重的关系,不止一次表现为打动了我的细致体贴,分明是真实而无保留的感情的征兆。

在我离开IHES之后那些紧张的年月里,IHES渐渐沉入遗忘,一如这一情节长期不为我所理解的教诲。于是又有十多年,我的朋友在我这里(想当然地)依然是我数学上的特权对话者;或者更确切地说,1970到1981年间,他是我唯一的对话者(除一桩情节外):在我数学活动断断续续的时期,每当需要对话者的时候,我想到的就是他\dmark{}\origpage{247}。

也是把他当作离我最近的数学家,在最初几次机缘里(1975至1978年间),我同样自然而然地向他求助,请他为与我一起工作的学生提供帮助、担保或支持。其中第一次,是1975年Sinh~夫人的学位论文答辩——她在越南,在异常艰难的条件下写成了论文。他是我为论文委员会联络的第一个人。他推却了,话里带着意思:这只能是一篇冒牌的论文,他绝无意为之作保。(不过我到底有办法说动Cartan、Schwartz、德尼与齐斯曼的好心,请他们为这出“把戏”助一臂之力——答辩在关切与温暖的同情的气氛中举行了。)在其后三年里,又经过三四次同类经验,我才终于明白:在我那位显赫而有影响的朋友身上,有一种蓄意的对抗姿态,对着我的“1970年以后”的学生,也对着只带有我的影响之印记的工作(至少“1970年以后”着手的那些)。我不知道,我在好几次这样的场合所觉察到的那种公然的轻蔑,是否也多多少少见于他对另一些数学家的关系——那些他视为远在自己之下的人。他引以为荣的那套无所节制的精英主义的精神本身,会使我作肯定的猜想。无论如何,自1978年起,我不再为任何事情向他开口。这并未妨碍他那份令人气馁的力量,又找到机会有效地施展。

大约也是在那一年,头一批征兆——起初很隐约——显现出一种对我自己的数学活动的轻蔑态度。第一次机缘,是我关于胞腔地图的反思,起于一个关于它们的发现,它曾使我如遭雷击(参见《纲领概要》第3节:“与儿童涂鸦相联的数域”)。这项发现(诚然“平凡”,本来毫无使我的显赫朋友动容乃至感兴趣之处)是另一个数学之\textbf{梦}的起点与第一份材料——那个梦,其规模可与Motif之梦相比,只是到三年之后(1981年1--6月)才随《穿越伽罗瓦理论的长征》开始成形。这些笔记与同期其他的笔记(手稿凡两千页)构成对那片“新大陆”的头一次巡视——一桩关于儿童涂鸦的平凡议论,使我隐约瞥见了它。

\dmark{}\origpage{248}在这项紧张的工作过程中,有两三次我给我的朋友写信,把我的一些\textbf{想法}告诉他,有时也提出一些技术性的问题。当他肯对我的问题开口时,他的评点始终同样清晰、同样中肯,见证着他年轻时便已打动过我的那同一副“手段”。但一份自满已经磨钝了那份曾经使我陶醉的求理解的渴望,也磨钝了那种经由“小”事把握大事的能力,那种倾听大小事物而把握或构想宏伟蓝图的能力。那种能力不属于理智、不属于单纯的“效率”、不属于对一门已成学科的“掌控”或对已知技术的“掌控”;它在理智的层面上,映照一件本质全然不同的事——孩子的那份\textbf{惊叹的天赋}。这份天赋在他身上仿佛已经熄灭,仿佛从来不曾有过。至少就他对我的关系而言是如此——而在此之前,对我的“1970年以后”的学生,已然如此。他已成了一个大人物;他对数学的进路,变成了不折不扣的那种“竞技”态度——一两个月前我才头一回考察过它,而且我自己对它也绝非陌路……

或许我本可以给自己找到一个理由,解释那在共同热情中的契合、那条曾经连接我们的深深纽带的显然缺席。我本会很满足——多半如此——在机会来临时,向朋友的巧思、向他关于数学事物世界的广博知识,提出或多或少技术性的问题,或者仅仅请教信息。但就在那一年(1981年),那份轻蔑的征兆突然变得如此粗暴\fnote{43}{(*)}{(5月28日)关于这第二个转折的新的照亮,另可参见“反常”一注,n\textsuperscript{o}~76。},以至于我对再同他交流数学问题——哪怕偶尔——全然失去了兴趣。($\Rightarrow$注67)

\subsection{白板}

\notehead{67}{}

(4月26日)正是在昨天写前面那些文字时,我把我们关系中的这个新转折,与1982年(也就是几乎恰好在这个苛刻的转折关头)Lecture~Notes那部“卓越文集”的出版联系了起来——它为我的数学下葬揭幕,没有鲜花,也没有花冠!既然我已被判定数学地“死去”,我的朋友仍在这里或那里继续回答一些归根结底已无从必要的数学问题,这算得上某种恩典……

\pmark \origpage{249}试着去倾听事件的意味,我感到,同样不是偶然:轻蔑、数学上的漠然(对着那些——他的“健全的数学本能”本该告诉他——灼热而多汁的东西!)的初次显现,至少就他对我的关系而言,差不多就落在“预葬文集”SGA~4½出版前后的时刻,即五年之前\fnote{44}{(*)}{关于此事,可参见前天“同伙”一注(n\textsuperscript{o}~63$'''$)。(6月5日)本注的反思,在本注及随后三则注(“白板”“格格不入者”“绿灯”“颠倒”)里得到继续;它们使人隐约看到“SGA~4½之役”的意味,及其与母研讨班SGA~5之“拆解”的联系。这一反思又在“我的学生们”的行列里、尤其在“我的学生们(1)--(7)”的续篇里重新拾起;在那里,一场对那个我的上同调学家学生们学艺的研讨班的真正屠杀的图景,渐渐显露。此外,整场战役里摊开着一种漫不经心的轻蔑——我对它在同一时期初次显现的那种“隐约的轻蔑”的知觉,在我朋友对我的关系里,只不过是一个非常苍白的反照。还有一桩联想,是一两个星期前来到的,关于我朋友对我的关系中那个“第一个转折”的时刻——1977年底或1978年间。正是在1978年,我的朋友得到了他那枚实至名归的“奖章”(为了Weil猜想的证明)。这一新头衔(系于一桩“困难尽人皆知的”猜想的证明)被我的朋友以何种方式内化,在“葬礼悼词”(关于我这亡故者的文章)及其对偶文章(关于他自己的)里,触目惊心;那两篇文章,诚然是整整五年之后,才在一个“重大场合”问世。对此可参见“葬礼悼词(1)——或恭维”一注,n\textsuperscript{o}~104。}。单是这一卷出版时的种种情势,就已足以见证一种蓄意的轻蔑姿态——既是隐约的,又是炫耀的。不劳烦征求我意见,甚至不告知我一声,就把我列为Deligne的“合作者”,而且小心提防不给我寄一册——单是这一点,在我看来,就比一篇演说更有说服力。更不必说,Deligne这部书,就其主体而言,本是要把我在十五年多以前发展的那些工作——那时我还没有听说过我这位光彩夺目的朋友的名字!——变得让更广大的读者更易亲近。一份轻蔑,以及继之而来的傲慢,必定由两头喂养:一头是我的缺席——它使我对一切毫无察觉,浑然不觉地“照单全收”;另一头是某种气候——它使这类张冠李戴得以“通行”,似乎引不起最轻微的评注。总之,关于这一卷,以及关于他筹备的那部下葬文集,我没有收到任何人(尤其在我自以为在数学家世界里还拥有的一众朋友中间)的只字回音。

\pmark \origpage{250}在导言里,作者更是开门见山地亮明底色。该卷的目的,是让“非专家”免于“求诸SGA~4与SGA~5那些芜杂的报告”,“剪除多余的细节”,“让使用者得以忘掉SGA~5——人们可以把它看作一系列枝蔓,其中一些还非常有趣”(对这些“枝蔓”倒还算客气!)。“SGA~4½的存在,不久将使人得以照原样出版SGA~5”——一句神秘的话,因为人们纳闷:这一出版(出的是一部他们劝人忘掉的东西)拖拖拉拉已有十二个年头,而且拥有一组自洽完美的结果(它们不曾等Deligne来提炼与证明),凭什么要以SGA~4½的存在为前提。在提出这个问题的同时,我也隐约看到一个简单的回答,以及那个可怜的SGA~5研讨班坎坷命运的也许一种解释(注68)(那是我于1965/66年间纵横发展起来的,比DeligneSGA~4½分卷的出版早了十一年\fnote{45}{(*)}{对这一“奥秘”的澄清,可见“绿灯”一注(n\textsuperscript{o}~68)的一条脚注(4月28日所写)。})。当人们被告知(第2页):SGA~5的原始版本里,“Lefschetz--Verdier公式只是猜想式地确立的”(这对Verdier未免刻薄——他本该会证明自己的定理,而那定理比SGA~5更早\fnote{46}{(**)}{(6月10日)对此的更精确说明,见“屠杀”一注(n\textsuperscript{o}~87)的子注n\textsuperscript{o}~87(原书如此,未印附号)。}),而且“此外,诸局部项在那里也未加计算”。这对非专家读者(本卷首先面向的读者)可能显得是一处令人遗憾的缺漏。略知门道的读者则很清楚:这些“局部项”,到今天也仍然没有“计算”出来;而且,倘若有人问那位光彩夺目、武断专行的作者,他在这句话里(就一般情形而言)所说的“计算”究竟何指,他自己也会窘迫万分\fnote{47}{(***)}{(6月10日)在一般的Lefschetz--Verdier公式里,对一个系数层与其自身的上同调对应,那些“局部项”(对应于不动点集的连通分支)在写出公式的一刻便毫无歧义地得到定义。这些局部项的“计算”问题,只在具体的个案里才取得精确的意义;其中最简单的之一,是Frobenius态射的情形——在那里,它们简单地由这些点上的纤维所诱导的自同态的寻常迹给出。这条公式,在口头研讨班里,已作为一个远更一般的公式的特例得到完整的证明。}(不过,显然没有人想到去问这句不客气的问题)。

一句含混的话——“这个研讨班(? )含有\dmark{}\origpage{251}另一个证明,这一份在Frobenius态射的特殊情形下是完整的”——似乎暗示:SGA~5终究没有给出(谁又能料到呢,在一部“枝蔓的文集”里!)它所宣布的那个主要“结果”——一条迹公式,从而蕴含Weil式的L函数有理性——的完整证明;幸好“这个研讨班”前来救场,亡羊补牢,犹未为晚……

在第4页,我们得知,“Arcata”诸报告的目的,是“给出平展上同调诸基本定理的证明,剥去包裹着它们的、SGA~4里肆虐的non-sense”\fnote{48}{(*)}{在“我那个年代”,英文里那个沿用已久的词“général non-sense”(泛泛之谈,意为:有时烦人却往往必要的通论)并没有贬义,而更多带着一点玩笑与憨直的意味。那个沿用已久的定语“général”在这里被“忘掉”,肯定不是偶然:这样一来就成了“non sensé”,用漂亮的法语说,不多不少正是“无意义”,还让人联想到“嗡嗡的空谈”(bombinage)和“胡扯”。}。他还有这份慈悲,不去细说SGA~4里肆虐的那份令人惋惜的non-sense(比如topos以及诸如此类的恐怖之物——读者可以庆幸自己福星高照:靠了这部光辉文集的天意降临,终于对那可惜的“矿渣”来了一场大扫除……)(注67$'$)(注67$_{1}$)。

我此刻把这一卷的导言与各章导言通读了一遍,把在我看来最明白地亮明底色的那些评语与自白抄录于此,另有两三句(路数是:枝蔓嘛,当然,不过“非常有趣”)在我看来主要是用来“把苦药裹上糖衣”的(而且糖衣果然顺利咽下)。比如,作者倒有这份诚实,开宗明义地说清楚:“若求完整的结果与详细的证明,SGA~4仍然不可或缺”。这一卷,无论其精神与动机多么暧昧,并不属于一场诈骗\fnote{49}{(**)}{(5月26日)不过,请看后天的“颠倒”一注(n\textsuperscript{o}~68$'$):在那里,我回到了这一印象——它被证明是仓促的。在反思的进程里,一场大面积的“SGA~4½--SGA~5”战役渐渐显露,它主要是为了Deligne的“收益”而做成,有我全部“上同调学家”学生的协助或默许。我所相信可以确认的那份“诚实”(依据导言第7行刚引用的那句话),在这里扮演的,是“对照线”的角色,用以障人眼目——是地道的“暂停!”风格。我的朋友从1968年起就在用这种风格(见“罐装的权与十二年的秘密”与“驱逐”两注,n\textsuperscript{o}~49与63)。另可参见“暂停!”与“中国皇帝的新衣”两注,n\textsuperscript{o}~77与77$'$。}式的战役。它的角色,在我看来,更近于一次试探——显然是结论性的:原来根本用不着这么多顾忌!

\dmark{}\origpage{252}有一种\textbf{荒谬性上的步步升级}(显然无人察觉!):从这一卷,到他筹备的那一卷(SGA~4½,以及LN~900)。在两处,人们都看到一个手段出众的人——一个生来就该去发现、去巡游、去探测浩瀚世界的人——执意去“重做”一位先行者的工作:先是我,然后是我的一位旧学生(Saavedra);而与此同时,他对这些先行者的工作——它们都是细心做出、并且深探到底的——拿不出任何本质的新东西。(他所能加进去的东西,总共恐怕二十页、三十页就写得下,依我看。)在第一个个案里,给出的理由还说得过去:让非专家的使用者得以“无泪地”亲近平展上同调\fnote{50}{(*)}{(6月10日)写这条注时,我才刚刚“登陆”,还没有感觉到“SGA~4½之役”的真正意味(及其与SGA~5坎坷命运的联系——对后者,我恰在那时才有了一次突然的预感)。后来我懂得:以那个骗人的名义SGA~4½发表的这堆五花八门的文本(见“颠倒”一注,n\textsuperscript{o}~68$'$),根本不是要充当SGA~4与SGA~5研讨班(它们构成我已刊数学著作的心脏)的一部“无泪”的普及读物,而是一次用它取代我已刊著作的行动:把我的著作呈现为一个多少有些含混的先声,把它自己呈现为关于平展上同调的真正杰作——仿佛这杰作当归功于Deligne。对这样一种冒名顶替的惊人表述(出自一支匿名的手笔),可在“SGA~4½”这次“试探”六年之后,“葬礼悼词(1)——或恭维”一注(n\textsuperscript{o}~104)里看到。},不必倚仗卷帙浩繁的SGA~4与SGA~5。(不过,人们头一回看到这位作者对常人世界如此体贴,竟压过了做数学的乐趣……)第二个个案里,这项工作实际上等于\textbf{照抄}Saavedra跟我做的学位论文!那篇论文本是一个完美的参考;其中一个陈述的证明有误、另一个陈述带有一条多余假设,这肯定不是重写整篇文章的理由。当然,对这样一件怪事,没有给出任何“理由”。

然而,我无需亲手捧起SGA~4½,就感到了这件表面上荒谬的事情的意味:Deligne在十年之后“重做”Saavedra的论文!它无疑同那件只稍逊荒谬的事的意味相同,而后者为它做了铺垫:Deligne(在十二年之后)对Grothendieck已刊著作的某一部分做一个“摘要”(边缘上还带几分居高临下)。恰恰是这一部分,倘若他还继续对代数簇的上同调(他始终无法摆脱它)感兴趣,他在任何情形下都不能装作可以没有。而Saavedra的论文,在所有已刊的、带着我的影响之印记的工作里,恰恰是他无论如何都绕不开的那一部——倘若他想把我发展的\dmark{}\origpage{253}MotifGalois群观念“记在自己名下”,并终于(十五年后!)开发利用这个显然至关重要的观念。经由先执笔SGA~4½,五年后又经由LN~900里那篇江河浩荡的Milne--Deligne文章(别名Saavedra),我的朋友沾沾自喜地给了自己一种虚幻的解脱感——解脱于他分明感到是痛苦负担的东西:不得不时时引证那个他正要取代、否认的人,哪怕仅仅引证某个引证他的人。

要抵达这一亲身的信念——关于这两个“荒谬”举动所共有的意味——我根本无需通读我这位多产朋友那(五十一项)出版物的全单;那份清单,十天前我才(头一回)收到。老实说,我甚至没有想到把我手里的四份抽印本再翻一遍\fnote{51}{(*)}{此外,还不算载于IHES《Publications Mathématiques》的那些工作——所长Nico~Kuiper,将近十五年来,一直好意地把它们寄给我。},为我相信知道的事找佐证。倘若今后我还会查阅我朋友的著作,那将是为了在里面找到别的东西——那些我在别处已经知道得足够多的东西之外的东西。那时我肯定会有学习美好数学的乐趣——那些数学,当年我曾有幸当面聆听、亲口从他那里学习,那乐趣还要更大!

\notehead{67$_{1}$}{}

\begingroup
\renewcommand{\labelenumi}{(\arabic{enumi})}
\begin{enumerate}
\item (6月14日)在SGA~4½里,我又记下了两桩微型的(细节上的)巧取。其一在“SGA~4,SGA~4½,SGA~5的引路线”(请欣赏这串富有暗示的排列!)里:作者写道(第2页),为了在平展上同调里确立“一个类似于凝聚对偶的对偶形式主义……Grothendieck利用了奇点消解与纯粹性猜想”,从而给人这样的印象:这一形式主义,最终是靠他Deligne本人才确立起来的——在(对许多应用已然足够)0维或1维正则概形上有限型概形的情形(见同一段)。他心里非常清楚:六种运算的形式主义(即整体对偶理论)是我不依赖任何“猜想”确立的;他所加的这条限制,只对双对偶定理(或“局部对偶”定理)才有依据——而那条定理,到头来,在SGA~5里(在Illusie的笔下)还改了姓,叫“Deligne定理”!

另一方面,第100页上有一节,题为“Nielsen--Wecken方法”——那正是我在代数几何里引入的方法,用来证明一条Nielsen--Wecken型的公式;这两位作者(在超越语境里)是用一种三角剖分技术证明它的,那技术在\dmark{}\origpage{254}代数语境里毫无用处。Deligne知道这一方法(连同Nielsen与Wecken二位先生的大名——那篇漂亮的德文文章,他根本无需去读!),是从我嘴里、在SGA~5那场“技术枝蔓”的研讨班上听来的——而SGA~4½正是要让人忘掉那个研讨班!在这一节里,既没有对SGA~5的暗示,也没有对我的暗示;关于这一方法的父权,读者可以自行选择:要么归于Nielsen--Wecken(倘若他消息极其不灵),要么归于这位光彩夺目又谦虚过人的作者。

有意思的是:整部书里,Verdier那“Woodshoie”式的证明(照录原书拼写)——关于一条迹公式,包括我所需要(为了Frobenius态射)的那个情形——只字未提。这一证明(显然已被人遗忘,让位给SGA~5里发展的那个更一般的方法)本是把我的L函数的上同调诠释完全立稳的缺失环节。显然,Deligne与Verdier之间有过(想必是心照不宣的)交易——Verdier把Weil猜想迹公式的信誉让给Deligne,作为交换,Verdier保有他在前一年(1976年)据为己有的SGA~5那一部分。(对此可参见“好的参考文献”一注,n\textsuperscript{o}~82.)另一项补偿:导出范畴与三角范畴的“第0状态”在SGA~4½里的发表——我的名字在那里同样缺席。况且四年之后,在Deligne笔下,代数几何里的平展对偶又得了“Verdier对偶”之名——Verdier这笔买卖做得不坏!(见“不义——或归来的意义”一注(n\textsuperscript{o}~75)末尾。)

\end{enumerate}
\endgroup

\subsection{格格不入者}

\notehead{67$'$}{}

(5月27日)\fnote{52}{(*)}{本注出自前一条注“白板”的一条脚注,是对它的补充,恰好一个月后的同一天写成。}所引的那些段落,连同围绕这部名叫SGA~4½的卓越文集之出版的全部情势,见证了我朋友的一种蓄意的嘲弄与轻蔑姿态——对着我的著作的核心部分:由SGA~4与SGA~5这两个血肉相连的研讨班构成的全体。从4月24日(见“同伙”一注,n\textsuperscript{o}~63$'''$)到5月18日(见“遗骸……”“……与身体”两注,n\textsuperscript{o}~88、89)的反思过程中显露出来的这些“情势”里,原初的SGA~5研讨班遭到劫掠——落实为1977年的那场屠杀式出版——绝非最轻的一项。(尤其见“屠杀”一注,n\textsuperscript{o}~87.)

我朋友这份蓄意的嘲弄,若想起以下事实,便取得全部意味:口头的SGA~5研讨班,正是年轻的Deligne\dmark{}\origpage{255}与概形、与上同调技术、尤其与对偶形式主义、与$\chi$-进上同调的初次接触——1965年,他21岁,在IHES登岸,目标十分明确:跟我学“代数几何”。正是在这场口头研讨班里,以及在两年前SGA~4的笔记里,他得以有幸第一时间学习那些直到今天仍主宰其著作的想法与技术\fnote{53}{(*)}{差不多同样的话,其实也适用于我的其他每一位上同调学家学生:Verdier、Illusie、Berthelot、Jouanolou——对此可参见“团结”一注及其后四注(n\textsuperscript{o}~85至89)。}。

“SGA~4½之役--SGA~5”、乃至我朋友皮埃尔对我个人之关系的这一本质的背景,显然在我写前一条注(“白板(1)”,n\textsuperscript{o}~67)时并不在场,在我之前关于葬礼的反思里也不在场。那个“年轻的Deligne”的记忆——他登上SGA~5研讨班,那时他还一切都等着学,而且确实(而且很快)学到了许多——只是在反思的最后阶段,仿佛违背我的本意,才浮上来。从我那个小小的数学“微观宇宙”里出现年轻的Deligne那一年起,我心中便立下一条蓄意的主意:不把他算进我的学生之列(仿佛把他算进去,就会对一位天资如此夺目的人有失谦逊);这一主意使我还一并低估了——或者不如说,直到最近这几星期才完全忽视——一个分明而具体的实在:人们通常用“师-生”这双重的称呼(而这称呼正是我一向拒斥的)所表达的那个实在\fnote{54}{(**)}{(6月14日)这一蓄意的主意,分明见于我四个月前在“耶稣与十二使徒”一注(n\textsuperscript{o}~19)里终于决定如何谈论他的那种方式——仿佛这样一来,我就在对那个乐于与我划清界限的人,违背了一项缄默或谦虚的义务……}。我曾乐于忘掉、忽视:确实有过某种从我到他身上的“传递”——某种对我、对他都具有巨大\textbf{价值}的东西,尽管对其价值所在,他与我的领会想必大不相同。在他与我之间那四年密切的数学接触中,我所传递的,是某种倾注了我自己最好部分的东西,一件由我的力量与我的爱滋养的东西——一件(我相信)我毫无保留地赠出、却不曾掂量、甚至或许不曾真正感到其分量的东西。

可以肯定,我所给予的,喂养了他身上那种求知的热情\dmark{}\origpage{256}——它与我身上跳动的那份热情同调——但也喂养了\textbf{别的东西};那东西我很久之后才感觉到,而且当时不曾把它与那桩已经发生、却乐于被我忽视的“传递”联系起来。换个说法:我所给予的,在另一个对我隐蔽的层面上,\textbf{也}被接受了——不是作为探测一个迷人的、永不枯竭的未知的工具,而是作为(起初)取而代之、(后来)确立支配、确立对他人毫不容情的“优势”的\textbf{工具}。

不去区分我朋友身上归于那个渴望发现的“孩子”的份额、与归于那个渴望取代、渴望支配(乃至碾压)的“老板”的份额,而只从表面看某些观念、技术、工具在一部著作中所占的份额——即便如此,最近六个星期里的一项意外发现仍然触目惊心:我朋友的著作,从我们相识那年起就振翅起飞,而它直到今天,仍在由我传递给他的东西滋养着。十五年前退出数学舞台时,我曾想象:我传给我这位“非学生”的朋友的那“一点点”(尽管我分明看到它在一次惊人的最初冲刺中的作用),将是第一个跳板,一次起飞将把他带得远远越过起点,使他\textbf{远离}我的著作、远离我这个人。而实际发生的却是:我的朋友直到今天,仍系在那起点上,\textbf{系}在那部他既要否认、要么交付嘲弄或遗忘、又要“利用”的著作上。这是对一个与父亲或母亲的冲突性纽带的典型个案:那根纽带,无限期地把那个乐于在自身培养这份冲突、而不肯冲向世界的人,拴在他注定要离开、要超越的人的轨道上……

今天我看清:经由那份蓄意把我的年轻朋友当作“格格不入者”、而不简单当作一个有幸比其他学生更有才具的学生来对待的主意,经由那份在他面前低估或忘记我所传递之物的分量(以及由此放进他年轻双手的\textbf{权力}……)的主意——经由我的这些态度,我在不知不觉中喂养了他身上的自负与冲突;而这两者,都对我隐蔽。与此同时,我入了一场双人游戏——或者不如说,是一场配合默契的双人游戏,我实在说不好是谁“开的头”(假如这问题有意义的话):我以“谦虚”自命,说我的年轻朋友实在太耀眼,不该做任何人的学生,说我所能给他的一点点实在不值一提——而他,则(早在我离去之前)与我这个人、与我的著作划清界限,在我的\dmark{}\origpage{257}纵容目光之下,否认那分明滋养过他的土壤。

只有在写本注的这一刻,我才终于看清这场游戏;对它的散漫知觉,不过是一两个星期前才在场的。我也看清:我身上的这份“谦虚”或“谦卑”,是假谦虚、假谦卑——是一种不单纯:不肯简简单单地把事情看成它们本来的样子。在这场游戏里,有对我的年轻朋友的纵容——播下的种子滋蔓了百倍!——还有更隐微的,对我自己的纵容:为一桩“特权的关系”、一桩非凡之至的关系垒起一个高台\fnote{55}{(*)}{可与5月10日的“高升”一注(n\textsuperscript{o}~63$'$)对读:在那里,我头一次知觉到,在我与我的朋友皮埃尔的那段关系里,有纵容这一配料。这份知觉始终零散而残缺,直到今天,才在本注“格格不入者”所进行的反思中明确起来。}。(正如一切不单纯,说到底,多半是一种对自己的纵容……)

\subsection{绿灯}

\notehead{68}{}

(4月27日)老实说,我从不曾反思过SGA~5研讨班那离奇坎坷的命运背后的意味。它1965/66年的口头进行,并未引起什么特别的困难;而由一批批(常常掉链子的)志愿者所做的成文工作,却拖了\textbf{十一年}\fnote{56}{(**)}{以我为口头报告所作的详细笔记为底本,把整个研讨班写定,对我而言本不过是几个月的工作。}!直到1976年,Illusie才终于把事情接过来,着手写定尚付阙如的部分,并出版全帙。今天——在这个研讨班之后将近二十年——我才头一回明白“这里面有东西要理解”。也许只有我一个人……

我冒出的第一个念头是:在研讨班那些多少积极的听众中间——他们也多少熟悉此前的SGA~1至SGA~4——必定发生过一种\textbf{饱和}现象:对着那滚滚而来、如一场无可抗辩的怒潮般的“Grothendieck学问”\fnote{57}{(***)}{这与我那班“多少有点被震懵了”的学生的印象相合;那一印象,见“教学的失败(2)——或创造与自负”一注(n\textsuperscript{o}~44$'$)所引的信。}。显然,某些执笔者信仰不足,他们想必不太明白这一切要走向哪里,我为什么偏偏要那么固执,整整一年,要把平展上同调的那些\dmark{}\origpage{258}本质形式性质翻来覆去、直到完全驾驭,连同与之相系的那一整套新观念的武库。尤其能说明那种普遍冷淡的是:研讨班的收官报告——提出未决问题与猜想(据我所知从未发表)的那一场——和梳理各语境中Euler--Poincaré型公式与Lefschetz公式的开篇报告,竟都没有留下痕迹。我不记得当时曾知觉到这种冷淡(甚至后来也没有,直到今天\fnote{58}{(*)}{(5月26日)正是在我多少重新“泡进”SGA~5研讨班的语境之后,我才想起一种不安的印象:那是我翻阅刚收到的、已出版的研讨班文本时产生的(那应当是1977年,出版之年)。那份“残缺”的印象(当时只停留在弥散、未成形的形态)主要地——也许甚至完全地(我当时没有花多少工夫细看,尽管那很值得!)——是由于开篇与收官两份报告的缺席,而尤其(我相信)是由于宣告这一缺席时的那份漫不经心,仿佛这是理所当然——何必费那份力气把它们收进去呢!我在某个层面上必定“感到了什么”,但直到这个月(将近七年之后!)我才费心让它浮上来、加以检视——在“屠杀”一注,以及紧随其后的“遗骸……”“……与身体”两注里。}),因为我一头扎在当时的任务里。

SGA~5的命运——它原初拥有与我的任何其他研讨班同样强的\textbf{统一性},却在随后的十一年不成文的岁月里渐渐被\textbf{拆解}(注68$'$)——本可以让我看到:我以那样的执拗所推进的那些大计划,我曾有几年的时间找到人手辅佐,却根本没有成为一项共同的事业,而始终只属于我个人。我的纲领在这里或那里唤起了些应景的合作,却从不曾在我当时任何一个学生身上变成一种观念之力——一种会促使他去做比他跟我做学位论文时更长程、视野更开阔的工作的力量;那篇学位论文在他生命中的主要作用,是让他学会他选定的数学家这门手艺。

依我看,唯一一个从整体上把握了(且不说据为己有)某种总体视野的人——超越就某类问题、或为发展某类工具而进行的特定“合作”的框架——是Deligne。这想必就是为什么,我(无须把事情明言)在他身上看到的,与其说是一个“学生”,不如说是一位\textbf{天造地设的继承人}。“继承人”这个词,比先浮现在我眼前的“接续者”更好地表达我想说的——后者会暗示一部受制于所受遗产的作品。相反,我把这份“遗产”感受为一个单纯的\textbf{馈赠}:我有能力做出它\dmark{}\origpage{259},以助一种个人视野的铺展;那视野还将从许多别的馈赠中汲取养分(事实上,早在我离去之前就已经如此),并且注定将毫不费力地超越先行与滋养它的一切。

回到SGA~5的悲惨命运:昨天掠过我心头的念头是,这一命运与Deligne对我个人、对我的著作之关系的暧昧,恐怕不无联系——尤其考虑到他那强大的数学人格,必定对我的全体学生施加过一种支配力\fnote{59}{(*)}{(4月28日)这一支配力的一个具体而雄辩的征兆是:SGA~5的出版,直到Deligne认为合适、向Illusie示意让他积极张罗的那一刻,才终于完成——也就是说,恰好在他自己需要把它当作其SGA~4½“摘要”的底本、以便取而代之的那个\textbf{时刻}。(对此可参见Illusie执笔的SGA~5导言的结尾。)这一下照亮了、并使前天我在“白板”一注(n\textsuperscript{o}~67)里还称之为“神秘的”那句声明取得全部意味:“SGA~4½的存在,不久将使人得以照原样出版SGA~5。”“照原样”在这里是一记尖刻,我恐怕是唯一(从前天起)感觉到并赏识其分量的人!(相对于原初研讨班,已刊版本代表的那场“拆解”,人们有目共睹。)}。他必定在这份研讨班笔记的坎坷命运中,在其被剥除那构成口头研讨班之统一与冲劲的一切之后,在内心深处得到了满足。然而反思之后,有一点很清楚:这些坎坷的第一性的、本质的原因,\textbf{不在}与会者中任何一个人的性情里。这一原因虽然我尚未清晰地辨明,但毫无疑问,它首先关乎我自己的\textbf{这个人},\textbf{也}关乎那些在65/66年间假装承担研讨班成文工作的人们。它肯定就在他们对我的关系之中,或者也许,在他们对我为他们所化身的某种做数学的方式(或某个纲领,或某种看待事情的视野)的关系之中。SGA~5的命运,如今在我眼前,是一种雄辩而执拗的\textbf{显影剂},显影某种我迄今不曾费心检视的东西——只因为我甚至不曾察觉它;直到此刻,我也只是隐约瞥见\fnote{60}{(**)}{(5月26日)这正是前面倒数第二条脚注里所说的那个“某种东西”;它在过去这几个星期的反思中、尤其从那个时刻(5月12日)起,终于浮出了水面——那天,自它1977年问世以来,我头一回费心看一看,“一场辉煌的研讨班”落到我的上同调学家学生们手里、经他们十一年后做成屠杀版之后,变成了什么。}。也许,这几行文字会促使这场集体厄运的某位当事人,把他对这件事的切身印象告诉我。

\pmark \origpage{260}然而,也许有一课——至少暂时地——我现在就能从SGA~5这一情节里吸取;这一情节先预演、随后例证了我所投身的那个著名“纲领”在我离去之后几乎全线的那种壮观的\textbf{停摆}。与我在六十年代那些兴高采烈的年月里或多或少相信的相反(那时我正为自己终于找到好心人辅佐而喜不自胜!),今天在我看来:经由顽强而缜密的工作使一个宏阔的个人视野取得肉身,本性上不可能是一场\textbf{集体}的冒险或事业。或者不如说:倘若真有“集体事业”,那也不是那种围绕同一个人进行十年、二十年(甚至三十年)的工作。只要这一视野注定要成为人人共享的遗产,它就会仅仅在需要的压力下、在此处或彼处,经由这个或那个人日复一日的工作而取得肉身——那个人对这位先行者,也许只知道名字(充其量如此!);而那位先行者的视野太过宏阔,凭他一己之臂,无法使之成形\fnote{61}{(*)}{(4月28日)也许,“我的一己之臂”本足以实现我在六十年代末所设想的那项宏阔的工作纲领——条件是,在随后的二三十年里,我让自己成为这一纲领的独家仆人。我今天感到庆幸的是,我没有走那条路——它本可能是我的路,而我如今清楚地看到它的陷阱与危险。}

\subsection{颠倒}

\notehead{68$'$}{}

(4月28日)作为这场拆解的一个例子(此外还有许多\fnote{62}{(**)}{(5月28日)我只是到5月12日的反思里,才下决心盘点这场“拆解”——见那则名字更贴切的注,“屠杀”(n\textsuperscript{o}~87)。}),我重新想起了SGA~5一场关键报告的命运:它最终由Deligne本人(我相信他从1965年起就揽下了这件事,为的是十一年后才“履行”自己的承诺……)依据我的口头报告写成,然后不经任何手续,直接收进了SGA~4½!说的是:与正则概形上的代数闭链相联系的上同调类的形式主义——它经由过渡到所论闭链之支撑上的“具支撑”上同调,而顺畅地展开。几乎与平展上同调中的一切构造一样(它们在许多别的语境里同样有用,并且在那里已成惯例),这一构造是我在五十年代末在凝聚上同调的框架里发展的(这里是Hodge上同调与De~Rham上同调;在“抽象”代数几何的框架里,它们头一回得到研究,是在我最早的几场Bourbaki报告之一里)。它如此\dmark{}\origpage{261}自然,以至于显然蕴含与上积的惯常相容性\fnote{63}{(*)}{(5月28日)凝聚框架里,见我的Bourbaki报告n\textsuperscript{o}~49(1957年5月),§40(原书如此)。在5月8日的“好的参考文献”一注(n\textsuperscript{o}~82)里,我发现:这些想法,连同我在同一场SGA~5研讨班里为闭链所配的同调类而发展的那些想法(以及许多别的),已被 J.L.~Verdier据为己有,对SGA~5研讨班的存在、对我的这个人,都不置一词。这场战役发生在1976年,即“SGA~4½之役”前一年(在我看来,它与后者紧密相连),而且是在1965/66年母研讨班SGA~5的全体往届听众与与会者的眼皮底下、在他们全都知道的情况下进行的。}。

写这些文字时,我发觉这场戏法(把这场关键报告偷运进SGA~4½)竟成就了这样一个光辉结果:Deligne——他明明参加了65/66年的SGA~5研讨班\fnote{64}{(**)}{(5月28日)而且,正是在那里,他头一回听说了他在那部海盗文集SGA~4½里陈述得如此光辉的那些事情!对此可参见昨天的“格格不入者”一注(n\textsuperscript{o}~67$'$)。相对于他朋友Verdier前一年的手法、以及他自己在别的场合使用过的手法,我的朋友在这里倒是守在公然劫掠的界限之内:因为他把我列为那场闭链报告的作者(诚然,附带那项光辉的成果:把我列为他的“合作者”),而且尚未装作全然不知道:平展上同调理论、迹公式等等里面,有我的一份。至于在这一方向上的决定性进展,可参见“葬礼悼词(1)——或恭维”一注(n\textsuperscript{o}~104)。}——\textbf{不在}封面我的“合作者”之列(这件事昨天翻阅已出版的Lecture~Notes n\textsuperscript{o}~589时就已打动了我);而我,反倒有资格(在研讨班十一年之后)以“Deligne的合作者”的身份露面。不得不说,这是一次相当绝妙的\textbf{颠倒}!在SGA~4½出版之时——我就这样在不知情中与之“合作”——我停止一切公开的数学活动已经七年;以至于我甚至从不曾过问那个可怜的SGA~5的出版,它对我来说,属于一段已被我抛在身后的过去……

(4月30日)至于SGA~5,它如今显出这样一副面孔:一堆有点五花八门、没头没尾的文本(头尾都在半路丢失了!),只有参照SGA~4½的文本才“站得住”。有一件事引人注目,而我恰在此刻才注意到:SGA~4½这个名字本身,就分明暗示这份文本\textbf{先于SGA~5},而后者只有参照前者才存在\fnote{65}{(***)}{(5月28日)关于这场把SGA~4½“强行塞进”SGA~4与SGA~5这两个不可分割的部分——它们构成我已刊著作的心脏——之间的“暴行”的更深意味,见“遗骸……”一注(n\textsuperscript{o}~88)。}。倘若这份文本的作者\dmark{}\origpage{262}的心思不那么暧昧\fnote{66}{(*)}{(5月28日)“暧昧的心思”这个说法,在这里分明是一个委婉语!},而他出于情感的缘由,执意要把他的“摘要”(“外加若干新结果”)塞进他曾扮演过角色的SGA丛书,那么顺理成章的名字,当然是SGA~5~1/5。

我在这里看到第二场戏法;它使我掂量出:Deligne在SGA~5的命运中所占的份额,比我三天前所想的还要重。这也使我收回前一天表达的感想——SGA~4½并不属于一场诈骗式的战役。显然没有人(从Illusie开始——他的诚意当然无可怀疑\fnote{67}{(**)}{顺便也该趁此机会,感谢Luc~Illusie:他以那样的细心与忘我,把某些陷于困境的报告的成文工作、以及“全包”的出版,进行到底——而且是在肯定谈不上鼓舞人心的条件下;其中,我的全然缺席,肯定不是最轻的一项!(5月26日)在n\textsuperscript{o}~84至89各注、尤其“屠杀”一注所继续的后续反思之光里,我对Illusie的这些感谢,取得了巨大的、未曾预料的喜剧维度——写这些文字时,我远没有预感到!诚然,我是顶着我身上的一种抗拒写下它们的;那抗拒尤其表现为(已在计划中的)致谢在注的“正文”里被“遗忘”,以致我不得不靠一条脚注来“补上”。这份抗拒,无疑出于那种不安:从我头一回亲手捧起这部名叫SGA~5的文集起,我就感到了它(而且我相信,直到最近几个星期,我再也没有捧起过它);那份不安,我在前一条注“信号”(原书此处题名如此)的一条脚注(今日5月26日所写)里谈过。这份疏忽,恰好例证了在禅观中、对当下正发生于自身之事保持警醒的留意的重要性。缺乏这份警醒,这里的反思就停留在禅观之下,停留在一个表浅的层面——而留意到那份抗拒,本会引导我去探它的根,并由此也更近地看一看那场美好的研讨班变成了什么(这件事,我两周之后才做)。})——察觉这场“战役”,这无疑要归功于我已领教过的那份“支配力”,我想,还有我朋友这个人身上的魅力:二者都使他高悬于一切怀疑之上!

\subsection{化圆为方}

\notehead{69}{}

(4月27日)十一二岁的时候,我正被拘留在里欧克罗集中营(芒德附近);我发现了用圆规作图的游戏,尤其为六瓣的玫瑰花饰心醉神迷——把圆周用圆规的开度在圆周上截取六次,分成六等份,恰好落回起点。这一实验的发现使我确信:\dmark{}\origpage{263}圆周的长,恰好等于半径的\textbf{六}倍。后来(我想是在芒德的中学,我最终还是上了学),我在教科书里看到,这个关系据说要复杂得多:人们有 $\ell=2\pi R$,其中 $\pi=3.14\ldots$;我坚信书错了,书的作者们(以及自古代以来想必走在他们前面的所有人!)肯定从来不曾做过这个极简单的图——它明明白白显示,人们干脆就有 $\pi=3$。典型的是,我发觉自己的错误(把弧的长与连接其两端点的弦的长混为一谈)是在我把我的惊讶——惊讶于先行者们的无知——向另一个人倾诉的时候(一位被拘留者,玛丽亚,她志愿给我上过几节数学和法语的私人课),恰在我准备向她演示为什么应当有 $\ell=6R$ 的一刻。

一个孩子对自身理解力所怀有的这份信心——宁可信任自己的能力,也不把在学校学到的、书里读到的东西照单全收——是一件珍贵的东西。然而它不断遭到周围人的挫抑。许多人会在我这里讲述的经验里,看到孩童僭妄的例子:它终究得在习得的知识面前低头——事实终于戳穿了某种可笑。但一如我亲身经历的这一情节,其中丝毫没有挫折感、没有可笑之感,而分明是又一次发现(在我草率地把 $\pi=3$ 这一错误公式当作解释的那次发现之后):一次错误的发现,以及同时,人们应当有 $\pi>3$ 的发现——因为显然,弧的长\textbf{大于}连接其两端点的弦的长。这条不等式,恰好又向着被驳回的 $\pi=3.14\ldots$ 那个公式的方向;于是它显得讲得通了,同时我那时想必隐约察觉:也许有一些并非那么愚蠢的人,曾经钻研过这个问题。此刻我的好奇心已经满足;我不记得当时还想对这个数——它想必重要到人们为它单配了一个字母——的来龙去脉知道得更多\fnote{68}{(*)}{(4月28日)前面的追述,唤起了别的回忆;它们表明,那个大名鼎鼎的 $\pi$,比我起初记得的更让我着迷。在一本书里(也许就是同一本)找到的近似值$344/133$(原书如此)打动了我——它那么漂亮,我简直难以相信它只是近似的!那时我除了分数之外不知道别的数,我为表达 $\pi$ 的那个既约分数的分子与分母会是什么模样而着迷——它们必定是些很不寻常的数!不必说,在这些关于化圆为方的孩童反思里,我没有走多远。}。

\dmark{}\origpage{264}这一经验,无疑是头几课之一:当我自己的理解似乎与公认的知识相抵触时,要有某种审慎——这样的局面可能值得仔细检视。审慎是经验的果实;它贴合并补足(而不改变)对自身认识与发现能力的自发信赖,以及认识到我们本就具有这种能力所带来的笃定。

\subsection{葬礼}

\notehead{70}{}

(4月28日)昨晚重新想到SGA~4½封面上那一幕——我在不知情中,以我显赫前学生的“合作者”的身份露面——这事在我看来难以置信到这种地步,以致我疑心自己的记忆背叛了我,怀疑我其实分明被征求过意见、并且没多想就点了头。但这一假定,与我直到去年还抱持的态度背道而驰到如此程度——即我根本无意再发表数学(更不用说以某人的“合作者”的身份,而且是某位与我的关系当时在我看来已带着深深的暧昧的人)——以致它比它本该“解释”的那件事还要“难以置信”得多;而那件事,归根结底,对我没有任何神秘或不可解释之处!为了求个心安,我还是查了我朋友1976年至今的信件(为数不多,很快查完),当然没有找到任何关于SGA~4½出版的暗示。我还是给当事人本人写了几行,请他对这出我不太欣赏的“闹剧”给我些解释……\fnote{69}{(*)}{(5月26日)我的朋友到底赏了我一份回答,它驱散了最后一丝疑云。他让我以“合作者”的身份露面,确确实实是因为那场SGA~5报告——他把它写成并收进了SGA~4½;而他不认为有必要为这次移交、或为以“合作者”的身份露面征得我的同意,也不认为有必要给我寄一册这本我曾“鼎力合作”的文集,鉴于“我已经七年不做数学了”。(6月5日)我刚收到(迟做总比不做好!)Contou-Carrère的一封信(日期为5月30日),回复我4月14日的一封信——我在信中问他(为求心安):在我的书里,他是否曾见过一份SGA~4½的副本。看来确有这样一份副本,Contou-Carrère留在了自己手里(除非是他买过而忘了?)。另一方面,Deligne的回答似乎确实证实:他不认为有必要寄我一份:“给你寄一份4½的副本,本确实可能是个好主意;我无疑料想你那时不会看出它的趣味”(5月15日的信)。}

三天前,当我在反思里提到三年前我与我朋友皮埃尔的关系中发生的那个转折——我失去了\textbf{兴趣}\dmark{}\origpage{265},不愿再同他交流数学问题(见“两个转折”一注,即注66)——我想起了一份当时强烈在场的印象。要把它安放好,我得先说明:在流逝的那十年里,当我的朋友事实上扮演着我唯一的数学对话者的角色时,我曾指望(一如我指望他扮演那个角色一样想当然):他会成为我向他吐露的数学反思与\textbf{想法}的中继站,把它们转达给可能感兴趣的数学家们。如我在别处解释过的(见第50节“一段过去的重量”),正是拥有这样一位对话者-中继站的感觉,赋予我那些断断续续的数学活动时期以比解馋更深的意味:把它们联系到一项超越我个人的共同冒险。无疑也是这份感觉,使得在那么长的时间里,我从未感到丝毫发表我所发现的东西的欲望,更从未感到丝毫对自己退出数学舞台的懊悔。(这样的懊悔,说到底,从未出现过;我是毫无蓄意地、“重新出现”在那个“舞台”上的,甚至在我自己察觉之前!)

至于我的朋友在多大程度上回应了这一期待,我说不好——很有可能,只要他对我还保持着那份数学上的热忱——它由好奇与亲切的同情共同推动,正是它使他在我与数学家世界的关系里(也在某种意义上,在我与数学本身的关系里)所扮演的那个特殊角色成为可能、并显得自然而然——他就一直扮演着那个角色。一两天前,当我向自己提出前面那个问题时,我收到了(仿佛作为部分的即答!)Larry~Breen的一封信,寄来1974与1975年间几宗通信的抄件,其中有Deligne1974年写给他的一张两行短笺,附着我刚写给Deligne的一封关于Picard叠形式主义的信的抄件,征询Breen对我那封信的意见。他在其中以“大师”一词指涉我;我从中听出一种半开玩笑、半亲切的声调。我不记得还有别的场合,有人把我1970年离去之后向我朋友吐露的事情,经他人之口回传给我。很可能有过,而我忘了;更不必说,即便在我数学活动的那些片段里,需要请教朋友的时候也相对少见;而且直到1977或1978年,我向他吐露的那些反思,分量都有限。所以,\dmark{}\origpage{266}严格说来,在那一刻之前,并没有多少东西可供“中继”\fnote{70}{(*)}{我可以算作例外的,是我关于分层结构的剖视理论的最初反思——七十年代初,我想必向Deligne提起过。他以一种纵容的同情接待了我在这方面的期望,有点像人们对一个什么都不怀疑的大孩子的态度。(他在与我的关系里常常是这种心绪,而且肯定常常有理!)我朋友的怀疑——它由他所熟知的某些野性现象的知识所推动,而那些现象是我所不知道的——当时并没有说服我;毋宁说,他向我指出的那些事实,使我从那时起就疑心:人们为“做拓扑”而通行的“拓扑空间”语境,不足以柔顺地表达某些我觉得本质的拓扑直觉,比如“管状邻域”。在其后十年里,我几乎没有机会回到这些反思,我的那些“疑点”被我淡忘了些许;它们经由我1981年12月--1982年1月的反思重新变得现实(并成为亲身的信念)——那场反思,是被为“Teichmüller塔”建立一套“剖视”理论的需要所激发的。(对此可参照《纲领概要》第5、6节。)(6月5日)作为另一个例外,我可以算上我关于虚拟相对概形与虚拟Motif(在一般基底概形之上)的反思——我记得曾向Deligne提起过。由于这些东西与他已决意埋葬的一部瑜伽(直到1982年的开棺出土)紧密相连,他没有表现出对我向他解释的那些想法有任何共鸣,这并不奇怪——那些想法当然使我心醉神迷;关于它们的几点指示,见注n\textsuperscript{o}~46$_{9}$。}。

事情在1977年起变化:那一年,自六十年代以来头一回,我强烈地“咬住”了一种异常丰饶的实体。那是我关于地图的反思的开端,以及(大约同时)顺藤摸瓜地,关于正多面体的一条新进路的开端(见《纲领概要》第3、4节)。也是从那一刻起,对我很清楚:我刚刚把手指按上去的那些事实,开启了无法预料的视野,其广度与深度,可与Motif观念诞生时所隐约瞥见的(以及后来远不止隐约瞥见的)相比。

奇怪的是,那一回,我竟然还是带着他将成为回声的期待去找我的朋友——尽管“Motif”这个名字周围已沉默了七八年,那沉默本已足够雄辩,足以教会我:我的期待是幻象!这份惊人的不察,恰好例证了我心中那条蓄意的主意(甚至在一两年前发现禅观之后):对我与数学、与数学家们的关系,一概不予留意——它们本该属于一段遥远而彻底翻篇的过去!我头一条向着这个\dmark{}\origpage{267}方向的反思\fnote{71}{(*)}{关于这场反思,见“扫兴的老板——或高压锅”(第43节)。}恰恰发生在1981年——那是我与我朋友关系中的第二个“转折”之年,我已有机会谈到它。但即便在那场延续了几个月的禅观里,与其他数学家的关系也只被轻轻拂过;而与其中那个无疑最亲近的人的关系(至少在我们共同热情的层面上)——据我所记,连拂都没有拂过。然而那本会非常有用!

但带着事后的认识、经由此刻的反思,有一点变得清楚:那一刻所发生的事——它曾那样使我惊讶、使我失望(一次隐约的轻蔑突然现身,而我本以为可以分享一次给了我深刻印象的发现那犹带体温的喜悦)——正是注定要发生的事。恰恰是我要传达的东西的\textbf{分量}——它曾激起我对一种与我的兴趣合调的兴趣的期待——应当在我朋友身上,头一回激发出\textbf{使我气馁}的反射。这一反射理应更加强烈,因为从SGA~4½出版那一刻起,我就已经被“预葬”了。三年后我再次进攻时——那时我的朋友(手持他关于绝对Hodge闭链的漂亮定理)正准备把下葬办得像模像样,随即有次年问世的“值得纪念的文集”\fnote{72}{(**)}{即Lecture~Notes~900那一卷;见“一个梦的回忆——或Motif的诞生”一注(n\textsuperscript{o}~51)。}——同一个反射又起了作用,只是粗暴得多。(这一情节当时断绝了我们数学层面的交流,但并没有“使我气馁”……)

在两种情形里,那份漠然显然是真诚的,一如它在别的情形里、对着我之外的人表达时那样。这已不是我头一回在他身上(或在别人身上)看到:与求知的渴望异质的力量,把那渴望中和掉,并取代数学家的嗅觉。

正是在这两次机缘——1978年,然后1981年——我头一回如电光石火般瞥见了我朋友身上那桩矛盾的“\textbf{代价}”;那桩矛盾,多年以来我就知道,但作为他著作中、他对数学事物理解中的一种阻碍与限制的分量,此前从不曾向我清楚显形。但只是在这一我已持续一个月的禅观——关于自我离去以来阴险地进行着的某场\textbf{下葬}——之中,这一\dmark{}\origpage{268}分量才终于渐渐显于全光之下。

在显见的层面上,这些天与这些星期里我所发现的下葬——几年前我已有预感,却从未想到把特定角色派给任何人——首先是对\textbf{我的数学著作}的下葬,经由它、并且首先,是对\textbf{我的这个人}的下葬。当然,在所有人中最有资格对这场下葬伸出援手(许多人在内心深处都在呼唤它)、并主持这场无名葬礼的,是那位不久前还在众人眼中俨然合法继承人的朋友。倘若他主持了葬礼,他肯定不是独自参加了葬礼!但更深一层:我的朋友在这些漫长的十二年间所悄悄下葬的,不是别人,正是\textbf{他自己};或者说,是他身上那样一种东西——那件不打动任何人、精微而难以捉摸、如一朵花或一枚果实的香气、无价的东西。($\Rightarrow$注71)

\subsection{墓}

\notehead{71}{}

但顺着联想的线索,我已离开本题——本题是要唤起某一记“强印象”;三天来,它的回忆执拗地回到我这里。这一印象,定位于我与朋友关系中那个“转折”的时刻:当我面对一种蓄意鄙夷姿态的征兆(既含蓄又粗暴地显明)——那些征兆使我在数学层面上结束了我们的关系。我当时懂得:时刻已到,延续这样一种关系,我再也无所期待;而“决定”自己做出了,没有分裂,没有懊悔,作为这份迟到的(而且非常不完全的)理解的第一个果实。

我心中没有愤怒,更没有苦涩。(我记得在我们交往期间,从未对我的朋友动过怒,也从未有过苦涩——除了我离开IHES的那桩情节的时刻;而且当时被包括在内的,也不止他一个。)但有一种悲伤:在同一个继续为我所珍的人的关系里,翻过这一页——尽管曾经把我和他系在一起的最强的纽带,已经干枯、已经死去。而作为此后几年里仍在的一根刺,还留着那份未被消化的挫折:我曾带来想要分享的喜悦,带给那个在我看来最亲近、最有资格分享它的人,却撞在一份自满的紧闭之门上。这份挫折,我看,最终是被我此刻继续着的禅观消解的。就在今天,它还在向我显示:发生在我身上的,正是注定要发生的;\dmark{}\origpage{269}而这份挫折的第一责任人,不是别人,正是我自己——是我自作主张,耽溺于某一实在的幻象图景,而不肯使用我健全的能力、用醒着的眼睛去看这实在。

正是在这份悲伤的底色上,也在这份期待的挫折的底色上,出现了那奇怪的印象——它当时不是作为一场反思(当时并没有发生)的果实或终点而来,而是作为一份直接的、无可置疑的直觉。那就是:我能在数学层面上对我朋友说的一切,以及多年来我已对他说过的一切,我是托付给了一座\textbf{墓}——或者已经托付给了一座墓。虽然我从不曾对任何人谈起这份印象,也从未在后来的反思里把它白纸黑字记下,我却清楚地记得:当时在场的,正是那座\textbf{墓}的意象,以及表达它的那个词(用法语),也就是我刚刚写下的那个词。这份“印象”或意象,想必在那一刻,作为某种理解的(不妨这样说的)视觉表达而涌现——那种理解,在某个层面上,想必早已成形、早已在场,是长年累月间必定发生过的一整组知觉的果实,尽管注意力不曾留住它们,记忆不曾登记它们;一些想必十分简单、十分分明的知觉,但我没有“留住”它们,因为它们对于我身上那位常常有权随意拣择的人……显得不合意。无论在当时还是后来,这份断然的意象,都未曾与任何具体的、可触摸的“事件”的记忆相连——一件朝着这意象方向去的事件,本可能在我心中唤起它;对这一骤然浮现的意象的记忆,此后想必也很少掠过我;今天是我头一回稍稍停留在它上面。

倘若当时没有浮现任何记忆或联想,那肯定是因为我连接待它们的最低限度的余裕都没有。奇怪的是,那时(如果我把时刻定位得准\fnote{73}{(*)}{(6月11日)一些旁证向我确认:确实如此。这个“第二个转折”,位于1981年的下半年。})我正投入一场关于我与数学之关系的禅观,而这一情节——它毕竟相当有力地,经由一个现在,向我诉说某段过去——竟没有使我想起中断反思的“线头”,好把方才发生之事的来龙去脉纳入反思;那件事,在我的一生里并非无关紧要。

\dmark{}\origpage{270}此刻(在唤起这一意象、并说出它一出现就与一切记忆或联想脱节……之后)浮现的第一个(老实说,也是唯一一个)联想,是我的Motif之“梦”所遭的命运——在数学诸视野之中最为我所珍的那一个。倘若那段过去,也许仍对我保有一种隐秘的支配,那正是经由这个梦——而这份隐秘的支配(我相信我在写下这些文字的时刻隐约瞥见它)本身,有着超越言词的、梦的力量。倘若在过往投资的遗产里——在一场对数学的满怀热情的投资的遗产里——曾在流逝的十年间出现一种未说出口而深刻的挫折,那正是看到:一片死一般的沉默,围绕着我视为活的东西,而我曾把它们当作活生生、精力充沛、只待冲向天光的东西,托付给我的朋友!我一走,能够而且有使命看护这次绽放的,是他,唯有他;把唯有他(与我一起)才亲密感到的东西,交到人人手中。而我虽然从不曾以这样或那样的字眼向自己说清这件事——甚至从不曾(据我所记)停下哪怕一念的工夫,想一想我留下之物的命运——但随着岁月流逝,在我身上的某个地方,我必定已经懂得:这个始终为我所珍的梦,我是把它托付给了一座“墓”。

于是,随着这一唤起、以及它在我身上激起的第一个联想,我看见一串别的联想循着它的尾流涌来,向我揭示:我分明刚刚触到一个命门——那也许正是头号要紧之处,我的数学家过去的(久被忽视的)重量由此施加。

但此刻,我想,还不是循着这些联想走的时候——反思的这一“最终”阶段已经拖得很长了。关于我的朋友皮埃尔,关于Motif,我想在这场反思里已经说得够多——肯定,对许多人的口味来说,已经太多!我相信,就这些注而言,是时候以一场对这场双重下葬之反思眼下教给我的东西的\textbf{盘点},来收束它们了。

\section{VI.~事物的回归——或全体一致的赞同}

\subsection{一脚踏进旋转木马}

\notehead{72}{}

(4月29日)\dotfill\fnote{74}{(*)}{我认为最好在这里免掉读者整整一页关于禅观的一般考虑——那是一种顾左右而言他的方式,是不肯进入正题的抵抗的征兆。}

\pmark \origpage{271}我觉得,就那桩占据着我的事情而言,应当做的描述与澄清的工作,就其“局部的图景”而言,主体已经完成。(况且显而易见:眼下这些注定要发表的注,只是实际工作的一幅缩影——在这里把汇成这一幅或那一幅“局部图景”的所有要素一一摆明,根本无从谈起。)当然,经由这同一项工作,某种总体图景也必定已经成形——还模糊,它等待着被表述,以便取得形态与生命,并向我说出它要对我说的话。自我昨天的反思以来,我感到它已蓄势待绽,推着我借给它一个声音。

老实说,昨天的反思(我刚刚重读)所主要教给我的,\textbf{不关乎任何别人,只关乎我自己}。我带着某种如释重负,看到反思回到对我自己的反思这一坚实土地——因为一个星期以来,它屡屡给我这样的感觉:牵连别人的这个人,多于牵连我自己。昨天的反思终于向我揭示了一件分明显而易见的事:即我对一段过去——我的“数学家的过去”——的执着之力,以及那部大名鼎鼎的Motif之“梦”在其中所扮演的特殊角色。

话一旦终于说出口,它的显而易见便跃然眼前——最新近、也许最清晰的征兆,便是那阵情绪:两年后才发现某桩“事件”之时所引爆的情绪,即Motif在我的前“学生”兼朋友的引导下,那场“悄悄的”(而且迟到的)“重返”数学兽栏!这阵情绪立刻表现为一场自以为已告终结的反思的重新开工——重新开工立刻落实为五十页回顾反思的洪流!于是(而这一确认,在这场不合时宜的重启中已数次来到我面前)似乎我还没有像一两个月前、在阶段性终点与一阵(绝非幻象的)解脱感的高涨中所以为的那样,“下了旋转木马”\dmark{}\origpage{272}——带着这样的教诲:“我并不比别人高明”,“学生超过老师,我无须惊讶”\fnote{75}{(*)}{见“打住这木马!”一节,n\textsuperscript{o}~41。}。然而这教诲并未阻止我惊讶——只要“学生”在一个我全然未曾预料的方向上超过我,就足够了!不过,倘若这教诲不曾阻止“我惊讶”,它在流逝的反思中却不止一次对我弥足珍贵,使我避开那些惯常的陷阱(或至少,避开这些陷阱中的\textbf{某一些})。

回到这份“支配”之力、这份我对Motif之梦的执着之力:这份执着在本书的许多别处已经显现——既在《收获与播种》里(其中Motif被不止一次谈起,而且用词已然雄辩之至),也在《纲领概要》里(在那里,“客观地”说本无须谈到Motif),或在《主题概要》里(在那里,Motif在一群活力充沛的雏鸡中间,有点像没有孵过的蛋)。在后一部文本里——它上溯到十二年之前,显然是怀着疏离的心境写就的——关于Motif的最后一段,据我看,是唯一能让人忽然感到一股暖意流过的段落……

引人注目的是:自我离去以来的这十四年里,这份执着从不曾向我显现;直到昨天,我才终于隐约瞥见这一显然的事实,直到今天才终于把它表述出来。在将近三年前的那场禅观里(1981年7月至12月),我终于确认了第一个显然的事实:即我身上对数学的那份热情的持存——它在流逝的岁月里,曾以十分雄辩的方式表达过自己。但我对一段过去的执着,据我所记,当时未被察觉,直到今天仍是如此。

然而,我必定是随着“一段过去的重量”那场反思才开始隐约瞥见它的——那场反思,是在对数学家过去的禅观似乎已经收官之时、出于求个心安而到来的(只除了一点:我还不曾知觉到这段过去的\textbf{重量}!)。而且我写它的时候分明感到,自己还停留在事情的表面,并未真正穿透。我随后不得不补加的那些注(先是注46、47)在当时把我引向一个方向——有好一阵,它使我离开我自己:先是把注意力系于一部数学著作(以及在我看来它最“重要”的那些方面),继而是那部著作的坎坷、他人在其中的角色,而不是我自己。

\dmark{}\origpage{273}我刚刚重读了“一段过去的重量”那场反思(第50节)。临近其结尾,我确实开始隐约瞥见:那股“翻转之力”(向着一种不止是偶尔的数学投入)可能是“对一段过去的执着”——对数学家的过去——所使;但更是对“最近十年的过去”、对“1970年以后”的过去,而不是对那些已白纸黑字写下的事物的过去、做成了的事物的过去、1970年以前的过去。几行之后,我却回忆起——但只是“顺带地”——在“我那时眼前的宏阔纲领……只有一小部分得到实现”。写下这几行时,我所想的,主要是“宏阔纲领”中那些直接可实现的部分;而它们的动机之力(!)却远不及“Motif之梦”所代表的那种。(它的论证——但绝不是它的表述——在当时,是“地平线上”的大任务之一……)

显然,我对“Motif之梦”的执着(一如想必一切执着),首先(即便不是完全)具有自我本位的性质。那是不止\textbf{贡献}于一项共同事业的欲望,也是看到这一贡献\textbf{得到承认}的欲望。假定“Motif的宏阔图景”确实已按我自六十年代末所见的那般幅度被勾勒出来,但我在这一视野的绽放中的那一份却被隐瞒——我的不快,想必不会小于(甚至可能大于?)我得知那部“值得纪念的文集”时所感到的(在其中,我分明看到被我提炼并带到光天化日之下的某些观念与想法被重新拾起,却(至少我是这样感到的)被剥去了那曾令我深深心醉的呼吸与炽烈的生命)\fnote{76}{(*)}{(6月14日)这份“不快”,在我看来,首先出于那份无耻之感:一种对某条纽带的蓄意蔑视,装作无视它、把它看作无足轻重。当一个人发现的观念或结果被他人重新发现时——这是常有的事——情形全然不同。}。只要那要看“得到承认”的自我本位欲望——它对着我的数学过去的这些或那些或远或近的事物——尚未烧尽,我自称“下了旋转木马”就还为时过早。数学的“木马”已不再\textbf{包容}我,一如它当年包容我、一如它包容我的某些朋友。但我肯定还在上面留着一只脚;而且我疑心,只要我还掺和做数学,这只脚就会留在上面!

\subsection{事物的回归(或一脚踩进菜盘)}

\notehead{73}{}

(4月30日)我方才重新想起了SGA~5研讨班的命运,以及这一命运与SGA~4½的出版相系的方式。一个曾经混乱的局面——我只是最近这些天、并以\dmark{}\origpage{274}顺带的一瞥检视过它——如今在我眼前显得非常清楚。我刚就此给三天前的反思补加了一条脚注\fnote{77}{(*)}{这条长得犯规的脚注,变成了一则独立的注:“颠倒”(n\textsuperscript{o}~68$'$)。}(见“信号”一注,即注68(原书此处题名如此)),而且我觉得,加上我前天已有的评注(同样以脚注的形式)与此前一天的反思(“白板”,注67),我已表述得足够清楚,无须再为如今已显出足够雄辩的局面再做一幅总结的全景\fnote{78}{(**)}{然而5月9日及其后几天,我还要回到它;见注n\textsuperscript{o}~84--89。}。

行文至此,有一点很重要,必须确认:SGA~5所遭“悲惨命运”的第一责任人、那份“被弃局面”之利用的第一责任人,不是别人,正是我自己。倘若各位“志愿者”(他们揽下了自己并不真想做的成文工作)显然对自己心里没数,那我也同样没有看清自己:我执意不听从一个已然雄辩的局面给出的教训,把事情托付给三心二意的“合作者”,而不肯亲自接手、亲自去做那从那时起就落在我头上的成文工作。毕竟,从口头研讨班结束,到我离开数学世界的那一刻,整整三年过去了(而后者立刻表现为:在其后的十四年里,我对我已刊著作几乎全然的漠然)。诚然,那三年里我完全被别的任务占据:SGA研讨班的延续(SGA~6与SGA~7),EGA的写作,对日复一日涌现的、常常多汁的问题的反思——而在这些问题之中,Motif总体视野的渐渐成熟……被这些任务缠住,我选择了对一场已成过去的研讨班的命运闭上眼睛——那场研讨班(连同前一年的SGA~4),是我就已完成的工作而言所能做出的最深刻的数学贡献,而且多半也是影响最广的一份。

在我一去不返地离开之后,局面只可能更加恶化,使我那些前学生中最显赫的一位得逞:他把那部大名鼎鼎的SGA~4½,插进SGA~4与SGA~5那non-sense与多余细节的矿渣之间,上演了那出绝妙的战役;并让我荣任那份以中心关键文本自居之物的“合作者”——那文本的使命(如他以那种使他生辉的天真所说)是慈悲地让人“忘掉”包裹着它的沉重矿渣……

\dmark{}\origpage{275}总之,我早在离去之前、并经由离去本身所做的种种选择,蕴含着对我已刊著作(或就SGA~5而言,待刊著作)命运的影响,同样也影响我“著作”中仍停留在梦的形态的那一部分——而且是\textbf{未发表}的梦。我不后悔我的选择;当我今天确认这些选择的某些不合我意的后果时,也轮不到我抱怨!倒是轮到我去检视这些后果(而且越不合我意,越要检视!),对事实\fnote{79}{(*)}{(5月28日)这里应读作“我所知道的那些事实”。从后天起,一些全然意外的新事实,将使关于葬礼的反思重新启动,并使我把它相关的注的篇幅增至三倍。}做一幅总体图景(这件事已经做成),并从中汲取它们能给我的教益。这是留给我要做的;今天的反思,至少,也许是在这个方向上的第一步。最近这几天里,我身上已做成了某些比照,我想首先把它们白纸黑字写下来。

在六十年代初的那段时期,推动我一般地在我的学生们身上投资的主动力,是找到“\textbf{人手}”来实现我的本能指认为紧迫而重要的“\textbf{任务}”(至少在我所特有的数学眼光里)的欲望。这份“重要性”肯定不是纯主观的,不是一个“趣味与颜色”的问题;而且常常(我相信),把我的某项任务当作己任的学生,分明感到它“够分量”,也许还感到它在更宏阔的蓝图中的位置。

然而,就这份推动力、这份催促我实现任务的内在动力而言,所关的并不是某种“客观的”重要性——Fermat猜想、Riemann假设或Poincaré猜想的“重要性”,曾使我全然无感,我并不真正“感到”它们。把这些任务与一切其他任务区分开来的,在我与它们的关系里,是它们是\textbf{我的}任务;是我感到了它们,并把它们变成我自己的。我很清楚:感到它们,曾是一项精微而深刻的工作、一项创造性的工作的终点,它曾使人得以把握作为某项任务之对象的关键观念与问题。它们曾经是——而且无疑(在很大程度上)至今仍是我这个人的一部分。把我和它们系在一起(或今天仍系在一起)的那条纽带,在我把某项任务托付给一个学生时,绝没有被剪断——恰恰相反,这条纽带获得了新的生命、新的活力!这条纽带无须说出\dmark{}\origpage{276}(而我在这里把它“说出”,哪怕只对自己,也是头一回)。这条纽带,对那个选择了与我一道工作、并就其自选的某项任务而工作的学生来说,是明摆着的;对我,以及(我确信)对任何别人,也一样。这是构想了一件事的人与这件事之间深深的纽带——它不被后来那些也把这件事变成“他们的”、并向它贡献出自己最好部分的人所削弱,而(依我看)恰恰被加强。这条纽带,我从不曾细心检视过。它在我看来,深深扎根于“自我”的本性之中,具有普遍的性质。人们有时装作无视这条纽带,仿佛自己高踞于这类小家子气之上——甚至可能,我也曾一度入过这种装模作样\fnote{80}{(*)}{可以肯定的是:我遵循着那种“得体的做派”——即对这类事情视而不见,因为它们有悖于时兴的形象!(5月30日)关于这条纽带,见“……与身体”一注,n\textsuperscript{o}~89。}。但最近这些年里(或最近这些天与这些星期里),那寥寥几次:当他人摆出无视这条纽带——他明明知道这条纽带——的态度,而这条纽带把我系于某项已由他人(或由我自己)完成或仅仅指定的任务时,我被触到了痛处。人们可以把这个痛处叫作“虚荣”或“自负”,再给它扣上别的字眼——我并不声称这些字眼在这里不合用;但无论人们给它什么名字,我谈论它毫无羞耻,也不因自己如今的样子而羞耻;而且我知道,我所谈的这件事,是天底下最普遍的事!当然,一个人对其“作品”的这份执着,其力量并不千人一律。在我的生活里——“做”自童年以来就是我几番重大精力投资的恒常焦点——这条纽带曾是强的,今天依然是强的。

于是我可以这样说:推动我与学生之间关系的主动力,是我把他们看作可欢迎的“人手”,用以实现“我的”任务。这一表述可以显得犬儒,而它只是表达一个明摆着的实在——我的学生们想必同我一样分明感到了它。这些任务是“我的”,丝毫不妨碍他们也把它们变成“他们的”——正是他们在自己身上对任务的这种认同,动员了完成任务所需的能量;一如对这个任务的认同,曾在我身上动员了使它诞生并成形的能量,并且继续动员我持续投注于这一主题的能量。这份能量是不可或缺的,否则我甚至无法作为“老师”而“运转”——也就是作为把一门手艺(它同时也是一门艺术)传授下去的长者;而这件事,离开可观的能量的动员,无从进行。在我做教师的全部过去里,我从不曾感到矛盾\dmark{}\origpage{277}:同一个任务,对我这里工作的学生而言深刻地是“他的”,同时又深刻地保持为“我的”。我不认为这一局面带丝毫冲突的性质,也不认为它曾给冲突的意图提供过攀附的机会\fnote{81}{(*)}{倘若,受到某种气候的怂恿,我的一个学生曾想在某项与我合作的工作里藏掖我曾有过的角色,那么事情发生在一个他早已不在学生之位的时候。}。在这种对同一任务同时投资、并认同于它的局面里,学生和我自己(依我看)都各得其所,处于一种全然清晰的工作关系之中;而就其自身而言(依我看依然如此),它不含有任何冲突性的成分。但在纯属个人的层面上,这份关系却停留在表面——这丝毫不妨碍它亲切,甚至友好,有时甚至充满感情。

对我的任务的投资,并\textbf{经由它们}对我的学生们-合作者们的投资,是(我说过)自我本位性质的(想必一如一切投资)。这些任务的实现,对“自我”来说,当然首先是一种自我扩张的手段:经由实现一项宏阔规模的总工程——“我的一己之臂”本无法完成到底的工程。从我数学生涯的某个时刻起,就有这份恒常的暧昧:一种同居、一种紧密的相互渗透,一边是“\textbf{孩子}”——他的求知与发现的渴望,他对隐约瞥见之事、对细察之事的惊叹;另一边是\textbf{自我}、那个“\textbf{老板}”——为自己的作品而欣喜,渴望扩张,渴望经由作品的倍增、或经由一项宏大规模的总营造的执拗而不懈的追求,来增添自己的荣耀!在这份暧昧里,我看到一种分裂,它继续压在我的生活上,给它打上深深的印记——一种也许将伴我终生的分裂。这样的分裂,当然不为我这个人所独有;但也许,在我那被“最好”与“最坏”一并充满的生活里,这份分裂采取了比在别人身上更极端的形式。

于是我可以这样说:对那个贪得无厌、渴望扩张的“自我”而言(它在场上的并非孤身,但确实在!),我的学生们首先是可欢迎的“合作者”,甚至不妨说“工具”——可欢迎的“人手”,用以营造一座会诉说“我的”荣耀的巍峨工程!\fnote{82}{(**)}{写下这句话时,我略有踌躇,字斟句酌,明知有人会把它夺走,当作那个可怕的学阀终于撕下面具的犬儒自白!但我也知道:我拦不住那个一心想淹死一条碍事的鱼的人,他自会从容行事。这不会阻止我继续我的话题:发现并说出明摆着的事情——包括上面写下的那份寒微的真相;它只会让那种从不曾费心朝自己内部看一眼的人感到意外。}

在我看来,这一点已经相当清楚地\dmark{}\origpage{278}在三年前那场关于我与数学之关系的禅观里(进而,关于一般的“做”)显形过,尽管后来我也许把它淡忘了些许。正是这件事,在这最后几天里在我的思想中在场,使我得以同另一件值得注意的事实比照:恰恰是我那时的一位学生(带引号的学生,姑且如此!);而且在所有人中与我最亲近的、也是唯一能不费力气地、就其整体地“感到”我身上那些宏阔蓝图的人——那些蓝图仿佛不容分说地催促我去实现它们——恰恰是他,在我离去之后(而且在内心深处,想必更早……)经年累月地实施了那场\textbf{葬礼}(Enterrement)——规模与著作(Œuvre)本身相当的葬礼(这两个法文词在这里用大写并不过分!),并且最终“主持了葬仪”(Obsèques,再多一个大写,好压得住秤!)。

在这份局面里,打动人的,是那种乌布王式的、巨大的、不可抗拒的\textbf{喜剧性}!最近几天里,我想必隐约感到了这份喜剧;但它只是在刚刚,在我为我自己那场庄严的葬礼放上最后一个大写的瞬间,才向我显出它的真面目——在一阵突然而不可抗拒的大笑里!恰恰是这份\textbf{笑},在反思的这个号称“最终”的阶段里,直到此刻一直缺席;在这里,主导的音色,宁可说是那位“体面先生”的愁容——他的正当期望落了空(甚至可以说被卑劣地欺骗),当愁容没有让位于那些抖得干净利落的嘲讽评点的时候(会说话的人自然会说话!)。我分明感到,经过这场长时间的、忧伤音调的枝蔓(这个词让我想起了什么……),我重新走上了正路。

而此刻,一个顺理成章的名字浮现出来,用以命名这条“注”(人们已不再清楚一条注是关于什么的注,但无所谓),是收束它的时候了。它将是:“\textbf{事物的回归}”。($\Rightarrow$注74)

\subsection{全体一致的赞同}

\notehead{74}{}

我终于感到——呼!——我触到了这一“最终阶段”的终点;它绵延了十二天,其中每一天都(一如从前)以“最后一天”的面目出现。也许,终场的话,就在几分钟前已经说出。我的(象征性的)下葬,是一场\textbf{事物的回归},一场由我亲手播种的收获。(而我那有血有肉\dmark{}\origpage{279}的下葬——倘若我有这份幸福死去,身后留下能为我送葬的活着的男男女女——也将是一次回归:回到我出生时就离开的某物之中……\fnote{83}{(*)}{(5月28日)这一与我自己的死亡的突然联想,带着力量来到了。我一度想把它推开,随后又想删去这突如其来的括号——它仿佛饭里的头发,来得不是地方。我克制住了,出于某种敬意。奇怪的是,第二天我得知:就在我继续反思的那个4月30日晚上,在我居住的市镇,一位朋友的重病缠身的姊妹去世了。我头一回见到德尼丝,就在当天,在她的灵床上。次日5月2日,我与我的朋友、与许多别的活着的男男女女一起,在一个明媚的春日里,把她送入土中……})。我觉得,还可能要补加的,多半只会是\textbf{尾声}方面的了。

那位大名鼎鼎的“最亲爱的学生”,并不是我亲爱的学生们中间唯一一个兴冲冲为我送葬的人;而那些分明伸过手的人们,在他们中间,恐怕也不是唯一出席葬礼而不以为苦的!但说到底,谁是张三谁是李四,于我何干!(关于这件事知道得更多——如果只是为此——不会教给我任何新东西。)我终于懂得了这场“事物的回归”;懂了,我也就领受了它的益处。

然而,我还没有取尽这份益处为我储备的全部实质。我还不清楚地辨明:究竟是我身上的\textbf{哪样东西},使得某些前学生觉得,下葬与葬礼于他们有利可图。难道仅仅是我谈过的那份“贪求”?它(依我看)与别的“老板”们的贪求相去不远;而当年他们与我初试身手时,曾不费力地与它相安(而且无疑甚至不曾察觉,至少在意识的层面上)?那么,是“机缘”(我的离去等等)“造就了贼”,并且\textbf{显影}了他们身上——一如在“众学生之首”身上——的一种\textbf{普遍的倾向}:一旦情势有利,就埋葬自己的“老师”或“父亲”?也许,我的“老师”(或“父亲”……)形象还超出了通常的尺度,而这一情势,曾有助于以完美的齐奏引爆这套“下葬综合征”?此刻我看不清!也许,我将收集到的回声(但愿如此)会使我看得更清,更好地消化此刻摆在我面前桌上的这份意外的食粮。

参加下葬与葬礼的,并非只有学生们,尽管据我所知,没有一个非前学生在其中扮演过突出的角色。显然,我的许多旧朋友从中各得其所。这回,这件事在我看来并不太神秘。

\dmark{}\origpage{280}如我已有机会顺带说过的,我不止一次确认:我不合时宜地退出数学舞台,在我昔日的朋友们身上造成了深深的困扰。那是这样一类东西所激起的困扰:人们在其中隐约感到一种对深刻的重新审视、对革新的\textbf{挑衅}。在这个案子里,自然,数学家中间的这份困扰在我的朋友们身上最烈——在那些认识我、并能感到我曾对那些至今仍是他们的价值投入过多大力量的人中间;更不必说,这些朋友中的每一位,自己也曾在这些价值中、并在这些价值所提供的丰厚“回报”中,做出并且继续做出一份力量相当的投资。早在“生存与生活”(Survivre et Vivre)时期一开始,我就已有充分的机会在别的科学工作者中间观察到这样的困扰。但这从未妨碍:每一次,当我在某位昔日朋友的身上——同情的纽带依然把我和他系在一起——确认出疏远的、有时是敌意的毫不含糊的征兆时,我都感到意外。使我的“弃绝”对某些人变得格外难以容忍的,恰恰是:我本该是他们中间的“最好者”之一,是他们最不可能疑心会给他们来这一手的人!(而且,我确实好几次以为感到,数学界的某位旧朋友身上有一丝\textbf{怨恨}的音调。)于是,很自然,他们在一场时尚里得到了满足:那场时尚宣告,所有这些“Grothendieck学问”,说到底不过是一大摞纸换不来多少东西,等等等等。一个人,无论多么显赫,不足以成一场时尚——要发起的时尚,还得回应许多别人心里的期待、暗藏的欲望,然后才成为共识、才取得支配地位\fnote{84}{(*)}{(5月28日)同一方向上,可参见5月14日的“掘墓人——或全体教众”一注,n\textsuperscript{o}~97。}。

在我离去之后的这十四年里,我也许一直倾向于低估这一离去在“上流社会”里造成的困扰——而对我而言,1970年6月的这次离去进行得如此自然,以至连一项“决定”都无须做出:新的任务一夜之间接了旧任务的班;旧任务忽然退后,仿佛被一段遥远的过去吸收殆尽!(诚然,在蒙彼利埃大学的同事们中间,我也不曾面对这样的困扰;他们构成一个与我离开的那个环境全然不同的世界。)也许,我同样低估了这份困扰在我的“1970年以前”的前学生们中间所能起的作用——他们中的许多属于同一个环境,并在其数学投资中“倾其所有”。很有可能\dmark{}\origpage{281},这份困扰在他们身上起的作用,不亚于在我以为还拥有于同一环境中的别的朋友们身上。无论如何,(我的某位旧朋友或学生与我之间的)每一种情形,都是独一无二的、不同于所有其他情形的个案;我能做的一般推测,只有非常有限而临时的分量。

回到更坚实的个案的地面,有一件事打动了我:我得以确认曾积极参与亲爱的老师下葬的那两位前学生,恰恰也正是最初以轻蔑的姿态、以使人气馁的意愿引起我注意的那两位——对着更年轻的数学家们,他们是“1970年以后的学生”,或者我的观念与数学进路的影响在他们身上清晰可见。这一巧合当然毫不奇怪(这当然不妨碍事件每一次都使我惊讶!)。另一桩有趣的巧合是:两人都是个人关系最友好、甚至最充满感情的那一类(对其中一位,这份关系以同样的音调一直延续到今天)。这与那条一般的确认相合:恰恰是最亲近的关系,最有力吸引并固定冲突的力量。

还有一桩巧合打动了我。在我将近二十五年来有过的学生中间,有两位,在我看来,既以其出众的“手段”、又以其与这手段相称的数学投资,区别于所有其他学生。(一项与我在我生命的二十五年里所做的那种力量相当的投资。)而且,对两位,我都曾有所顾忌,不肯把他们算进我的学生之列——尽管确实,他们都在与我的接触中学到了对他们有用的东西\fnote{85}{(*)}{(5月28日)这是一个委婉语,正如我后来身不由己地确认的那样!对此可参见昨天的“格格不入者”一注,n\textsuperscript{o}~67$'$。}。事情的本性是:两人都各自发现了他们自己的任务,无须我从我曾有(或仍有)的储备中向他们提出任务;两人的学位论文工作,都是在与我这个人无涉的情况下完成的\fnote{86}{(**)}{(5月28日)这并不完全准确。两人都在其工作中,本质上使用了我锻造、并且是他们在与我的接触中学会使用的工具。在这一角色之外,构成其学位论文的工作里的Hodge--Deligne理论(“Hodge理论II”,《Publications Mathématiques》n\textsuperscript{o}~40, 1972, 第5--57页),直接出自他得自我的Motif瑜伽——“混合Hodge结构”是对(在Motif的眼光里同样“明摆着的”)问题——“把复数域上不必半单的Motif观念‘翻译’成(适当意义上的)Hodge‘结构’”的“明摆着的”回答。在一次出色完成的“翻译练习”之外,这项工作中当然有“与我这个人无关”的原创而深刻的想法。但同样清楚的是:倘若没有我引入数学、并且他们在与我的接触中率先获得的那些想法与工具,那套Hodge--Deligne理论今天就无从存在(而且恐怕Deligne或我的另一个学生的几乎全部著作,也无从存在)。}。如此多的共同点\dmark{}\origpage{282}!作为相异点,我要说:两人中较年轻的(如果我没弄错)今天“荣誉等身”(其详尽清单,我替读者、也替当事人那份众所周知的谦虚省去了),而且是影响力最大的数学家之一,也就是说,也是最有权势者之一;另一位,眼下是代理助理,暂代一个职位,而该职位的正式任职者明年就要回来。还有别的相异点,它们在某种程度上解释了这一命运的悬殊——一如也有别的相似点,在这里不值得一一铺陈。倘还要举一条,那就是:在我有过的全部学生中,个人关系最亲近、最友好的,恰恰也是与这两人;而共同的热情,当时一下子就在他们每一位与我之间建立了强的纽带。我此刻想说的这份\textbf{巧合}是:据我所知,他们也是仅有的两位学生(带引号,自不待言!),对着那个“上流社会”,曾尽其可能地淡化和抹去这条通向我的、如此简单而分明的纽带。

这是一份真正触目的巧合;写下这些文字的时刻,它的意味仍从我这里滑脱。对两人,我都可以援引各自的时势缘由,彼此不同。而且很有可能,甚至可以说大概:在两人身上,在某个大概已不再是充分自觉的意图的层面上,这样一种缘由(一方的自负,另一方的谨慎)起过作用。但我怀疑,这个现成的解释,无论在这个个案还是那个个案里,能提供对事情的理解。肯定,在更深处,别的力量——那些真正的力量——曾在自负或怯懦的熟悉表象背后起过作用。肯定,表达这些力量的行为,有某种重要的东西要向他们每一位诉说。但肯定同样:在两个如此不同的人身上出现同样的行为——仿佛他们串通好了(鉴于二人命运的悬殊,这当然不可思议)——也有某种重要的东西要向我诉说,而且不关乎别人,恰恰关乎我自己。难道又恰恰是那没完没了的“\textbf{对父亲的拒斥}”的再现?然而,这种拒斥\dmark{}\origpage{283}本有许多可以选择的表达途径!抑或,正因为无意识那如此可靠的本能——它(当涉及“触碰”时)总能不偏不倚地直取最敏感、最不设防之处——竟使得两人不约而同地落在了\textbf{同一个}地方?实际上,我倾向于这样想。但这是一件推论出来的事,而不是一件看见的事;而由于缺乏那双有幸看得清楚而深的眼睛,我感到自己有点像个盲人,在黑暗中尽力摸索,尽力用他的手、他的耳朵、他的皮肤去“看”——可它们生来不是用来看的。

然而,为了不以这条\textbf{困惑}的音调收束(它有损于我的名声),而以一条令一位善意的、或许会有的读者欢欣的音调收束,我将只说出方才浮现的那个一锤定音的名字;在我看来,它分明表达了这场(关于一场下葬的反思之)\textbf{尾声}的各条考虑所共有的内容,即:

\begin{center}\textbf{全体一致的赞同!}\end{center}

\chapter{C)~上流社会}

\section{VII.~研讨会——或Mebkhout的层与反常性}

\subsection{不义——或归来的意义}

\notehead{75}{}

\dmark{}\origpage{285}(5月2日)看来我学无止境!我刚得知两份文本,它们(至少对我而言)以料想不到的光,照亮了前文已谈过的那场“藏掖”(对Mebkhout著作的藏掖)(“当值无名之辈与好上帝定理”一注,即注48$'$)。文本所揭示的,正是那两位显赫的同行兼前学生所扮演的角色——我曾确认他们心存轻蔑的漠然,却不曾怀疑其职业诚意。两份文本都属于\textbf{吕米尼研讨会}(1981年7月6日至11日)的会议录;研讨会题为\textbf{《奇异空间上的分析与拓扑》},刊于Astérisque~n\textsuperscript{o}~100(1982)。

第一份文本是研讨会的导言,署名\textbf{B.~Teissier}与\textbf{J.L.~Verdier}(即那位曾充当 Z.~Mebkhout正式的论文导师的同一位)。这份一页半的文本,开篇先解释某条“所谓Riemann--Hilbert对应”,它分明注定要在研讨会中扮演头等重要的角色(而它不是别的,正是“好上帝定理”,又名Mebkhout定理)。在这一对应里(这正是它的魅力与深度所在,也是必须引入导出范畴的原因):一个正则完整\textbf{模}(即集中在零次的一个正则完整复形),对应着一个由$\mathbb{C}$-向量空间的层构成的可构造复形;据称,它可由纯拓扑的性质来刻画,而这些性质,对一个任意域上定义的、不必光滑的簇上的平展层的可构造复形,同样保有意义。据解释,这就是研讨会“主旋律”——“\textbf{反常性、相交复形、纯粹性}”——的出发点:所谓“\textbf{反常}”的(\dmark{}\origpage{286}层)复形\fnote{1}{(*)}{(5月4日)关于这一奇怪的应用,见注n\textsuperscript{o}~76“反常性”。}不是别的,正是那些在“实质上”对应于(“Mebkhout式地”)最简单的正则完整微分算子复形——后者借一个单独的$\mathcal{D}$-模来表达——的层复形。

第二份文本,是\textbf{A.A.~Beilinson、J.~Bernstein与P.~Deligne}关于反常层的长文的一部分\fnote{2}{(**)}{(5月4日)我此后收到了全文;它确认了我手头这一部分已经向我表明的东西。};导言里称之为研讨会的中心工作。如同目录与我手头的其他页码所见证的,这篇文章宣告了导出范畴与三角范畴以骤然的强力重返公共广场——追随着Mebkhout那些鲜为人知的著作,与那道著名的“所谓Riemann--Hilbert”定理。

难以置信,却千真万确:两份文本里,Z.~Mebkhout的名字都缺席,正如它在参考文献里也缺席。我要说明:不仅 J.L.~Verdier对Mebkhout的工作了如指掌(这是理所当然的!),Deligne也同样了如指掌;而且,对这样一个对数学时事如此灵通的人——况且这还是与他切身相关的主题——甚至难以设想会是别样\fnote{3}{(***)}{我尤其提请注意:Mebkhout的著作与其“好上帝定理”,相对于Deligne(1969年的)较早的工作,是一个决定性的进展——而Deligne没有发表那些工作。对此可参见已引用的注n\textsuperscript{o}~48$'$。}。

B.~Teissier\fnote{4}{(****)}{(6月12日)B.~Teissier久已关注Mebkhout的工作,并因此是极少数对他持鼓励态度的人之一。因此,他对这桩巧取了如指掌,而且是明知故犯地给予了配合。他曾向Mebkhout辩白说,反正他“什么也改变不了”。}以及吕米尼研讨会的其他与会者情况如何——尤其是与Deligne共同署名前引文章的那两位\fnotevariant{5}{(****)}{(*****)}{(5月28日)我此后得知:A.A.~Beilinson与 J.~Bernstein是通过 P.~Deligne(1980年10月)与Mebkhout本人(1980年11月,在莫斯科的一场报告里,内容极为详尽)得知Mebkhout的结果的。早在1981年6月(原书此处如此)的吕米尼研讨会之前,这两位作者就在他们关于著名的Kazhdan--Lusztig猜想的证明里,本质上使用了好上帝定理。可与“一种不义与无力的感觉”一注(注n\textsuperscript{o}~44$''$)里所引的Zoghman~Mebkhout的信对读。(6月3日)关于研讨会全体与会者之团结的更多细节,见下一条注“研讨会”,n\textsuperscript{o}~75$'$。}——我不知道。似乎与会者中没有一个人曾如此好奇,想去了解把他们动员起来的那些观念与关键定理的父系渊源。

我料想,事情不言自明——有点(很有点)像次年将为Motif重返同一“公共广场”揭幕的Lecture~Notes的LN~900那样\fnote{6}{(******)}{对此可参见注n\textsuperscript{o}~51、52、59。}:父权属于发起并主持研讨会的那些光彩夺目的数学家中最光彩的那一位。对所有人来说,可以肯定的至少是:那绝不是Riemann,也不是Hilbert——否则,那场光辉的研讨会\dmark{}\origpage{287}就该在1900年召开,而不是1981年,即在无名的学生于Jean-Louis~Verdier门下答辩论文的两年之后。

我在这里得以确认的这类战役,如今也许已是家常便饭\fnote{7}{(*)}{我还想到另外两场同方向的“战役”,它们落实为LN~900(参见前一条脚注)的出版,以及五年前的SGA~4½(对此可参见注n\textsuperscript{o}~67、67$'$、68、68$'$)。(5月9日)关于与前面几场紧密相连的第三场这样的战役,见“好的参考文献”一注(n\textsuperscript{o}~82)——这次是 J.L.~Verdier手笔的另一篇“值得纪念的文章”。},并且完全得到认可,只要实践它的是占据高位的数学家,而付出代价的那位只是个无名之辈(尽管人们好心地邀他来,让他高兴)。这些实践者中的一位,倘若凭其手段也凭其作品,俨然一位大数学家(这使他立刻高悬于一切怀疑之上),也丝毫不改变事情的本性。肯定,我太守旧了——在我那个年代,这类战役叫作\textbf{巧取};而这一桩,在我看来,是容忍它的那一代数学家的一场\textbf{耻辱}。

天才的光辉并不减轻这样的耻辱;它给它添上一个新的维度——在我们的科学史上或许是独一无二的维度\fnote{8}{(**)}{在数学之外的任何一门科学或艺术的历史里,我也从未听说过类似的事情。}。它\dmark{}\origpage{288}可以使人隐约看到:在这一举动的荒谬与无端的表象背后(它出自一个命运已对他奢侈到无边、却仍然耽于巧取的人……),或许有别种力量在起作用,而不只是出风头的欲望,或羞辱、或使那个自觉无援无声的人绝望的无谓欲望。

既然我分明已深陷“风俗画”之中,我顺带指出(几乎当作理所当然):我的名字,在前引文本里同样缺席。然而,我能欣然确认:前引文章里没有一页(在我手头的那些页码里\fnote{9}{(*)}{(5月4日)还有别的页码,我此后也知道了。})不是深深植根于我的著作、带着它的印记的——直到我引入的那些记号,直到那些随处可见的观念所用的名字——那些名字,是我在它们尚未得名时就认识它们的时候,给它们取的。当然,合乎规矩的调整是有的:比如我在五十年代提炼的双对偶定理\fnote{10}{(**)}{平展对偶理论同样如此:在其慷慨的朋友Deligne的笔下,它也改姓“Verdier对偶”!},就为这场合改名为“Verdier对偶”——还是那个Verdier,没有搞错……\fnote{11}{(***)}{(5月5日)可与注n\textsuperscript{o}~48$'$、63$''$对读。在这场绵延将近十五年的漫长的葬礼中,也在这场刚刚于上个月、由主要的“被预葬者”所发现的葬礼中, J.L.~Verdier分明始终与他显赫的朋友不可分割:在这可悲的场合,他不计其数地为他撒下那些必不可少的万千花束。}

不过,我的名字不可能连间接出现都没有:经由对几部至今无可替代的文本(尽管有SGA~4½——它还不太胜任其使命)的偶尔引证,即EGA与SGA。(在“SGA~=~Bois~Marie代数几何研讨班”这个缩写的解释里,我的名字当然不出现;但在EGA里,人诚实不诚实总藏不住:全称是给出的,作者名单里包括我……)另一个打动我的细节,它见证了下葬综合征的强迫之力——在一位分明毫无“强迫症倾向”的人身上:我看到的仅有的两处对SGA的引证,每次都恪尽职守地特别写明“SGA~4里M.~Artin的定理……”,生怕哪位多心的读者会想:该定理说不定出自那位被小心隐去的人——尽管天知道,那场报告分明是做过了的,而且出自一位叫得出名字的作者!(注77)

\dmark{}\origpage{289}这一切,想必人们会认为,在今天的“上流社会”里,本是兵家常事。这场小规模交战虽然不能让我高兴(它也不是为此而设的……),却并不真正损害那位被预葬者:他那象征性的遗骸,就这样被交付给这场抢夺集市的种种偶然——而我不过两个星期前,才惊叹地发现这场集市。它并不以无力之中遭受的“不义”之感啃噬我的生活。它没有折断把我推向数学事物、推向周遭世界之事物的喜悦与冲劲;它没有在我身上烧掉这些事物的精微的美。我可以自认为幸运——而我也\textbf{确实}幸运……

我也为我那意味曾向我滑脱的意外的“归来”而喜悦。倘若这归来只需教给我这些天里已经学到的东西,它就不曾虚度——它已经使我满盈。($\Rightarrow$注76)

\subsection{研讨会}

\notehead{75$'$}{}

(6月3日)关于研讨会的其他与会者,我得知了一些细节,它们驱散了一切疑问。尽管官方日程里没有安排Mebkhout的任何报告,Verdier却不得不当场、在最后一刻请他做一个报告,以补某场官方报告的缺漏(那场报告曾托付给布日林斯基——他对D-模理论不甚了解)。Mebkhout由此得以陈述其观念与结果,尤其是好上帝定理,其陈述方式使人们对这条定理——以及与之相随的那套哲学——的父系渊源不留任何疑问;而正是它们,使得代数簇上同调壮观地重新启动,并落实为这场研讨会。于是,\textbf{研讨会的全体与会者,经由这场报告,都被告知了这一父系渊源}。我料想,所有人无一例外,此后也都读到了研讨会的会议录,尤其是导言与前引的Beilinson--Bernstein-Deligne文章。显然,没有一个人觉得有什么不对——或者即便觉得,也没有流露半点。Zoghman~Mebkhout在这方面的回音,一无所获。于是,研讨会的全体与会者,都满可以被视为这场在研讨会期间做成的欺瞒的共谋。

这场集体的欺瞒,早在研讨会当时就已分明,因为当Deligne在其关于所谓“反常”层的口头报告里不提Mebkhout的名字时,没有任何人觉得有什么不对。报告人只限于陈述好上帝定理,并说\dmark{}\origpage{290}他不会在报告中证明它。另一方面,他倒也(带着他那惯常的谦虚)把话说得很足:猜出他所称“反常”的层那些卓越的、事先难以预测的性质,“没有任何功劳可言”——那些性质,由他方才谈到的“Riemann--Hilbert对应”明摆着提示出来\fnote{12}{(*)}{可与所引文章的第10、11页对读。(6月7日)关于藏掖之技艺的细节,见下一条注“变戏法的人”,n\textsuperscript{o}~75$''$。}。人人都觉得正常:他避免点名那位有幸发现这条天赐对应的“有功之人”,并摆出一副父权在己的模样——而恰恰是他们刚被告知、或将在随后几天里被告知:事情绝非如此。人们想必认定:一个研讨会上的无名龙套,竟是一条如此卓越定理的作者,这必是某种不可容忍的错配;于是人人尽了一份力,拨正船头,确立一个共识,把父权归于那位分明是天造地设的人——那位\textbf{本应}是作者的人\fnote{13}{(**)}{(6月5日)况且一切都环环相扣!“学生”行列(继“研讨会”行列之后)中继续进行的反思,以及某种声调(尤其还见于最近与Deligne的一次简短的通信往来——见“葬礼”一注(n\textsuperscript{o}~70)的第一条脚注),向我表明:对Deligne、对我的其他上同调学家学生们来说,久已分明:平展上同调的发现与其掌控,那位本该的作者,同样是Deligne;而在某个层面上(即支配行为与态度的层面),他们深深相信:\textbf{归根结底}那是他,而我在旁边,不过扮演一个毛糙而笨拙的助手的角色——与其说有用,不如说妨碍一套理论(通向“Deligne定理——原Weil猜想”)的和谐展开,妨碍一种令所有相关者满意的角色的分配……}。

有一处很有代表性的细节:\textbf{Mebkhout的报告没有出现在研讨会的会议录里}。Verdier曾请Mebkhout不要成文他的报告,理由是:研讨会旨在陈述新结果,而Mebkhout的结果早已发表两年有余。

当人们不肯把自己关在技术话语里,而去看这场光辉研讨会里真正发生了什么——看推动各方的力量与胃口——人们会以为在看一部电影:某座遥远\dmark{}\origpage{291}的巨型都市的暗巢里,黑手党的天下。然而这是一幅地道我们自家的风俗画,演员们是法国与国际科学最高贵的翘楚。那位对一切发号施令的大主厨,不是别人,正是不久前还之于我扮演着谦和微笑的精神之子——或至少合法继承人(同样谦和,同样微笑)的那一位。至于那可差遣、可宰割的“软柿子”——在满是寸步不让的“硬汉”的世界里——由于一个我尚未完全参透其意味的奇怪“偶然”,也与我这个人紧密相连。他是我的“学生”,一如大主厨(而且同他一样,是带引号的“学生”……)——那是在我已被宣告死亡并下葬多年之后,进到我的门下来学艺的……

\subsection{变戏法的人}

\notehead{75$''$}{}

(6月7日)人们尽可以在那篇“值得纪念的文章”(前两条注里谈到的)里,欣赏那套漫不经心的藏掖的炉火纯青。作为全部工作之本质动机的那个范畴等价,头一回被引进,是在导言第4页(第10页,第9--15行)的一句话的转角处,不给它任何名字,随即转入“反常”层观念的一连串推论(第10、11页)。此后直到第16页末,它再未被提起;在那里我们读到\fnote{14}{(*)}{下面的引文里,着重号是我加的。}:

\begin{quote}
“应当指出,在以下几点上——\textbf{它们本应在这份笔记中占有一席之地}——我们未能完成任务。

——反常层与完整模的关系。如本导言所述,它曾起过重要的启发作用。关键陈述是4.1.9(\textbf{此处未证明})……”
\end{quote}

(随即转入别的“本应占有一席之地”的要点。)

我赶紧去看:作者们没肯劳神纳入其工作(或至少没肯纳入其证明)的这个“关键陈述”究竟是什么。找n\textsuperscript{o}~4.1.9……我碰到一条“评注4.1.9”——想必不是它;我找的是一个“关键陈述”,一个正式表述的定理或附注,并带有一个出处,表明作者们\textbf{在}哪里证明过或将证明它——既然他们不\textbf{在}这里证……可我怎么找,也没有“定理4.1.9”的踪迹——只有一个段落编着4.1.9的号。我姑且把这条“评注”读下去(不抱什么信念\dmark{}\origpage{292}——想必是编号出了错……):里面说“4.1.1在复上同调中的类似是真的……”糟糕,难道我还得回到4.1.1去看这说的是什么?我跳过去,接着往下读——瞧,我本已不再指望了:十一行之后,有一句话以“人们知道……”开头,以“诱导范畴……与反常层的范畴之间的一个等价”收尾。

呼——原来就是它!可我再怎么往下找,也找不到只字片语,来澄清这句谜语般的“人们知道……”。那位本不曾“知道”的读者,只觉得自己全然配不上这局面。对他来说,反正有一点是清楚的(且不论他配不配):这个“本应在那份笔记中占有一席之地”的结果——人们在此借一条技术评注的转角处“提醒”读者应当知道的东西——分明归于那“笔记”的作者们,或其中一位;也许是其中最显赫、并执笔了这篇文章的那位(有一种“家传的笔法”是不会错的……),也就是做了那场口头报告、而其人人皆知的谦虚当然不许他说“是我!”的那位——但人人都懂了,无须他说……

这立刻让我想起最近几星期反思里的几段记忆。头一段,是Deligne1968年的第一篇文章:在4月22日的“驱逐”一注(n\textsuperscript{o}~63)里——那是在LN~900谜底揭穿之后三天——我(十六年后)终于费心把它细看了一遍。我在这里又见到同一套笔法,想必经过十三年的“磨合”而略有变体。在那篇1968年的文章里——其主要灵感来自我——他在文章结尾处顺带、谜语式地提了我的名字,权作“手续完备”。在这里,他连这份小心都不再费——经验久已教他:完全用不着了!相反,在他年轻时那篇文章里,既然他感到非提我不可,他就以完全藏掖该文的最初动机作为补偿(连权的瑜伽一起,以便六年后以一个替补的父权把它端出来,再过八年,又等着Motif的开棺出土……)。无论如何,即便藏起(并留作独家收益……)这篇文章本质的算术动机,这篇文章也“站得住”,完全可读,配得上其作者把事情做得尽善尽美\dmark{}\origpage{293}的名声。而在这里,他所发展的理论,若没有那条启发性动机,便不可理解。于是他给出了动机,却以“关键陈述”这个定语来指涉它,同时把它从腿底下打发掉——不赐给它一个名字,不赐给它一个成形而冠以定理或命题之名的陈述;连“对应”(所谓Riemann--Hilbert的)都没有——他把这份差事留给了他的朋友Verdier与Teissier。他无须给它命名(鉴于其分量之轻\fnote{15}{(*)}{(6月14日)为安放这个“轻”,我提请注意:Deligne曾在IHES主持一场研讨班,试图发展可构造离散系数向连续系数的翻译,而未获令人满意的结果。对此可参见“当值无名之辈与好上帝定理”一注,n\textsuperscript{o}~48$'$。}——他肯定五分钟就能证明它!),也无须提名任何人——别人自会代劳,并且完全合他的心意。分明有一部瑜伽、一套哲学,作者以完美的娴熟与权威挥洒自如,而无须命名任何东西——那个他装作不屑的“轻”(“本应在这份笔记中占有一席之地”),他心里清楚:只要他知道在该沉默处沉默、肯等待,它自会外加送到。他头一回成功地玩这一手,那个“轻”是“权方面的若干考虑”,借一条谜语式评注的转角处影射了一下(且等六年后,再把权的哲学大张旗鼓地端出来)。据我所知,第二回是在我1970年离去之时——那个“轻”是“Motif之梦”,它在十二年间都不配得到一言的礼遇(想想吧——不过是个梦,而且还是个亡故者的梦,外加未经发表!),且待这一回发现真正的motifs(以及拿它们能做什么),再把那无可争议的父权——一如既往地谦虚——戴在自己头上\fnotevariant{16}{(***)}{(**)}{关于这套“以轻蔑来据有”的技术的其他评注,见次日之注,n\textsuperscript{o}~59$'$。}。

\subsection{反常性}

\notehead{76}{}

(5月4日)我清楚地记得,头一回听到“反常层”这个名字(译按:法语pervers还有“邪恶”的意思)——那应当是两三年前——它曾令我不适地一震;它在我身上激起一种不安之感。这份感觉,在此后两三次再听到这个不寻常的名字时,又出现了。那是一种内心的“却步”,停留在意识的表层;倘若当时我停下来检视它,它想必会表达为类似这样的话:给一件数学的东西起这样的名字,是什么主意!甚至给\dmark{}\origpage{294}别的任何东西或生灵起这种名字——严格说,唯有对人可以破例——因为明摆着,在宇宙间的一切“物”之中,唯我们人类,是这个字眼有时可以施用的对象……

我相当有把握(不敢说全然有把握):头一回对我谈起所谓“反常”层的,不是别人,正是Deligne本人——那是在吕米尼研讨会之后,他顺路到我家来的时候\fnote{17}{(*)}{倘若确实如此(如今我已深信不疑),那就得归功于我朋友的谦虚:因为(至少在意识的层面上)我并不曾疑心,引入并命名它们的人不是别人,正是他。是读了那篇“值得纪念的文章”,我才明白过来。(5月28日)老实说,文章本身对此也同样不置一词——正如它不置一词说Deligne是Riemann--Hilbert对应的父亲。然而,对“反常层”这一命名的父权,我从未有过怀疑;而且这一父权,此后确实得到了证实。}。那甚至想必是我们之间最后几场数学谈话之一——他到我家来之后,再没有过别的。恰恰是在那次来访时,显现了那个“征兆”;几个星期或几个月之后(当这一征兆在那次会面之后的数学通信往来中得到再次确认时),正是它使我断绝了数学层面上的交流\fnote{18}{(**)}{在纯属个人的层面上,这份关系仍以过去那种亲切友好的音调延续着,表面毫无变化。我的朋友习惯大约每两年,趁某次远足之便,来看我一次。去年夏天他又来过,那是一个可喜的机缘,使我也得以认识他的妻子列娜与他们尚且幼小的女儿娜塔莎。我记得那是从又一场吕米尼研讨会回来的路上;关于那场研讨会,我几乎没有听到什么回音(只有Mebkhout的几句阴郁而含糊的暗示——人们又一次好心地邀请了他,而他也想不出比再度入局更好的事……)。他们在我家住了两三天,接触自始至终都极好。}。(关于这一情节,见“两个转折”一注,n\textsuperscript{o}~66.)

回到那些所谓(我以为是错称的)“反常”层:显然,“正常地”说,这些层本该叫作“Mebkhout的层”,那才算公道。(不止一次,我曾把一些我提炼并研究过的数学观念,冠以先行者或同行的名字——他们与观念的联系,远不及Mebkhout与这个美好观念的联系来得近——何况这观念,依我看,其音色与其说“反常”,不如说“崇高”!)Deligne在发现并命名这个出自Mebkhout著作的观念之际的心境——他自己分明早已“奢侈到无边”,却仍准备巧取他人——这份心境,满可以正当地称之为“反常的”。肯定,我的朋友自己,在其内心深处,在某个不会为所炫示的门面所欺的层面上,也曾感受到这一点。在这个名字的授予里(乍一看,它显得悖谬),我感到一记\textbf{挑衅}之举,一种在如此完全的权力之中的醉意——它甚至可以炫耀(象征性地,借一个挑衅性的名字的陈列——其分明之极的真意,\textbf{没有一个人}肯去读!)它那巧取他人的“反常”本性。

\dmark{}\origpage{295}在我看来,绝非不可能:在某个深层,我曾知觉到我朋友身上这份心境的音色,而这曾促成我方才谈到的那份不安\fnote{19}{(*)}{我甚至倾向于认为确实如此。我不止一次在我身上确认:对事情的深层知觉,其精微与敏锐,与掠过意识层面、浮于意识表层的东西,全然不可同日而语。那个完全“醒着”的人,想必正是这些知觉被持续地整合进有意识的视野与有意识的经历之中的人——也就是那个完全按其真实的手段生活、而不只是靠这些手段中微不足道的一小部分生活的人。}。这份不安,尤其表现为:对他想必曾向我作的种种解释,我一概心不在焉;而我不相信,在这次会面之前,曾有哪次我没有持续专注地听他说话——尤其当事关数学的时候。我心中对这个被(天知道为什么)称作“反常”的观念,有某种阻塞——我并不真想听人谈它,尽管它与我一直(且在某种程度上依然)十分贴近的一些问题紧密相连。

干脆说吧:Deligne等人的这整篇文章,就是典型的、如假包换的“\textbf{Grothendieck学问}”,它们满可以同样出自我手(唯独主要观念的名字除外)!这大体上就是我在前一条注(n\textsuperscript{o}~75)第二部分里已表达过的;也是我通读所引文章的那一刻就已感到的——只是那份弥散的感觉,还不曾落实为这一触目的确认——它是我刚刚才做的。这份确认,以一种惊人的方式,使我重新对这桩深层的矛盾敏感起来:那个(在某种意义上)不能自已地复制着他所否认的人、并同化于他的人——那个他要埋葬的人,同时\dmark{}\origpage{296}也正是他\textbf{想成为}的人,而且(在某种意义上)他\textbf{就是}的人。

前天,在写前一条注(“不义——或归来的意义”)时,我已被这一巧合打动:我与朋友之间关系的这一转折——它骤然失去了一段共同热情中的契合,而那契合曾是这段关系之所以存在、并最强韧的内驱——恰逢我的朋友从那场值得纪念的研讨会归来之时;而那场研讨会的意味,刚刚向我显露。1981年7月我们相会时——在某个层面上,它与我们会面的其他场合同样友好、同样亲切——令我瞠目的,是那个“征兆”:论声调与神色,隐约,却又粗暴地分明,一份蓄意的轻蔑。它像是我的朋友预先支取的一笔\textbf{定金}——这一回,是在个人关系的层面上——对着那份隐含的、同样“隐约”(而且同样“粗暴地分明”)的轻蔑:那是他刚刚在吕米尼研讨会上、在当时一场当红名家间的技术精湛的盛大演出之中,公开对作为公众人物的我表达的。也是同一份“轻蔑”,刚刚对着那个胆敢(哪怕只有一点点)自称源出我门下的人表达出来(只是这一回,带着一种全然不同、“反常”得多的粗暴);而那人因此注定在我朋友皮埃尔眼里——至少在某个层面上——只成了“另一个Grothendieck”\fnote{20}{(*)}{在我们的个人关系里,我的朋友以我名字亚历山大的爱称(源于俄语)称呼我;这也是自我童年起,我的亲人、最亲近的朋友们对我的称呼。}——而如今,当务之急是不惜一切代价碾碎他……

\subsection{暂停!(Pouce!)}

\notehead{77}{}

(5月5日)翻阅这篇值得纪念的文章\fnote{21}{(**)}{关于“值得纪念的文章”,见注n\textsuperscript{o}~75。}——它(据说)主宰了1981年6月那场同样值得纪念的吕米尼研讨会——另一个细节打动了我。最后一章,顶着“从$\mathbb{F}$到$\mathbb{C}$”这个意味深长的名目,连篇累牍地描述一条卓越的原理——我把它引入代数几何,恐怕已有二十年了;那应当还在Motif观念诞生之前(Motif经由原Weil猜想,给出它最\dmark{}\origpage{297}深刻的例证)。这条原理断言:对涉及一域上有限型概形的某些类型的命题,只需在有限基底域上(亦即在一种“算术性质”的局面里)证明它们,即可推出其在任意域上、尤其在复数域上的有效性——在后一情形,所论的代数-几何结果有时可以经由超越途径改述(比如以整系数或有理系数上同调的语言,或以Hodge结构的语言,等等)\fnote{22}{(*)}{(5月6日)我记得,运用这样一条原理的第一个例子,见于Lazard关于仿射空间$\underline{E}$上代数群律的幂零性的定理(在任意域上)。它的证明曾深深打动我;我由它取得启发,处理了许多别的命题,并把它做成一套“哲学”,它主宰了我关于Motif理论的反思。}。我的朋友知道它,不是从别人、恰恰是从我嘴里,经年累月,就着许多例子\fnote{23}{(**)}{其中一个例子,见“驱逐”一注(n\textsuperscript{o}~63)。}。这条原理的父权(其初等形式甚至在EGA~IV里有明文——请别问我哪一节哪一条……)更是众所周知的\fnote{24}{(***)}{(6月5日)我自称这条原理的“父亲”,或许言过其实:据我所知,它的第一次应用当归Lazard(见前一条脚注)。一如在别的场合,我的角色,是感到他人一个想法的一般性,并把它系统化到使之成为一种“反射”或“第二天性”的地步。在权与Motif瑜伽的框架里,头一个使用这条原理的,很可能是Serre(而非我),凭他虚拟Betti数的那个想法——正是它,把我引上一部一般的权与Motif瑜伽之路。(Serre的那个想法,见注n\textsuperscript{o}~46\textsuperscript{9}。)诚然,流行的惯例是:一条已成通例的推理“原理”,其父权不归于可找到其最初踪迹的作者,而归于头一个知觉到其一般分量、把它系统化并推广的人。在这个意义上,可以说 N.~Katz的更正(下面这句话里谈到)把这条原理的父权归于我,是有道理的。}。以至于1978年赫尔辛基大会上为我这位光彩夺目的朋友授予菲尔兹奖章时, N.~Katz禁不住在颂扬 P.~Deligne的演说里顺带提了它一句,从而(不动声色地)更正了其显赫的受奖人那有点难堪的系统性“遗忘”。我是几天前才读到这篇演说的,与那篇“值得纪念的文章”同时。

无论如何,在这篇文章里,“算术”向“几何”过渡的哲学,被以这样的措辞呈献:任何一个不知情的读者都不可能怀疑,那位光彩夺目的主笔作者(请原谅\dmark{}\origpage{298}我的失言……)刚刚才发现这条分量如此之大的奇妙原理。

诚然,我没有为这条方法申请专利,而我光彩夺目的朋友也从未在任何地方说他是它那神奇的发明者;正如他并未明言他是那条著名的“对应”的父亲(请欣赏这个词,它散发着好一股十九世纪的馨香!)——那对应被谦虚地记在Riemann与Hilbert名下(这两位,倒也配得上给一位如此显赫的后继者的孩子们当教父);正如他在那部“值得纪念的文集”(LN~900,原书此处作LH~900)里也未指明:发明Motif、MotifGalois群以及与之相随的一整套哲学(他至今才掏出其中一段)的,分明是他;SGA~4½那部名著也同样无话可说——在那里,人们甚至再度不吝赐我以该卷“合作者”的荣衔:该卷从零开始如此光辉地发展平展上同调,还肯赏脸(尽管那两卷带着可惜的“多余细节的矿渣”等等)求助于SGA~4与SGA~5这两卷注定被遗忘的卫星卷,并且宽宏大量地承认后者的功劳:提供若干补充与技术枝蔓(其中有些甚至“非常有趣”)\fnote{25}{(*)}{关于“SGA~AT之役”(原书如此)的细节,见“白板”“格格不入者”“绿灯”“颠倒”四注(n\textsuperscript{o}~67、67$'$、68、68$'$)。}。

在所有这些情形里——也在最近五六年里我得以确认的许多别的微型个案里,尽管我从没想到过要\textbf{辨明我的不安}、给我所目睹或参与的事情起个名字\fnote{26}{(**)}{在一个个案里“辨明我的不安”的第一步,是不到三个月前在《收获与播种》里迈出的:那场(事后证明十分费力的——而且不无原因!)反思“注记,或新伦理”(第33节)。这场反思在那则注“年轻人的势利,或纯洁的捍卫者”(注n\textsuperscript{o}~27)里得到继续;然后,不到两周之前(在前一天发现“值得纪念的文集”(LN~900)所带来的冲击之下),又随注n\textsuperscript{o}~59“新伦理(2)——或抢夺的集市”重新开始。写那一注时,我心中对动用“抢夺的集市”这个相当粗鲁的字眼,还留有一丝踌躇。此后接连而来的发现向我表明:任何踌躇都不合时宜。}——在所有这些情形里,我认出同一套\textbf{做派}。我的朋友永远、彻头彻尾地处于“\textbf{暂停}”的豁免状态——他可以从容取用,带着那份由同侪与下侪的仰慕(无论如何是有根有据的)所给予的坦然良心;那仰慕,担保着完全的豁免。

\subsection{中国皇帝的新衣}

\notehead{77$'$}{}

\dmark{}\origpage{299}(5月7日)当然,那些看着我朋友Deligne行事、并且多少“在局中”的人——我是说那些并非刚下船、并非刚刚从当事者本人或其他当红名家(即便不都金光闪闪)的出版物里才学到“时下数学”的人,亦即那些了解来龙去脉的人——这些同行(毕竟,他们也还没有那么稀少!),在\textbf{某个层面上},是明白发生了什么的。在那几桩“有点扎眼”的个案里,他们必定感到了那份特殊的小小不安——那是我自己在这些“微型个案”面前、比起他们的个案微小百倍的情形面前,也不止一次感到过的。但他们所感到的是如此\textbf{庞大}、如此\textbf{不可思议},以至于它想必从未浮出水面——一如它最终开始在我身上浮出水面那样,经由一项\textbf{工作}:它表达为围绕着一个微型个案的两份文本(前一条脚注里谈到的那桩个案)。我确实没有听说,在我们的科学或任何别的科学的历史里,有过堪与比肩的事。对某些人来说,它不曾“浮出水面”,反而想必\textbf{成了风气},或至少被视为\textbf{正常}——只要一位分明是天才的、为万人所仰慕的人物,以天底下最自然的神色实践它,在众目睽睽之下,而这事从来不曾(据我所知)引起过只言片语的评注。

最近这几天里,我禁不住多次重新想起那个故事:“中国皇帝的新衣”。那位皇帝,受了无耻骗子的欺、也受了自己虚荣的欺,传谕说他将举行盛大的出巡,身穿举世未见的最华美的衣裳——那衣裳,是那些号称裁缝艺术家的骗子们斥巨资为他赶制的。而当他出巡,满朝臣子盛装华彩、前呼后拥,那些“艺术家”们打躬作揖,皇室倾巢而出——无论在仪仗里,还是在聚拢来瞻仰这第七大奇迹的人群里,竟没有一个人敢相信自己眼睛的见证;人人恪尽职守地赞叹,并为这身衣裳那无与伦比的华彩锦上添花——皇帝就披着它!直到一个混在人群里走迷了路的小孩失声喊道:“可是皇帝什么也没穿呀!”——刹那间,所有的人仿佛同一个声音,跟着这孩子喊道:“皇帝是赤身裸体的!”

我觉得自己就像那个孩子:他相信自己眼睛的见证,哪怕他所看见的,是如此闻所未闻、前所未见、为众人忽视并否认的东西。

\dmark{}\origpage{300}至于孩子的声音是否足以使某些人回到他们健全官能的谦卑见证,那是另一回事。故事就是故事:它向我们诉说着实在的某样东西——但它并不是实在\fnote{27}{(*)}{(6月14日)写完这条注之后,“中国皇帝的新衣”这个名字,向我显出是这场葬礼的一个天然的副题,表达它一个格外触目的方面。随后,反思转而面向我的全体学生,乃至数学界建制的“全体教众”,这一副题显得不再那么切合。然而,我终于明白:那则首先在我想到我的朋友Deligne时来到心头的寓言,同样适用于葬礼的各个侧面与各个插曲——它们每走一步,都在不可思议之中抵达乌布王式的荒诞(那不可思议,人人都恪尽职守地装作不见),而它却是真的。这个方向上的反思,尤其可参见“进步是挡不住的!”“研讨会”“受害者——或两种沉默”“玩笑——或权复形”“欺瞒”“掘墓人——或全体教众”各注(n\textsuperscript{o}~50、75$'$、83、85$'$、97)——没有一条特别关涉我的朋友皮埃尔。}。

\subsection{墓外相会}

\notehead{78}{}

(5月6日)才过去五天,我终于得到了那份慷慨的一包文献——来自我的朋友Zoghman~Mebkhout;其中尤其是那两份已被检视过的“值得纪念的研讨会”文本——一场围绕着一桩纪念碑式的\textbf{欺瞒}搭起来的研讨会!“不义——或归来的意义”那条注——我在其中尽力消化这桩新“事件”那令人难以置信的意味——就写在我收到这包文献的当天(5月1日的次日),犹带着发现的余温\fnote{28}{(**)}{连同“注记——或新伦理(1)”一节,这条注是仅有的不得不重写数遍的两篇注或节之一:因为头一稿(乃至第二稿)所“流出”的,还带着我一贯看待事物的方式所带有的全部惰性,远在有待检视的实在之下。}。

自4月19日——我终于得知Lecture~Notes那部“值得纪念的文集”(LN~900——见注51、52)——以来,这是关于宏大葬礼之种种盛典的第三项大发现;它在我看来分量也最大:既因它照亮了一些与我关系极近之人的作为,也因其作为一个时代(显然独一无二的时代——诚然,我对历史所知寥寥……)的“风俗画”所具有的含义。

第二项发现紧随第一项而来:即“Motif”的开棺\dmark{}\origpage{301}出土——它们被埋了整整十二年。继“值得纪念的文集”之后,轮到“值得纪念的研讨班”——那场从未举行过的“研讨班”,顶着假名字(无论SGA还是4½这个号),还附赠一篇幽灵论文的“第0状态”,更不必说(真正的)SGA~5研讨班的一场中心报告(它以晚出的面目出现,尽管比它早十二年);一场为战役的需要而被“借走”、不经任何手续的报告。这场光辉的战役,以及它在降临到可怜的SGA~5研讨班(从头、从尾、从中间被拆散!)身上的那些离奇坎坷中所扮演的角色,是在4月24日至30日间持续的一场反思中渐渐显露的。(对此可参见“同伙”“白板”“格格不入者”“信号”“颠倒”五注,n\textsuperscript{o}~63$'''$、67、67$'$、68、68$'$。)

那项发现刚消化完——与此同时,我“我的朋友皮埃尔”的回顾反思也接近尾声——而4月30日,我刚以“尾声——或全体一致的赞同”这条名字双重欢腾的“终注”,在这场无休无止的葬礼下面,骄傲地画上最终的句点(这回是肯定的——我总算走到了头!)——就在这时,这包倒霉的文献到了,把句点、尾声、版式、编号……统统推翻。对文献以及随附的评注与书信匆匆一瞥,便明明白白:我的句点完蛋了,那场头等葬礼的堂皇部署——我正准备修饰其最后细节——也完蛋了;我只好重新套上司仪的挽具……

不过天知道,我的朋友Zoghman有的是时间把局势告诉我!局势以潜伏的形态,想必已持续十年;以“急性形态”(这还是委婉语)至少三年——自那场研讨会起。在会上,他分明感到了风向,而无须等到次年、以其显赫的前老板兼庇护人为教父的那份高度官方的会议录问世。

他的论文答辩(1979年2月)之后几个月,他曾到过我住了六年的那个村子,想送我一份论文。不巧,几天前我刚离开(从此没有回去过,除路过之外……),退居孤独。他只遇到我的女儿,论文是她后来转交给我的。是第二年,我想\dmark{}\origpage{302},我们才终于正式相识,在蒙彼利埃大学;我们聊了一两个钟头。那时我对数学并不怎么上心,想必再不会记得一篇我只花了几分钟翻过的论文,也记不得其作者的名字。但这并没有妨碍接触的热忱。我清楚地记得一股即刻的相互好感。数学谈得不多(据我所记),谈的更多是或多或少的个人事情。Zoghman后来告诉我(此事我已忘却):他还是向我大致解释了一点D-模的“哲学”,而且他对这次会面很满意,感到我在从他那里获悉新事物时哪怕微微“震动”了一下,虽然这些新事物(在某种意义上)又是“意料之中”的。我印象最深的,是他的为人给我的印象——一种执拗而平静的力量之感,一个“闯将”的力量之感。在那个时刻,比起去年我们相会时、或其后的通信中,我更强烈地感到性情的深切投合——尤其在“闯将”这一点上。但两次会面之间的两三年,似乎已颇损耗了他……

我不记得在我们第一次短暂的会面时,Zoghman曾向我谈起他所工作的孤立,谈起那些曾是我的学生的“权威们”没有给他任何鼓励。倘若他有所暗示,想必不曾坚持。在当时,这件事本没有什么可让我惊讶的\fnote{29}{(*)}{(5月30日)这并不完全真——我把后来产生的失望情绪投射到了过去。我还记得,去年夏天与Zoghman相会时,我曾惊讶:我的上同调学家学生们(尤其Deligne、Verdier、Berthelot、Illusie)竟没有一个在Zoghman的工作中扶持他。这份惊讶,十来天后Deligne来我家时又一度重演(我想必向他提了Zoghman一句,没有遇到回音),后来又经与Illusie的一次电话再次重演。(对此可参见“欺瞒”一注,n\textsuperscript{o}~85$'$。)}。我说不好那是在1981年6月吕米尼研讨会之前还是之后\fnote{30}{(**)}{(6月3日)那是在之前——1980年2月,他论文答辩后一年。}。倘若在之后,他心里想必还压着滚烫的事情——而他全然没有那副模样。他更像一个知道自己想做什么、知道自己要什么的人,安静地走自己的路,不寻衅,也无人寻他的衅。

\dmark{}\origpage{303}后来我们没有继续通信。但我清楚地记得他;去年年初,我给他去了个短信,姑且一问:他是否有空着手去做一项为“温和拓扑”而设的宏伟奠基工作——它(依我看)只等着一个有他这般才具的人来做。Zoghman起初没有向我明说:他对这一前景并不真感兴趣;不过他似乎很高兴借这个机会再见一面。那时我离局太深,没有看清局势:我以为D-模理论已然功成、就此收官,一如凝聚对偶理论那样(注78$_{1}$),而Mebkhout也许正缺“大任务”。直到去年夏天我们会面,我才明白:就在他亲手开创的理论里,“大任务”并不缺——有些甚至还未动工,因为连被人看见都还没有!

无论如何,那是一个现成的机会,来第二次会面——这一回不是像第一次那样旋风式的。去年夏天,6月里,Zoghman在我家大概住了一个星期。在数学层面上,这次会面主要使我得以多少弄懂D-模的瑜伽。我“解冻”得很慢:与我旧日的上同调之恋已有些失联,而且尤其一头扎在《追寻栈》的写作里——它处在相当不同的音区。Zoghman看到我一面心不在焉地听,并不气馁;他不倦地卷土重来,带着动人的耐心。我想,当我终于明白那些大名鼎鼎的D-模不是别的,正是我很久以前所称的\textbf{模的晶体}、并因此于奇异空间上仍保有意义的那一刻,我被触发了。于是,我看见我那晶体-微分过去的整片直觉之网从被遗忘的深处浮起,我那“六种运算”过去有点生锈的反射,重新咬合……

于是轮到Zoghman有点跟不上;或者不如说,是事后他才决定,不肯把手指伸进那副齿轮(一如我的朋友皮埃尔不肯——尽管只要我还在场,他曾满腔热火……)。($\Rightarrow$注78$'$)

\notehead{78$_{1}$}{}

\dmark{}\origpage{304}然而,凝聚对偶确有一定数量的“精细”结果——尤其关于“对偶化微分模”的结构、它们与“朴素”微分模的关系、以及非光滑平坦情形的迹映射与留数映射——是我在五十年代末发展的,而且据我所知从未发表。尽管如此,就其实质而言,凝聚对偶理论(至少在概形的框架里),一如平展对偶理论(及其对局部紧空间的离散系数上同调的变体——由Verdier依平展的模型发展)、又如线性代数或一般拓扑,都显现为实质上\textbf{已完成}\fnote{31}{(*)}{(6月12日)就平展对偶而言,只要纯粹性猜想与“双对偶定理”尚未在全部一般性下得到证明,这就不完全真。}的理论,因而其本性是打磨停当、随时可用的\textbf{工具},而不是多少不为人知、有待探究与消化的\textbf{实体}。

\subsection{受害者——或两种沉默}

\notehead{78$'$}{}

我们的会面,是在一种友好信赖与亲切的气氛里进行的。然而这份气氛并未兑现其许诺。我如今明白:从那一刻起,我朋友心中的信赖就远非完全。那是那场大名鼎鼎的研讨会之后两年、会议录在Astérisque里刊出之后一年\fnote{32}{(**)}{(10月9日)Zoghman提醒我:这些“会议录”实际上迟至1984年初才问世。}——也就是他正付出一场骇人听闻的巧取的代价的时刻。可他直到仅仅四天前,才肯把这件事告诉我!去年他来的时候,是从又一场吕米尼研讨会\fnote{33}{(***)}{(5月7日)这里的记忆有一点点混乱——我倒觉得,他是正准备去参加那场研讨会。当然,从那时起,我所忆起的那些“辛酸而含糊的字眼”,缘由并不缺。但这辛酸,又因他在我家小住之后途经吕米尼而再度加剧。他从吕米尼回来后给我打了个电话,我由此得了些回声。从那一刻起,我就有一种十分清晰的直觉:他赶到吕米尼去,是为了享受被“那些人”(我没有细问是哪些人)怠慢的滋味——那些人慷慨地邀请他,为的是享受把他当无足轻重之辈来对待的乐趣。我想必把这话对他说了,或至少暗示了;这想必无助于改善当时我的朋友对我的心绪。}\dmark{}\origpage{305}回来的(这回干脆就是以D-模为主题的)——人们又一次慷慨地邀请了他,而他也赶紧应召而去。他谈起它,措辞既辛酸又含糊,话里带着意思:如今他“火中取栗”已毕,“一切都是别人做的”。我完全可以想象那幅图景——尤其Verdier,忽然想起三角范畴(连导出范畴也算上,反正一不做二不休!)的父权——那父权他搁置了十来十五年,勉强容忍他的“学生”Mebkhout在工作中使用它们……(注81)。

尽管他当时不肯把事情说清,Zoghman心里对Verdier想必积怨已深——这完全可以理解,看看他前老板那副不甚鼓励的样子。然而,我的其他上同调学家学生——Deligne、Berthelot、Illusie——也同样不曾屈尊对他所做的事情发生兴趣、多少扶持他一把。但人们几乎要说:对Zoghman而言,这本是理所应当——(人们会说)他从不曾在长辈们中间体验过别的态度。倘若他当时怨谁,那唯一而且不折不扣的,是Verdier。

从Zoghman的暗示(他分明无意细说)里,我懂了:“人们”在系统地把他的工作轻描淡写——一笔带过而已。这毕竟是人世间最寻常的事。对一件事的重要性的评判,在很大程度上是主观的;给自己的工作、给自己伙伴与盟友的工作,比对别人的工作——尤其对那些由于这种或那种缘故而人们有意贬低者的工作——记上更多的功劳与重要性,是寻常而几乎普遍的事。(而在这里,那个“缘故”对我而言实在算不得什么谜!)没有什么能使我疑心:远在这些寻常态度之外,这里有一场纯粹的巧取战役——其中根本谈不上“轻描淡写”,而分明是要把Mebkhout对那些使停滞之地重获生命的观念与结果的父系渊源\textbf{整个藏掖}起来……

然而,倘若天底下有一个人,我的朋友向他倾诉是再自然不过的,那正是我:在那几年执拗的工作里——有时辛酸,逆着时下的风尚——是我的著作给了他灵感;是我,在家里亲切地接待他,并且在尽力学习他乐于教给我的东西之时,轮到自己多少\dmark{}\origpage{306}做了他的学生\fnote{34}{(*)}{而且,Zoghman不但不对我谈他自己的下葬,也不对我谈我的下葬——尽管他置身头排,注视那场下葬的进行,已然将近十年!老实说,他的那些“庇护人”(边缘上还带着几分勉强)甚至曾让他亲手抬过装着我的遗骸的棺材的一角;但他们没有原谅他:他是满座宾客中唯一敢偶尔说出那个众人皆讳莫如深的名字的人!于是,我的朋友想必在与我的关系里自觉处境尴尬,而他未能在自己身上找到那份单纯,去承担一段(像我的一样)满载暧昧的过去,并不绕弯、清清楚楚地说话。谈他的下葬,也就谈我的下葬、谈他本人在其中扮演的角色……无论如何,我最终发现这场大名鼎鼎的葬礼的全部盛况,是顶着一种“沉默的共谋”才做到的;这一共谋既包括我的朋友Zoghman,也包括我的朋友皮埃尔,想必还包括我在数学“上流社会”里的大多数朋友。(6月3日)关于更多的细节,见随后的注n\textsuperscript{o}~78$''$。}。

我朋友怀着亲切的热忱来访之后,紧接着就有了一次“反弹”。我有这样的印象:他决意把八到十年间——在他那些曾是我学生的某些人所给的冷漠与轻蔑的刺痛下——积攒起来的疑忌与辛酸,移栽到我的身上。在其后的几个月里,我们之间的通信始终没有离开酸中带甜的音域——它最终停在一张贺年卡上,没有得到回音。

直到3月底,我才重新联络Zoghman,给他寄去“一段过去的重量”以及我当时为该节补加的各注(n\textsuperscript{o}~45、46、47、50)。我是想问他:是否同意我像已经做的那样,让他出现在关于我著作的那段简短反思里(在“我的孤儿们”一注,即n\textsuperscript{o}~46中)——尽管对人人都会一清二楚:我使用的是他给我的信息,而他或许认为那些信息是保密的。我全然没有把握:我的朋友会不会(如在他之前的别人那样)宁可“把自己压扁,也不愿招人嫌”。倘若真是那样,我会难过。

等他的回信,我觉得日子漫长——十天后才收到。我多少料到它还会是模棱两可的——但\dmark{}\origpage{307}这一回,它分明是热忱的。他毫无保留地表示同意,甚至为我在其中谈到他的那些措辞而动容。

在他那封长信(八页)的第6页,他才顺带地、就其定理那“令人印象深刻的”应用之众多(“无论在平展拓扑的框架还是超越的框架里”)提了一句:这条定理在文献里始终顶着“Riemann--Hilbert对应”之名\fnote{35}{(*)}{其信件的引文,见“一种不义与无力的感觉”一注,n\textsuperscript{o}~44$''$。}。他说得那么附带,而且那一手字迹潦草得仿佛故意,险些整个从我眼皮底下溜走!但我到底记住了;那真是一件怪事。怪到几乎难以置信;再说,也许我的朋友言过其实,分明他对人人都怀恨在心,包括我在内——尽管我心无恶意,这一点本是够清楚的。于是我补写了一条注(好个Zoghman,我还以为已经收尾了呢!),名之为“当值无名之辈与好上帝定理”,外加另外两条:“本能与时尚——或强者的法则”(写它时,在别人之中我也想到了他)与“罐装的权与十二年的秘密”。这条“当值无名之辈”的注,我起初写得并无完全的信念;Zoghman在我眼里纠结之至、充满矛盾,以致我自问:我这是要驶向哪里——我只做他的回声,对事实本身并没有多少第一手了解。我压根没有想到可能有巧取,更没有想到Verdier或Deligne本人牵连在内。Zoghman对我说过的话里,没有任何东西能提示这个……

然而,这两人都与那条好上帝定理系得太紧:不经过他们至少心照不宣的同意,其父系渊源几乎无从藏掖。这事想必在其后的几天里在我身上发酵。我想起Deligne曾对那个(十年后)由Zoghman解决的问题苦苦思索过;至于Verdier,他毕竟担任过研究导师;即便他没为学生费多少力气、宁可用冷淡与气馁来打发他,他至少应当知道这项工作里两条主要定理是什么——Zoghman肯定向他解释过,在那几场大名鼎鼎的、Verdier\dmark{}\origpage{308}“恩准过”的“会谈”里!于是我给这条注增补了一则评语,论Mebkhout的工作与Deligne一次较早尝试的关系,外加一条论Verdier角色的脚注。这同时也是对我朋友Zoghman的一次试探……

人们也许会想:这样一来,Zoghman必会抓住时机,亮出他藏了三年的炮群,让分明的真相轰然炸响,让受压迫者的大业凯旋!全然不是!十五天的沉默之后,来了一封信,谈数学里的一切,唯独不谈好上帝定理——或者不如说,关于它,他只限于按我的请求,给我其论文中的精确出处。(我到底想知道,这条我如此斩钉截铁地担保过的大名鼎鼎的定理,在哪里得到证明!)

我不得不在回信里,就“对我的著作的宏大巧取”——我刚发现的那桩(随“值得纪念的文集”LN~900,外加“预支了许多乐趣”:不日将在系图书馆里结识SGA~4½)——说上几句,才在又一个十天的沉默之后,使我的朋友终于引爆!

这一回他总算“倾囊而出”——而且是一\textbf{大}包:挑选得当的文献,使我(一个不常光顾图书馆、甚至连我在大学办公室里越堆越高的抽印本也不常翻的人……)得以对一种“气候”形成均衡的印象;在那气候里,许许多多的人,至今仍未参与我那漫长而庄严的葬礼\fnote{36}{(*)}{(6月12日)不过,Katz、Manin、Langlands似乎不在其列……(1985年3月)关于Katz的另一种说法,然而可参见“把一个个 i 点清”一注(n\textsuperscript{o}~164(II5))与“调动”(n\textsuperscript{o}~169)“插曲2”。(1985年4月)Langlands亦然,见“预先出土(2)”一注,n\textsuperscript{o}~175$'_{1}$。}。除了主要的“铁证”(那场大名鼎鼎研讨会的两篇文章,它们使那桩不可思议的欺瞒原形毕露)与另一篇“值得纪念的文章”(这一回出自Verdier的手笔\fnote{37}{(**)}{关于这篇文章,见“好的参考文献”一注,n\textsuperscript{o}~82。})之外,还有 N.~Katz论“菲尔兹奖得主”Deligne的演说,外加Langlands与Manin在同一届1978年赫尔辛基大会上的报告;然后是Deligne1970年尼斯大会上的“Hodge理论I”(其中第3行又一次影射“Grothendieck的猜想性Motif理论”(注78$'_{1}$))\dmark{}\origpage{309},以及同一位Deligne1974年温哥华大会上的“代数簇上同调中的权”(其中我的名字未被提及(注78$'_{2}$));最后还有与 A.~Borel(又一位老哥们——我同时得知他已回苏黎世……)的通信,以及Mebkhout在法国科学院报告上的两条短文,其中1980年那篇是其学位论文第五章(前一年答辩)的摘要,把好上帝定理更突出了一点\fnote{38}{(*)}{这条注的精确出处——Mebkhout的论文与好上帝定理——见“大部头与上流社会(或:鱼目与珍珠……)”一注,n\textsuperscript{o}~80。}。还不算一份文件——嘘!——以保密之名传达,我在这里一个字也不多说……

两封信随这份厚实的包裹而来(4月27日与29日),一封很长,两封都很有分量。如今他终于把“底”(这一回是真的!)亮出来了,Zoghman却仍旧像自我重新联络他以来那样,力劝我极端谨慎。倘若听他的,我会小心不把我的反思笔记公开,它们将只在我们二人之间保持绝对的秘密——至少牵连任何人的那一部分,因为“他们”“大权在握”,“人人都站在他们一边”\fnote{39}{(**)}{(5月30日)我兴之所至,在这里说得有点过头。Zoghman从未在任何时刻建议我不要发表我笔记的这一部分或那一部分。近来他甚至坚持:这些笔记理应以书的形式问世,以造福“后世”——一份限量印制、类似预印本的东西,在他看来有点“竹篮打水”。}!然而,我分明告诫过Zoghman:这些笔记——我寄给他的只是与他相关的摘录——注定要公之于众,而且越快越好。

一切要素似乎终于齐备,足以使受压迫者的公义大业凯旋;但那位“受害者”似乎仍在竭尽全力,存心把水搅浑——仿佛带着一种(人们会说的)事后的懊悔,后悔亮出了那道大名鼎鼎的“底”——直到命中注定的5月2日为止,Zoghman想必是它唯一的持有人。这份暧昧,几乎每一行都透出来(我几乎没有夸张),直到我刚收到的最后几封信——包括最后那封:他以一种阴郁的凯旋神色,把“值得纪念的文章”全份寄来(而先前那“大包”里,他到底只舍得与这份关键铁证的头二十页分手\fnote{40}{(***)}{(10月9日)Zoghman向我说明:其实他手头起初并没有全文的复印件,是后来才复制的。}。)

\dmark{}\origpage{310}至于皮埃尔朋友——我是说Deligne(对世人而言,他既非皮埃尔,也非“朋友”……)——他简直恨不能含着热泪唱起赞歌——这样一来,倒仿佛不再是Zoghman是“受害者”,而是可怜的Deligne:他被周围的人带坏了,毒害至深;唯一的恶人、把他带坏的人,是Verdier(而且……不妨顺着我的目光看……):分明是我“一定得罪过Verdier”,以致他恶意到只为害人而害人;更不必说,做他导师的也是我,授予他博士头衔、荣耀以及其他一切的也是我——总而言之,“绝对权力”的手段全在我!\fnote{41}{(*)}{这已不是我头一回听到“绝对权力”这样的调门,人们想借此让自己相信自己的无能,并为之辩护。倘若有人曾把对自己本人的“绝对权力”托付给什么人,那么对Zoghman而言,那不是别人,正是Zoghman自己!}

显然,倘若我的朋友怨谁,他真正怨的,并不是他那位显赫的前老板——那个人,他十年里有幸见面“会谈”总共不过三次(照他最近写信对我说的来理解)——一个遥远得令人眩晕、全然够不着的人;他怨的,是那个他想来就来、可以同吃同住的人……\fnote{42}{(**)}{(5月8日)我朋友对我的关系中,冲突的毫不含糊的征兆,恰恰出现在那次小住之后——那次他“分享了我的面包与住处”,气氛毫无保留地亲切,消除了我们第一次短暂会面想必不曾完全抹去的“距离”感——这肯定不是偶然。我在这里遇到一个我由来已熟悉的局面;我曾在“敌父(1)、(2)”两注(第29、30节)里(以相对一般的措辞)表达过它。当年我写它们,作为对先前反思的评语时,我不曾料到:我在其中描绘的那个原型局面,竟会始终处于一场尚在到来的漫长反思的中心——而那时,我还自以为已近旅程的尽头!}

Zoghman每前进一步,披露一点新东西,使我更多地了解他身为受害者的那场巧取局面(并能多少有助于解开它),我都感到那像一次\textbf{挣脱}——一场筋疲力尽的内心斗争的终点。有一个\textbf{角色},他似乎已以身家性命相认同,死死攀住,如同攀住他最珍贵的财产——那\textbf{受害者}的角色\dmark{}\origpage{311};而要保持这个角色,他就必须围绕这个角色以及为之辩护的局面,保持最绝对的秘密\fnote{43}{(*)}{(5月30日)自这些文字写成(5月6日)以来,我朋友的态度已发生急遽的演变;近来我不再觉察到他对受害者角色的攀附。当然,下面的文字(一如上面的文字)所关涉的,是我朋友生活中的某些情节,绝不声称刻画一种性情、或描绘一种恒常的成见。}。而此刻,在他不情愿的配合之下——那配合仿佛是被一个局面的逻辑硬挤出来的;而那个局面,不是别人,正是我带着这些多事的、关于一场风平浪静的葬礼的反思造成的——这个秘密行将终结,连同也许还有那个他不知从何时起便乐于固守的角色;在这样的时刻,他确实可能痛苦欲裂,比以往任何时候都更怨我。

我朋友Zoghman的这场“下葬”,是由\textbf{两种沉默}的合力做成的:每一种回应着另一种、又反过来挑衅它,在一轮毫无罅隙的回旋里,一方的角色紧紧贴合一方的角色——巧取者们与被巧取者。倘若我不止一次被这样一个景象攫住:“下葬者”同时、而且更深地,是他自己的“被葬者”,那我同样被另一幅景象攫住:在另一位朋友的身上,一个“被葬者”同时、而且更深地,是他自己的“下葬者”——与那些他乐于为之充当自愿受害者的人,紧密共谋。而且我看得很清楚:他自己那场巧取的第一责任人,不是别人,正是我的朋友Zoghman本人——三年来,他经由沉默,默许那些对他恣意妄为的人对他的羞辱。他手里握有一切,可以一战;他却三年来选择忘记自己还有手,未曾抵抗便告败北\fnote{44}{(**)}{(5月30日)这当然是一个有斗争家气质的人的主观看法;那种气质,在这位身上似乎阙如。不过,看来这些文字写成之后,战斗的心弦在我朋友身上已经苏醒,他已决意与那桩让他付出代价的不义一战。(1985年4月18日)关于对我朋友心境的一种不同、也不那么“苛刻”的照亮,另可参见“根”一注(n\textsuperscript{o}~171$_{3}$)。}

\notehead{78$'_{1}$}{}

这份简短的预注,我从不曾亲手捧读过;我读到的只是更详尽的“Hodge理论II、III”——发表于《Publications Mathématiques》。正因如此,我一直有这样的印象:Deligne不曾认为有必要哪怕暗示一下,他关于Hodge理论的那些想法之生成中,\dmark{}\origpage{312}Motif理论所起过的角色。我对自己说:倘若他有意愿提及我可能在他身边扮演过的角色\fnote{45}{(*)}{(5月30日)而且直到几星期前,我还在系统地淡化这一角色。对此可参见5月27日的“格格不入者”一注(n\textsuperscript{o}~67$'$),我在那里头一回觉察到我身上的这份态度,并隐约瞥见其意味。},他无疑会借着构成其学位论文工作的“Hodge理论II”来做——那正是提及这类事情的天赐良机\fnote{46}{(**)}{(5月30日)我也不记得曾被邀参加其论文答辩委员会。葬礼那时已然进行得热火朝天……}。我刚刚看到:他经由这句碑文式的话\fnote{47}{(***)}{Serre也经由指涉记号[3],在同一行里含蓄地露面——好奇的读者,在“Hodge理论I”的参考文献里能找到他的名字。这条草草的对照线,想必是从1968年到今天之间唯一的一处:对它一口气提到的那些“源头”(Serre(别名[3])、motifs、Grothendieck……)——无论多么谜语式——的一点影射。(5月28日)然而此后我又撞见另一处这样的影射;鉴于那个场合之特殊,它格外有趣。对此可参见“葬礼悼词(1)——或恭维”一注(n\textsuperscript{o}~104),以及紧随其前的“掘墓人——或全体教众”一注(n\textsuperscript{o}~97)的结尾——后者安放了那个“特殊的场合”。}——影射“Grothendieck的猜想性Motif理论”——一劳永逸地尽了提及我的手续,外带一个出处(Demazure在Bourbaki研讨班上的报告)。

又是无话可说!他从未想到要指明:他学得这套理论(且别忘了,整个是猜想性的!)是通过\textbf{别的来源},而不是Demazure那份寒碜的文本——后者绝无可能给出一套如此丰富的理论(整个是猜想性的!)的图像;那套理论,以底纹的形式贯穿Deligne其后围绕权的瑜伽的全部著作——且待那部海盗文集LN~900的步步升级, MotifGalois群(这回连一行含着亡故者名字的简短出处都没有……)才终于(十五年后)在那里开棺出土。

反思之后,在这句简短的引文里,我认出同一套“暂停!”做派——一句纯礼节性的引用,为的是把手续办清;附带的那个出处,绝无照亮读者之效(在这个案子里,即照出与那些恰好需要藏掖的想法之间那些分明而深刻的关系\dmark{}\origpage{313}——而它们在随后的十二年里也确实保持着藏掖),却\textbf{有欺瞒之效}\fnote{48}{(*)}{写这些文字时,一个联想不由分说地来到:早两年发生的、围绕“权”的第一起显影事件;它在“罐装的权与十二年的秘密”一注(n\textsuperscript{o}~49)的开头被谈到,并在“驱逐”一注(n\textsuperscript{o}~63)开头处得到更详细的论述。至于一般的“暂停!”做派,见“暂停!”一注(n\textsuperscript{o}~76,原书如此,题名对应注77)的反思。这套做派,开始对我变得非常熟悉了!}。

\notehead{78$'_{2}$}{}

这份文本\fnote{49}{(**)}{“代数簇上同调中的权”,P.~Deligne,1974年温哥华大会,《会议文集》,第78--85页。}我无需亲手捧读(我是仅仅几个星期前才知道它存在的),就知道我的名字不在其中。Serre的名字同样不在——而最早隐约瞥见“权的哲学”的人,正是他;那套哲学,是我随后以极大的细节提炼出来的。

\subsection{老板}

\notehead{78$''$}{}

(6月3日)Zoghman向我解释:他只是逐渐地、起初还模模糊糊地,才意识到围绕我的著作所做的那场“巧取”。1975年Verdier给他的手稿(见“好的参考文献”一注,n\textsuperscript{o}~82)对他曾如同天赐,尤其向他引入了可构造性观念及其本质性质,以及双对偶定理——他从后者取得启发,做出了D-模语境里的双对偶定理(或“局部对偶”定理)。只是多年之后,在读SGA~5时(屠杀版诚然,但还没有被屠杀到足以瞒过一个像他那样留心的读者),他才开始觉察到什么。长年以来,他对那位遥远的长辈满怀钦佩与感激,深信他大量汲取的那些观念属于后者。甚至似乎有多年之久,他确实深信:那套所谓“Verdier对偶”的对偶理论确实归Verdier——或至少归“Serre--Verdier”;同样,他所称“Poincaré--Verdier对偶”的那个对偶想法,也确实归Verdier。直到1979年前后(他答辩之年),他才开始觉察有什么不对——但我料想,他必定小心地不对其显赫的“老板”流露半点,一如在我们1980年2月与1983年6月会面时,不对我也一样。直到“反常”研讨会,1981年6月,当他开始感觉到\dmark{}\origpage{314}自己本人的著作正在被人藏掖,他才同时开始更清楚地认识到:自己闯进的是怎样一个世界\fnote{50}{(*)}{Zoghman那时对他前老板的评价,已低落到这种地步:他甚至深信,Verdier六十年代所做的一切(我在注n\textsuperscript{o}~81“赊账的学位论文与全险”的一条脚注里评阅过),或多或少是我在口授、或至少在提示。}!肯定,对他而言,我必属于这个世界:在那里,我的旧学生们(或至少其中某些)占据高位,以对待亡师的同样漫不经心,劫掠着遗孤学生。倘若真有区别,那不过是:我是亡者,而他们活得好好的,并且以确凿的方式证明着这一点……

我可以想象,甚至在“反常”研讨会之后,Zoghman仍难以相信自身健全认识能力的见证——尽管那见证分明地教给了他发生了什么。那篇大名鼎鼎的、由B.~Teissier与其“老板(原文如此)”Verdier署名的研讨会会议录导言,他直到1984年1月才亲手拿到。在驳回那份明摆着的事实近三年之后,震动想必愈发剧烈——我自以为是这样理解的。两个月之后我才重新联络他,3月底给他寄去“我的孤儿们”与“拒受遗产——或矛盾的代价”各注;又过了一个月,他才终于决定向我“亮底”,把“反常层会议之欺瞒”告诉我。

\subsection{我的朋友们}

\notehead{79}{}

瞧,我正准备把这场反思收尾并公之于众——它将终结Zoghman自己在那场他付出代价的巧取周围保持的沉默;而那场巧取的阴暗收益,他也照单全收\fnote{51}{(**)}{(5月30日)我提请注意:这场反思所由生的,是我朋友身上的一些心境;它们如今看来已被超越。(可与注n\textsuperscript{o}~78$'$的5月30日两条脚注对读。)}。也许它会令他不快,一如也许会令我的朋友皮埃尔不快——一旦它完成、文本誊清并印出,我就会把它亲手送到他手里\fnote{52}{(***)}{我原不曾想到,在我前面尚余的年月里,还有机会回到首都住上几天。但十多年来,我的朋友皮埃尔频频移步,到我蛰居的乡野深处来看我;于是这一回,趁着这个例外场合,轮到我移步了——那同时也是对一个屡次发出、而我从未应邀的邀请的回应。}。我能献给我朋友Zoghman与朋友皮埃尔的\dmark{}\origpage{315}最好的东西,他们两位多半会把它当作最坏的东西来接受:当作一场灾祸,或一桩冒犯。更何况,我的见证是公开的——一如他们两人的沉默,都是公开的仪式性行为,使这一位和另一位都须承担责任。

拒绝或接纳我的见证,是他们的选择;对Jean-Louis亦然——我曾把他算作我的朋友,一如今天的Zoghman与皮埃尔。这些选择与我切身相关,却不是我的选择。我毫无预言它们的兴致。我不久就会知道;我怀着强烈的兴趣、一份悬念——却没有一丝焦虑的阴影——期待着到来的星期与月份将带给我什么。我唯一的关切、我唯一的责任,是:我所献出的,确实是我所能献出的最好的东西——也就是说:真。

也许有人会惊讶:我竟毫不容情地谈论我以“朋友”之名相称的人;并在这名字里看出一句套话,甚至看出其中并不存在的讥讽声调。当我把Zoghman~Mebkhout或Pierre~Deligne称作“朋友”时,那是对同情的、亲切的、敬重的感情的重申——它们在我写下这些文字的时刻,在我身上在场。敬重告诉我:我无须“容情”对待一位朋友,一如我无须“容情”对待自己——一如我,他有资格与那份寒微的真相相遇;一如我,他无须容情。

倘若我不以“朋友”之名指涉Jean-Louis~Verdier,那绝不是因为我认为他比我的朋友Zoghman与皮埃尔、或比我自己更不“好”、更不“配”,而是因为生活恰巧把我们分开了。十五年前、乃至更久以前把我和他系在一起的同情的、亲切的感情,已被时间或多或少磨蚀,不曾有机会经由哪怕一点个人的接触而复活。我为重建这种接触做过的寥寥几次尝试,没有得到回音;我不知道,这些反思的阅读,会不会使一段已然僵死的关系复活。但即便此刻他于我不再是“朋友”,我也不认为:我像对待自己、对待我的朋友们那样毫不容情地对待他,就是不敬;而且我深知,反其道而行,于他、于任何人,都无益。更不必说:无论他还是我的朋友皮埃尔,倘若他们真想“自卫”(或进攻)而不肯冒险看自己一眼,并不缺手段,也不缺后援。更不必说:在他们有条件\dmark{}\origpage{316}令人气馁、或碾压人的地方,这一位与那一位,都不止一次这样做过,毫不容情,也毫无怜悯。

\subsection{大部头与上流社会(或:鱼目与珍珠……)}

\notehead{80}{}

(5月9日)况且,也该是时候了:我到底给出那条大名鼎鼎的Riemann--Hilbert——(不具名的Deligne)——亚当与夏娃——上帝——(尤其不是Mebkhout)定理的出处:人人都在大量引用它(包括我自己),而显然还没有人想到要问它在哪里得到证明。我自以为从我的朋友Zoghman那里懂得了:那条“值得纪念的定理”就在他的论文里;我果然在其目录里找到了它,顶着(确实土气、确实配得上莽汉的)名字“一个范畴等价”,第三章,第3节,第75页。更倒霉的是,它甚至没资格叫“定理”,而是叫“命题3.3”(而且更糟的是,我的名字出现在同一页上,还带着着重号!)。我承认,由于没有读过前面那75页来弄清脉络,我甚至不完全有把握是不是它——Zoghman向我确认是的;我信他\fnote{53}{(*)}{(1985年4月17日)最终显明:人们通常使用的“好上帝定理”的形式,并非此处所引定理的形式,而是一个相近的形式,可以用同样的方法证明。见“一个视野的绽放——或闯入者”一注(n\textsuperscript{o}~171$_{1}$),尤其其中今日所写的那条脚注。}。证明(据说)是同一篇论文第五章的对象——论文于1979年2月15日在巴黎第七大学答辩,答辩委员会由 D.~Bertrand、R.~Godement、G.~Houzel、Le~Dung~Tráng、J.L.~Verdier组成。有兴趣的相关人士,倘若尚未收到作者寄赠的副本(凡他无论对错地疑心可能感兴趣的人,他都寄了论文),只消向他开口,他将乐于从命……他当然给每一位我的前上同调学家学生都寄了副本;没有一位给过活着的音讯。想必他们其间都改了行,真不巧……

得承认:Zoghman做买卖的本事,确实不怎么样,不会以清楚而诱人的方式呈献他的货色——这是学得来的手艺;而他不像我的前学生们有那份运气,得以跟一位行业高手、且从不吝惜时间的高手学徒。可他也无从抱怨:他那“三次会谈”是得过了;说不定哪天,哪位“权威”会想起来,甚至为他那份难啃的大部头回一句“已收到”。他自己想必也发觉了:这份大部头很难啃(即便它对Riemann与Hilbert都并非白费……):他在法国科学院报告上发了一条短文——那毕竟短些——以引起人们对其大名鼎鼎定理的注意;题目嘛,您猜一百次:《论Hilbert--Riemann问题》!我早就知道,我的朋友Pierre~Deligne的历史修养并不比我强:他只需把年代顺序拨正\dmark{}\origpage{317},再贡献出那个漂亮的民间传说式的称呼“对应”,戏法就完成了——Zoghman,这可是他自找的……这条短文是1980年3月3日的,Série~A,第415--417页。

Verdier想必是在他“恩准”其学生的那“三次会谈”之一里(或在答辩时)得知这条定理的;但倘若真要问,他想必什么也没察觉。Deligne呢,他到底何时察觉到了什么,我说不好;但可以肯定:1980年10月他已知情,而据他自己说,Bernstein与Beilinson也是。Mebkhout还亲自去了莫斯科,向Beilinson与Bernstein(唯恐他们读起来费劲)详详细细解释其结果。我不知道是他们还是Deligne读过那篇论文或其后那篇科学院短文;但想必他们到底弄懂了里面的东西,因为次年那场大名鼎鼎的吕米尼“研讨会”恰好围着它转——出于天底下最大的偶然。

总结一下,并把我情报部门最新惠予的讯息算在内:对局势完全知情、并参与了那场所谓“反常层会议”之欺瞒的,至少有五人,即(按演员姓名的字母顺序)A.A.~Beilinson、J.~Bernstein、P.~Deligne、J.L.~Verdier与 Z.~Mebkhout——外加整整一场研讨会的与会者(原文称他们为acultes,疑为adultes“成年人”):肯定还都是光彩夺目的数学家,他们分明不求别的,只求被欺瞒、把鱼目当珍珠\fnote{54}{(*)}{(6月3日)实际上,看来研讨会的全体与会者,无一例外,都已在现场被告知了局势。对此可参见今天写成的“研讨会”一注,n\textsuperscript{o}~75$'$。}。这再次证明:我们数学家——从显赫的奖章得主到无名的学生——并不比张三李四多一丁点聪明或明智。

\section{VIII.~学生——别名老板}

\subsection{赊账的学位论文与全险}

\notehead{81}{}

\dmark{}\origpage{319}(5月8日)我觉得是时候了,该对“幽灵论文”这桩事作更周详的表述了;此前我只在两条注里(注48与63$'''$)“顺带”谈过。一位不细心的或不怀好意的读者,也许会说,我同时责备我的前学生 J.L.~Verdier两件自相矛盾的事:既“埋葬”了导出范畴,又把它们“发表”了(在SGA~4½里)并自居其父;正如同一读者会说,我责备 P.~Deligne,既“埋葬”了Motif,又把它们开棺出土(在LN~900里)。因此,或许并非多余:对1960年至今的局势作一番回顾。

1960或1961年前后,我向Verdier提议一项可能的论文工作:发展同调代数的新奠基,基于导出范畴的形式主义——那是我在此前几年里,为概形语境中的凝聚对偶形式主义之需而提炼并使用的。当时已经说清楚:在我向他提出的纲领里,前景中并没有什么严肃的技术困难;主要是概念性的工作,其出发点已然在手,或许需要相当规模的展开——规模堪比Cartan--Eilenberg那部奠基之作。Verdier接受了这个题目。他的奠基工作进展顺利,于1963年落实为一份关于导出范畴与三角范畴的“第0状态”,由IHES油印。那是一份50页的文本,1977年作为附录复制进SGA~4½(如注63$'''$所说)\fnote{55}{(*)}{单是这份文本,对一位天资出众的青年研究者两三年的工作而言,似乎嫌瘦。但Verdier当时的大部分精力,用于修习同调代数与代数几何不可或缺的基础——尤其通过听我的研讨班,以及两两对谈的工作。他对对偶形式主义的贡献(见下文)来得较晚:在我与Artin于SGA(1963/64)里详细发展了平展对偶的形式主义之后,当时我建议他(在其导出范畴的奠基工作之旁),在“寻常”拓扑空间、及这类空间的可光滑化态射的框架里,发展同一形式主义。大约在我以SGA~1开创“代数几何研讨班”丛书之际(1960年),Verdier联络了我,与Jean~Giraud、Michel~Demazure同时——问我有没有工作给他们;而他们敲的正是对的门!一件曾打动我的巧合:打从我写“我的孤儿们”一注(n\textsuperscript{o}~46)起,我记得他们三位联络我的时候,恰好刚组成一个小研讨班,名为“孤儿研讨班”(主题:自守函数,计算进路(原文approche calculs à brin de zinc,其中à brin de zinc措辞待考)),因为他们的老板(或CNRS的教父?)不打招呼就走了一年,让他们意犹未尽,又有些茫然。这片空,很快就被填上了……}

\dmark{}\origpage{320}答辩之所以没有在1963年、而是1967年举行,是因为难以想象:这份50页、尚待来临的奠基工作的胚芽,能构成一篇国家博士论文——这个问题当然根本不会被提出。也是出于同一理由,1967年6月14日答辩时(评审委员会由 C.~Chevalley、R.~Godement与我——主席——组成),根本不打算把这项工作作为论文呈上。提交委员会的文本,17页(加参考文献),以一项正在撰写的大工程的\textbf{导言}的面目出现。它勾勒这项工作的主要想法,并把这些想法安置在其众多使用的语境之中。第10、11页,对这项奠基工作所计划的章与节作了周详的描述。

倘若 J.L.~Verdier的理学博士头衔,是凭这份17页、勾勒着连他自己都说并非出于己手的想法的文本授予的\fnote{56}{(*)}{论文开头写道:“本论文是在 A.~Grothendieck的指导下写成的。它所包含的本质观念出自他。没有他最初的启发、他始终如一的帮助、他富有成果的批评,我不可能把这项工作进行到底。愿他在这里找到我深深的感激的表达。\par 我感谢Claude~Chevalley肯担任我论文答辩委员会主席,并有耐心读完这份文本。\par 我感谢 R.~Godement与 N.~Bourbaki,把我领进数学之门。”——按:“本论文”一词,恐怕只能指那项着手了的奠基工作的全体——所提交的文本构成其导言;而严格说来,那项工作在答辩之时并未“进行到底”。(5月30日)这一不自洽,恰好映出那份局势的暧昧;而那份局势,我作为论文导师——并且(照我手头那份论文封面上的说法)作为委员会主席——是第一责任人。我对一位光彩夺目的学生不够“严格”,有所纵容;这一纵容的方向,与我对Deligne所表现的(见“格格不入者”一注,n\textsuperscript{o}~67$'$)如出一辙,并且出了一份力,使同样的果实结了出来。},那显然是委员会与他之间的一纸本着相互信任订立的契约\dmark{}\origpage{321}:他承诺把这项工作做到底、并公之于众——对这项工作,他呈献了一篇光辉的导言。这份契约,候选人没有履行\fnote{57}{(*)}{越发引人注目的是:1983年12月, J.L.~Verdier竟拒绝了我请他加入Contou-Carrère论文答辩委员会的提议——委员会里还有 J.~Giraud,而我任研究导师;他的理由是:那篇论文(分明已全部写成,并经 J.~Giraud细心读过)与那个委员会,不经\textbf{巴黎}各大学论文委员会(原文如此)的核准,不具足够的严肃性的担保。}:他宣布的那份文本——一份按照一个已被证明有效的新观点写成的同调代数奠基文本——从未发表。

显然,倘若Verdier1961至1967年间的工作,仅限于写下1963年那份骨架式的“第0状态”,委员会绝不会想到接受这份“赊账的论文”。当时,其工作的成文想必已有相当的进展,可望一两年内完成;而出于实际的缘由,让Verdier不必等到那项本该为之奠基的工作完成就取得头衔,显得合宜。

还得补上一句:1964至1967年间,Verdier对对偶形式主义做出过几项有趣的贡献(注81$_{1}$);这些贡献,连同他本该继续的那项奠基工作,可以佐证人们给他赊的这笔账。单是他的那些对偶贡献,严格说,本就足以构成一篇像样的博士论文。然而这样一篇论文,绝不合我素来命题的风格:那些题目,都是对一个我感到其需要与紧迫的理论,做系统的、贯彻到底的发展(注81$_{2}$,原书误作82$_{2}$)。我不记得Verdier曾想过提出这样一篇“凭既有成果立论”的论文;我怀疑我是否会接受——因为这样一篇论文,与他我之间立下的“契约”毫无对应:我把导出范畴这个美好的题目托付给他,责成他发展大面积的奠基。

我承认我全部的责任,作为 J.L.~Verdier的论文导师与委员会主席:我轻率地(与 C.~Chevalley、R.~Godement一同——他们信任我所给的担保)把博士头衔授予一项尚未做成的工作\fnote{58}{(**)}{在这份责任之外,我还该补上另一份:在其后的两年里(在我退出数学舞台之前),我未曾督察Verdier是否确实履行了他所立的契约。得承认,我的精力当时全部投在我自己承担的奠基工作上——更不必说Motif的与其他的种种反思——以致那桩提醒他人尽其义务的不快差事,我顾不上多想。Verdier放弃出版那项预定工作的决定,我是七十年代初才知道的——那时我已全然不再关心数学,“反应”的念头根本无从起。}。

\dmark{}\origpage{322}倘若我今天确认这份轻率的某些果实,我原无权抱怨。但这并不妨碍我公开地做这份确认,并且声明:我的前学生 J.L.~Verdier的所作所为,只关他自己的责任,不关任何别人。

不履行对我、对那位曾信任他的委员会所立的契约,是一种埋葬的方式:埋葬我引入、并由他承揽以一项大工程去奠基的导出范畴的观点。这项工作也许做过,却从未交到使用者手里。这是对一整组他自己也曾协助发展的想法“画一道叉”。

Mebkhout的著作重新拾起导出范畴的观念,没有从Verdier那里得到任何鼓励(同样也没有从我其他任何一位身为“上同调权威”的学生那里)。对导出范畴的事实抵制,在我看来,直到1981年前后都是彻底的\fnote{59}{(*)}{(5月30日)这些带几分存疑语气的措辞,其实并不合宜。Zoghman~Mebkhout向我确认了(他是付出代价才知道的):我对“Grothendieck风格”的同调代数所受待遇的存疑表述,与实情完全相符。}——就在那年,它们随着需求的突然推挤,在那场值得纪念的吕米尼“研讨会”里强力重返(见注75)。

然而,Verdier那份“论文”的“第0状态”,早在此四年之前,即1977年,就已作为附录刊于SGA~4½卷(见注n\textsuperscript{o}~63$'''$)——也就是其论文答辩十年之后,而在(据我所知)Mebkhout是唯一在其工作中使用导出范畴的人的时刻,逆着此前七年的风尚。如无差错,他一直是唯一的一个,直到那场已提到的研讨会上围绕那条大名鼎鼎的“Riemann--Hilbert对应”的大“抢滩”——在那里,Deligne,别名Riemann--Hilbert,以这条“对应”之父的面目出现——原文如此;而Verdier(连同他那被其慷慨朋友大量引用的天赐“第0状态”),以导出范畴与两千年风格同调代数之父的面目出现——我这个人则绝口不提,更不必提Mebkhout\fnote{60}{(**)}{可与“同伙”与“不义——或归来的意义”两注(n\textsuperscript{o}~63$'''$与75)中的评语对读。}。

\dmark{}\origpage{323}在这些事件的光照下,我想我懂得了这份“第0状态”出人意表的发表缘由——SGA~4½的导言(还是那位一以贯之的朋友执笔的)说它“已经无处可寻”,而当时除了(也许)Zoghman~Mebkhout,谁也不关心去“找”它\fnote{61}{(*)}{不过,恰恰是4月底翻阅我刚收到的一份 Z.~Mebkhout著作的参考文献时,我才知道这份“第0状态”业已发表——那份上了年纪的文本的存在,我甚至已经忘了……}。于是恰恰有这么一个倒霉人,在他自己的角落里,逆着所有人;执拗地使用着这些过了时的观念,却没人弄得清他想达到什么——他终究那么执拗,以至于一个疑心开始冒头:万一哪天这个家伙端出几样够分量的东西呢,谁知道……毕竟,他斗胆引为灵感之源(在大师的著作旁边)的那个人,当年可是用这一切证明过、发现过东西的;那些东西,人们无法统统装作忘记,即便忘了它们的作者——而且大师本人,Jean-Louis~Verdier亲自,不正是靠着那条“Lefschetz--Verdier公式”起家奔向荣耀的吗?而没有这些该进纸篓的观念,那条公式他连写下来都费劲,更遑论证明……

我那影响巨大的前学生在这将近十年里(自从他摆脱了那道烦人的手续……)\textbf{押注导出范畴没有前途}、并且还将这样押注到关键时刻(那场大名鼎鼎的研讨会)的钟点,他想必认定谨慎为上(谁知道呢……),抢先于可能到来的事件——总之,来一份“全险”:发表的(当然不是那项本该终有一日构成一篇论文的大工程,而是)一份“证人文本”、一件“以备不时之需……”的证据;一份能证明他对一个\textbf{孤儿}之父权的文本——那孤儿,他曾乐于厌弃,并在等待事件来临之际,继续否认\fnote{62}{(**)}{倘若 J.L.~Verdier真有让埋葬了七年之久的导出范畴瑜伽为人所知的愿望,他本该选择发表构成其论文的那篇导言文本,而不是一份无人问津、只有放在那部瑜伽及其众多使用的底色上才见趣味的技术文本。但人们懂得:他绝无意愿在那份50页的证人文本旁边,附上其论文的那17页——里面有如今叫人难堪的断语,论的是那个绝不可提名的人的角色……}。

\notehead{81$_{1}$}{}

\dmark{}\origpage{324}所说的贡献是:
\begingroup
\renewcommand{\labelenumi}{\arabic{enumi})}
\begin{enumerate}
\item 局部紧空间语境中对偶形式主义的奠基;
\item Galois模语境中的相应形式主义(与J.~Tate合作);
\item 那条所谓Lefschetz--Verdier的\textbf{不动点公式};
\item 局部紧空间中的对偶。
\end{enumerate}
\endgroup

贡献2)与3)相对已知的东西,构成一个“意外”。最重要的贡献,依我看是3)。它的证明,从对偶形式主义(无论“离散”还是“连续”系数)轻易得出;这并不妨碍它构成我们上同调武库中“万用”公式的重要一味。这条公式的存在,是Verdier发现的;对我而言,那是一次(愉快的!)意外\fnoteplainmark{63}{(*)}{(1985年4月19日)关于这条漂亮的公式、它的角色、以及它在葬礼中的离奇坎坷,我在《收获与播种》第四部分的“真正的数学……”“……与‘无意义’”“把戏与创造”三注(n\textsuperscript{o}~169$_{5}$、169$_{6}$、169$'_{6}$)里还要谈到。}。

局部紧空间语境中的对偶形式主义,就其实质而言,是对我在概形平展上同调语境里所做之事的“势在必行的”移植(而且没有那种一切尚待去做的局面所固有的困难)。不过它带来一个有趣的新想法:把函子$f^!$(不先对$f$作可光滑化处理)直接构造为$Rf_!$的右伴随,外带一个存在性定理。这一办法后来被Deligne在平展上同调里采用,使他不借可光滑化假设,在那个框架里定义$f^!$。

这些评语,我想,足以说明:到1967年,Verdier已证明了他做原创数学工作的能力;这当然是为他赊账的决定性因素。

\notehead{81$_{2}$}{}

作为另一个例子,我指出:在局部紧空间的语境里,依“六种运算”与导出范畴那套“万用”形式主义的精神,对偶形式主义的详尽发展——Verdier在Bourbaki研讨班上的报告,可算它的胚芽。即便只就拓扑\textbf{流形}的语境而言,据我所知,Poincaré对偶形式主义至今也没有一部令人满意的参考文本。

\dmark{}\origpage{325}(6月5日)还有两个方向,我遗憾地看到,Verdier不曾认为值得把工作做到底——那工作他起手的势头已足够强,本可\textbf{收取其信誉}(我是指:在离散系数与局部紧拓扑空间的语境里,开启一套对偶形式主义),尽管那些本质观念并非出自他,而且他(一如对导出范畴)全无心肝去做一项任务的\textbf{仆人}、把一套完整的形式主义交到使用者手里(一如我在SGA~4、SGA~5、SGA~7三场研讨班里竭力所做的)。我所预见并建议他发展的对偶纲领,置于一般拓扑空间(不必局部紧)的框架、以及其间“分离的”、局部“可光滑化”的态射(即局部地,源嵌入$Y\times\mathbb{R}^n$,$Y$为靶空间)之中。这显然是同\textbf{任意}概形的平展上同调框架的类比所明摆着提示的。Verdier看得出:在局部紧空间的框架里,态射局部可光滑化的假设是多余的(这曾是一件意外)。但这并不妨碍:局部紧空间的语境(从而排除并非局部紧的“参数空间”)分明碍手碍脚。一个更令人满意的语境,应当同时涵盖Verdier所选的情形和我预见的情形:拓扑空间(乃至topos?)多少任意,而态射$f:X\to Y$须满足两条限制:1)分离;2)“局部可紧化”,即$X$局部嵌入$Y\times K$,$K$紧。

在这个语境里,一个“被允许的”态射的纤维,是任意的局部紧空间。下一步,则是允许$X$与$Y$不是拓扑空间,而是“拓扑多重簇”(即“局部如同一个拓扑空间”的topos),乃至任意的topos,同时对态射作合宜的限制(有待明确),使纤维成为\textbf{局部紧的多重簇},必要时附加补充条件(也许近于佐竹的$G$-流形的观点),比如(而且在万不得已时!)局部地取$(X,G)$之形,$X$为带一个\textbf{有限}算子群$G$的紧空间。

据我所知,连“寻常的”\dmark{}\origpage{326}Poincaré对偶,也还不曾在紧光滑拓扑多重簇(光滑:即局部如同一个拓扑流形)的情形里得到发展。有限群的分类空间的情形似乎表明:人们几乎只能指望在模去挠性的意义下得到一条(绝对整体的)对偶定理;更确切地说,须取一个$\mathbb{Q}$-代数作为系数环。除这一限制外,倘若Poincaré对偶(“六种运算”风格)在这个语境里照原样成立,我不会惊讶。从来无人看过它一眼(除了那些不可救药的微分几何学家——他们装模作样地看叶状结构的“叶空间”的上同调),这并不奇怪:看看我的上同调学家学生们——Deligne与Verdier领衔——对多重簇观念本身所设置的普遍抵制。

老实说,缺的是这样一场奠基的反思:在任意topos及其上的“离散”系数层的语境里,描述(倘能做到)topos态射的“真性”“光滑性”“局部真性”“分离性”诸观念,从而提炼出一个“被允许的”topos态射$f:X\to Y$的观念——对其,$Rf_!$与$Lf^!$两种运算(前者为后者的左伴随)各有意义,从而取得六种运算形式主义的惯常性质。这里,topos被视为不带环层的,或也许带着环层(必要时假定为常值或局部常值);先假定(至少起初)带环层的topos态射$f:(X,\mathcal{A})\to(Y,\mathcal{B})$满足$f^{-1}(\mathcal{B})\to\mathcal{A}$\textbf{是一个同构}(注81$_{3}$)。上面的反思提示:当我们把自己限于特征零的系数环层(即$\mathbb{Q}$-代数)时(译按:特征零本身不足以保证环为$\mathbb{Q}$-代数,例如$\mathbb{Z}$),“被允许的态射”的观念可以明显放宽,以涵盖这样的“纤维”:比如多重簇(拓扑的或概形的),而非寻常的“空间”(拓扑的或概形的)。

这个方向的最初尝试(除我处理过、随后Verdier依同一模型处理过的情形之外),归Tate与Verdier:在离散群或投射有限群的语境里。这一初步尝试的记忆,曾鼓励我去年继续这个方向上的反思:在小范畴(它们推广离散群、充当同伦模型)的语境里。虽然没有走远,这场反思仍足以使我相信:必须存在一套完整的六种运算形式主义,在小范畴的范畴(Cat)的语境里。(对此可参见\dmark{}\origpage{327}《追寻栈》,第VII章,第136、137节。)这样一种理论在(Cat)、乃至在Pro(Cat)里的发展,一如这类理论在拓扑或概形的空间与多重簇的语境里的发展,对我而言,其主要趣味在于:它是朝着在一般topos的语境里更好地理解“离散对偶”迈出的一步。

去年Illusie向我透露:他在半单纯空间(或概形)的情形里,曾同对偶理论的种种疑难苦斗。那听上去分明还是同一回事——在一个个案里察觉到一套六种运算形式主义的存在,并理解它。但似乎,一场奠基反思的前景本身,就有本事使我的前学生们——至少我的上同调学家学生们——人人浑身冰凉。我当年同他们费过力气,是怀着这样的信念:他们不会(在概念工作的层面上)恰好在与我同行的那个地点戛然而止;每当一个新的局面显出他们与其伙伴同我做的工作不够用时,不会只是搓手发愁。人们所做的概念工作,从长远看\textbf{总是}不够用的;数学的进步,恰恰在于重新拾起它、超越它——除此之外别无他途。1955到1970年间,年复一年,我发现前些年所做的已不敷需要,便立刻重新伏案——至少当别人(比如Mike~Artin,以他自己意义下的“代数簇”观点)还没有抢先把这件事做掉的时候。但似乎,我的学生们把我给他们做的榜样也一并埋葬了,连同我这个人、我的著作。

\notehead{81$_{3}$}{}

我记得,在(比如说)平展上同调的六种运算形式主义里,充当系数的环层局部常值的假设是多余的——本质的假设是:它们是素于各剩余特征的挠层,\textbf{并且}$f^{-1}(\mathcal{B})\to\mathcal{A}$是一个同构。一旦放弃后一假设,人们便进入一门(据我所知从未明言的)理论:它把“空间离散”对偶与“凝聚”对偶(关涉系数环层及其同态)混合起来。于是,人们设想:在概形(或更一般的topos)$X$、$Y$上,把系数环层$\mathcal{A}$、$\mathcal{B}$换成$X$、$Y$上的相对概形(不必仿射)$X'$、$Y'$,把带环层的topos态射\dmark{}\origpage{328}$(X,\mathcal{A})\to(Y,\mathcal{B})$换成如下类型的交换图表
\[
\begin{tikzcd}[row sep=2em, column sep=3em]
X' \arrow[r] \arrow[d] & X \arrow[d] \\
Y' \arrow[r] & Y
\end{tikzcd}
\]
并带一套这一类型语境中的“六种运算”形式主义。当$X$、$Y$等为点的topos时,人们应当找回惯常的凝聚对偶。

\subsection{好的参考文献}

\notehead{82}{}

(5月8日)说的是 J.L.~Verdier的文章“与一个闭链相配的同调类”,1976年刊于Astérisque~n\textsuperscript{o}~36(SMF),第101--151页。从某种意义上说,这篇相当不可思议的文章(尽管如今已没有什么应当再使我惊讶……)与Deligne等人的“反常文章”互为表里。除一点保留外,它实际上等于花五十页,在一个略有不同的语境里,\textbf{照抄}我在十年或十五年前纵横发展过的观念、构造与论证——术语、记号,应有尽有,一字不差!我几乎以为自己回到了1965/66年间SGA~5研讨班的一场:那些东西在那里被讲明过(显然,详尽到与会者吃不消\fnote{64}{(*)}{这一方向上的评语,见注n\textsuperscript{o}~68、68$'$“信号”与“颠倒”,我在其中检视这场研讨班成文工作的离奇坎坷,以及这些坎坷与Deligne“SGA~4½之役”的关系。随后的反思,向我显出这些坎坷、以及母研讨班经Verdier与Deligne合谋肢解的又一个未曾预料的侧面。两人的出版物为这场肢解揭幕,分别刊于1976与1977年——它们构成给Illusie的“绿灯”,让他去筹备(十一年之后……)SGA~5的出版(那一卷,如Deligne在SGA~4½里所说,“可以看作一系列枝蔓,其中一些还非常有趣”)。})整整一年。至少在那场研讨班之后,对所有多少“在局中”的人而言\fnote{65}{(**)}{关于我回到这份“仓促”印象的一场反思,见“沉默”一注(n\textsuperscript{o}~84)。},这一切都属于“众所周知”的领域。Verdier当然到场了,一如Deligne(唯一一个从不掉队的人——尽管那是他头一回把脚迈\dmark{}\origpage{329}进我的研讨班\fnote{66}{(*)}{这场研讨班的那一年(我想),正是我结识Deligne的一年;他当时该有十九岁。他很快“入了局”,甚至揽下成文我前一年平展对偶报告的差事(那些报告,他想必经由我的讲解与我的笔记而知),以及关于与一个闭链相配的上同调类的那场报告——后文在引过的“颠倒”一注(n\textsuperscript{o}~68$'$)里谈到,在本注里也还要谈到一点。以他那般手段、对主题的那份完全驾驭,他竟等了十一年才做成文,然后不经告知,把它收进其SGA~4½——如今回顾,这向我表明:早在1966年(而不仅是我曾以为的1968年——见“驱逐”一注,n\textsuperscript{o}~63)——即我们相识的头一年,我的朋友对我的关系中,便已存在一种深层的暧昧,而且从那一刻起就以全然清晰的方式表达着;而我一直回避去面对它,直到今天!}——总得有人这样……)。倒是有一件事:到1976年,那场大名鼎鼎研讨班的“成文(原文如此)”,由那些早已厌烦的“志愿者(原文如此)”拖着,已经拖了十年——我现在明白:其中一位“志愿者”到底还是按他自己的方式揽下了“成文”的差事,早在SGA~5于1977年出版之前!看来这位倒霉研讨班的坎坷,方便的还不止Deligne一人——他以自己的方式,从一盘散沙的局面里渔利。不过在那个时刻,Deligne倒还小心:一边把SGA~5的一场关键报告拆下来、名正言顺地并入其SGA~4½,一边到底还在其成文(关于与一个闭链相配的上同调类)里提一句“依Grothendieck的一次报告”。(诚然,他从中得到一份补偿:得以据此把我列为他的“合作者”!——见“颠倒”一注,n\textsuperscript{o}~68$'$。)

回到与一个闭链相配的\textbf{同调}类(别混淆!)——按标题,它是Verdier文章的对象——这个形式主义,我在口头研讨班上,分几次报告,以穷尽的细节发展过;而且面对的听众分明在求饶(照例只有Deligne一人,始终神采奕奕、毫无倦容……)。那是那一年我在平展框架的对偶形式主义上发展过的无数“长练习”之一;我感到需要完全驾驭一切在我看来必须彻底领会的要点。这里的趣味在于:拥有一个在不必正则的环绕概形上成立的形式主义——到正则情形的\textbf{上同调}类的过渡,以及与我那套使用具支撑上同调的老构造的联系(它立即给出与上积的相容性),都是直接的。我还发现:\dmark{}\origpage{330}研讨班的这一部分,属于未收入已刊版本的那一堆——想必Illusie(出一版拿得出手的(哼)本子的全部准备工作,到头来都压在他身上)该庆幸Verdier揽下了这件事,依样画葫芦(此处即:一字不改!)。

按照如今已经功德圆满的套话,“几乎无须说”:我的名字既不在正文里,也不在参考文献里(除了那条永恒的SGA~4引证——它含蓄地带着我,尽管人们早晚得找到替代……)。没有任何暗示提到一场与SGA~5这个缩写相应的“代数几何研讨班”——作者想必听说过它——尽管我分明记得见过他,殷勤地、懂事地做着笔记(像所有人一样;当然,除了Deligne……)。

况且,我说我的名字不见于正文,只夸张了一丁点:它有一次露面,神秘而碑文式,在第38页,第3.5节“基本的上同调类,相交”(人们终于抵达问题的死结了!)。那个出处是一句谜语式的话,其意我承认参不透:“系统地使用权复形(???又是这些该死的权!)的想法归Grothendieck,并经Deligne整理成型”——对这些神秘的“权复形”——据说我出过想法、而我此刻是头一回听说——再无任何解释。此后全书再未提及(此前37页里也从未提及)。谁懂谁懂!至于该节的内容,它一字不差地抄自十年前举行的那场SGA~5研讨班(而那时这一构造已有五六年的历史,见注n\textsuperscript{o}~68$'$);那场研讨班,他自然绝口不引。对Deligne的引证(他竟“整理成型”一条想法——当我的朋友还是中学生时,那条想法就已然现成!)是一朵“花”:这一想法想必之所以来到作者心头,是因为年轻的新来者Deligne确实揽下了成文我那场报告的差事(而且一搁十一年,为了人们知道的那些收益,见前引注)。这朵“花”,属于这对形影不离的朋友之间礼尚往来的交换。

不过,文章里有一个(想必)新颖而极有趣的结果(定理3.3.1,第9页):解析可构造离散层经由一个解析真态射的高阶直像的稳定性。Verdier约十五年前从我嘴里学得各种语境中的可构造性观念,以及\dmark{}\origpage{331}那条稳定性猜想——我在五十年代末提出过它(并对凡愿听的人谈起),那时我还没有结识他的荣幸。读这篇文章,一个不知情的读者(但这样的人已开始稀少……我恐怕又在重复自己)绝不会想到:作者并不是正在趁热端上他刚刚才发现的观念与陈述。他无须说那是他——不言自明。这就是那套大名鼎鼎的“暂停”做派,分明已经蔚然成风。

除了这个细节(依我之见,它合乎行业的新规范)之外,这篇文章想必总有十页上下(五十页里),围绕这个有趣的结果,呈献作者本人的工作。比例放到相应的大小,在Verdier身上、一如在Deligne身上,最打动我的,是他完全有能力做出漂亮的数学。即使在这篇令人痛心的文章里,前引定理也透出这一点。但他(一如其朋友)把自己保持在掘墓人的心境里,便与他显赫的朋友一样,只靠其手段中微不足道的一部分运转。一个(曾令我瞠目的)表面平庸的征兆,见于一位分明给出过机巧与嗅觉之证明的数学家:对他的“学生(原文如此)”Mebkhout的工作的分量,全无直觉——他乐于以自己的伟大俯视他,却从不曾自己做出深度与原创性堪比的工作\fnote{67}{(*)}{同一个令人瞠目的嗅觉迟钝,在这同一场合也见于Deligne:他仿佛到1980年才“感到风向”(Mebkhout那些想法的重要性),而Mebkhout在这一方向上自1974年起就工作着。我的朋友身上,自然嗅觉被自满所堵塞,我已有不止一次机会观察到;尤其自1977年(或1978年)起——那一年似乎构成第一个“转折”(对此可参见“两个转折”与“葬礼”两注,n\textsuperscript{o}~66、70)。}。这不是说他也许没有能力做到Mebkhout或我的地步。而是他从不曾给自己留一点机会去做大事——也就是说,放开缰绳,让一股热情牵引,而不是把数学及其天赋当作\textbf{工具},用以炫目、支配或碾压。直到此刻,他一直满足于照原样拾取那些现成而丰产的观念与观点。确实,他似乎已全然丧失了何为\textbf{数学创造}的感觉。

然而我记得,他与我一起工作的时候,这份感觉还在。没有任何外在于他的东西,妨碍这份感觉\dmark{}\origpage{332}重新浮出。一如其朋友——在那人身上,我常常感到那件精微而活泼之物的同样的隐没,被同样的自负所堵塞。

这篇50页的不可思议的文章,刊于一份有分量的杂志,为我照亮了“注记——或新伦理”那桩事故(第33节)新的一面:一篇\textbf{几页}的科学院短文——概括一项扎实而\textbf{原创}的工作,题材重要(依我浅见),是一位天资极高的青年数学家\textbf{两年工作}的果实——竟被两位权威以“毫无趣味”为由拒稿\fnote{68}{(*)}{细节见“棺材4——或没有鲜花也没有花冠的topos”一注,n\textsuperscript{o}~96。}。其中一位权威,不是别人,正是Pierre~Deligne——同一个Deligne,却肯一字不差地、亲自照抄我的一个学生那份不起眼的博士论文(他还恪尽职守地引用它)。(这份复制件,经一项显赫的署名拔高,成了那部同样显赫的丛书里“值得纪念的文集”LN~900中最厚重的一篇!对此可参见注52、67的末尾。)

分明,这幅“风俗画”日甚一日地丰满起来;而我为此甚至无须走出隐退、走街串巷去混迹“上流社会”。在几篇挑选得当的“大文本”里翻阅的零零星星几个钟头,已足以使我大开眼界……

\subsection{玩笑——或“权复形”}

\notehead{83}{}

(5月8--9日)我重新想到了那个“权复形”——它出现于Verdier值得纪念的文章里那条“出处——暂停”\fnote{69}{(**)}{见前一条注“好的参考文献”。}之中——一条俨然滑稽戏、纯粹无意义的出处。这条离奇的出处刚一入眼,一个联想便来到心头,并且一直在脑际萦回。我面对看似离奇、似乎挑战一切合理解释的东西——而其意味却又清清楚楚、并且被清晰地知觉到,只是在一个不同于常规逻辑的层面上——这已远不是头一回。常规逻辑,是我一生中几乎全部时间在意识层面运转的唯一逻辑——结果是:我不断地被“离奇”的、不可理解的事件甩在身后——它们那不可还原的离奇,令人焦虑!我的生活变化很大,自那一刻起(不到十年之前):我开始在我官能的一个更宽广的音区上生活。我懂得了:一切离奇、一切所谓的“无意义”都有\textbf{意味}——单单\textbf{知道}这一点\dmark{}\origpage{333},从此对无意义背后的意味保持好奇,就常常向我敞开其明摆着的意义。

在这“权复形”的无意义里,我感到一记与“反常层”的命名同类的\textbf{挑衅}\fnote{70}{(*)}{见“反常性”一注,n\textsuperscript{o}~76。}——此处的乐趣在于向自己证明:在一份有分量的杂志里、在一个自命为标准参考文献的文本里\fnote{71}{(**)}{而且,这个文本看来今天确实是一份标准参考文献——至少多年来,它一直是Zoghman的枕边书之一(他最近寄给了我)。正是在那里,他学得了可构造性观念(它在其定理里起本质作用);而且长年以来,他深信Verdier是这条对他至关重要的观念的神奇发明者。},人们\textbf{可以放胆}说一句与此相类的离奇话,而\textbf{没有一个人}会想到哪怕提出一个问题!而且我确信:这一挑衅里所含的赌注——文章发表八年以来的赌注——\textbf{赢了},直到今天:今天我是头一个向作者提出那个天真的问题的人。

当然,一桩离奇出现的时间(或地点)——在这里,恰是全文唯一一次指涉我个人的当口——绝非偶然;其形态也绝非偶然:这里借对一类观念——“权”——的影射(它与整篇文章的主题全然无关),并借对一个从未存在过的复合观念“权复形”的即兴,来实现。立时来到我心头的那个联想,很可能提供这桩离奇——在挑衅、权力示威之外——的更精确意味的钥匙。那是与另一处谜语式、同样纯礼节性的影射(但尚不带离奇的附加维度!)的联想——在注49开头所引的Deligne文章里\fnote{72}{(***)}{即“罐装的权与十二年的秘密”一注。从我们此处关心的角度对这篇Deligne文章作更周详的检视,见后文所引的“驱逐”一注(n\textsuperscript{o}~63)。}。那正是一处晦暗的影射,在一篇“权”一词严格缺席的文章里——除Serre与我,谁也无从看出——影射某些“权方面的考虑”,它们使我(文中特意写明:以一个欠一般的形式)猜出了该文的主结果。如我在更详尽的“驱逐”一注(n\textsuperscript{o}~63)里所解释的:在这份纯礼节性的影射背后,透出\textbf{藏掖}的意图——既藏我的角色,也藏那些他打算独占其收益的观念(关于“权”及其与一般上同调、尤其与Hodge上同调的关系)。这一意图,Verdier想必领会得格外透,因为他自己\dmark{}\origpage{334}就在同一个音高上“运转”(至少就他对我的关系而言;而依我看,这正是那对形影不离的朋友之间主要的黏合剂)。在两种情形里,诚实的做法本应是:开篇就明白指出主要观念、或促成文章的那个或那些问题的源头。

把这一点说清之后,这桩表面无意义的象征语言背后,我所知觉的意味便是:我可以毫无顾忌地当众陈列一桩明摆着的\textbf{无意义},同时借这个无意义、借对“权复形”那句荒谬的影射-出处,表达我真正的意图:我对Grothendieck在这项工作中的角色,无意透露分毫,一如Deligne当年对“权方面的\textbf{考虑}”的凭空影射——那条影射,对当时的读者,其意味不多一分,恰如我此刻发明的“\textbf{权}复形”,为着事情的需要、也随我的高兴!

我刚刚把昨天写的这条注誊清——其间我被Verdier的一个电话打断:那天我本想找他,正是为了当面问这个问题。我向他解释:我正在晚年勉力学一点上同调——这门学问我从来没弄懂过,他知道得很清楚——Mebkhout为供我学习,给了我一篇Verdier的旧文章;Mebkhout曾长期拿它当枕边书。我如今正尽力读它,但那里有这条谜语式的出处——他当然好意引我——可我全然不懂他指的是什么。

他非常高兴,甚至有点受宠若惊——可不是嘛,宽大的笑容从一副慈父般的欢容后面透出来——我竟然在我这把年纪,靠他那页旧纸学起上同调来了。我没有料到:当我说他分明知道我对上同调从未入门时,他会想到反驳我——那分明是久已有之的定论……至于那大名鼎鼎的“权复形”,我在听筒里再次感到他的宽大笑容(人们会说我编故事!),为终于有人(而且正是收信人本人)注意到那件被放过了这么久的事而欣喜。同时,又带着一丝窘迫——我想,那不是不知如何作答的窘迫,而是没能把一份乐趣\dmark{}\origpage{335}藏住的窘迫(一如人们对一桩略有荤味的故事所感到的乐趣……)。掉队如我,他实在无须在这一头费心!他不假思索地把话岔到Deligne身上(我没有提这个名字):Deligne在一篇文章里做过一个证明,其中还引了我——他记不太清在哪里;总之那里谈的是权,可不是嘛——当然他有点忘了——但那确实不是算术的权;这一点我全对,不是一回事……

声调欢容而不容置辩,他示意已经为我花了不少时间——略显匆忙,却不失那副从容、略带庇护的口吻。我为用这样一个有点愚蠢的问题打扰他表示歉意,并为他的解释道谢。我的歉意是真诚的,我的道谢也是——他确实把我想知道的一切都教给了我\fnote{73}{(*)}{即便摆出那副掉队的神气,我也并不真觉得自己在演戏(我没有这份天分);那全然自然——说实话,对这套我已近十五年不曾操持的玩意,我确实有点掉队!但我相信,即便老糊涂到该上灵车,我仍闻得出空核桃与满核桃的区别……}。

\section{IX.~我的学生们}

\subsection{沉默}

\notehead{84}{}

\dmark{}(5月9日)昨天我也许写得急了点:我在“好的参考文献”(见注82)里写道,作者兼前学生毫无顾忌地照抄的内容,“属于多少‘在局中’的人‘众所周知’的领域”。我曾试着为自己弄清,这些“多少在局中的人”究竟是谁——结论是:\textbf{不多不少,正是1965/66年间SGA~5那场研讨班的亲爱的听众们}——而且,如我已有机会说过的,这些听众多半或多或少掉队;而从这场研讨班的成文工作落在一群志愿者手里、我不肯觉察其信念不足的坎坷来判断,常常是“多”于“少”(当然,照例是Deligne例外)。事实上,只要SGA~5尚未成文出版、以便人们经由阅读而“入局”,就绝不会有别的“局中人”!这场研讨班实际上(天意弄人!)是在我最亲爱的两位学生与战友的两大“值得纪念的出版物”之后才出版的:其一,Verdier1976年的那篇文章——其中对由他执笔首次发表的观念之来源只字不提;那些观念,已被大量谈过\fnote{74}{(*)}{尤其见注n\textsuperscript{o}~67、67$'$、58、68$'$。};其二,Deligne随SGA~4½\origpage{338}。此后,人们才好心地邀请Illusie,来张罗其余部分的出版!

我已记不太清那场研讨班的与会者——比如Artin在不在场。我想,无论如何,我第一期的学生多半都在——当然Sinh~夫人与Saavedra除外(当时我尚未结识他们),或许还有Hakim夫人。此外还有Bucur(已故)、Houzel、Ferrand——我不算Serre——他对上同调的重型装备素来没有胃口,只是远远地、谨慎地偶尔露一脚。虽然除Deligne外,也许没有人十分清楚这一切将通向哪里,但我想,总该有十来位听众(不算很投入)至少跟得足够紧,可以被算作“在局中”\dmark{}。从昨天起在我脑际萦回的念头是:在所有这些“在局中”的人中间——他们因此俨然以上同调行家自居(即便不都像Illusie与Berthelot那样以“上同调”论文成名——那些论文分明够分量)——即便把Verdier与Deligne除外,想必仍有相当一批人,亲手捧读过Verdier这篇文章!Verdier身上某种神气,使我深信:从没有人向他暗示过,或许有什么地方不对劲。我也清楚:同样从没有人向我提过这件事——我是5月2日,即恰好一周前的今天,经由Mebkhout才知道这篇文章的存在的;而他对那场巧取,当然已知情多年。

这给十天前“全体一致的赞同”(为埋葬我寒微的身量而达成的)那句兴高采烈的确认,赋予一个十分具体的意味(注74)!这份赞同,把我“1970年以前”的学生中的许多人(乃至全体)——也就是今天在数学界引领风骚的许多人——都包括在内;它还把(或曾把)我的朋友Zoghman本人包括在内:他被当作上流社会的灰姑娘,逆着所有人死死攀住某种“对我的著作的忠诚”(借用他自己的说法\fnote{75}{(*)}{(6月7日)最近来访时通读关于葬礼的各注,Zoghman向我指出:他那个“对我的著作的忠诚”的说法,并不真正表达他的想法。他心里想的,毋宁是对自己的判断力与数学本能的信任——它们向他表明,我的著作为他带来了他所需要的某些想法。因此,那是一种对\textbf{他自己}的忠诚;而要做真正革新的工作,这确实是本质的东西。}),并有胆量与执拗,时不时以此自命,后果则有目共睹。谁能懂,就去懂吧!

总之,我暗示那份有分量的杂志刊登了一篇照抄“众所周知”的空壳文章,是我错了。作者当着(虽非所有、却)众多见证者的面照抄的东西,既未发表、也非“众所周知”(唯有一个闭链在凝聚框架里的上同调类除外——那个我久已发表);而那些观念,我再没脸去贬低:我并不曾认为,花一年工夫、在众多听众面前在一个研讨班里发展这些观念与其他观念,是虚掷光阴。想必Verdier的文章,是我所发展的观念与技术的一小部分的“摘要”,有用而做得好:恰好使它们得以进入“众所周知”的领域,成为把上同调(或同调)用于那些多少配得上“簇”之名的对象者的日课。\dmark{}\origpage{339}从这个观点看,Verdier做了有用的事\fnote{76}{(*)}{他做了,诚然,以原初SGA~5研讨班的“拆解”为代价——那场拆解,他与Deligne是主要的操办者与“受益者”。\par (6月7日)三天后(5月12日)的反思(见“屠杀”一注,n\textsuperscript{o}~87)显明:与这场拆解相比,所发生的更是一场“屠杀”,而Illusie的关联比Verdier还要直接——即便他并非“受益者”,而是代人行事。},我最终并无不满之处。然而,根据我从我的前学生、今天仍是朋友的人那里——在电话里——以及从他这个人身上我得以感到的许多别的事情(其中最“重”的、或至少最“轰动”的,是反常层会议之欺瞒)所感到的——我分明感到\textbf{有什么地方不对劲}。那场值得纪念的研讨会,从数学上说,想必在许多方面都极其辉煌。而“不对劲”之处,处在全然另一个层面。我可以试着以言词勾勒它,但我分明感到,那没有多大意义。在那个不感到这场研讨会不对劲的人那里——也在许多别的研讨会里,毫无欺瞒——当我做了这场“勾勒”的努力、哪怕完全成功地做到令我满意——他也不会多感到一丁点……

对我仍然悬而未决的问题是:这个“征兆”——它表现为一件如今想必相对寻常的琐事(一位作者把他人的未发表观念呈献为自己的)——这个征兆,是世风普遍败坏的征兆,因而只是今天数学界某种“时代精神”的典型征兆呢,还是它更想就我个人、就那个曾经是我、如今又回到我身上的人(经由昔日学生们对我的态度)教给我点什么。

两种可能的意义绝不互相排斥。倘若没有某种世风的状况怂恿,我的前学生们对我的关系,本无从循这条路径表达。况且,早在这个“征兆”之前,我已看到别的许多征兆;在我看来,它们在“风俗画”的层面上更雄辩。而在这个征兆里打动我的,是它区别于所有其他征兆的特殊之处:它似乎\textbf{同时牵连我的大部分昔日学生}。

这样的情势绝非偶然。把它干脆记在“世风败坏”(千真万确的败坏)的账上,便是回避其更个人的意味——那意味牵连我,一如牵连我的每一位前学生。倘若我说“每一位”,似乎超出了这个征兆的实际幅度,那是字斟句酌。因为这个征兆适时地提醒我:我的某一位昔日学生,竟不曾至少面对过这类局面,几乎是不可设想的。多年来,我感到一种关于我的某种“风向”,吹拂在我离开的数学家世界里(如今我分明看清这风的来源与缘由,我想)。不可能他们中的任何一位,从不曾感到这阵风的吹拂——无论是经由像这篇掘墓文章的发表这样的“事故”,还是经由任何别的场合。无论当事人愿意与否,这样一场遭逢,必然向他提出(或重新提出)他与我的关系的问题——我曾教他手艺。而我确认的征兆,在把我引到这件事的那一个之外,是:\textbf{那些曾是我的学生的人里,没有一个就此事给过我任何回音}\fnote{77}{(*)}{(5月31日)有意思的是,唯一一个曾向我暗示存在一场下葬的人,是一位非洲朋友——十年前,他跟我做过一篇第三阶段博士论文(即“1970年以后的学生”,身份寒微);我一直与他保持友好的交往。他有所暗示的那封信,应当是两三年前;在那个时刻,这毫无使我惊讶之处。我当时没有就他的印象问细节;他只是最近,才重新谈起它们。}。这是一桩“巧合”,其意味仍向我滑脱——但它不可能没有意味(注84$_{1}$)。

天开始破晓——我感到是停笔的时候了。我拿不准:在《收获与播种》里,此刻此地继续深挖这桩触目巧合的意味,是否合宜。这也许是留给别的明天的收成——倘若我今夜的反思,在我昔日学生中的某一位那里遇到回音。($\Rightarrow$注85)

\notehead{84$_{1}$}{}

(5月16日)我的昔日学生们之间的这份完美的协调——在对我的全然沉默之中——与其他征兆指向同一方向。其一是迎接“外来者”情节(见第24节)的全然同样的沉默——关于那份沉默,我在注n\textsuperscript{o}~23v里已多少追问过。另一方面,除Berthelot(他给我寄过许多抽印本)、Deligne(寄过四份,在其约五十项出版物里)与Illusie的一份之外,我的任何一位前学生的抽印本,我都没有收到过。关于他们对我之关系的两价性,这说明了许多。寄送抽印本——即便人们怀疑我是否会在自己的工作中使用它们\fnote{78}{(*)}{(5月31日)直到1976年,这甚至想必可以算是被排除的:七十年代初,我曾相当明确地说过,我不打算再重操数学。1976年在IHES的那场关于带除幂的De~Rham复形的报告,则相当清楚地表明我仍在关注数学。}——本会是最明显不过的方式,向那位教了他们手艺的人表明:这门手艺在他们手里并未僵死,而是活生生、在运转。但同样真确的是:至少对其中某些人,他们的出版物也见证了他们参与一场心照不宣的下葬;对那场下葬,最好还是别告知被预葬者——手艺也罢,非手艺也罢……相反,我却收到过好几位从事晶体上同调的作者的许多抽印本\fnote{79}{(**)}{(5月31日)说的是我不认识的年轻作者;我料想他们效法了Berthelot——在他们眼里,他想必是长辈。这里稍显奇怪的是:至少两年来(自1982年9月6--10日的吕米尼研讨会起),Berthelot正亲力亲为地埋葬我(对此可参见“共同继承人……”一注(n\textsuperscript{o}~91)5月22日的脚注)——莫非他对我的关系近来起了转折?我不记得收到过那篇晶体上同调及其盟友的综述文章的抽印本——文中对我的名字讳莫如深;他肯定小心没寄给我!},甚至收到过许多我只闻其名的分析家同行们的抽印本——当他们的工作重提(有时解决)我在三十年前或更早提出的问题之时;尽管明摆着我不会回到那个我已告别的题目,而且从“实用”的角度看,那些抽印本是白寄了。但这些同行想必感到了某种我的学生们不屑于感到的东西。——当然,在六十年代,我的学生们在我的一切出版物里都是头等的受赠者:无论我的文章,还是EGA与SGA两大系列;他们每一位(除Sinh~夫人,或许还有Saavedra)都应当拥有我1955至1970年间出版的全部著作(想必近一万页)。

诚然,我的前学生们有很好的同伴:我在“上流”数学界的任何一位旧友——包括那些著作与我的著作血肉相连的人、或在六十年代我的工作纲领的展开中扮演过角色的人——都不曾认为值得在我离开共同环境之后继续给我寄抽印本\fnote{80}{(***)}{(5月31日)当然,促使他们给我寄的心理缘由,远不及对我的学生们那样强;但,人们也许天真地想,总该比我的分析家同行们、甚至比我收到过抽印本的众多代数几何学家——我并不认识或只略识面——强些。分明,自我离开共同环境之后,“曾是朋友”这件事,在我数学界的旧友们身上造成或强化了我已有机会观察到的那种拒斥的自动机制。(关于这些态度——《收获与播种》里此处彼处顺带影射过——见5月24日的“掘墓人——或全体教众”一注,n\textsuperscript{o}~97.)}。最近,在我寄出《纲领概要》的十五二十位旧友(包括几位学生)之中——那份概要(除别的以外)向他们宣告:中断十四年之后,紧张的研究活动重新开始,且研究主题与我们当年共同推进的血脉相连——只有两位(Malgrange与Demazure)费心寄来几行致谢。我所收到的较周详(而且热忱)的些许回音,来自我近来结识的青年数学家,以及我的老朋友Nico~Kuiper——尽管他对我所做的这一类事情全然外行。他是经第三者得知这份文本的,并为我那不期然的“重返”由衷高兴\fnote{81}{(*)}{(5月31日)这几乎是唯一的、来自我的一位旧友(或前学生)的、对“重返”表示首肯的回音。这当然毫无奇怪之处:亡者现身,总归不合时宜地打断葬礼的顺畅进行……\par (6月17日)然而最近,我荣幸地收到Mumford一封热忱的信:他自称对《概要》里勾勒的想法感到“激动不已”(thrilled)且“非常兴奋”(very excited),并向我确认:我为其Teichmüller塔的组合描述所需要的那个关键技术结果,确实已得到证明。这是自1978年以来,头一回有我的一位旧朋友对我的“阿那贝尔式”想法产生共鸣——那想法的非凡分量(堪比Motif瑜伽),在我看来,自始便是明摆着的。\par (1985年3月28日)在这些文字写成之后,我还收到 I.M.~Gelfand一封非常热忱的信(注1984年9月3日),回复《概要》。}。

\subsection{团结}

\notehead{85}{}

(5月11日)那场倒霉的SGA~5研讨班的事,继续在我脑际萦回。“好的参考文献”\fnote{82}{(**)}{见注n\textsuperscript{o}~82。}分明给这段往事投下一束新的光,并同时赋予那场光辉的“SGA~4½之役”一个新的意味。

越想,SGA~5的事在我眼里分量越\textbf{重}。我的第一印象——那是仅仅几个星期前我刚“下船”时(见注68、68$'$)——是:65/66年那场研讨班的穷听众中间的一盘散沙,被我的朋友皮埃尔以其大名鼎鼎的战役加以利用;而在那场战役里,别人都无关。至于SGA~5的诸种不幸,既不怪他也不怪任何人,而宁可怪我自己:我没能——唉!——点燃我的志愿者-执笔者听众的热情,也没能替他们做那项他们执拗着不做、嘴上却说着“马上就动手”\dmark{}\origpage{343}的工作。然而近几天显明:毕竟有一位,其热情十年后才苏醒,把他在那里乐于取用的东西发表出来(对研讨班不置一词),从而在别的“志愿者”们还没下决心动手的时刻,为自己立了一份“好的参考文献”。

自昨天起,对我越来越清楚的,是:直接卷入这场研讨班之藏掖的,不止两个“恶人”,而是\textbf{我的每一位“上同调学家”学生}。如无差错,他们每一位都听过这场研讨班——即(按我的“上同调学家”学生们出现的先后为序):Verdier、Berthelot、Illusie、Deligne、Jouanolou。(Jean~Giraud不算;他在与SGA~5或其前驱SGA~4的主题相当不同的音区上工作。)

这场研讨班,是我首先\textbf{为了我的学生们的收益}而做的——即便有时他们哀求饶命——\textbf{我认为它不是一堆狗屎}。那一年里,他们每一位都学到了一大包“使用上同调的数学家”这门手艺!我在平展框架里、以周详得多的方式,重做我最初在凝聚框架里发展的那些观念——这些东西,除了这场专为他们的收益而办的研讨班,他们在别处无处可寻:在我之前,没有人费心做过;除我之外,没有人甚至感到有什么要做、为什么做。(照例,Deligne除外:他在这个研讨班本身里经月习得,比别人悟性快。)正是听过这场研讨班(及其前一场)并在自己身上多少消化了它,而没有别的,才使他们从此对对偶形式主义“在局中”——而且他们是\textbf{仅有的}在局者。这份\textbf{特权},依我看,为他们造成一份\textbf{义务}:即留意使这份特权不致滞留在他们独家手里;使那些从我嘴里学到、并在他们直到今天的全部后续工作中构成不可或缺行装的东西,在合理而惯常的期限内——至多一年,宽限两年——交到人人手中。

\dmark{}\origpage{344}人们会不无道理地说:首先轮到我来督察这件事。但我当年是真心接受学生及其他听众自荐承担成文工作的(对肯认真投入的人,成文只会带给他们最大的益处)——并非为了享受一项本属于我的工作由别人代劳、而我袖手旁观的好处。我继续(有Dieudonné的协助与其他人——其中1966/67年间还有Berthelot与Illusie)发展另一些在我看来同样紧迫的奠基文本,而那些,当年没有别人能代我做、或在没有我的协助下做\fnote{83}{(*)}{1960到1970年间,我必须以平均每年约一千页文本(EGA、SGA、文章)的节奏工作,其中几乎全部都成了通行参考文献(这一点,我在写它们时、或鼓励某位合作者在我的协助下写它们时,就已经很清楚)。}。这些文本本身,也成了不可或缺的参考文献——包括对我的“上同调学家学生们”:他们同所有人一样,很高兴在需要时找到现成的它们。

凭着经由与我的接触、经由听我的或参与我的研讨班而获得的对上同调观念与技术的驾驭,合力成文这场研讨班,与对“大名鼎鼎的数学共同体”所提供的服务相比——也许还有,稍后,对我的一份忠诚义务相比——是一件微不足道的小事。我说过,对我这个有手感的人,写定整场研讨班是几个月的活。五人分工,加上他们在那些年里各自积累的成文经验,手头又有我的详细手写笔记,每人的投入至多一两个月。他们比别的执笔者(比如Bucur——他分明巴不得把一件明显超出他能力的差事,交给更年轻、动机更直接的手)装备好得多。

只要我还在场(也就是其后的三年里),我理解“靠我”的反射可能起作用:理应由我统筹一切、跟“志愿者们”周旋。很可能,倘若我当时要求他们每人限期做两三场报告,我同样承担,以便了结此事,他们不至于推却。\dmark{}\origpage{345}自我退出数学世界的那一刻起,局面全然改观。他们发现自己成了\textbf{某份遗产的独家保管人}:既是默示的(没有遗嘱),又非常具体。诚然,从实际看,我的离去等于一场\textbf{消失}——我确实“亡故”了,意思是:除他们之外,无人知道这份遗产,无人能用它,无人(无论为善还是为恶……)关心它的命运。

倘若在我离去后的七年里,这份遗产始终隐没(除1976年的“好的参考文献”!),那是因为\textbf{我的学生们不曾在整段时光里让它变成公共的}。按比例说,这局面在我看来,与“Motif瑜伽”的局面相去不远:那部瑜伽,除我之外,唯有Deligne洞悉入微,而他判定把它留在自己怀里、独享其益。若说乍看有区别,那是:此处的“受益人”是五个而非一个;而且一个人所藏之物的深度,与五人共同藏起之物的深度,不可同日而语。

每个人的深层动机,我当然不知道——即便是Deligne的情形,我的领会仍模糊,而且想必仍将模糊。但在“实际”的层面上,Deligne的棋(连同“SGA~4½之役”以及其余)再清楚不过。同样清楚的,是这些战役\textbf{没有全体的团结便无以进行}。依我看,Jouanolou却不太在局中——他不俨然“权威”,我的印象是他久已离开上同调的泥沼(注85$_{1}$)。但我难以想象,Illusie与Berthelot不曾亲手捧读过SGA~4½与“好的参考文献”两者;他们跟我一样识字,也不比我笨。

倘若Illusie突然张罗起SGA~5的出版——恰在Verdier取用、Deligne取用、Deligne需要一座后勤基地以成其大名鼎鼎的SGA~4½(从中顺手牵走孕育了这份文本与其全部著作的那两场研讨班)之时——而Illusie手里已有十年功夫可以做这件事,这肯定不是偶然。倘若我1966年所做的收官报告——论未决问题与猜想——“不幸没有成文,其非常漂亮的开篇报告也没有[sic]:后者巡视了各语境(拓扑的、复解析的、代数的)中的Euler--Poincaré公式与Lefschetz公式”,这同样肯定不是偶然——而那是一场我不在行的下葬。同样不是偶然:把研讨班的一场关键报告截肢、不经任何手续送进SGA~4½,在Illusie与Deligne看来竟如此自然(而且只配在“细节变更”里顺带一提)\fnote{84}{(*)}{(5月16日)实际上,正如次日我就发现的(见注n\textsuperscript{o}~87):对母研讨班(或父研讨班!)SGA~5,在Verdier、Deligne与Illusie的手里,发生了一场真正的“屠杀”。}。

Luc~Illusie——我对他怀有如对Pierre~Deligne一样的感情,而他(一如Deligne)在我面前始终极其亲切\fnote{85}{(**)}{甚至在我1970年离去之后,Illusie对我仍有过细致的体贴——比如,长年以来,岁末节庆之际,他都寄给我非常漂亮的贺年卡。恐怕我不常回信致谢、报个平安——这些忠实的友情的信使,来自一个显得无限遥远的过去,而我已经与它失联。\par (5月16日)相反,Illusie从无继续或重建数学层面接触的念头;去年(我为数学问题联络他时),我感到他的迟疑。自我离去以来的十四年里,我只收到过他一份、唯一的一份抽印本,日期1979年。}——他的(有意识的与无意识的)意图,我不知道。但我确认:他在Deligne身边,成了一场\textbf{毫不脸红的欺瞒}的共同演员:那场欺瞒,把1965/66年的母研讨班SGA~5——恰恰是Deligne头一回听说概形、平展上同调、对偶及其他“枝蔓”的地方——说成一份顶着欺眼之名SGA~4½的文本集(晚八年写成,却佯装在前:既凭标题里的编号,也凭Lecture~Notes里的出版号,更凭作者那句不寻常的评语“它的存在(SGA~4½的)不久将使人得以\textbf{照原样}出版SGA~5”——着重号是我加的)的一种无形而近乎滑稽的附录——并且还摆出对那份寒碜的合集所自出的全部工作的毫不掩饰的轻蔑。

没有这些被如此潇洒地轻慢的工作,奠定其应得名望的Deligne的大部头著作,眼下没有一部写得出来——倘若真要问,一百年后也写不出来(Illusie与我的其他上同调学家学生们想必同样)。在这场“SGA~4½之役”的精神里,有一种\textbf{无耻}\dmark{}\origpage{347};Illusie(想必连自己都不觉察)为它作保;而它之所以能如此招摇,只因有一个\textbf{共识}的心照不宣的首肯。除Deligne本人之外,卷入这一共识的头一批人,恰恰曾是我的学生、某份遗产的主要受益人:那份遗产,在他们眼前,被交付给抢夺集市的风险与轻蔑。

还有那份武断的自满神气,那份慈父般的庇护神气——前天在我们通电话时,我还在我的前学生身上领教过\fnote{86}{(*)}{关于这次通话,见“玩笑——或‘权复形’”一注(n\textsuperscript{o}~83)。};还有我的朋友皮埃尔身上那份较隐约的居高临下——自“SGA~4½--SGA~5”双重大捷的翌日起我就领教过(而当时,乃至其后七年,我对它连一丝疑心都没有)——这些神气,绝不是孤独的产物,而分明仍是一个\textbf{从未被质疑过的共识}的征兆。这些神气,告诉我的不仅是Verdier与Deligne,还有我曾经的学生们全体;而首当其冲,是那些(就其工作主题与日常操持的工具而言)最直接相关的人。

“欺瞒”这个词——它不召自来——适时地使我想起另一场欺瞒,其中摊开着同样的犬儒:即所谓“反常”研讨会的那一场。如今,这两场在我眼前\textbf{血肉相连、不可分解——使两者皆成为可能的是同一个精神}。也许除Jouanolou(他已不太混迹“上流社会”)之外,我认为这些前学生是这场耻辱的共同责任人、团结一致的上同调学家。至于Berthelot与Illusie,没有任何东西允许我预设恶意或不诚(在Verdier与Deligne的个案里,那是不容置疑的)。但我至少确认一种盲目、一种健全官能之使用的堵塞;其深层缘由当然向我滑脱。倘若他们心中没有那份冷漠与轻蔑的蓄意,那么Zoghman~Mebkhout——七十年代唯一公开自称出自我的著作的人,而且题材恰与两人切身相关(尽管他们不屑于觉察)——本会享有最低限度的“善意推定”,使他们至少稍微看一看他所做的事情,从而觉察他自1974年起所投入的方向的趣味;那趣味是\textbf{明摆着的}\dmark{}\origpage{348}!两人都不屑于从一个装作还出自Grothendieck的无名小卒那里觉察任何东西。无名小卒的论文,是他亲手寄给他们的;我不知道他们开没开封,或读没读那些更短、更易消化的说明文字——总之,他们不屑于哪怕寄一张收条(Deligne亦然——分明是他定调子的\fnote{87}{(*)}{(6月12日)其间我得知:两人都不曾参加那场研讨会(1981年6月的吕米尼)。不过,可参见“欺瞒”一注,n\textsuperscript{o}~85$'$。})。

这当然不曾妨碍:与那场值得纪念研讨会的其他与会者一起,他们带着兴趣了解了那条卓越的“Riemann--Hilbert对应”,却想不到就其来源或父权提出丝毫问题,或至少(作为扎实的数学家)问一问它在哪里得到证明(注85$'$)。不过,在这一点上,我信得过Deligne:他自会乐于优雅地向他们讲解那份证明——对像他们那样的人,想必明白之至——恰恰是那种借Hironaka式奇点消解的证明,他们久已学得,而且不是从别人、正是从我这里(注85$_{2}$)。Riemann--Hilbert,Hironaka,阿布拉卡达布拉——戏法告成!

显然,一如Verdier、一如Deligne,他们已全然忘了何为\textbf{数学创造}:一个经月历年缓缓澄清的视野,揭示那件“明摆着”却无人看见的东西,在一个无人想到过的“明摆着”的陈述中成形(而在本案里,一位Deligne曾徒劳地试过整整一年……)——随后随便什么人都可以用现成的技术在五分钟内证明它,靠着他在某场遥远研讨班的长凳上有幸学得、如今却不屑(或不记得)想起的技术……

倘若我谈论Berthelot与Illusie时毫不容情,那并不是我尤其想给他们抹黑(在与他们的两位朋友清账之后)。我知道:他们既不比他们亲爱的同行们“更坏”、更蠢,也不比我更坏更蠢;而我在他们身上、在这一场合所察觉的那份嗅觉与健全判断的欠缺(有时,还有那份对他人必要的尊重的欠缺……)绝非痼疾,而是一个\textbf{选择}的效果。无疑,这一选择曾给他们带来中意的\textbf{收益}\dmark{}\origpage{349};而随着我的反思而来的另一份“收益”,对这一位或那一位,也许不合时宜。倘若如此,那只说明:它再度重演同一个选择——即在自身官能的微不足道的一小部分上运转的选择,不惜把鱼目当珍珠、把珍珠当鱼目,并无望地把空核桃(好哥们的)与满核桃(无名外人的)混为一谈。各人自便!($\Rightarrow$注86、87)

\notehead{85$_{1}$}{}

Jouanolou是我学生中除Verdier外唯一一个不曾把发表其论文放在心上的人。在我看来,这是对他所发展的奠基工作——即导出范畴观点下的$\chi$-进上同调——冷淡的征兆。由于他这一主题的工作,大部分是在我离去\textbf{之后}进行的——即在同学们以Deligne与Verdier为首、对我引入同调代数的观念(尤其导出范畴观念)之普遍冷淡已然亮出信号的时刻——气候并不鼓励Jouanolou认同于自己的工作、给它(应得的)发表之荣。而当这同一些Deligne与Verdier,追随着Zoghman~Mebkhout(别名Verdier的无名学生、别名Grothendieck的遗孤学生)的著作,终于(在大张旗鼓、互相吹捧之中)发现导出范畴的重要性时(见注n\textsuperscript{o}~75、77、81),那份遭冷落的Jouanolou论文,自“反常”研讨会起,重新恢复了全部的现实性;倘若1970年我离去之后,概形上同调理论的发展照常进行,那份现实性本不曾中断。一个触目的细节,例证Deligne在我离去之后其取舍的一场急遽“转向”:恰恰是Deligne本人(他分明深知在三角范畴框架里发展$\chi$-进上同调形式主义的重要性)向Jouanolou提供了一条技术关键想法,供一个有待研究的、成形的$\chi$-adic三角范畴的定义之用;那条想法在论文中得到发展。(对此可参见我1969年论Deligne工作的“报告”,第8节。)

(5月30日)关于Jouanolou的工作,另见“共同继承人……”一注,n\textsuperscript{o}~91。

\notehead{85$_{2}$}{}

一份有意味的“巧合”:恰恰是在同一场SGA~5研讨班里,这一整批人学得了那条证明原理——它既用于证明平展上同调里的双对偶定理(在拥有奇点消解的情形),也用于诸$R^if_*$无真性假设的\textbf{有限性}\dmark{}\origpage{350}定理,对$R\mathcal{H}om$、$Lf^!$亦然。(这些有限性定理,同样被从SGA~5的已刊版本里藏掖走、并进了SGA~4½;Illusie甚至不认为有必要在其导言里点明——我是写这些文字时才发觉的!)Zoghman没有听过那场研讨班的福气(他得到的是“好的参考文献”的替代品),是在另一处我使用过该办法的地方(为$\mathbb{C}$上光滑概形的De~Rham定理)学得它的。他其实也能在“好的参考文献”里学到:在那里,我的证明在解析框架里被照抄,以确立我的SGA~5学生们从此乐于称作“Verdier对偶”的东西(那条对偶,在我结识他的荣幸之前,我就已经知道了)。分明,一切都环环相扣!\textbf{同一份证明}(连同陈述,一并抄自我)被Verdier用作一种对偶的父权凭证——那种对偶,他除了在这场被肢解、被交付轻蔑的SGA~5研讨班里,无处学得;而它又被用来\textbf{对付}Mebkhout,成为(正因其“明摆着”)心照不宣的借口与手段,毫不脸红地巧取一项重要发现的信誉。

(5月30日)我记得,我头一回使用Hironaka式奇点消解、并领会消解作为证明工具的非凡力量,是为了对Grauert--Remmert定理做“三勺汤”式的证明——该定理描述某类复解析空间的有限覆盖上的复解析结构;以及在$\mathbb{C}$上有限型概形的情形里的类似陈述。(并不排除:恰恰在那次,原理是Serre向我提示的。)后一结果,是证明平展上同调与寻常上同调之比较定理的主要配料(其余,仰仗$Rf_!$形式主义归约;再稍加消解,从$Rf_!$过渡到$Rf_*$……)。

\subsection{欺瞒}

\notehead{85$'$}{}

(6月3日)实际上,我得知:他们本无须为这份父权提出问题,因为Berthelot与Illusie都是从Mebkhout嘴里学得好上帝定理的——前者在1982年2月,后者早在1979年(Mebkhout论文答辩之年)。虽然两人都不曾参加那场研讨会,他们却是那场研讨会之欺瞒的同谋,因为\dmark{}\origpage{351}他们不可能不知道那场围绕Mebkhout对好上帝定理尤其之父权的藏掖。我还可以想象:他们与研讨会的全体与会者一起,争先恐后地做了那场集体欺瞒的头一批受骗者——那场欺瞒,由他们的朋友Verdier与Deligne悉心操办(我的五位上同调学家学生中,有四位俨然是它的同谋)。至少就Illusie而言:去年夏天Mebkhout来我家之后,与Illusie通电话时,他对Mebkhout显然满不在乎,令我吃惊——他十分惊讶(几乎为他老师难过:他本指望老师有更好的判断力……)看到我在代数簇上同调理论的重启中,把头等角色派给Mebkhout。一些力量可观的共识,早已决定把Mebkhout归入无名小卒之列;我的朋友Illusie与这份三重矛盾活得逍遥自在,不问一个为什么:好上帝定理及其哲学的头等角色;围绕这些事情之父系的藏掖(他本人也在大帮人之中参与其中);以及他对Mebkhout的器量与角色的轻视——尽管他分明知道:Mebkhout是这些东西从未被点名的作者,而这些东西革新了数学的一个领域——在那个领域里,Illusie本人俨然权威。

我在这里重新找到那种健全意识与健全判断的完全堵塞——甚至在一桩表面上如对科学问题的评判般不涉私人的事情里;这种堵塞我已有不止一次机会影射,而每一次都重新令我错愕。而我在Illusie(想必还有许多别人)对Mebkhout——我的“遗孤学生”——的关系里确认的这份矛盾,肯定不是别的,正是他与我的关系中那桩更关键的矛盾的众多效果之一。正是这份矛盾——在他身上尤其,也在我的其他学生身上——在本书这支由我昔日学生们组成的葬礼行列的各注所进行的反思中,越来越清楚地显形……

\subsection{亡故者}

\notehead{86}{}

(5月11日)一如常有的情形,我是带着几分迟疑开始这场新反思的,主题“SGA~5--SGA~4½--反常性”——它似乎已被检视再检视到腻:“这会给一位早已听腻的读者留下糟糕的印象;再进这样的细节毫无雅致可言,SGA~5呀SGA~4½呀,这些都是过去了的事,不值得再饶舌……”

\dmark{}幸好我没有被这类熟悉的调门吓退——它想以“分明不值得”、“由它去吧”为口实,阻止我探一件事的底(至少就我当下所能探到的深度)。\origpage{352}倘若我曾发现过我认为有用而重要的东西,那从来是在这样的时刻:我懂得不听那份冒充“理性”、甚至“体面”的声音,而顺从我身上那份不体面的欲望——去看看甚至被认为“没意思”、寒碜,乃至粗劣、不体面的东西。我一生中不记得有哪一次,我曾后悔把某样东西看得更近一点——对着那些恨不得阻止我的积习反射。在《收获与播种》里,这些抑制反射比别的场合更强,因为这场反思注定要公之于众,这就立刻带来分寸(当我牵连第三方)与简洁(顾念读者)的约束。但我最终并不觉得,这些约束曾有任何一刻阻止我触及我想触及的东西、或把它深入到我感到欲望的地步。在曾一度显得棘手的临界情形里,我径直前行,带着这样的把握:必要时,我总还留着不把那场不客气的反思所“流出”的东西收进《收获与播种》的余裕。这些“边缘个案”,只在我犹豫是否牵连他人时出现;牵连我自己,则从未犹豫。但即使在第一种情形,事情竟是(这对我是一桩意外):我从不曾动用过那份“余裕”:《收获与播种》的文本,是我反思的全本——至少,是这场反思中找到通往书写之路的那一部分。

我感到:随着前一条注\fnote{88}{(*)}{即同日的“团结”一注,n\textsuperscript{o}~85。}那场简短的反思,局面大大澄清了。我是说:一个曾经故意弄得混乱的局面的某个本质侧面——我刚以“主题”的三重名目(SGA~5--SGA~4½--反常性)唤起过它——以全光显现在我眼前:即一种“团结”、一种“默契”;它们此前只被模糊地知觉过。这绝不意味着我自以为探明并理解了一个复杂局面的全部机关与来龙去脉——它至少直接而分明显然地牵连七个人:Zoghman~Mebkhout(在某种意义上充当某一局面的“显影剂”)、我的五位前上同调学家学生\dmark{}\origpage{353},以及我自己。我甚至不敢自称:对“SGA~5等等”这一局面里、在我自己身上起过作用的全部机关与动机——那场“倒霉的研讨班”距今已近二十年!——我都有所知觉。但我感到,比起昨天(或甚至今天早晨),我的状态好得多,足以理解并安放我希望将从这个或那个主要当事人那里到来的回声。

摆在我面前的主要问题(它想必在反思的另一个阶段已在场,如今带着新的力量重现)是(依我看)这个:经由我的学生们(差不多全员)进行的这场葬礼所发生的一切,是一件全然\textbf{非典型}的事,系于我的个人与我的独特命运的某些特点(如近十五年前我退出数学舞台,及围绕它的种种情势,等等)?抑或相反,是一件“再自然不过”的事,出于机缘的简单巧合——按“机会造就窃贼”的原理?我犹豫相信后者,尽管此刻辨不出、甚至隐约瞥不见:我这个人身上的哪个特殊方面,竟有这种本领,在我昔日学生们中间造就一份如此完美、如此一致的\textbf{赞同},去埋葬“老师”,并埋葬那些自称出自他、或著作分明带着他的印记(却并不因而是“他们的”)的人。是我这个人周围那圈“父亲”的光环吗——我已有机会谈过?抑或,我的离去这一件事本身,对他们的每一位置疑?此刻,我全然说不出来,因为缺乏那双会看的眼睛……也许到来的月份会就此教给我点什么\fnote{89}{(*)}{(5月30日)这个方向上的反思,见“掘墓人——或全体教众”一注,n\textsuperscript{o}~97。}。

最近三个星期里,我不止一次想到另一桩奇怪的“巧合”:葬礼“在其全部光辉中”的发现(经历LN~900--SGA~4½--SGA~5--反常层会议四幕,然后回到SGA~5与SGA~4½)——这项发现恰在此时完成:我刚结束一场对我的数学家过去、我与学生们之关系的深入反思。也就是说:恰在我刚刚——尽我官能之所能、并在当时已知事实的限度内,照它们经由常常雾蒙蒙的回忆所还原的样子——就这段过去“与自己算清账”的时刻\dmark{}\origpage{354}。换个说法:那恰是我终于\textbf{准备好}得知这件事、并从中获益的时刻。

“偶然”把事情安排得如此之好,以至禅观中连一个中断都没有。那场反思——它以对那些(依我所感)最重要的观念之命运的一场简短回顾开始\fnote{90}{(*)}{见3月31日的“我的孤儿们”与“拒受遗产——或矛盾的代价”两注(n\textsuperscript{o}~46、47)。}(那场反思还停留在某种模糊之中,只有某种基调执拗地浮出……)——这场反思在4月19日那个星期四,全然自然地继续下去。诚然,那还是带着“无耻”——借用方才的词,它也恰当地描摹了我当时感到的一桩事——的印象所激起的情绪:读那部“值得纪念的文集”LN~900之时。

在这“同一场”反思的新起点上,主发动机是“老板”——我的自尊、我的体面之感受了伤;写下我的情绪,我便在某种程度上从中解脱。分明是“自我”、“老板”在随后十天里领舞——那些日子,没有笑容也没有笑声,只有密不透风的严肃。无疑我必须经过那里、经过那十天的弯路,反思才回到它所离开的中心——我自己的身上。我仍记得那次回归的如释重负——像出了隧道重见天光!就在那时,我找回了笑与微笑,仿佛从未分离。那是4月29日。次日30日,当月最后一天,我欢欢喜喜地为反思的这一最终阶段画上句点。

那肯定也是我完全准备好接收下一个“包裹”的时刻——这一回由我的朋友Zoghman惠寄——“研讨会”包裹,两天后到达。到今天,我消化那份包裹的实质已进入第十天。而在这一阶段,尽管我为这场没完没了反弹的反弹而咬着嚼子,笑容却一天也不曾离我而去。而今天,我真的相信(诚然,已是第无数回!),是终于画句点的日子。

五天前,我就有过这同样的抵达终点之感,觉得剩下的只是总务活:这里那里补几条脚注,把涂改过多的页誊清\dmark{}\origpage{355}(那每一回都是思路尚欠清晰的征兆,它要求经由这看似机械的功夫各就其位;而文本经过誊清,总是带着一副新面目出来……)。那是在我刚写完如今那则“我的朋友们”注(n\textsuperscript{o}~79)之时;它自然而然地续成了几段“终曲”。但我终于把这些终曲与注的开头分开了。因为事实证明:那桩大名鼎鼎的总务活爆了:那些不打行距敲出的“脚注”,变成了真正的(而\textbf{非}脚注的)漂亮条注,颇具规模,不得不隔行重新打出,然后想法安置在此处或彼处。又过了几天,我才向明摆着的事实低头:继名为“研讨会”的行列之后,另一个行列正在成形,要加入送葬的队伍——而压轴的行列,将不是我心中已定的那场研讨会,而将由\textbf{学生}领队。而且就在今天,当第一个行列——已缩成一条注——刚添上第二条(“一种不义与无力的感觉”)时,我也知道了将由谁领队:是“遗孤学生”。于是,这支由一个小写的“学生”(遗孤的,合乎其寒微的身份)开道、又由一个大写的“学生”(这回毫不寒微)压阵的队伍,我看,终于齐备了!

现在,我想,是时候了——在第一次“假到站”之后——回到最后一曲《自深处》(De Profundis)的和声;它们今天比五天前更贴切。以下是我当时记下的;它们同样表达我此刻的心情。

(5月31日)最终,这又是又一次“假到站”——“终曲”这一回还是提前了!二十天过去,其间“总务活”不断爆成对被忽视的这一侧面或那一侧面的反思的重启。又有六条注加入本该压阵的“学生”行列。灵车在“学生”身后现身,载着四口棺材,掘墓人随行。分明还缺些东西,好给一支似乎谁也没在送葬的队伍以实体与意味。

经验使我变得谨慎;我等着事件到来,此刻不冒险预言:队伍是否终于齐备,或者是否还有一个被遗忘的行列,会在最后\dmark{}\origpage{356}一分钟钻进来,不肯错过那终极的仪式\fnote{91}{(*)}{(6月12日)谨慎确有其必要:一个新的行列“我的学生们”已从最初称为“学生”的行列里分离出来——后者改称“学生——别名老板”。}。

\subsection{屠杀}

\notehead{87}{}

(5月12日)\fnote{92}{(**)}{本注接续前一天的“团结”一注(n\textsuperscript{o}~85)的反思。}为了多少算上同调学家的读者的教益,也尤其为了我自己,我想把这场对一个辉煌研讨班的按部就班的劫掠——它落在我的两位前上同调学家学生手里,而在其余人的善意的目光之下\fnote{93}{(***)}{反思的后续将显明:这些“其余人”中的一位,曾为这场战役有效地出手,代人行事。}——的细节巡视一遍:正是在那场研讨班里,他们比全世界早十二年、从工匠本人手里,学得了使他们成名的手艺的基础与精微之处。

我的两场口头报告,从未以任何形态交到公众手里。其一是收官报告,论未决问题与猜想——它“不幸没有成文”,鉴于其分量之轻;而屠杀版的导言作者也确实认为,连提一提所论的是\textbf{哪些}未决问题与猜想,都是多余的。他又何必费那份心呢——那不过是些问题(人人尽可随意自设!)与猜想(连证明都没有!)(注87$_{1}$)。其二是开篇报告,它把研讨班一举安放进一个更宽广的语境(拓扑的、复解析的、代数的),并巡视了Euler--Poincaré、Lefschetz、Nielsen--Wecken各型公式——其中一些,构成研讨班的主要应用之一。导言作者用“……也如同……”一笔带过这场报告的消失,那种措辞对一份\textbf{漫不经心}的心境说得很多;在当时——研讨班的作者已从流通中消失七年之际——这种心境分明是理所当然。

有整整一组报告,是我做的:关于与一个闭链相配的同调类与上同调类的形式主义(上同调情形取正则环绕概形)\fnote{94}{(****)}{细节见注n\textsuperscript{o}~82“好的参考文献”。}。它们被公平地瓜分了:上同调归Deligne,同调归Verdier——后者还是稍微越界到上同调;作为回报,则不吝以那大名鼎鼎的“权复形”向Deligne道个小小礼\dmark{}\origpage{357}\fnote{95}{(*)}{见注83“玩笑——或权复形”。}。(更不必说他还把$R\mathcal{H}om$的有限性定理与双对偶定理据为己有——一字不差地抄自研讨班;反正,狮子的份额归Deligne,那是理所当然……)导言作者认为连提一提同调的报告也无须。确实无须:前一年,他的朋友Verdier已揽下补足那份缺失的“好的参考文献”的差事(对一场研讨班、对我,均不置一词)。

有关于$R^if_*$(f非真)运算之有限性定理的口头报告,以及作为推论,关于$R\mathcal{H}om_*$与$Lf^!$运算的报告。关键定理以Hironaka式奇点消解的技术证明(因而只在拥有消解的情形有效)。我所使用的这些论证,自研讨班以降已成通例(见注85$_{2}$)。Deligne得以在另一些更便利于应用、如今在大多数应用里得到验证的假设下,证明这些有限性定理,以及双对偶定理。人们本可以指望:他会要求把这些改进纳入那场他有幸学得平展上同调、学得其后来全部著作之基础的观念与技术的研讨班。但这一情势,却被拿来当作“理由”,把研讨班的这一部分截去。至于双对偶定理,从此在Illusie的笔下(在概形框架里)成了“Deligne双对偶定理”(报告I的导言)。这也算公道:解析情形里,Verdier前一年就早已把父权判给了自己(甚至无须费心另找一份证明)。

有一场报告,发展一条“泛Künneth公式”,它曾由Illusie成文。此前没有人想到过提炼这一类陈述;其灵感在于这样的直觉:“一般地”——即在基的一般点附近——一个相对概形表现得如一个“局部平凡纤维”(在拓扑语境里)。经由一条与前面提到的证明相仿的优雅证明,Deligne得以消除我所做的奇点消解假设。判给了——报告删除,以对同一位Illusie在那场号称“在前”的研讨班SGA~4½里的报告的一条引证“取而代之”。

\dmark{}\origpage{358}有一组报告,论非交换迹的形式主义;它是作为手段发展的,用以在从未被处理过的情形里显豁Lefschetz--Verdier公式的局部项。据说,这些报告最终由Bucur成文,其手稿“在一次搬家中天赐般地丢失了”——这简直成了通俗笑剧!\fnote{96}{(*)}{想必正是这一情势,当场启发了Deligne对SGA~5的那条光辉批评:说Lefschetz--Verdier公式(别忘了——它“还停留在猜想阶段”!!!)的局部项在那里竟连计算都没有!(关于这条批评的离奇——在知情的读者看来,它与Verdier前一年那条大名鼎鼎的“权复形”难分伯仲——见“白板”一注,n\textsuperscript{o}~67。如此一来,成风的是Verdier!)}而在Illusie执笔的SGA~5导言里,这些报告居然变成了“Grothendieck的\textbf{交换}迹理论,以光辉方式推广了Stallings的理论”(而它们分明是\textbf{非}交换的!)。这个口误\fnote{97}{(**)}{即把“非交换”迹理论之父权派给我的那个口误——那顶王冠,倒也没人等我来戴。它竟然一直保留到已刊版本里,越发引人注目:Illusie是我学生中工作上最一丝不苟、直到最后细节的一位,想必。}只能出自一位秘书欠佳(或太好……)的灵感;她想必与我朋友Ionel~Bucur的搬家公司过从甚密。(“光辉地”一词是我的笔插进去的,为的是更好地还原这个同样天赐的口误所万无一失地暗示的思想。)

我无从抱怨,既然Illusie把重做的活揽在了自己身上(而且,照他的说法,还是个“更精致的”版本,因为这按层论的方式调了味——不过我记得,Illusie,在我那个年代,你做过的“精致”创新,可比这一手多……)。想必他费了不少时日——如果我没记错,当年我把这台机器调试好,花过几个星期;说不定我的手稿也在同一场天赐的搬家里丢了,天知道哪位被我的滔滔不绝淹没了的亲爱听众,至少记下过能读懂的笔记……

有一件事值得注意,我此前不曾留意:他没有把这场报告插回它原定的报告XI的位置(那想必也是它在口头研讨班里的位置),而宁肯在那个位置留下一个大窟窿,把他的报告做成一场伪托的报告,称为“局部项的计算”。这标题倒分明对应我记得自己在口头研讨班里做过的东西——奇怪。但就在他为这场报告所写导言的第1行,作者赶紧为我们\dmark{}\origpage{359}澄清误会:“这份报告,成文于1977年1月,\textbf{不对应于研讨班的任何一场口头报告}”。随后接以Lefschetz--Verdier公式(这名字分明对我说过什么;而我记得自己分明纵横发展过一套非交换迹理论,正是为了在某些情形里计算“局部项”……),继以一条Langlands公式,以及一份1967年的Artin--Verdier证明(那毕竟是口头研讨班“终曲”之后的一年;它不可能不影响到这两位作者——其中至少一位,或许两位都听过研讨班)。到这一页的末尾,我们才顺带得知——与开头的宣告相反——还有“这份报告的第二部分,性质技术得多”(这种语言我在哪里读到过……),它(请欣赏这个措辞)“\textbf{受Grothendieck为建立曲线上某些上同调对应的Lefschetz公式所用的方法的启发}”,附一处对同一场研讨班报告XII的引证,尤其对不可或缺的SGA~4½的引证;分明,为了这么点东西,没有理由把它插回大窟窿的位置——方才那份“更精致的版本”想必已尽善尽美。Illusie与Deligne还肯以“启发”之源的礼遇引我,可谓客气;而他们朋友Verdier前一年的榜样,分明早已表明:这份顾忌,完全用不着了。

我回到Illusie为顶着SGA~5之名的那一卷所写的导言。我们在那里再度得知——一如Deligne在其SGA~4½导言里已宣告的——研讨班之所以终于出版,全靠\textbf{了他的朋友}:

\begin{quote}
“我感谢 P.~Deligne说服我,在报告III的新版里,写出Lefschetz--Verdier公式的一份证明,\textbf{从而排除了这份研讨班出版的障碍之一}”。
\end{quote}

我们又一次深陷闹剧——那位顺从的Illusie,在SGA~4½的导言里照单全收!倘若研讨班十多年不曾出版,那是因为(难就难在想得到)没有一个人(在1977年救场的Deligne之前)想到过:为那条(不无道理地)称为“Lefschetz--Verdier”的公式写一份证明,也许是个好主意——而那份公式的父权,除了他那位\dmark{}\origpage{360}形影不离的朋友、我的前学生Verdier亲自,\textbf{至少从}1964年起(注87$_{2}$)就骄傲地佩戴着;也就是说,至少早在我那场研讨班结束之前两年——它当时只待好心的志愿者,便能交到人人手中!

最后,作为研讨班的又一场(最后的?)残损:Serre就“(Serre--)Swan模”所做的漂亮报告消失了——那份报告题为“Brauer理论导引”。所幸Serre,眼见事态的走向,明智地把那份报告收进了他的书《有限群的线性表示》(Hermann,1971),把它交到数学公众手里。(注87$_{3}$)

这一回,我想,这幅画我巡视完了。这幅画,关乎一场研讨班的命运——我曾把自己的最好部分倾注其中(注88)\fnote{98}{(*)}{关于“自己的最好部分”这一说法的意味,见随后的“遗骸……”“……与身体”两注,n\textsuperscript{o}~88、89。其中第一条,把SGA~5研讨班——连同与之密不可分的SGA~4——安放为我“彻底完成”的著作部分的主干。};二十年后重逢,它已面目全非,被那些曾是它独家受益人的人屠杀——或者至少,被其中三人,并获得所有其他与会者的首肯。

我并不后悔再一次费心把渐渐强加于我注意力的东西追到底。我曾在对我的一位前学生的关系的长篇回顾之末,确认这场“事物的回归”\fnote{99}{(**)}{见同名注(4月30日,n\textsuperscript{o}~73)。},当时已隐约预感到:他不是唯一一个“兴冲冲为我送葬”的人——而只是到了此刻,我才得知它的气息、它的“气味”(借用当时我的一场梦里出现的说法)——一场\textbf{暴力}的气息。这份气息,既被那种呈献高度技术性实体的(表面超然、不动声色的)话语所掩藏,又被它所泄露\fnote{100}{(***)}{尤其指伴随SGA~5(Illusie执笔)与SGA~4½(Deligne执笔)的那些导言性文字里的话语。}。这场暴力所瞄准的,经由一副任人宰割的“遗骸”,恰恰是那位曾是“老师”、“父亲”的人本身——而与此同时,“学生们”久已占据了他那令人艳羡的位置,不曾遇到任何抵抗;并且\dmark{}\origpage{361}久已在他们中间选出新的“父亲”,注定取代旧者、统治他们。

我感到这份气息,然而它于我仍是异己的、未被理解的。要“理解”它,想必这份气息得在我身上活着,或曾在我身上活过。但四年前,我头一回感到并掂量了我生活中一桩我从未想过、总以为理所当然的事:即我童年时代对父亲的认同,\textbf{不曾}带有冲突的印记——在我童年的任何时刻,\textbf{我既不怕我的父亲,也不嫉妒他},同时对他怀着毫无保留的爱。这层关系——或许是给我的一生打下最深刻印记的关系(在这场四年前的禅观之前,我甚至不曾觉察它)——在我的童年,如同对另一个既是强力的又是慈爱的自我的关系——这层关系,不曾打上分裂与冲突的印记。倘若在我那常常撕裂的一生里,对我身上所怀之力量的认识始终活着;倘若在我那绝不免于恐惧的一生里,我既不惧怕任何一个人,也不惧怕任何一件事——我把它归功于这份寒微的情势;直到我年过五十,它仍不为人知。这份情势是一件无价的特权,因为对自身之内创造之力的亲密认识——它\textbf{就是}那股力量——正使力量得按其本性自由表达:经由创造,经由一种创造的生活。

而这份特权——它使我免于冲突最深的那几道印记之一——此刻却像一副枷锁,像我生活经验里的一个“\textbf{空缺}”。一个难以填补的空缺,在那里,许多别人拥有情感、意象、联想的丰富织体,向他们提供道路(只要他们好奇肯走),既深刻理解他人、也深刻理解自己;面对那些境况——我靠反复与旁证才能勉强把握,却仍像一个异乡人——我身上求知的欲望,忍饥挨饿。

\notehead{87$_{1}$}{}

(5月31日)这场收官报告——与开篇报告一起,想必是最有趣、最实的一份——显然并非对所有人都白费,如我读到MacPherson的文章“Chern classes for singular algebraic varieties”(Annals of Math. (2) 100, 1974, 第423--432页;1973年4月收到)时所看到的。我在其中,以“Deligne--Grothendieck猜想”之名,重新找到我在这场报告里、于概形框架中引入的主要猜想之一\dmark{}\origpage{362}。MacPherson在复数域上代数簇的超越框架里重述它,Chow环换作同调群。Deligne1966年在我那场报告里学得这条猜想\fnote{101}{(*)}{(6月6日)诚然,形态略有不同;见本注今日所写的后续。\par (1985年3月)Deligne本人提供的细节,见“把一个个 i 点清”一注,n\textsuperscript{o}~164(II~1)。}——那恰是他在那场研讨班里初次露面、开始熟悉概形语言与上同调技术的一年(见“格格不入者”一注,n\textsuperscript{o}~67$'$)——几年之后,把我算进猜想的命名里,想必就不再合宜了……

(6月6日)我借此机会在此写明:我在研讨班里、于概形框架中所陈述的猜想是什么;当时我肯定还指出了它在复解析框架(乃至刚解析框架)中的明摆着变形。我把它构想为一个“Riemann--Roch”型的定理,只是取离散系数而非凝聚系数。(Zoghman~Mebkhout还对我说过:他的D-模观点应当使人可以把这两个Riemann--Roch定理都看作包含在同一条晶状Riemann--Roch定理之中;后者因此在特征零,代表我在数学里引入的两条Riemann--Roch定理——一条在1957年,另一条在1966年——的自然综合。)取定一个系数环$\Lambda$(不必交换,但为简单设为诺特的;且为平展上同调之需,其挠性与所论概形的特征互素……)。对一个概形$X$,以
$$K.(X,\Lambda)$$
记由$\Lambda$-模的平展可构造层构成的Grothendieck群。利用函子$Rf_!$,该群函子地依赖于$X$(对诺特概形$X$与分离有限型的概形态射)。对正则$X$,我假设存在一个典范的群同态,扮演凝聚RR定理中“Chern特征”的角色\dmark{}\origpage{363}:
\begin{equation}\tag{15.1}\label{eq:discrete-chern-map}
ch_X: K.(X,\Lambda)\to A(X)\otimes_{\mathbb{Z}} K.(\Lambda),
\end{equation}
其中$A(X)$是$X$的Chow环,$K.(\Lambda)$是由有限型$\Lambda$-模构成的Grothendieck群。这个同态,应当由“离散Riemann--Roch公式”对正则概形间一个\textbf{真}态射$f:X\to Y$的成立而唯一确定;该公式写作与凝聚Riemann--Roch公式同样的样子,只是把Todd“乘子”换成相对全Chern类:
\begin{equation}\tag{15.2}\label{eq:discrete-rr-formula}
ch_Y(f_!(x)) = f_*(ch_X(x)\, c(f)),
\end{equation}
其中$c(f)\in A(X)$是$f$的全Chern类。不难看出:在一个拥有Hironaka强形式奇点消解的语境里,RR公式确实唯一确定诸$ch_X$。

当然,这预设我们处在Chow环有定义的语境里。(我不知道有谁哪怕尝试过,为并非有限型于一个域上的正则概形写一套Chow环理论。)否则,也可以在凝聚语境里惯常的“Grothendieck”环$K^{\circ}(X)$所配的分次环里工作,按惯常方式过滤(见SGA~6)。还可以把$A(X)$换成偶数次的$\ell$-进上同调环,即诸$H^{2i}(X,\underline{Z}_\ell(i))$的直和。这有一个缺点:引入一个人为参数$\ell$,并给出较“纯数字”的公式;而Chow环的魅力恰在其连续的结构——过渡到上同调时被摧毁。

早在$X$是代数闭域上光滑代数曲线的情形,$ch_X$的计算就牵涉Artin--Serre--Swan型精致的局部不变量。这就是说,那条一般的猜想是一条深刻的猜想;对其的追求,系于对这些不变量的高维类似物的理解。

评注。同样以$K^.(X,\Lambda)$记由Tor维数有限的$\Lambda$-平展层可构造复形构成的“Grothendieck环”(当$\Lambda$交换时,该环作用于$K.(X,\Lambda)$……),人们同样应有一个同态
$$ch_X: K^.(X,\Lambda)\to A(X)\otimes_{\mathbb{Z}} K^.(\Lambda),$$
它同样(mutatis mutandis)引出同一条Riemann--Roch(RR)公式。

\dmark{}现在设$\mathrm{cons}(X)$为$X$上可构造整值函数的环。人们以或多或少许套套逻辑的方式定义典范同态
$$K.(X,\Lambda)\to \mathrm{Cons}(X)\otimes_{\mathbb{Z}} K.(\Lambda),\qquad K^.(X,\Lambda)\to\mathrm{Cons}(X)\otimes_{\mathbb{Z}} K^.(\Lambda).$$
倘若现在我们只限于\textbf{特征零}的概形,那么(利用具真支撑的Euler--Poincaré特征)可以看出:群$\mathrm{Cons}(X)$对诺特概形的有限型态射是共变函子(此外,作为环函子,它是反变的——这与特征无关),而前面的套套逻辑态射是函子的。(这对应于一个“众所周知”的事实——我想,在SGA~5口头研讨班里并不曾证明——即\textbf{特征零}时,对代数概形$X$上一个局部常值的$\Lambda$-模层$\mathcal{F}$,其经由$f_!:K.(X,\Lambda)\to K.(e,\Lambda)\simeq K.(\Lambda)$的像等于$d\chi(X)$,其中$d$是$\mathcal{F}$的秩,$e=\mathrm{Spec}(k)$,$k$为假定代数闭的基域……)这立刻提示:Chern同态(1)与(1$'$)应当可以从套套逻辑同态(2)与(2$'$)出发、复合一个“普适”Chern同态(不依赖任何系数环$\Lambda$)
$$ch_X: \mathrm{Cons}(X)\to A(X)$$
而导出;于是,RR公式的两个“带$\Lambda$系数”的版本,便在形式上显现为包含于一条在可构造函数层面上的RR公式之中——后者总以同样的形态书写。

当我们对固定基域(重新允许任意特征)上的概形工作,或更一般地对固定的\textbf{正则}基底概形$S$(比如$S=\mathrm{Spec}(\mathbb{Z})$)上的概形工作时,与惯常写法(自1957年以来熟悉的凝聚框架)最贴合的RR公式形态,经由引入乘积
$$ch_X(x)\, c(X/S) = c_{X/S}(x)$$
而得到(其中$x$在$K.(X,\Lambda)$或$K^.(X,\Lambda)$中皆可);人们可以称之为$x$相对于基$S$的\dmark{}\origpage{365}\textbf{Chern类}。当$x$是$K^.(X,\Lambda)$的单位元、即取值$\Lambda$的常值层的类时,我们得到$X$相对于$S$的相对全Chern类在典范同态$A(X)\to A(X)\otimes K^.(\Lambda)$下的像。这么规定之后,RR公式等价于:对$S$上(有限型的)变动的正则概形$X$、固定$S$,这些相对Chern类的作法
$$c_{X/S}: K.(X,\Lambda)\to A(X)\otimes K.(\Lambda)$$
对真态射是函子的,(5$'$)的变形亦然。特征零时,这化归为相应映射
$$c_{X/S}: \mathrm{Cons}(X)\to A(X)$$
对真态射的函子性。正是在“绝对Chern类”映射(6)的存在与唯一这一形态下——在$S=\mathrm{Spec}(\mathbb{C})$的情形——该猜想呈现于MacPherson的工作之中;相关条件(这里与特征零的一般情形一样)是:a)(6)对真态射的函子性;b)$c_{X/S}(1)=c(X/S)$(在这里,即“绝对”全Chern类)。相对于我最初的猜想,MacPherson所呈献并证明的形态,有两点不同。其一是“减”:他不置于Chow环、而置于整系数上同调环——更确切地说,经由超越途径定义的整系数同调群。其二是“加”——也许正是在这里,Deligne对我最初的猜想做出了一份贡献(除非这份贡献出于MacPherson本人\fnote{102}{(*)}{(1985年3月)确实如此;见前一条脚注所引的注n\textsuperscript{o}~164。}):即对映射(6)的存在与唯一,无须把自己限于正则概形$X$,只要把$A(X)$换成整系数同调群。由此很可能:在一般情形亦然,只要以诺特概形$X$的\textbf{Chow群}(一般不再是环)充当$A(X)$(或更好:$A.(X)$)。或者换个说法:虽然$ch_X(x)$的启发式定义($x$在$K.(X,\Lambda)$或$K^.(X,\Lambda)$中)本质上地利用了环绕概形正则的假设,可一旦把它乘上“乘子”$c(X/S)$(当概形$X$有限型于固定正则概形$S$之上),所得乘积(4)似乎无须对$X$的任何正则性假设仍保有意义\dmark{}\origpage{366}——作为张量积
$$A.(X)\otimes K.(\Lambda) \quad\text{或}\quad A.(X)\otimes K^.(\Lambda)$$
中的元素,其中$A.(X)$记$X$的Chow群。MacPherson的证明精神(它不使用奇点消解)提示:(5$'$)的同态有可能有一个显式的“演算式”构造:“将就着”$X$的奇点本身,连同系数层$\mathcal{F}$(其类为$x$)的奇点,在$X$上“收取”一个取值于$K^.(\Lambda)$的闭链。这也合乎我1957年随凝聚Riemann--Roch定理引入的观念的精神:我在其中尤其做自交计算,小心不让所论的闭链“动”。一个显明的初步归约(经由把$X$嵌入一个$S$-概形)是$X$为正则概形$S$的闭子概形的情形……

此外,应当可能发展一条\textbf{奇异的}(凝聚)Riemann--Roch定理——这个想法我早已熟悉,说不出从何时起;只是我从不曾认真检验过它。某种程度上,正是这个想法(除去与“上同调、同调、卡积”形式主义的类比)曾在SGA~6里(1966/67)使我系统地引入$K.(X)$、$K^.(X)$与$A.(X)$、$A^.(X)$,而不满足于只用$K^.(X)$工作。我不记得1966年在SGA~5研讨班里是否也想过这类东西、是否在口头报告里暗示过。由于我的手写笔记已散失(也许在一次搬家里?),我恐怕永远无从知道了……

(6月7日)翻阅MacPherson的文章,有一件事打动了我:“Riemann--Roch”一词在其中未曾出现——这也正是我没有立刻认出我1966年在SGA~5研讨班里所做那条猜想的缘由;那条猜想,对我来说(现在仍是),是一个“Riemann--Roch”型的定理。写文章之际,MacPherson似乎甚至不曾觉察这层显然的亲缘。我料想缘由在于:Deligne——我离去之后,他以自己乐意的形态使这条猜想流通——曾尽力“抹去”它与Riemann--Roch--Grothendieck定理的显然亲缘。我自以为感到他如此行事的动机。一边,这削弱了这条猜想与我个人的联系,并使它目前流通的那个名称“Deligne--Grothendieck猜想”更可信。(附记:我不知道它在概形情形是否也在流通;倘若在,我很想知道顶着什么名称。)但我看更深的缘由,在于他那个执拗的念头:尽一切可能否认并摧毁我的著作、我的数学视野的根基处的统一性\fnote{103}{(*)}{可与“遗骸”一注(n\textsuperscript{o}~88)里关于“SGA~4½之役”之深层意味的评语对读:那场战役同样意在,经由外来文本SGA~4½对SGA~4与SGA~5这两个不可分割的部分——它们正发展着我的著作——的“强行塞入”,把我的著作围绕平展上同调的深层统一,炸成一堆无定形的“技术枝蔓”。}。这是一个触目的例子,说明在一个手段出众的数学家身上,一个与一切数学动机全然无关的执念,竟能使那种我称之为数学的“健全本能”的东西昏暗(乃至全然堵塞)。这份本能不可能不察觉一个“同一条”Riemann--Roch定理的“连续”与“离散”两个陈述之间的类比——我当年在口头报告里当然也强调过。如我昨天所写,这份亲缘想必不久将由一个成形的陈述(Zoghman~Mebkhout所猜)确认——至少在复解析情形——使人得以从一条共同的陈述导出两者。显然,在Deligne对着Riemann--Roch定理的那份“掘墓”心境里\fnote{104}{(**)}{这些心境——恰对着Riemann--Roch--Grothendieck定理——在“葬礼悼词”里表现得格外清楚;见“葬礼悼词(1)——或恭维”一注,n\textsuperscript{o}~104。},他根本不可能发现那个在解析框架里联结两者的唯一陈述,更无从在一般概形框架里提出类似陈述的问题。一如他在这样的心境里,无从提炼D-模这个丰产的观点——它以过于自然的方式,出自他正要埋葬的观念——甚至在多年之间,无从承认Mebkhout那项在他失败之处成功的丰产工作。

\notehead{87$_{2}$}{}

(5月31日)那是1957年——我的Bourbaki报告论L函数有理性之年;我在其中启发式地使用Verdier的结果(???)(尤其该个案里局部项的预期形态),而不等Illusie十三年后——应Deligne之邀——才肯去证明它。而且我觉得,当年Verdier把他那条超\dmark{}\origpage{368}一般的公式拿给我看——它来得如一场意外——时,他是用“六种运算”形式主义在几行里证明它的:这类公式,(几乎)写出来就是证明了它!倘若有什么“困难”,至多在核验一两处相容性的层面上\fnote{105}{(*)}{(6月6日)此外似乎:经由双对偶定理(其间已晋升为“Deligne定理”),Lefschetz--Verdier公式的初始证明曾依赖一条奇点消解假设;Deligne得以在一域上有限型概形的情形摆脱它。这正是浑水摸鱼的大好机会,给人这样的印象:SGA~5竟从属于那个“在先(原文如此)”的“研讨班SGA~4½”(而且它确实先出版了!)。}。况且,无论Illusie还是Deligne,都一清二楚:我在研讨班里为各种显式迹公式所给的证明\textbf{是完整的};它们绝不依赖Verdier那条一般公式——后者只扮演了“扳机”的角色,怂恿人们在尽可能一般的情形里显豁并证明迹公式。两人在这里的不诚,昭然若揭。对Deligne,我写“白板”一注(n\textsuperscript{o}~67)时就已清楚——但对不知情的读者,当然不清楚;同样不清楚的,还有那种放弃使用自己健全官能的知情读者。

(6月6日)至于Illusie,他完全入了其朋友的戏,试图搅浑水,摆出一副超技术的口头研讨班的样子——仿佛它连全部结果的完整证明都给不出,尤其迹公式。然而那些公式在65/66年间\textbf{确实}得到证明(而且是头一回),恰恰是在那里,他与Deligne有幸学得它们,以及与之相随的一整套精微技术\fnote{106}{(**)}{在顶着SGA~5之名出版的那一卷的导言第二段里,Illusie把围绕平展上同调Lefschetz公式的三场报告III、III~B、XII当作“研讨班的心脏”;而人们看到,在报告III~B的导言里,他小心注明(与事实相反)“这份报告不对应于研讨班的任何一场口头报告”;在报告III与III~B的导言里,他尽力给人这样的印象:它们从属于SGA~4½,而报告III被呈献为“猜想性的”!!实际上,SGA~5研讨班的全体,在技术上独立于报告III(Lefschetz--Verdier公式)——后者扮演启发性动机的角色;而报告III~B不是别的,正是Bucur那次搬家所造成的“窟窿”(报告XI);那窟窿,成了这场进一步肢解的可喜口实。为使人相信那场“技术枝蔓”研讨班的版本(其朋友Deligne的授意),Illusie小心砍掉了开篇报告——我在其中为将在这场研讨班里发展的各大主题画了一幅初步的全景;迹公式在其中只占一小角(因其朝Weil猜想方向的算术含义而取得特殊重要性)。关于这些“大主题”的概观,见后面的子注n\textsuperscript{o}~87$_{4}$(原书误作87$_{5}$)。}。

\dmark{}这使我想起:当然,我曾费心在研讨班里证明Lefschetz--Verdier公式\origpage{369}——那是最起码的,也是我打算发展的局部与整体对偶形式主义的一次格外触目的应用。这几天问题来到我心头:既然分明有十来份报告的成文工作,在我亲爱的学生们手里陷于困境,Deligne与Illusie要为其SGA~5出版的“障碍(原文如此)”命名,本有挑选的余地——为什么偏偏选中他们的好哥们Verdier的定理——他当时正把父权当作应得之物佩戴着,一如他佩戴导出范畴与三角范畴的父权,而那些他也从不曾费心成文(或至少交到公众手里)。这里有一种\textbf{荒谬之中的挑衅}(或我的前上同调学家学生群体里的某种集体犬儒——我认为他们在那场屠杀战役里全体同谋),它使我想起Verdier前一年光辉地发明的“权复形”(见同名注,n\textsuperscript{o}~83),或(在不义的音区里)Deligne给本该叫“Mebkhout的层”的那些层所取的名字“反常”(见“反常性”一注,n\textsuperscript{o}~76)。我在这样的发明里感到,是对整个数学共同体的支配与轻蔑之举——同时是一份\textbf{赌注};它分明一直赢到亡者不期现身之时——他俨然像一群沉睡者中间唯一醒着的人……

\notehead{87$_{3}$}{}

(6月5日)一场屠杀盘点之后,人们自会掂量Illusie在其顶着SGA~5之名的卷子的导言第2行里的这份自白:

\begin{quote}
“相对于最初的版本,唯一重要的变更涉及报告II[泛Künneth公式],未予复制;以及报告III[Lefschetz--Verdier公式],已全部重写,并增补一个编号III~B的附录\fnote{107}{(*)}{它被呈献为“研讨班的心脏”的一部分!(见前一条脚注。)}。除若干\dmark{}\origpage{370}细节修改与脚注增补之外,其余报告一仍\textbf{其旧}”(着重号是我加的)。
\end{quote}

这里,Illusie又一次甘当其那位罄竹难书的朋友另一句俏皮话的应声虫,即:SGA~4½的存在“不久将使人得以\textbf{照原样}出版SGA~5”(见“白板”一注,n\textsuperscript{o}~67)——而Illusie在其报告与导言里,处处尽力为这一冒名张目——即:SGA~5(他与朋友学艺之地)竟依赖于那部海盗文集SGA~4½,由其后十二年间拾得或劫掠的破烂与断枝拼成——每页转弯处都不吝对SGA~4½的大量引证……

最后的字,照例归Deligne:一个月前(5月3日)他给我写信,回复一份简短的信息之请(见“葬礼”一注,n\textsuperscript{o}~70的开头):

\begin{quote}
“总之,倘若SGA~4½这份文本问世时,你已七年不做数学[?!],那只是相应于SGA~5出版[?]的漫长迟延,\textbf{它太不完整,无法照原样有益地出版}。

希望这些解释令你满意。”
\end{quote}

它们即便不曾“令我满意”,至少令我大开眼界……

\notehead{87$_{4}$}{}

(6月6日)或许是时候指出:口头研讨班里发展过的主要主题有哪些;已刊文本,只有靠旁证才能对其形成概念。

I)对偶理论的局部侧面;其本质技术配料是(一如凝聚情形)双对偶定理(补以一条“上同调纯粹性”定理)。我的印象是:这后一条定理的几何意味——作为一条局部Poincaré对偶定理——我当年在口头研讨班里分明解释过,却被我的学生们从此全然忘却\fnote{108}{(*)}{核实之后:这一几何诠释,至少在Illusie的成文里得到了保留。}。

II)迹公式,包括比惯常迹公式(其两边为整数,或一般为某系数环——如$\underline{Z}/n\underline{Z}$或$\ell$-进环$\mathbb{Z}_\ell$、乃至$\mathbb{Q}_\ell$——的元素)更精微的“非交换”迹公式;它们置于一作用在所论概形上的\textbf{有限}群的代数之中,系数取于一个合宜的环(如前一括号里所设想者)。这一推广来得十分自然:因为即便在惯常型的Lefschetz公式里,对“扭转的”系数层,人们也被引至以一个Galois覆盖(一般分歧)替换原概形,以“解扭”系数,而Galois群作用于其上。如此,“Nielsen--Wecken”型的公式,便自然地进入概形语境。

III)Euler--Poincaré公式。一边,借助Serre--Swan模,对代数曲线的一条“绝对”公式作了周详研究(推广了驯分歧系数的情形;该情形给出较朴素的Ogg--Chafarévitch--Grothendieck公式)。另一边,有未发表的、深刻的“离散”Riemann--Roch型猜想;其中一条,七年后以一种杂糅的版本重现,顶着“Deligne--Grothendieck猜想”之名,由MacPherson经超越途径证明(见注n\textsuperscript{o}~87$_{1}$)。

我对这两大主题(Lefschetz公式、Euler--Poincaré公式)之间深刻关系所能做的评语,也一并散失,未留痕迹。(一如我的习惯,我把全部手写笔记留给了“志愿者们-执笔者(原文如此)”;口头研讨班的手写笔记,我本来有一整套——尽管有些简略——如今片纸无存。)

IV)与一个闭链相配的同调类与上同调类形式主义的详论,自然出自一般对偶形式主义与那条关键想法:在所论闭链上做“具支撑”的上同调,利用上同调纯粹性定理。

V)有限性定理(包括泛型有限性定理)与任意支撑上同调的泛Künneth定理。

研讨班还发展了一套从挠系数过渡到$\ell$-进系数的技术(报告V与VI)。那是研讨班最技术的部分;研讨班一般用挠系数工作,事后再“取极限”以导出相应的$\ell$-进结果。这一观点是一个临时的\dmark{}\origpage{372}权宜之计,且待Jouanolou的论文(至今未出版!)直接在$\ell$-进框架里给出所需的形式主义。

我不把对几个经典概形的一些计算、以及Chern类的上同调理论,算进“主要主题”;Illusie在其导言里把它们捧为研讨班“最有趣的”之一。由于日程已满,我在口头研讨班里不曾认为有必要在这些计算与这一构造上流连:只需把十年前我为Riemann--Roch定理、在Chow环的语境里给出的论证,几乎一字不差地搬来即可。另一方面,明摆着应当把它们纳入成文的研讨班,以便给平展上同调的使用者一份便利的参考。Jouanolou揽下了这项工作(报告VIII);他本该把它看作一项美差——既服务于数学共同体,又为自己学到一些日后必用的基本技术——却把它当作苦役:其成文一拖数年\fnote{109}{(*)}{(6月12日)翻阅那份报告,我进而确信:Jouanolou与我的其他上同调学家学生之间,有完美的默契。}。想必他的论文也如出一辙:至今仍是一份幽灵参考,一如Verdier的那份……“极限过渡”那一部分也不该算作研讨班的“主要主题”之一,因为它不与某个特定的几何想法相连。毋宁说,它反映平展上同调语境(区别于超越语境)特有的一项技术复杂:平展上同调的主要定理首先关涉挠系数(与剩余特征互素);要有对应于特征零系数环的理论——Weil猜想所必需——必须对$\mathbb{Z}/\ell^n\mathbb{Z}$型系数环取极限,以获得“$\ell$-进”结果。

把这一切说清之后,口头研讨班五个主要主题里,唯一一个似乎以完整形态出现在已刊文本里的,是主题I。主题IV与V干脆消失,被SGA~4½吸收,附带的收益是可以大量引证它、给人这样的印象:SGA~5依赖一份自居在先的Deligne文本。主题II与III以残缺形态出现在已刊卷里,并始终维持对SGA~4½文本的同一冒名(那份文本,实际上整个出自母研讨班SGA~4、SGA~5)。

\subsection{遗骸……}

\notehead{88}{}

\dmark{}(5月16日)SGA~4与SGA~5这两场前后相续的研讨班——对我而言,它们如\textbf{同一场}“研讨班”——从虚无中同时发展出:topos的\textbf{语言}——强大的综合与发现的工具;以及平展上同调——一件打磨停当、效率完美的\textbf{工具};其本质形式性质,彼时甚至比寻常空间的上同调理论理解得还透\fnote{110}{(*)}{即便把自己限于最贴近“簇”的空间,如可三角剖分的空间。}\origpage{373}。这一整体,代表我在数学里、就彻底完成的工作而言所能做出的最深刻、最革新的贡献。同时——它并无意于此,虽然一切始终以事情之明摆着的自然进行——这项工作代表我作为数学家在其著作中所完成的规模最大的技术壮举\fnote{111}{(**)}{若干困难或意外的结果,是由别人(Artin、Verdier、Giraud、Deligne)得到的;工作的某些部分是与别人合作完成的。这丝毫不减损(至少在我心里)我对这项工作在我全部著作中之地位的评判的分量。况且,我想在《主题概要》的一篇附录里,更周详地回到这一点,在分明已属必要之处把点补在i上。}。这两场研讨班,对我是密不可分的。它们在其统一之中,同时代表\textbf{视野}与\textbf{工具}——topos,与一套完整的平展上同调形式主义。

视野至今仍遭拒斥,而工具近二十年来深刻革新了代数几何其最令我心醉的一面——“算术”的一面:经由一种直觉、一套“几何”性质的概念与技术行装所把握的算术。

肯定不只是为了暗示其上同调“摘要”先于SGA~5部分,Deligne才给它戴上SGA~4½这个欺眼之名——反正要戴,大可以叫SGA~3½!在“SGA~4½之役”里,我感到一种意图:把那部他的全部著作皆出于斯的著作(那部他始终无法摆脱的著作!)——其分明而深层的统一,在SGA~4与(真正的)SGA~5两场研讨班的整体里昭然可见——呈献为一副\textbf{分裂}的面目(一如他自己也是分裂的……),被一份外来的、轻蔑的文本的强行塞入\textbf{一分为二};那份文本,妄想显出是活的心脏、思想与\dmark{}\origpage{374}视野的精髓——而它在那视野里不曾有过任何份额\fnote{112}{(*)}{这份思想——就其主导观念与本质结果而言——早在年轻的Deligne登台、来与我接触学代数几何与上同调技术(1965至1969年)之前,就已完全成熟。(5月30日)对此可参见“格格不入者”一注,n\textsuperscript{o}~67$'$。};而环绕它的两个“断片”,则成了近乎怪诞的附录,一堆对那部自居为中心与本质的著作——出自Deligne的笔下,而我寒微的身量则被客气地接纳(在彻底下葬之前)跻身“合作者”之列——的“枝蔓”与“技术补充”\fnote{113}{(**)}{见“绿灯”与“颠倒”两注,n\textsuperscript{o}~68、68$'$。}。

“偶然”把事情安排得妙极。这副“任人宰割的遗骸”——那场被“执笔者们”一直弃置不顾的“倒霉研讨班”,我离去时落在我的上同调学家学生们手里、听凭其处置的——那可不是老师著作中\textbf{随便}哪一部分!既不是SGA~1与SGA~2(我在其中独自发展、尚不自知那些即将成为未来主作“起飞”的两件不可或缺的技术辅助工具的东西),也不是SGA~3(我的贡献主要是为使概形的“四面八方”技术纯熟而做的无穷音阶与琶音——有时颇费力气),也不是SGA~6(系统发展我围绕Riemann--Roch定理与相交形式主义的十年旧观念),乃至SGA~7(按一场反思的内在逻辑,它从对中心工具的持有、对上同调的掌控中流出)。分明是我著作的\textbf{主干部分}——其成文工作尚未完成(而且是经他们的手……)——我把它,至少部分地,留在了我的上同调学家学生们手里。他们选择屠杀的,正是著作的这一主干;他们据为己有的碎块,忘了那使其获得意味、美与创造之德的那份统一(注90)。

同样不是偶然:他们手执七拼八凑的工具、否认使其从虚无中诞生的精神与视野,竟无一人能在革新之著作重生之处——逆着他们的冷漠与轻蔑——辨认出它。也无从解释:六年之后,当那件新工具终于被Deligne把握,他们竟全体一致地埋葬了那位在孤独中创造它的人——Zoghman~Mebkhout\dmark{}\origpage{375},被否认的老师之遗孤学生!同样不是偶然:在Deligne最初的冲劲回落之后(那冲劲,在几年间把他带向一套新Hodge理论的强力启动,与Weil猜想的证明),尽管他手段惊人、我的上同调学家学生们手段出众,我今天仍确认这场“阴郁停滞”——在一个无比丰饶、一切似乎尚待去做的领域里。这无须惊讶:近十五年来,灵感的主要来源与某些“大问题”\fnote{114}{(*)}{这份“灵感的主要来源”,当然是“Motif瑜伽”。它只在Deligne一人身上起过作用——他把它留在怀里独享“收益”,且以一种被剥除其大半力量的狭隘形态,拒斥这部瑜伽的某些本质方面。在由它激发、却遭忽视或被悄悄抹黑的“大问题”之中,我如今(虽不折不扣是局外人)看到:标准猜想;以及面向一切惯常类型的系数——或多或少贴近“Motif”本身(它们对后者扮演“普适”系数的角色:生出所有其他系数的那种)——的“六种运算”形式主义的发展。可与“我的孤儿们”一注(n\textsuperscript{o}~46)里的评语对读。},即便它们在场、人们每走一步都与它们遭逢,却被小心绕开、藏掖起来,如那位十五年来人们不断要埋葬的人的信使一般。

\subsection{……与身体}

\notehead{89}{}

(5月17日)那份曾在我身上活着的念头、那种事物的视野——我曾以为已把它传递出去——我看作一个活的身体:健全而和谐,带着活物更新的力量、构想与生育的力量。瞧,这个活的身体变成了一副\textbf{遗骸},在人们中间被瓜分——这条腿或这半边身子,被妥善制成标本,在某一家充当战利品;另一条,被剁碎,在另一家充当击头的棍棒或回旋镖;还有一条,天知道,原样下到自家的汤锅里(我们已经不在乎这个了!)——其余的,统统可以烂在垃圾堆里……

这就是——形象化的措辞诚然,但在我看来,恰当地表达事情的某种实在——终于向我显现的图景。那根棍棒,至多会在这里或那里敲裂一两颗头颅\fnote{115}{(**)}{(5月31日)而且它还满可以用来证明某条“困难尽人皆知的”定理!};但这些四散的碎块——战利品也罢、棍棒也罢、家常汤也罢——将永远不具那个活的身体那份如此简单、如此分明的力量:那创造一个新生命的、爱的拥抱的力量……

(5月18日)这个活的身体、这副碎块随风四散的“遗骸”的意象,想必是在刚过去的一星期里在我身上渐渐成形的。\dmark{}它在我的打字机笔下呈现的那副滑稽形态,丝毫不意味着这意象是半点\textbf{虚构}、一场顺着话头即兴而来的略带阴森的滑稽表演。这意象表达一份\textbf{实在};当它经书面表述取得物质形态的那一刻,我被深深地触动。\origpage{376}这份实在,想必早在我“离去”以来的十四年里、甚至更早,已在此处彼处以断片的方式被我获知。一些信息的断片,起初由一份心不在焉、倾注别处的注意力,登记在一个表浅的层面——但它们全都指向同一方向,想必已在更深的层面,聚成某种意象——一幅未曾表述的意象,我不曾费心去认识它,那时我有更要紧的事。这份意象,在3月底以来持续进行的反思中、即这六七个星期里,大大丰富并明确了。更确切地说:零散的信息要素,终于经一份全然在场的、有意识的留心的检视,渐渐聚成\textbf{另一幅}意象——在检视与探查的思维这个较表浅的层面上;经由一项工作,它看似独立于第一幅意象在更深层的在场。这项有意识的工作,六天前在“屠杀”之发生的骤然视野中达到顶点——那时我感到了一份\textbf{暴力}的“气息”、“气味”;我想,这是整场反思里的头一回\fnote{116}{(*)}{(6月12日)近些年里,我曾有过这样的时刻:感到我的某些前学生对另一些与我“同葬者”怀有暴力的意图;但从不曾感到一种出自集体意志(在这里结成五人)、指向我这个人、经由我的著作而行的暴力。}。也是在这样的时刻,一个健全和谐之躯分明遭“屠杀”的感觉,想必已在临近表面的层次出现;而那幅更深的弥散意象,想必也已开始浮出水面,为那幅成形中的意象带来一层血肉——一种单凭思维无力给出的“气味”。

这份“血肉”的侧面,昨夜又在一场梦里显现——正是受这场梦的推动,我此刻回到昨天写下的文字。梦里,我身体的几处地方被割了相当深的口子。起初是嘴唇与口腔里的割口,大量出血;我对着镜子用大量清水漱口,漱口水被血染得通红。然后是腹部的伤口,也大量出血,尤其其中一处,血一鼓一鼓地涌出,仿佛是动脉(做梦者无心顾及解剖学的写实)。我甚至想到:倘若照这样流下去,我兴许会躺倒在地的;我用手按住伤口,蜷起身子止血——血确实止住了奔流,渐渐结成一个血块与一大片痂。后来,我小心地揭开那片痂:细微的愈合已经开始。我的一个手指也被割了,上面缠着厚厚一团绷带……

我无意对这个梦做更精微、更周详的描述,也不打算在这里(或别处)深探它。这个梦“照原样”已经以惊人的力量向我显示:昨天我谈的那个“身体”——我在写它时,把它看作与我分离的东西,像一个我或许孕育、生出、如今已走向世界走自己路的——这个身体,至今仍是我这个人亲切的一部分:那是\textbf{我的}身体,由血肉与一种生命的力量构成,这力量使它经得起深的创伤、并自我再生。而我的身体,想必也是世上与我系得最深、最不可解的东西……

做梦者没有让我停留于“屠杀”与瓜分遗骸的意象。那份意象要还原的,是我曾强烈知觉到的、\textbf{他人}身上的一份意图与心境的实在,而不是我自己经历这场侵凌、这场残损——经由一件我仍与之紧紧相连的东西——的方式。我与之相系到何等地步,做梦者刚刚让我隐约瞥见。这与我在“事物的回归(或一脚踩进菜盘)”一注(n\textsuperscript{o}~73)的反思里(以较弱的力度)所知觉的东西相合:在那里,我试着稍稍勾勒“构想了一件事的人与这件事之间的深深纽带”的那份感觉——它出现在那天的反思里。在那场4月30日的反思(仅仅三周之前)之前,在我的一生里,我曾装作无视这条纽带,或至少淡化它,顺着时兴老生常谈那条划好的坡道。操心一件已脱手的作品的命运——尤其操心我们的名字是否还多少附着于它——被感受为一种小家子气、一种悭吝;然而,人人却觉得自然:一个我们养大的(并且自以为爱过的)亲生孩子,倘若选择否认他出生时所得的名字,我们会深受触动。

\subsection{继承人}

\notehead{90}{}

\dmark{}(5月18日)我不知道六十年代里,可曾有哪个学生(Deligne除外)感到了这份本质的统一——超越他与我一起进行的有限工作。也许某些人模糊地感到过;而这份知觉,在我离去之后的年月一过,便永无挽回地失落了。可以肯定的却是:从1965年我们初次接触起,Deligne就预感到了这份活的统一。\origpage{378}正是这份对一个宏阔蓝图之统一旨趣的精微知觉,无疑是他身上那份强烈兴趣的主要兴奋剂——对着我有可传授、可传递的一切。这份兴趣,在1965到1969年这四年持续的数学接触里,从不曾衰减\fnote{117}{(*)}{这一时期包括五年,其中一年(1966),我的朋友在比利时服兵役。}。它赋予我们之间的数学交流以我谈过的那份卓越的品质;与别的数学家朋友,我只是在寥寥的瞬间知晓过它。正是这份对本质之物的知觉,以及它在他身上激起的那份炽热的兴趣,使他得以仿佛游戏般学得我能教他的一切:既有技术手段(锌丝概形的技术,Riemann--Roch与相交的瑜伽,上同调形式主义,平展上同调,topos的语言),也有使其统一的总体\textbf{视野},最后是\textbf{Motif瑜伽}——它当时是这份视野的主要果实,是我到那时为止所发现的最强的灵感之源。

清楚的是:Deligne是我的学生中直到今天仍唯一的一个,曾在某个时刻(我想自1968年起)把我可传递的一切,就其本质的统一与其手段的多样性,完全消化并据为己有\fnote{118}{(**)}{当我谈“全部”时,应理解为:就视野与手段中的本质事物而言。这当然不意味着,不曾有一些未发表的想法与结果,我从未想过向他谈起。相反,我不认为1965--69年间的任何一场数学反思,是我没有“趁热”向我的朋友谈过的——始终怀着乐趣,也始终有所获益。}。正是这份情势——想必人人都感到了——使他俨然我著作天造地设的“合法继承人”。显然,这份遗产既不使他累赘,也不限制他——它不是一副重担,而给他添翼;我是说:它以其活力滋养着他生而有之的那些“翅膀”\dmark{}\origpage{379}——一如别的视野、别的遗产(诚然没那么个人……)也将滋养它们……

这份他曾在成长与起飞的关键年月里赖以为养的遗产,以及使其具有美与创造之德的那份统一——他曾那么懂得去感觉、它已成为他自身一部分的东西——我的朋友此后\fnote{119}{(*)}{奇怪的是,这份分裂,想必在我们相识的头一年就已在场(它已经表现为对SGA~5研讨班的一份暧昧态度——那是他与概形、与Grothendieck风格的上同调技术、与平展上同调的初次接触),最迟在1968年、以毫不含糊的形态(见“驱逐”一注,n\textsuperscript{o}~63)——也就是数学交流完美无缺、他的数学思想的起飞在我看来尚未打上冲突印记的时刻。他当时(“顺带地”)做出过许多有趣的贡献(我在SGA~4的导言里极乐于把它们捧出来),题材却是他在我离去之后竭力要埋葬的。}予以否认,不懈地藏掖遗产,否认并摧毁那份作为其灵魂的创造性的统一。在我的学生中间,他是第一个做出榜样的人:一面据有工具、据有“碎块”,一面拼命肢解它们所自出的统一、那个活的身体。他自己的创造冲劲,由此被他身上这份深层的分裂所阻滞、吸收、最终肢解;这分裂驱使他否认并摧毁恰恰构成其力量、滋养其冲劲的东西。

我看这份分裂经由三个互相支援、密不可分的效果表达。其一是能量涣散的效果:在否认、肢解、取代、藏掖的努力中四散。其二在于对某些观念与手段的拒斥——它们对本是他选作中心主题的那个主题的“自然”发展本质不可少\fnote{120}{(**)}{这份拒斥,尤其表现为:导出范畴与三角范畴被埋葬(直到1981年),六种变换的形式主义被埋葬(直到今天),topos的语言被埋葬(亦然);以及对同调代数与同伦代数那套宏阔奠基纲领的一种“以轻蔑来封锁”——那纲领,我现在(二十年后)试图以《追寻栈》给出一个草图,而他当然同样不曾不感到其需要。最后,即便他从Motif瑜伽(被埋葬直到1982年)汲取灵感,这部瑜伽仍被削去其一部分力量:它与六种变换的形式主义——其一个本质的形式侧面——脱了钩。依我看,这一侧面,也被严格地逐出Hodge--Deligne理论之外。}。其三,是执著于那个首当其冲的主题:在其中,要取代、排挤一位每走一步都在场的老师、并无休止地抹去他——恰恰是这个主题,被注入最强烈的那份矛盾,它主宰其数学家的一生\dmark{}\origpage{380}。

我第一手所知的东西,加上一份从不曾骗过我的素朴本能或嗅觉,使我清楚:倘若Deligne不曾被这份其工作内部的深层矛盾撕裂,今天的数学便不会是如今这副样子\fnote{121}{(*)}{写“今天的数学”这些文字时,我想的不只是我们今天对数学事物或多或少的认识。背景里还有对一个特定\textbf{精神}的念头——在数学家的世界里,尤其在被人们(不带讥讽或嘲弄的口吻)称为“上流”数学界的东西里:那个“定调子”、裁定何为“重要”乃至“正当”、何为不然的精神;那个控制信息手段、并在很大程度上控制前程的精神。也许我夸大了身处旗手之位的一个人对给定环境在给定时代的“时代精神”所能有的影响。Deligne的影响,依我看,堪比(无论好坏)二十年前Weil在接待我的那个环境里所显得具有的影响——那个环境,我曾二十年认同于它。(5月31日)可与“掘墓人——或全体教众”一注(n\textsuperscript{o}~97)的(互补)反思对读。}:它本会在其若干本质的部分,经历宏阔的革新——一如我曾亲身充当其主要工具的那一场革新——恰恰是同一位Deligne拼命对抗、使之改道的那一场!\fnote{122}{(**)}{(6月16日)我深信:单凭我引入数学的主导观念正常发展这一点——沿着六十年代已获得的冲劲(那冲劲,被“电锯效应”猛然切断,下两条注将谈到它……)——今天的数学,在我离去十五年后,本会在其若干本质的部分,不同于如今的样子……}

同样毫无疑问:他天造地设,本该成为一所强力几何学派的灵魂——延续那所曾围绕我形成的学派:一所汲取其母体活力、以及接我班者之创造力的学派。但这所曾围绕我形成的学派、这座曾环绕那些紧张的受业年月的滋养母体,在我离去后的翌日便已解体。所以如此,恰恰是因为在分明接我班的那个人身上\fnote{123}{(***)}{这份\textbf{既成事实的接班},以毫不含糊的具体征兆表达:他在IHES接了我的班(我是在他入所的次年离开的——见“驱逐”一注,n\textsuperscript{o}~63);并且他接手了代数簇上同调这一中心主题,用的正是我为这一目的、在十五年间(1955至1970年)发展的手段。},找不到那个也将是一场共同冒险所结合的群体的灵魂的人——为一项其规模超出每个人的手段的任务。

我的印象是:我离去之后,我的每一位学生都退回自己的角落——工作自然有的是,数学里从不缺工作——但这个“角落”不再嵌入一个整体,这份“工作”不再被一股潮流、一个更宏阔的旨趣所承载。肯定,自我离去起——甚至更早——我的大部分学生或前学生的目光,都转向那位天造地设的“继承者”:他们中间最光彩、也离我最近的一位。在这个敏感的时刻,我的朋友想必平生头一回感到那种对他人的权力突然握在自己手里:那种对一个学派的命运生杀予夺的权力——那学派,他出身于斯;与他同学四年的朋友们,想必期待他保证其延续。局面整个在他手里,是他来定调子……他也确实定了——经由摧毁遗产;而首先摧毁的,是那份信任与那份期待\fnote{124}{(*)}{(5月26日)在反思的进程中,我还察觉对我的默示继承人的另一份全然不同的“期待”——这一回不出自我的学生们,而出自“全体教众”——对此可参见“掘墓人——或全体教众”一注(n\textsuperscript{o}~97)的结尾。我几乎不怀疑:这两份方向相反的期待——一份系于一个非常特定的时刻,另一份贯穿葬礼的十四年——都是实在的。更有甚者,我倾向于认为:在我的不止一位昔日学生身上,两份期待想必同时在场:一面是期待在他们的最光彩者身上,也找到那保证一所学派、一部著作——他们在其中有位置、有份额——之延续的人;一面是期待看到(倘若可能)那个人的任何痕迹被抹去——他的离去,曾在划好的坡道的安稳之中,以如此的力量,突然向他们发问……}——那些曾与他同学于同一位老师的人,本不会不带给他的……

当然,许多人为Deligne的著作所折服,并非没有理由。但我也清楚:这部著作,在其令人惊叹的最初冲劲(以Weil猜想的证明告终)之外,离使出“全副本领”还很远。诚然,它见证一种少有的技术驾驭与从容,把他置于“最好者”之列。但它没有我年少时在他身上知觉到的那份寒微的德性——革新的德性。他身上曾怀有的这份德性、那份幼童的鲜活或天真,久已深深埋没、遭否认。我本要写:凭着这份“德性”与其不凡的天赋\dmark{}\origpage{382},也凭着他的天赋之舒展所曾受惠的那些特殊情势,Deligne本被召唤去“支配”我们时代的数学,一如Riemann、一如Hilbert,各支配其时代的数学。一些扎根于日常语言的积习思路,向我提示了“支配”这一意象;但它给出的,是对实在的一幅失真的把握。这些伟人当然曾完完全全“把握”“消化”“据为己有”其时代的已知数学;这无疑也给了他们一份对技术手段的出众驾驭。但倘若他们名副其实地显得“伟大”,那不是凭其技术的壮举——从一块顽石上“硬凿”艰难的证明。而是凭各自在数学的几个重要部分所带来的革新,经由简单而丰产的“想法”——也就是说:因为曾把目光投向那些简单而本质的事物——那些在他们之前无人肯垂顾的事物。那种\textbf{看见}简单而本质之事物的孩童般的能力——无论它们多么寒微、多么为众人所不屑——革新的力量、每个人身上创造的力量,正寓于其\textbf{中}。这份力量,在我认识的那个年轻人身上——一个无人识得、谦卑而热忱地爱着数学的人——以一种罕见的程度在场。岁月流转,这份寒微的“力量”似乎已从那位受人仰慕与畏惧、无碍地享有其威望、以及威望赋予他的(有时是独断的)对他人之权力的数学家身上,消失了。

这份在我朋友身上对一件极精微、极活泼、为众人所忽视而有创造之力的东西的\textbf{窒息},自我离去以来我已感到过多次;近年来尤甚。但需要这些星期的发现,以及我自3月底以来持续的反思(循着《收获与播种》的冲劲),才开始在其全部幅度上感到这份窒息在我朋友生活中的毁灭性效果——也在许多我曾亲近认识的其他人身上。不只对我“1970以后”的学生(及类似者)——他们领教过他的恶意(某些情形里,也许是无意识的),它对着每一个人施展,沉重地压在三人身上;而且——我现在隐约看到——也在我“1970以前”的学生中间:经由一道\textbf{旨趣上的连续性}的毁灭,以及整体感、统一感的毁灭——那份本赋予其工作以比一堆署着他们名字的抽印本的堆积更深、更广的意味的东西(注91)\fnote{125}{(*)}{(6月16日)这第二个侧面,只是在葬礼的反思中才向我显形。倘若我有幸见到一位显赫的数学家动用“令人气馁的权力”,那恰恰是在不久前还俨然我天造地设继承人的那位身上。写“令人气馁的权力”一节时,我想到他有过许多(在反思转回我自己之前),却还不曾有一丝疑心(至少在意识的层面上):这份权力,竟在那些人中间找到施展的机会——对他们,他本该(一如当年对我)扮演完美数学家的典范……}。

最近七年里,不止一次;最近这些星期、这些天里,又是不止一次,我在一种于某个层面上被感受为巨大的\textbf{糟蹋}的东西面前,感到一阵悲伤——当人身上与别人身上最珍贵的东西,仿佛刻意地被挥霍或窒息。然而,我终究也懂得:这样的“糟蹋”是人的境况的一个基调,它以这样或那样的形态,无处不在——在最寒微与最显赫的人们的生活里,一如在各族各邦的生活里。这份“糟蹋”——它不是别的,正是冲突、分裂在每个人的生活中的行动——是一种我刚刚才开始探其边的丰富而深刻的实体,一份属于我去“吃”去消化的食粮。由此,这份糟蹋,以及我在每一步遇到的别的糟蹋、以及每在路转弯处来到我面前、常常不合时宜的任何事物——这份糟蹋与这些不合时宜之物,身上都带着一份\textbf{益处}。倘若禅观有意义、有革新之力,那正是在这样的限度内:它使我得以领受那份(经我积习的折光)显为“为害”之物的益处,使我得以\textbf{以}那看似为毁灭而造的东西为食。

以自己的经历为食,任它更新自己而不总是绕开它——这就是充分承担自己的生活。我身上有这份力量;每一刻,用它或弃之于闲置,都由我。对我的朋友皮埃尔,对曾是我学生的每一个人,亦然——同我一样自由,以这份“糟蹋”为食;这份糟蹋,我在这场漫长禅观的最后几天里,正巡视完毕。对读到这些为读者而写的文字的人,亦然。

\subsection{共同继承人}

\notehead{91}{}

\dmark{}(5月19日)从四面八方传到我这里的、关于我昔日学生们的回音,稀薄之至。\origpage{384}我离去之后,几乎没有人肯给我一点活着的音讯——哪怕寄一份抽印本\fnote{126}{(*)}{(5月31日)对此可参见“沉默”一注之后的注n\textsuperscript{o}~84\textsuperscript{1}。}。然而,把我得到的这点东西聚拢起来,我可以形成一个概念——诚然极其粗略。也许在随后的几个月里它会明确起来,倘若这场反思促使他们中某些人现身。

我已有机会确认Deligne著作在我离去之后的深刻断裂,尽管在若干方面,他违背本意地显出接班人的样子,因而显出某种连续性。我曾感到这道断裂,必定深深波及我其余所有学生的工作。这正是我想稍稍接近地勾勒的印象。

这些学生中唯一一个其工作显然(至少乍看)延续了他与我一起所做的工作的,似乎是Berthelot\fnote{127}{(**)}{从Verdier在我离去后几年间所追随的对偶主题——在一个与我发展它的语境相近的解析语境——看,有一种如Berthelot情形的连续性印象。但我觉得,那多少是一种“例行公事的连续性”;而我主要寻找其征兆(或征兆之缺席)的,是一种创造性的连续——在未知中延续一股最初的冲劲……}。他也是长年来唯一不断给我寄许多抽印本的——甚至也许全部抽印本。它们都处在晶体上同调这个艰难的题目里;其系统的启动,是他论文的对象。然而我觉得,一如我其他“上同调学家”(交换的)学生,其著作同样带着对我引入的某些主要观念的冷淡:导出范畴(及Verdier提炼的三角范畴),六种运算的形式主义,topos(注91\textsuperscript{1})。如Zoghman~Mebkhout本人所说,他自己的著作——题材与Berthelot的如此贴近(注91\textsuperscript{2})——正处在这条线上:这些观念,连同Sato学派的观念。倘若它们不曾被我的上同调学家学生们——Deligne与Verdier领衔——否认,那么很可能早在七十年代初,Mebkhout的晶体理论(他1975年才着手发展,而且顶着这同一批学生的漠然)就已抵达六种\dmark{}\origpage{385}运算形式主义的完全成熟——而它至今仍未达到\fnote{128}{(*)}{(6月7日)贸然下这一评判,我曾有一番踌躇:它可能被解读为贬低Mebkhout理论的原创性。那与我的想法绝不相符;更何况,我对我的每一位上同调学家学生的手段都有极佳的评价(当他们不被外在于健全数学意识的成见所堵塞时)。我的朋友Zoghman本人,消除了我可能有的这份顾忌:他自己深信,“正常地”说,早在七十年代初,本该是我的学生们来发展他的理论。在某个层面上,他们自己想必是最先深信这一点的:那理论的本该的作者,是他们,或Deligne——而世风的普遍败坏再一助阵,便足够他们表现得俨然自己(或Deligne)确实就是作者!对此可参见“研讨会”与“欺瞒”两注,n\textsuperscript{o}~75$'$与85$'$。}。

我倒是记得曾向Verdier谈起那个使我着迷的问题:离散可构造系数与连续系数的关系;而毫无引起他兴趣的样子。这想必随后使Deligne着了迷,因为他用一整年(1969年)主持一场研讨班,以建立一部词典;那词典想必不曾使他满意,因为他随后把它弃之不顾。(见“当值无名之辈与好上帝定理”一注,n\textsuperscript{o}~48$'$。)况且此后,他被自己的下葬综合征“堵”到这种地步:直到1980年10月,他都未曾知觉Mebkhout工作的重要性——而当他终于觉察时,是以人们所知道的那副掘墓心境(见注n\textsuperscript{o}~75至76)。

据我所知,Verdier自答辩以来的著作,其实质仅限于在解析语境(它有时呈现额外的技术困难)里重做我在概形-凝聚框架里所做的事,不引入任何新想法。这甚至相当离奇:凭他分明发展出的反射、凭他消息如此灵通,他竟没有靠自己摇动手柄而撞上Mebkhout的理论——而且竟至少不曾认出:他的“学生”正在做些依我看有趣的事情,而且是他所错过的(一如Deligne所错过的)。

老实说,尽管我对离散系数与连续系数之关系的问题着迷,我却当真不曾疑心Mebkhout的晶体理论——它将在我离去后的十年间绽放。相反,有一个宏阔的主题,出自我五十年代(1955--1960)交换与非交换\dmark{}\origpage{386}上同调的反思;它仅仅被着手过(在“交换”的侧面,即以加性范畴的语言):在Verdier的工作里——六十年代初启动,答辩之后被弃置(见注n\textsuperscript{o}~81)。非交换的侧面,更晚在Giraud的论文里着手:他发展一种几何语言,以topos上的1-叠,表述维数$\leq 2$的非交换上同调。自六十年代下半叶起,这两道起跑线的不足就已显而易见:一边,“三角范畴”观念(Verdier提炼的)不足以表达与一个导出范畴相系的丰富的结构(该观念注定被\textbf{求导者}这个丰富得多的观念取代);另一边,需要为任意维数的非交换上同调发展一套几何语言,以topos上的n-叠与$\infty$-\textbf{叠}表述。人们(或我)感到需要这两条进路的综合;它将为同调代数与同伦代数充当共同的概念基础。这样的工作,也与Illusie的论文工作直接连续——两个侧面在其中都有代表。

经由求导者观念(非交换与交换框架皆有效),Bousfield--Kan论同伦极限的奠基工作(Lecture~Notes n\textsuperscript{o}~304,1972年出版),也处在这条弥散的纲领的脉络里——至少自1967年起,它只待人手来发展。去年1月——我尚不知道一个月后我将投身《追寻栈》——我把关于同伦型的“积分”的反思(同伦论者熟悉其名:“同伦(归纳)极限”)呈给Illusie;那时我还全然不知道Bousfield与Kan的工作的存在,也不知道这类运算已被别人考察过。显明的是,Illusie同样不知——尽管他被认为自1970年我“逝世”以来,始终泡在同调-同伦的水域里!由此可见,他与某些自然地属于奠基反思、属于他六十年代亲身推进的那条路线的实在,似乎失联到了何等地步\fnote{129}{(*)}{同伦型的“积分”这一观念,是在分层结构的剖视的语境里、于1981年底重新强加于我的。}。他想必给自己挖了个小窝,不再常出来……

\dmark{}\origpage{387}随着对topos观念本身、对全部“范畴的无意义”的轻蔑,毫不奇怪:Giraud如今对他头一项大工作主题已全然冷淡。诚然,Deligne两年前开棺出土Motif之际,曾装作忽然发现非交换上同调武库——gerbe、lien等——的趣味,仿佛是他刚刚引入它们,连同Motif与MotifGalois群\fnote{130}{(*)}{见“一个梦的回忆……——或Motif的诞生”一注,n\textsuperscript{o}~51。}。这种马戏,多半无力重燃他自己拼命扑灭的火焰……去年2月,我曾把一封二十页的信的抄件寄给Giraud;它成了《追寻栈》的开卷第一章。那是一场绝无技术性的反思;在其中,我得以“一跃而过”当年曾使Giraud(以及许多别人)却步的“炼狱”——即“非严格”n-范畴(我如今称之为n-叠)的观念:它始终停留在启发的层面,却分明是基本的。那就是《追寻栈》的起点。去年12月,在Contou-Carrère答辩之际,我们相遇(双方心境全然友好);我从Giraud那里得知:他连读这封信的好奇心都不曾有!我得到的印象是:他对这类东西已经一笔勾销。在那里可能有一片丰饶的实体——在他久已放弃的方向上——这个念头,似乎连掠过他心头都不曾。我试图让他明白——恐怕未能成功——那里有一项多汁而宏阔的工作,近二十年来等着有人做;而我到了暮年才终于着手,以至少在总体上勾出一个草图——听事物本身的口授——那片丰饶的实体;我这个“亡者”仍强烈地感到它,而我的学生们早已把它忘了。

Jouanolou同样放弃了一条研究方向——他刚以论文起手的方向。这一方向,成了由恰恰向他提供其选题之关键技术钥匙的那个人所倡导的时尚的轻蔑对象。三年前“反常”研讨会上围绕三角范畴的“抢滩”之际,同一位Deligne忽然(面不改色地)装作发现眼前那项宏大的奠基工作——其缺失忽然处处可感,而他正是十年来头一个劝阻它的人\dmark{}\origpage{388}——这样一项工作的需要,对我来说,早在1963/64年平展上同调起步时就已分明;对Deligne,也丝毫不差,自他开始听说$\ell$-进上同调与三角范畴之时——即次年他登上我的研讨班之时。问题在于:在$\mathbb{Z}_\ell$上(比如说,一个基底概形之上)构造“可构造三角范畴”、并在这一框架里发展“六种运算”的形式主义(这一步,我记得在Jouanolou的论文里已告完成)之外,还要做一项类似的工作:把基底环$\mathbb{Z}_\ell$换成一个多少任意的诺特$\mathbb{Z}_\ell$-代数,比如$\mathbb{Q}_\ell$或$\mathbb{Q}_\ell$的一个(代数?)扩张。这属于那些二十年来时机已熟、却始终等着有人去做的事情——且待吹在它们身上的轻蔑之风平息……

Raynaud夫人的工作的自然延续(平展上同调中的若干弱Lefschetz定理,以1-叠的语言),本会置于一个严格禁忌的$\infty$-叠语境里——不必提它!Sinh~夫人的工作亦然——1968年开始,1975年才完成——其自然延续,本会是一个所谓“monomiale”范畴的包裹$\infty$-范畴观念,或这类范畴的三角变形\fnote{131}{(*)}{见“一个梦的回忆……——或Motif的诞生”一注,n\textsuperscript{o}~51。}——不必想它!另一个是把她的工作转译成topos上叠的语言——多可怕!至于Monique~Hakim,她也不幸把论文做在一个题目上——自我不合时宜的离去以来,这题目如今在边缘上显得有点可笑:局部带环的topos上的相对概形,您说说看!她关于这一主题的小书,刊于《Grundlehren》(Springer),想必每年卖三四册——难怪我在那家名声不好,他们对接受我可能推荐的文本不再热心。对我而言,那是为一切“绝对”的“簇”观念(代数的、解析的等等)向一般“基底”做“相对化”的第一步试验——其需要,在我看来是明摆着的(注91\textsuperscript{3})。人们会说:直至今日,没有它也过得很好。但同样真确的是:人在这世上两百多万年,没有做数学也过得很好。总之,Monique~Hakim——她做论文的动机与我为她提议该题的动机并不相同——肯定绝无念头与那样一个主题保持哪怕一点接触;那主题,一旦脱离有利共识的语境、或脱离一股顶风冒雨坚持一种执拗而可靠视野的思想,便不再能对她有任何意味。

至于Neantro~Saavedra~Rivano,他似乎已全然从流通中消失——在世界数学家人名册(以及一切官方的名册)上,我甚至找不到他名字的踪迹。可以肯定的是:他那过于范畴家的论文题目,在那些裁定何为严肃、何为不然的先生们那里,想必不会有好名声。依我看,这篇论文最自然的延续,不多不少,正是那幅“Motif的宏阔图景”——对这个学生较为谦逊的抱负而言,那主题分明嫌大。然而他到底得了那份意外的荣幸:不过两年前,亲眼见自己的论文被其中一位大人物从头到尾重做。(对此可参见“下葬——或新父亲”与“白板”两注,n\textsuperscript{o}~52与67.)

最终,在我的十二位“1970年以前”的学生中,我唯一不太清楚其工作中是否有一道或多或少的苛烈或深刻的\textbf{断裂}——相对于他们在我的接触中所追随的工作——的,是Michel~Demazure与Michel~Raynaud(注91\textsuperscript{4})。我所知道的只是:他们继续做数学,而且(就其出众的手段而言,想必在意料之中)属于我方才所称的“上流”数学界。

前面这段从有时极其单薄的材料出发的简短反思,当然在很大程度上是假设性的、极其粗略的。我希望被提到的人肯原谅我或许粗疏的判断之误;倘若他们肯就此向我示意,我将乐于更正。在这里,我再次意识到:每个人的情形,肯定不同于所有其他人,代表一个远比一个如我般遥远的人所能合理把握的实在复杂得多的实在,更遑论以几行文字表达。所有这些保留说尽,我却仍感到这场反思并非无用——至少对我而言:以若干具体事实,稍稍安放昨天显现的那份仍属弥散的印象(它想必在未成形的层面,多年以来已在场):即一道\textbf{断裂}——它在我离去后的翌日降临于我的许多学生身上,并且在个人的层面上,映照一所“学派”的突然消失——一夜之间;在那道学派受业岁月的关键年月里,他们必定感到自己是它的一部分\dmark{}\origpage{390}。

\notehead{91\textsuperscript{1}}{}

(5月22日)我刚读到 P.~Berthelot的一篇综述文章,论CIRM吕米尼研讨会(1982年9月6--10日)“p-adic分析及其应用”;题为“刚几何与特征p代数簇的上同调”(24页),勾勒综合Dwork--Monsky--Washnitzer上同调与晶体上同调的主要想法。晶体上同调(受Monsky--Washnitzer启发)的出发点观念(连同其名字本身),以及经由引入由刚解析空间构成的site使之完备的观念——这些我在六十年代引入的想法,已成了所有从事这一主题者的日课,以Berthelot为首;其论文正在于发展并充实这些出发观念中的某些。这并不妨碍:我的名字,无论正文还是参考文献,都严格缺席。这样,第四位殡仪人学生清晰在案。下一个轮到谁?

(6月7日)有一件事值得注意:自我引入晶体上同调的起步观念十五年有余、自Berthelot确立该理论对真光滑概形确是“对的”理论的论文十有余年之后,人们仍未抵达我所称的晶体上同调的“掌控”局面——堪比为平展上同调在SGA~4与SGA~5研讨班里发展的那一种。对一个含对偶现象的上同调形式主义的(第一阶)“掌控”,我指的不多不少,正是对一套六种运算形式主义的完全占有。虽然我还不够“在局中”、无以赏识晶体语境特有的困难,但我不会惊讶:这一相对停滞的主因,在于Berthelot及其他人对这套形式主义观念本身的冷淡;这使他们(一如Deligne对其仍处童年状态的Hodge理论)忽视那头一个本质的“台阶”——通往一个全然“成年”的上同调形式主义所必达的台阶\dmark{}\origpage{391}。想必正是同一类心境,也使他不曾察觉Mebkhout的观点对他自己研究的趣味。

附记:当我在这里谈“晶体上同调”,而语境是放弃真性的假设时(对一个“全然成年”的形式主义,这是必需的),不言而喻:人们所用的晶状site,其对象是并非纯无穷小、而是“合宜的”(带除幂的)拓扑代数型的“增厚”(带除幂)。对晶状site的这一原初形态之推广的需要(对我来说,原初形态只是“对的”晶体理论的第一次逼近),我一开始就很清楚;而Berthelot是(连同出发观念一起)从我这里学得它的。这一关联的书面影射,见《主题概要》第5条。

\notehead{91\textsuperscript{2}}{}

相当离奇的是:除我之外,似乎没有人觉察,未具名的Mebkhout的理论是晶体理论的一个本质的新侧面。我——一个近十五年来已全然“脱钩”于上同调的人——却觉察到了,一旦Mebkhout去年费心向我多少解释了他所做的事情。而当我在Illusie面前(当作理所当然地)提起这一点时,他的神气,仿佛把这看作把两件分明毫无关系的东西(D-模与晶体)做了一点“边缘离奇”的拉扯。然而我第一手知道:他有数学家的嗅觉;我的其他学生(在这个案里即上同调学家们,从Deligne起)也有——但我确认:在某些局面里,这嗅觉对他们不再有任何用处……我越想,越觉得离奇:在这样的气候里,Mebkhout到底还是做成了他的工作,没让他的数学嗅觉被其长辈们的全然不理解——他们高踞其上……——缴了械。

\notehead{91\textsuperscript{3}}{}

主要是从五十年代末我在Cartan研讨班上关于复解析空间理论之基础、以及“Teichmüller式水平模簇”之精确几何诠释的报告起,我才懂得了对迄今所用的“簇”的通行观念(代数的、实或复解析的、微分的——或此后,“温和拓扑”中的变形)做双重推广的重要性。其一是放宽定义,以容许任意的“奇异性”与“标量函数”结构层中的幂零元——以我的概形奠基工作为模型。\dmark{}\origpage{392}另一项推广,是向合宜的局部带环topos之上的“相对化”(诸“绝对”观念,经由取一个点的topos为基而得到)。这一概念工作,成熟已逾二十五年,在Monique~Hakim的论文里着手,至今仍等着有人接手。一个格外有趣的个案,是相对刚解析空间的观念:它使人得以把寻常的复解析空间与变剩余特征的局部域上的刚解析空间,看作同一个相对刚解析空间的“纤维”;一如相对概形观念(它终于入了俗)使人得以把定义在不同特征域上的代数簇相互联结。

\notehead{91\textsuperscript{4}}{}

Demazure的论文工作,一如Raynaud的,本质上使用他们在与我的接触中学得的一套纯熟的概形技术;但各自工作的本质观念,不属于“Grothendieck式”的行装——这使他们的工作有别于我第一期其他学生的工作。很可能这一情势的后果,是其著作中的一份连续性,免于“老师下葬综合征”效果造成的断裂。这未必意味着这一综合征不曾以别的方式触及这一位或那一位。三年前,Raynaud对Contou-Carrère相对局部Jacobi簇工作的态度,曾使我吃惊。所宣布的结果深刻、艰难,而且美不胜收,远超对“众所周知”之事的简单推广;与Cartan典型曲线理论有出乎意料的联系;有漂亮的显式公式——全然在Raynaud的射程之内(也在我的射程之内)。他接待的那份冷淡,想必对Contou-Carrère的战略撤退起了决定性的分量:他放弃了那个他曾毫无保留地投入、而看来只会给他带来麻烦的题目……\fnote{132}{(*)}{细节见“棺材3——或有点太相对的Jacobi簇”一注(n\textsuperscript{o}~95)的子注n\textsuperscript{o}~95\textsuperscript{1}。}。我那封向他表示我(痛心地)惊讶于他对这些结果之美无动于衷的信,没有得到回音。

\subsection{……与电锯}

\notehead{92}{}

\dmark{}将近四年前,当我来到这一带安家时,离我家不远有一片漂亮的樱桃园。\origpage{393}我散步时常绕过去看看。我喜欢看那些茁壮的樱桃树——正当盛年,树干粗壮,仿佛自古就与这片土地长成一体;野草在树下自由滋蔓。它们想必没见过化肥,也没见过农药;樱桃时节,我就到那里去摘有味道的樱桃。树,想必有二三十棵。

有一天我再去,看见所有的树干都在人高的地方被锯断,树冠瘫倒在树干旁的地上,茬口朝天——一幅大屠杀的景象。用一把好电锯,想必顷刻而成,至多个把钟头。这样的场面我从未见过——人们锯树,一般总肯弯下腰,贴着地面锯。樱桃滞销,这我同意;这片园子想必也产不出几吨,这没问题——但这些树茬诉说的,可不是滞销与收成……

昨天,我又有了这同样的感觉:一根强劲的树干,根系深厚,树汁丰沛,枝条繁壮、延续着它的冲劲——被齐人高、干净利落地锯断,仿佛纯粹为了取乐。正是费心把主枝一根一根看过、看清每一根都被锯断,才终于使我看清发生了什么。那本该循着一股冲劲、循着一种根系深处内在的必然而舒展的东西,被一刀切断,不加手续,好让它在众人眼前被指定为笑柄。

这使我想起Zoghman谈起的那个“误会”:它想必发生在我的学生们(Deligne除外)与我之间。分明的是:冲劲与视野,都不曾从我传到我的任何一位学生身上(Deligne除外——他分明“除外”,名副其实!)。每一位都消化了一套技术行装——对他所选的题目做出色的工作有用(甚至不可或缺),日后兴许还用得上。是否有别的什么东西的起头——超越于此——我说不好。即便有起头,它面对电锯也绝无机会:电锯顷刻把它刈平……

我知道,只要还有人做数学\dmark{}\origpage{394}——而且除非彻底放弃人们做了两个多千年的那一类数学——他们早晚会不由自主地让我眼前这些僵卧的树枝一一复活。其中一些,已经被我的电锯朋友据为己有;倘若上天借他余年,很可能他还会照样处置另一些,甚至全部。但大多数,已全然不合他的风格。不过,也许他终究会厌倦于不断取代别人——那肯定极其累人,而且无利可图——而满足于做他自己(这已经不错了)。

% 第16章 D) 被埋葬者们;原书p397-412

\chapter{D)~被埋葬者们}

\section{X.~灵车}

\subsection{棺材1——或知恩的D-模}

\notehead{93}{}

\dmark{}(5月21日)我的反思在那些“根正苗红”的学生们——“1970以前”的——身上,已滞留两个星期。每一天,这场反思都摆出一副“最后一项补充”的面目,出于求个心安,补向一场(几乎)已告终结的反思。不止一次,是一条不起眼的脚注——贸然接上前一天或前两天的反思——越拖越长,直到一份独立“注”的规模。每一回,它都很快得名,把自己区别于所有其他的注,并恰好插进其葬礼行列的合适位置,仿佛它从来就在那里!每隔两天,我都得(每次都怀着乐趣)至少重做目录的结尾:目录仿佛已收尾,却随之添上两三位新的入列者,有时甚至一整个新的行列……

这支队伍的规模,渐渐变得令人不安:谁也不会肯把这全部读完!但它这样加长,老实说,并不是为了一个假想读者的可疑收益,而首先是为我自己的收益——一如我做数学之时。这些“最后一项补充”——我每一回都仿佛违背本意地登船——从不曾让我后悔。借着一次次的最后一项补充,我学到了许多别处无从学到的东西,倘若不这样“从材料出发”反思,我便无从学到这些。而这些东西,一件一件聚成一幅色彩鲜明的图景,宏阔而有多个侧面。此刻我仍看到它尚未全部完成——还有两处,似乎在讨最后一笔。

我觉得,谈过那些“根正苗红”的学生之后,如今也该稍稍谈谈\textbf{被埋葬者们}了——那些“与我一同,有权领受这场以沉默与轻蔑进行的下葬之殊荣”的人。一如我、一如那些兴冲冲下葬的人,这些被埋葬者也不是圣人,也没有殉道的使命。我相信,他们中没有一个不曾暗暗怪我\dmark{}\origpage{396}给他招来麻烦——那些我全然无意给他招来的麻烦(只因为他曾失于轻率,把赌注押在我身上,押在某种做数学的进路与某种做派上……)——或者,至少不曾试图与我划清界限,一旦认清赌注分明已输\fnote{1}{(*)}{(1985年2月)自我到蒙彼利埃以来,在校外,我总共得知七八份(简短的)出版物,呈献(以摘要的形态)一项与我一起做、受我启发的工作。我的名字,在所有这些里缺席。}\origpage{397}。何况我可以确认:那是白费力气——一旦被盯上,就完了;划清界限,只会喂养一份轻蔑、给它一份心照不宣的辩解,而不能解除它。同样不止一次、以多种方式,我看到下葬者与被葬者的角色相互毗邻、相互混淆\fnote{2}{(**)}{(9月2日)以各不相同的方式,他们每一位都曾在某个时刻,终于把对其工作的轻蔑内化并据为己有,默许那份藏掖其工作或把它归为“无趣”的共识。}。想必正是这些暧昧的侧面,造成我身上一份长久的迟疑:不肯比此处彼处顺带的影射更周详地谈“被埋葬者们”。很可能,除Zoghman之外,另外三位我认识的人里,没有一个会感激我在这里替他做“广告”——仿佛我给他招的麻烦还不够多似的。

一如《收获与播种》里的许多时刻,我终于越过我身上这份迟疑。我对自己说:即便对着那些曾因我而受苦的人——经由他们在某个时刻做出的选择;在那选择里,他们出于这样或那样的缘由,分明各得其所,尽管他们同我一样,不曾料到其选择所附带的种种不利——即便对着他们,我的角色,也不是帮他们回避一个实实在在的局面;无论愿意与否,他们都身陷其中,而这局面分明有其意味,哪怕它带着严重的麻烦。

在接上我那些亲爱的共同亡故者、共同被埋葬者的四口棺材的黑色系列之前,我或许该以一条不那么丧气的注,给读者换换心情。首先,在与我们大学数学研究所的“地方”层面的关系里,我完全没有体验到:我为一个职位的候选人所说的话、或候选人是我的学生(不言而喻,是“1970以后”的)、或其著作受我著作的影响,就必然于他不利。那种系统性抵制的态度,仅仅是“上流”数学界对我这个人的关系——推而广之,对那些显得与我“1970以后”相连的人——的特征。据我所知,这份抵制,自我离去以来的十四年里几乎毫无破绽,只除两桩寒微的例外。其一涉及一个学生:他开局颇有希望,本该跟我做一篇题材极为诱人的国家博士论文;他申请我们大学的助理教授职位,却被专家委员会驳回。他在全国一级被“捞了回来”,有Demazure的帮助——我曾就这位学生的工作给Demazure写过信\fnote{3}{(*)}{就晋升或取得职位与身份的“实际”层面而言,我1970年以来的教学活动的全部账目,归结为两桩取得职位并附带身份的事:一回是助理教授,另一回是助教。说来奇怪的讽刺:两回,这晋升都是当事人一切研究活动戛然而止、彻底停摆的信号。}。另一方面,有两次,《Topology》杂志接受了出自我的学生的文章:一篇“Stein因子分解与剖分”,出自Jean~Malgoire与Christine~Voisin;以及Yves~Ladegaillerie一篇将刊的文章,包含其1976年论文的中心结果(见注n\textsuperscript{o}~94)。

我尤其已有机会谈到过Zoghman~Mebkhout,这里只“存照”再谈\fnote{4}{(**)}{除《导言(6)》(葬礼)之外,谈到Mebkhout的还有:“我的孤儿们”“当值无名之辈与好上帝定理”“不义——或归来的意义”“反常性”“墓外相会”“受害者——或两种沉默”“大部头与上流社会”“赊账的学位论文与全险”各注(n\textsuperscript{o}~46、48$'$、75、76、78、78$'$、80、81)。}。Mebkhout想必自1974年起开始从我的著作汲取灵感,并且顶风冒雨、一直汲取到今天。我不知道我的哪一位“官方”学生做出了分量可比的著作——尽管Mebkhout的著作,必然带着其被迫在逆境中推进的印记。如我在《导言(6)》里所说,四年来,Mebkhout的观念与结果,人人都在用,而他的名字却被小心地藏掖\fnote{5}{(***)}{在那场下葬里充当掘墓人的,大有人在;1981年6月的吕米尼研讨会,几乎全员参与。除我的上同调学家学生们(对此可参见“我的学生们(2):团结”一注,n\textsuperscript{o}~85)之外,其职业诚意在此直接且严重地成为问题、而我有知情的,是 J.L.~Verdier、B.~Teissier、P.~Deligne、A.A.~Beilinson、J.~Bernstein。}。我的朋友何以能在承受轻蔑、继而承受不义——如一场不可抗拒的宿命——的同时继续做数学,这对我是谜——那份宿命,经由他想必(并且想必仍然)感到高悬其上、令人眩晕的人而来\fnote{6}{(*)}{当然,Zoghman~Mebkhout并不比我傻,而且他足够“在局中”,对我每一位上同调学家学生的工作都有精确的把握,能掂量其分量与限度,毫无理想化的倾向。但这并不妨碍:一些力量可观的抑制,使他连那样一个念头都不曾有——即公开质疑他们中的任何一个,哪怕在恶意分明之处。}\dmark{}\origpage{398}:那些人,他想必头一回听说时,俨然“赛场诸神”;而那时,他(一如当年的我)是一个生计拮据的寒微留学生。1979年答辩之时,他在奥尔良有一个助教职位。当时他竭力设法进CNRS,三次进攻——到第三回(1982年10月),人们终于肯给他一个研究员职位(相当于大学的助理教授或副教授)。这给了他,即便不是身份的担保,至少一份相对的安全。

在我所知的四位“共同被埋葬”的数学家之中,Mebkhout是唯一一个逆着所有人继续其工作的:他信任自己的数学本能,不被那份无情时尚所能提示的谨慎与机会主义的考虑绊住。他身上——他并非好斗的天性——有一种对自己判断的素朴的\textbf{信念},它同时也是一份\textbf{宽宏};而它(远甚于脑力“手段”)是做出革新而深刻的工作的首要条件。我对他的著作的看法,肯定仍不完全。就我所知其著作的主干部分而言,我觉得:以他那出众的手段,倘若处身于一份热忱而起作用的同情的气候里,他本可以在三年或四年里(而非十年)完成它,并把它带向更大的成熟——在喜悦而非辛酸之中。但三年也罢十年也罢,“成熟”与否,值得注意的,是革新的著作出现了,而且竟能在这样的条件下出现。

\subsection{棺材2——或电锯剖切}

\notehead{94}{}

Yves~Ladegaillerie1974年开始与我工作。那是“姑且一试”,在他生涯的一个低谷——我当时向他提出了几条关于1维拓扑复形在曲面中的嵌入的天真反思;那时我对曲面一无所知(除亏格观念),他懂得更少。这有点像Grothendieck学问(反正在我这里,一切总是这样开始……);它在他那里或多或少咬住了钩,直到有一天终于“翻了盘”——我已记不清何时何故。也许是在一个分明多汁的问题正在成形之时:关于紧致带边定向曲面中1维紧复形的同痕类的判定的一条关键猜想。真——假?悬念持续了六个月、一年;其间伊夫熟悉了(也顺带使我熟悉了)曲面理论的关键定理,同时推进其工作的“奠基”部分。已知的结果使猜想相当可信,却分明远远不够——因为那猜想蕴含巴厄与Epstein的那些难啃的结果,以及另一些有着离奇甚至可疑面目的东西。1975年夏,他终于证明了那条关键猜想。它实质上等价于:以基本群的语言,对一个可三角剖分的紧空间嵌入紧致带边定向曲面的全部同痕类,给出完整的代数描述\fnote{7}{(*)}{非定向情形的“类似”陈述是假的——分明这是一个精微的结果,被小心地“剖”在一组同样“可信”却仍为假的假设-结论之中。关于Ladegaillerie工作的其他评语,见《纲领概要》,尤其第3节开头。}。

\dmark{}从伊夫“咬钩”那一刻起,他一年、至多一年半做完了论文:结果、成文,一应俱全,而且几乎苛刻地完美。那是一篇光辉的论文,比我多数学生的论文薄,却与其余十一篇中的任何一篇同样实。答辩在1976年5月举行。

论文至今未出版。它虽不厚,据说还是太厚,不宜出版——这是人们给我的许多条极好的理由之一。我在“进步是挡不住的!”一注(n\textsuperscript{o}~50)里列举了几条。我为“安放”这篇倒霉论文所作努力的经过——它是我有幸激发的最好论文之一——可以写一本小书;它肯定富有教益,而我放弃写它。在那些有极好的理由忘掉去了解这些结果、并闭着眼睛把这一切埋葬的昔日好友之中,有(按登场的顺序)Norbert~A.~Campo、Barry~Mazur、Valentin~Poenaru、Pierre~Deligne——还没算经Springer公司转手的 B.~Eckmann\fnote{8}{(*)}{我与Eckmann并不相识;为使伊夫的论文在Lecture~Notes出版而做的通信,是与Springer主管LN的Peters博士进行的。我想,经由六十年代由我(SGA尤其是)或由学生(论文)出版的十五卷LN,我曾是凭我的担保为这套丛书尚在草创之际的前所未有的信誉与成功出过力的人之一。拒绝我所推荐的工作所给的理由(他们不出版论文),是句玩笑话。通信事务上的“新面貌”,我的头一份经验也出自这一情节:怀着一套真正令人叹为观止的默契, A.~Campo、B.~Mazur、 V.~Poenaru与Peters博士,对我第二封信,都竟免于赏我一复——当时我天真(我的悟性慢……)地卷土重来,而他们那封迟疑的回信已表明:他们没有费心了解Ladegaillerie工作导言里所陈述的结果。}\dmark{}\origpage{400}。中心结果终于要发表了——九或十年之后,削减到只剩骨头,以《Topology》一篇短文的形式(嘘——我在那份可敬杂志的编委会里有一位同谋……)。工作的其余部分,一边证明了人人久已不证而用的事情(当然,没有证明也过得很好!),一边发展典型的Grothendieck学问,全然背逆时下的惯例与良俗。我很清楚:倘若我的朋友Deligne不在未来十年里大张旗鼓地负责“发现”它们,别人在三五十年内也忍不住会重做它们——因为我健全的本能告诉我:那些是基本的东西。它们曾是我阿那贝尔式沉思中宝贵的引线;倘若上天借我余年,在《数学反思》发展阿那贝尔式代数几何瑜伽的那一部分里,我会有充分的机会指涉它们。

这段经历对我是一次显影,同类中的头一次——对某样东西的显影;那东西,直到葬礼的反思,我才充分认识。况且,自从我的心神另有所系,我倾向于把它忘了。Yves~Ladegaillerie——我有过的最光彩的学生之一——则从那一刻起就懂得了:要在今天的数学界被接纳,仅仅全力投入、并做出合乎一切卓越要求的工作,是不够的。他另有几手,七年来他投身于更接地气的任务与较少问题的回报。他运气好:早在与我那场不幸相遇之前,他就有助理教授的终身职位\dmark{}\origpage{401},给了他一份其厄运不曾危及的保障。去年,一点数学的火星似乎重新燃起,在一个与我近年来关心的问题十分贴近的主题上——瑟斯顿式的双曲几何及其与Teichmüller群的关系。甚至可能我们还会一起走一小段路;或者他走他个人的散步,纯粹为了乐趣,不期待数学本身所能给的回报之外的任何回报。他清楚:倘若他期待别的,他最好换一个对话者或同路人(甚至,换一段过去……)。

\subsection{棺材3——或有点太相对的Jacobi簇}

\notehead{95}{}

我与Carlos~Contou-Carrère的最初相遇,是在1973年我到蒙彼利埃之后不几天,数学所的走廊里。他会把我堵在某个昏暗的角落,在我来得及客气告退之前,把一股数学解释的洪流倾倒在我头上。他连珠炮般倾倒给我的东西,全从我头顶掠过;而他丝毫不曾察觉,或在我怯怯暗示之时有丝毫窘迫。他迫切需要对话者,而我并非他唯一的“违心对话者”。那又是一个我全然不沾数学的时期。有一两年,我一看见他(一眼就能认出)的身影在走廊尽头出现,就躲。直到林登——他曾在蒙彼利埃做过一年客座教授——向我透风:Contou-Carrère手段不凡,却眼看要触礁,因为不知道怎样使用它们。在那之前,Contou-Carrère倾倒给我的东西站不站得住、他有没有手段,连掠过我心头都不曾——一切那么遥远。也许林登的提示来得恰逢其时:我开始对数学问题重新发生一点兴趣。总之,我横下心来,问Contou-Carrère可肯给我解释一件他做过的事,好让我能懂。我疑心:我是头一个向他提这种要求的人——至少在他来法国已有好几年之后。要他显豁一件事不容易,却绝非不可能,而且值得。我很快发觉:林登没有看错——Contou-Carrère满脑子观念,只等着被提炼、被细心发展;而且在我们可能向他提交的几乎一切数学局面里,他都有一种即时的、非常可靠的直觉。凭这份直觉的敏捷与可靠——哪怕在他毫不熟悉的事情上——他超过我、令我惊叹;在别的学生身上,我以可比的程度体会到它的,只有Deligne\fnote{9}{(*)}{我不确信在其他数学家身上遇到过它——除Pierre~Cartier(他年轻时以这份卓越的能力深深打动过我)与Olivier~Leroy(下一条注将谈到他)。}。相反,他对书写有一道几乎完全的阻塞!不可思议:他做数学\textbf{不写}——天知道他怎么能哪怕做出那么一点,更不必说与他人的交流:在那方面,“触礁”是彻底的(见上文)。

倘若我有什么紧急而有用的事情要教给Contou-Carrère,那是书写的艺术——或更直截了当地,哪怕只是让他明白:数学是\textbf{写着}做的。我费了两年、也许三年,直到1976或1977年\fnote{10}{(**)}{(6月7日)核实之后:直到1978年2月。},不敢说是否全然成功。他第一部完全白纸黑字写成的大工作,是他的Schubert闭链论文——迟至去年12月(1983年)才答辩\fnote{11}{(***)}{那是一项长工作(我没有读过):他细心发展一些与我毫无关系的想法,其中给出一切“Schubert”型闭链的显式奇点消解——此前无人能做到。难得他有一次成形的写作,人们却责备它太详(更不必说他的陈述太一般……)!至于我,倘若我有批评,那将朝着相反的方向:虽然Contou-Carrère声称其方法适用于一切类型的半单群与Schubert闭链,他却只在一般线性群的情形做了工作——因此他在这件明确的事情上没有做到底:普适Schubert闭链的等变奇点消解\textbf{以及}这些Schubert闭链的奇点轨迹的描述。这份欠缺,在我看来,是那道对“从材料出发”的扎实工作、对书写的“阻塞”的一份遗产;那阻塞,长久以来是他的主要缺陷。}。1978年至今,我们的关系极为零星;我的角色,实际仅限于在他职业生涯中屡屡以这样或那样方式陷入困境之时——那职业,始终悬在最不牢靠的代理助教职位上——尽力扶持他。

有两三年,我曾试图向Contou-Carrère提供一套精确而柔顺的数学语言的基础,外加几条系统化的原则。带着这份行装,连同他的手段与观念的丰富,接续什么,他有真正的挑选余地。他没有从他自己的出发观念入手,而是接上了相对局部与整体Jacobi簇的理论——我曾把它作为可能的论文题目向他提过。一旦我放手让他自己干,他在仅仅一年之间做出了一项很漂亮的工作,其中一部分在一篇科学院报告短文里已有预告(注95\textsuperscript{1})。把这个矿脉做到底,本会代表几年令人着迷、强烈激励他的工作——足以同时学得概形技术的全部精微。那时我毫无疑心——明摆着,Cartier、Deligne、Raynaud三人都会热烈接待这项已完成的工作:深刻、艰难,且在好几个方面出人意料。Cartier确实满心欢喜地看到他的某些旧观念取得新的现实性。相反,是Raynaud的冷漠,一如Deligne的——后者把完整手稿在抽屉里压了六个月,不屑给一点活着的音讯\fnote{12}{(*)}{Contou-Carrère倒也先防了一手:在其短文里,对我的这个人——曾给了他出发纲领的人——只字不提。那也是白费:任他往里加多少自己的东西,有一种“做派”不会错,它系于某些主题;谁愿在今天的数学里混出前程,最好避开这些主题。(6月7日)向当事人查询之后,我发现:我在这里把围绕Contou-CarrèreJacobi工作的两个不同情节弄混了。细节与精确出处,见下一条注(n\textsuperscript{o}~95\textsuperscript{1})。}。这是二对一——足以感到风向。那有点太相对的Jacobi簇,被无限期弃置,任其自生自灭。电锯干得很出色……

这并没有为Contou-Carrère挡住厄运;为其写一份周详的报告,本可以再写一本小书,我由衷放弃。我想正是在这时——自我1970年离开那个机构以来(在那机构尚无自己的屋顶的年月里,我曾是四年间1958--62唯一在“实地”代表它、使它取信于人的人\dmark{}\origpage{404}),唯一的一次:我担当起为某人推荐一份邀请(为期一年)的责任,正当Contou-Carrère眼看要没有职位、流落街头。我知道:我推荐的人——一如当年的Hironaka、Artin或Deligne被我热情接待时那样无名——会像他们一样,为接待他的机构增光。当然,我不曾忘了这么说。对Contou-Carrère幸运的是:他的代理助教职位(诚然配不上一份如此挑剔机构的邀请之荣)最终得以续聘\fnote{13}{(*)}{我无从抱怨:五六年之后,借去年IHES建所二十五周年庆典之机,人们倒真赏了我这个人一份邀请,甚至让我在部长致辞的庄严招待会与随后一周的IHES居留之间选择——还包食宿(人家郑重向我担保)。我对老朋友Nico~Kuiper说:想到我,非常客气;但我这把年纪,不再出门旅行了……}。

对这一情节我并不太惊讶:当时我已了解Deligne的心境;况且Nico~Kuiper曾提醒我,在这个案子里一切取决于Deligne。(我甚至没有想到向他暗示:这事兴许也关涉科学委员会的其他成员——恰恰因为这是个案……)然而,在Contou-Carrère(我的“门生”——Verdier曾在一封信里想当然地这样称呼他……)的一切厄运之中,最触动我的情节,发生在1981年10月:事涉他申请佩皮尼昂的一个教授职位。佩皮尼昂的同行们(他在那里任代理助教)想必乐于身边有一个在任何数学分支里都从容、可咨询的人。一次教授职位出缺,他们把他列为唯一候选人——罕见之至的事,分明表明:他们想看的正是他、而非别人。Contou-Carrère在其阿根廷Santalo门下的博士论文之外,出版物相对少:主要是科学院报告短文,预告结果(其中一些深刻),而无证明。还没有人向他暗示过:在如今的时势下,在安身之前,最好有附完整证明的文章作“证据”——这一点我这边当然向他念叨过不少,只是从较不功利的角度\fnote{14}{(**)}{前一年,Contou-Carrère申请雷恩的一个教授职位;他认识Berthelot与Larry~Breen。当时大学咨询委员会认为他的申请可受理;但职位给了另一位候选人。没有人费心提醒当事人:倘若他想有职位的希望,就得发表他所预告结果的详细证明。次年委员会的驳回,对Contou-Carrère、对他的佩皮尼昂同行们、对我,都来得完全意外。带着事后的认识、在这场反思的光照下,我进而怀疑:随着其论文的成文(如今已被宣告“照原样不可出版”)与答辩,局面是否真的变了;他是否有希望在法国找到一个教授职位。}。总之,Contou-Carrère的申请\dmark{}\origpage{405}被大学咨询委员会判定不予受理,材料退回。当时使我倒吸一口凉气的,是:委员会主席(做出决定的全国机构)不以委员会名义、任何成员也不以个人名义,具有这点最起码的尊重:给当事人Contou-Carrère本人、或至少给佩皮尼昂数学所所长写几行,解释这次表决的意味——在毫无解释的情况下,这次表决只能被领受为对佩皮尼昂同行们选择的响亮否定,对其唯一候选人能否体面胜任其被提名之职的否定。委员会里有我的三位旧学生,其中两位亲自认识Contou-Carrère。他们当然知道他同他们一样是我的学生——何况材料里还包括我一份对候选人工作格外称道的报告。他们中没有一个、委员会其他成员中也没有一个,想到过这记不经任何手续的铡刀式表决所代表的羞辱,想到过对一位不比他们中任何人逊色的数学家的按部就班的鱼雷攻击。

正是这一事故,在我做数学家的一生中,头一回使我感到那种我已在反思中不止一次谈过的“气息”。四年前,经由“外来者”的情节,我已感到过它\fnote{15}{(*)}{见“我的告别——或外来者”一节,第24节。}。但那不是在我曾经所属的世界内部,吹在\textbf{他们自己人}身上——吹在一个毫无保留地认同于那个世界的人身上。我好几个星期、也许好几个月为此如病如恙。为摆脱当时攫住我的一份焦虑——我当时无心去认识它\fnote{16}{(*)}{我是在次年一场漫长的禅观期间,才意识到这份焦虑的;在那场禅观里,我发现了焦虑在我生活中的角色——它的在场(1976年之前是慢性的,1976年之后是偶发的),曾是我一生的“世上守得最好的秘密”。有一套效率极高的机制,藏掖焦虑的一切公认的征兆;那焦虑,对我自己、对我的亲人,都始终不为人知。},我四处奔走,东一封信西一封信,还写了一篇三十来页的文本《大脑与轻蔑》,带着黑色幽默的笔调;我终于放弃了发表它\fnote{17}{(**)}{劝我别发表它的,恰恰是那些我正准备为之开战的人——幸亏我有先见之明,在试图让它面世之前,把我的文本给他们看过。}。带着事后的认识,我看清:那本是就所遇之事的意义展开\textbf{禅观}的绝好时机。最滑稽的是:当时“妨碍”我觉察需要一场深入禅观的,恰是一场我那时正投入的长禅观\fnote{18}{(***)}{见“扫兴的老板——或高压锅”一节,第43节。}——而且是一场关于我与数学之关系的禅观(即便还不是关于我数学家过去的)!它被一个情节搅乱:生活以强力向我发问——而我在奔忙中回避那发问,随后重新扎进“禅观”。带着事后的认识,我看清:那场“禅观”当时配不上这个名字;它缺少真正禅观的一个本质维度:对\textbf{当下时刻}的我自身的留意。我当时“禅观”某些较为遥远的事件的意味,却忽视一份被压抑的焦虑(诚然,由于长期练习这类控制,它被完美地控制着)——那是我拒认那股被驳回的“气息”所带来的讯息的征兆。

但我正在离开本题。那场鱼雷攻击,当然有其必有的效果。佩皮尼昂的同行们被申饬了一次,就够了。显然,他们那里连代理助教的职位也没有了——至少对Contou-Carrère没有。他在蒙彼利埃临时找到一个替补职位,就本学年,编制者明年回来。

不过我并不太为他的前程担忧:一段时间以来,Contou-Carrère已有智慧抢在命运的打击之前,转而投向计算机。以他那出众的手段,想必久已高屋建瓴地驾驭这门学问,同时在闲余做他喜欢的数学。他是两个孩子的父亲;在如今的时势下,带着紧追不放的过去,数学分明是碰运气的事,甚至凶险。他大可以做一个辉煌的计算机学家:没人会因他多少当过我的学生而对他严苛。

\notehead{95\textsuperscript{1}}{}

(6月7日)大约77年底,我向Contou-Carrère提交了一份为相对局部与整体Jacobi簇理论的周详工作计划;其中包括,在局部情形,建议“重新拧紧”Jacobi簇与Cartier的ind-群,以找到一个具有更美泛性质的“完备”Jacobi簇,它应当是“自对偶”的。我没有任何可提出的证明,1978年2月之后也不再过问其工作:我发觉我的在场抑制而非刺激其才能。他在随后的一年里到底“启动”了;其第一篇注“相对曲线的广义Jacobi簇:构造与分解的泛性质”(论整体情形)发表于1979年7月16日(CRAS t.289, Série A - 203)。

次月,他为局部Jacobi簇找到了决定性的结果,却一年半不发表;然后在1981年3月2日的一篇科学院报告里发表“一半”(寻常的、未与Cartier群重新拧紧的相对局部Jacobi簇的泛性质),顶着(乍看不太有说服力的)名字“相对几何局部类域论”(CRAS t.292, Série I - 481)。至于完备局部Jacobi簇的理论——依我看更有趣味——存在一份科学院报告的底稿,从未发表,题为:“局部Jacobi簇,普适Witt双向量群与tame符号”。当然,1979年起我就知道他的结果:那是对我向他提议的临时纲领的一个完整实现;为它,须克服可观的技术困难,要求大量的想象力与技术力量。(如无差错)我只读过第一篇注;我奇怪他不发表后续,即局部部分——他从不曾就此清楚说明——但他分明对第一篇注所受的接待感到失望。1980年雷恩申请失败之后,鉴于附在他申请材料里的我的支持信\dmark{}\origpage{408}论及整体与局部相对Jacobi簇上的卓越结果,他想必还是认定谨慎为上(为准备次年佩皮尼昂的申请),至少再发表一篇关于局部Jacobi簇的注,如非和盘托出。又是两个月之后,1981年5月,他把第三篇注的底稿寄给Deligne与Raynaud(Cartier想必久已知情),我想是先探探路。(我不相信:自1979年8月他手握结果以来的任何时刻,他找Cartan呈送这第三篇注会有丝毫困难。)Raynaud与Deligne都不给他一点音讯——但1982年3月,Deligne给他寄来一篇文章的手稿“tame符号评注”,作者是加藤和也,献给Deligne;它在基底域的情形做出Contou-Carrère的理论,并猜想其对任意基底环有效。Contou-Carrère当时跟我谈起,自认深信Deligne把他的结果(不提其名,也不给证明的任何指示)通报了加藤。当时这事在我看来如此不可思议,以至于我没有把Contou-Carrère当真——如今我才明白:那完全是我光彩夺目的朋友Deligne惯常的“暂停!”做派。Contou-Carrère一副真正被冒犯的神气:有人竟“斗胆猜想”他俨然视为某种私产的东西。然而他自己的猜想得自我的传授,却同样不认为有必要在三条注的任何一条里影射我这个人\fnote{19}{(*)}{关于我在这类局面里、与我的某些学生之间常常扮演的那份共谋的角色,见“暧昧”一注,n\textsuperscript{o}~63$''$。}!他对我说来,想必理所当然;而Deligne对他做的单纯的动念却冒犯了他——但他不敢向当事人吐一个字。(我曾力劝他向Deligne说清;他却小心没有做……)

这些年来,我料想,他必定以某种方式对自己施加了暴力,不发表那些极美的结果——做它们时,他想必全力投入。倘若他这样对自己施暴,那是出于时势的考虑:分明不利于这一类Grothendieck学问。近几天他十分惊讶地收到同一位Deligne的一封信:他(不动声色地!)惊讶他没有发表关于“总体”Jacobi簇的那篇注,并向他索要他关于这一主题、甚至别的主题的一切。Zoghman~Mebkhout几天前已告诉我:Deligne正在使用这些东西,甚至在这一语境里点了Contou-Carrère的名。分明,时候快到了:Contou-Carrère终于将认下一个属于他的孩子——一个他谨慎地埋了近五年的孩子。也许,谁知道呢,两位“学生-仇敌”和解的钟点也已敲响:我的学生中最光彩的两位,一位是受勋的院士,一位是代理助教;然而(无论和解与否),久已是一对\textbf{兄弟}。

\subsection{棺材4——或没有鲜花也没有花冠的topos}

\notehead{96}{}

(5月22日)倘若我说我从未见过Olivier~Leroy,我几乎没有夸张。可以肯定的是:从他听说我那一刻起,他就决定像躲瘟疫一样躲我。他的缘由,我承认,向我滑脱。也许一个本能告诉他:我只会给他招麻烦;也许Contou-Carrère(他长久以来与他过从甚密)向他提过醒——我恐怕永远无从知道。但我到底有幸与他有过两次有实质内容的交谈;我记得很清楚。

第一次应当在1976、77年:我和Contou-Carrère不打招呼地登门看他,想聊聊数学——是否心里另有盘算,我说不好。不过,想必有言在先:奥利维耶考虑读第三阶段博士,而我袖子里有的是题目。因为在Contou-Carrère家见过他一两次,加上Contou-Carrère自己向我透露的,我有一个印象:奥利维耶悟性快,而且不只在数学上。三人那晚令人难忘。我很快向奥利维耶提起一个纲领:一个topos基本群的理论、以及topos框架里的Van~Kampen型定理;他一副感兴趣的样子。经由Contou-Carrère的代数几何研讨班,他想必有一点topos的常识;他似乎乐于有机会在一个具体理论的例子上“练练手”,用topos的语言。足足两三个钟头,我把我所见当发展的理论的周详总施工图倾倒给他;它随我谈论而渐渐丰满,同时,代数几何与拓扑的一大群具体局面在我身上浮起——那些局面,要在topos框架里表述;而每一回,我都得先向一个头一回听说的人“复述”\dmark{}\origpage{410}。那晚不止一次,Contou-Carrère(他几乎无所不读、胃口极好)眼神茫然、掉了队——连他也觉得一次太多;不止一次,我谨慎地问奥利维耶,是否今天到此为止、改日再续。我原可省了这份心——分明奥利维耶神清气爽,眼神明亮,完全自如;我甚至笑起来,因为难以置信他竟不垮——可他全然不垮!那是个也许二十岁的年轻人,只有一点概形的常识、一点拓扑与topos;我相信他倒也摆弄过不少无限离散群……老实说,那是三脚猫的底子;而凭这个,他竟能填上所有的空白,毫不费力地“感到”我这个老资格、凭十五年与这一主题的熟悉、在两三个钟头里全速讲述的东西。我从未见过这样的事——至多在Deligne身上,也许还有Cartier;Cartier年轻时在这一条线上也相当出众。

总之,分明是拍板成交:奥利维耶将就上述题目做他的第三阶段博士论文。然而他想必不曾料到路的尽头有什么。总之,在其后他埋头苦干的两年里、乃至之后,我再也没有见过他。他的名义导师是Contou-Carrère,当然;但有机会同这样一个机灵的人聊聊,本会使我高兴。实际上,连答辩的通知都不曾发我,而且我相信从未收到过那份论文的副本——但我记得亲手捧读过一份,是在一个有资格得到它的人手里\fnote{20}{(*)}{这场捉迷藏,愈发离奇,因为毫无疑问:在整个朗格多克,能与Contou-Carrère一道、对Olivier~Leroy的工作多少懂点什么的人,只有我。不必说,Leroy那份科学院报告的底稿,我也同样不曾亲手捧读。也许我在自欺;但我觉得:倘若我不曾被如此苛刻地隔离、以致实际上无从干预,我总归能找到办法,让那篇倒霉的短文发表——必要时经由Cartan或Serre;他们虽不在局中,却肯信任我,只要我担保工作的严肃性。(6月7日)想必很久之后我才知道Leroy已答辩,而且我太忙,无暇自问为何连通知都没人给我。直到Contou-Carrère本人的答辩之后,我才“翻然醒悟”——那场答辩,我被认定是其论文导师(x)。他想方设法,使评审委员会成员中唯独我没有资格得到其论文的正式定本!今天我终于收到一份副本——他以为(他写道)得到它“对我没有兴趣”……(x)更确切地说:有一两年,Contou-Carrère谨慎地同时耍着两位“导师”(谁也说不准……),两人彼此不知有一位“平行”导师。Verdier的导师角色,我是在最后关头才知道的:1983年春,当Verdier分明终究要他的命变得清楚时,Contou-Carrère才最终转而投奔我。}。答辩是在那篇科学院报告“沉没”之前还是之后,我说不好。我说的是那次沉没:在第33节“注记——或新伦理(1)”里,我相当周详但不点名地谈到它。经手这次沉没的两位数学家,是Pierre~Cartier(正是那位:其惊人的直觉之敏捷,在我谈起他年轻的同行的敏捷时重回我心头;而他却如此客气、怀着全部的遗憾,沉没了那份敏捷)与Pierre~Deligne——带着他那句历史性的话:这些数学“不逗他乐”。(它们在他年轻时,可是“逗过乐”的……)我还该加上Contou-Carrère本人:他没有为自己的学生抬一根手指——那会冒触怒权势人物的危险。他想必向Olivier~Leroy暗示:最好忘掉那篇不合时宜的论文。无论如何,清楚的是:Leroy分明对这一情节画了粗粗的一道——即便日后有机会发表的不只是一篇科学院报告、而是其工作的全部,我深信他也不会用\fnote{21}{(*)}{这道“粗线”的一个雄辩征兆:Olivier~Leroy两年前蒙彼利埃一个助教职位出缺时提交的申请材料里,既不提其第三阶段论文的题目,也不提曾是其导师的Contou-Carrère的名字。他也不提任何一项个人的工作,无论什么。显然,他当时还没决定要不要那份职位——于是,尽管其天赋惊人,职位给了另一位候选人:材料扎实、意图毫无疑义的一位。}。这一回,电锯又干得漂亮\fnote{22}{(**)}{一个有趣的巧合:最近我得到回音,Cartier有意把其一场Bourbaki报告献给我(我相信这是头一遭),而且那份报告恰恰专论topos理论——正是这些topos,被同一位Cartier判定不配进入一篇科学院报告。这莫非是最近这几年时尚风向的转变?肯定不是;而且一切都环环相扣:那份报告谈的是topos在逻辑中的使用!我的朋友Cartier这份感人的题献,在我看来,与去年在一个重大场合所致的葬礼悼词同一血脉(见“葬礼悼词——或恭维”一注,n\textsuperscript{o}~104):在那里,“topos”一词被提到(在一串抖得利落的恭维之中),紧接着急忙补上(作为唯一而雄辩的评语)它们“如今用于逻辑”——别处没有,自不待言,至少只要我那些乐善好施恭维的朋友们、凭其手中的权力,还能阻止的话……(Cartier报告的出处:Catégories, logiques et faisceaux, modèles de la théorie des ensembles, Bourbaki研讨班n\textsuperscript{o}~513,1978年2月。)我在某些人(那些定调子的人,如Deligne、Cartier、Quillen……)对topos在几何中这类革新而深刻的观念的那份居高临下的(与抵制的……)态度里,感到一种惊人的狂妄。且假定他们中哪怕一位有那份材料(或那份天真……),能像我经由引入平展topos与晶状topos、从虚无中引出代数簇的一种新的拓扑视野(并由此,在拓扑到来之前,引出代数几何与算术深刻革新的手段)——毫无疑问,他乐于在自己身上培养、在他人身上煽动的这份轻蔑的态度本身,就会解除那视野与革新的力量,只为一份自负效劳。}。

\dmark{}\origpage{412}尽管有这场厄运,1981年初的那几个月里,我到底有幸常见Leroy。那是一场我当时主持的微研讨班,论Teichmüller塔的代数-算术理论(《纲领概要》里略有谈及)。真正意义的听众,只有Contou-Carrère与Leroy。即便对一个超选的巴黎听众(我懂我在说什么),一整间屋子也未必有三四个不掉队的。老实说,我办那场研讨班——当时Contou-Carrère正全力打磨其Schubert闭链的想法——是为了Leroy,想着他或许会对如此辉煌的题目咬钩。分明他“感觉得到”我在做什么;但他(我相信)事先已决定不“咬钩”。奇怪的是他甚至肯费心来——某种东西必定使他着迷,一如我着迷;而他自己对究竟要什么,也不甚了然。当我明白他不会咬钩,我便止损了。面对两位观众(无论多么光彩)继续独白,我无意。我与Leroy的第二场也是最后一场交谈,就在那时。我相信,此后我再没有见过他。

Leroy与我之间,除七年前的那个晚上,没有过真正的数学讨论——这解释了为何对他的工作我几乎一无所知,除那份不合时宜的topos工作。他的厄运想必不曾增进他对像我、甚至Contou-Carrère或“上流”数学界其他人等的信任。我听说他在文学院办研讨班——那里有一群和气的、彼此相处融洽的数学人。他们讲授组合拓扑的想法——一个近十年来全在我射程之内的题目。既然我天性谨慎(可不是嘛!),我没有问他讲些什么,也不知道他是否打算发表。至于境况,他过着一种最不合法的生活(却既非外国人、亦无身份问题):到处带习题课,报酬(嘘……)不知出自哪些小金库,背着司库与审计法院。我想,他是否最终走数学这一行,尚未下定决心;长远看,无论有没有审计法院,这想必不是个舒坦的处境。倘若我那幅富于教益的葬礼图景——他在其中充当第四口编外棺材——能帮他驱散困惑,这一回是在知根知底的情况下,我会高兴。

\subsection{掘墓人——或全体教众}

\notehead{97}{}

(5月24日)我是顶着我身上某种迟疑,才终于决定点名提及我的某些亲近的旧友与同行——在数学界;我得以看到他们充当“掘墓人”(或“电锯”)的角色:从一开始就斩断某些身份寒微或无保障的数学家的尝试——他们试图重提我的某些观念、循其固有的逻辑加以发展,或者(如Yves~Ladegaillerie的个案)仅仅追随一种带着我的影响之印记的进路与做派。如我一说再说,这类牵连他人的迟疑——或仅仅是在未曾\dmark{}\origpage{414}征询其意见的情况下点其名的迟疑\fnote{23}{(*)}{比如,为收进一条注(注n\textsuperscript{o}~19)——其中将点名提及所有跟我做完国家博士论文的学生——我就有过这样的迟疑。我身上这份犹豫,想必来自我的许多学生对与我这个人的名字相连的抗拒;那份抗拒,想必几年前我已在未成形的层面知觉到了。我的旧学生们(带或不带引号)中,当时与我划清界限的意愿已被我清晰知觉的,只有Contou-Carrère(我恰在那时才在他身上发现它)与Deligne(在那里,此事自1968年起已相当分明,尽管我不曾疑心这份意愿将把他带向何方)。在Deligne的个案里,我迟迟不肯以“多少算个学生”的名义点名他的迟疑,格外强烈:我不愿显得在拿这样光彩的学生夸耀,而他自己则不肯让那把他与我、与我的著作连在一起的纽带显露分毫。我的反思,还使我进而明白:这条纽带,在我年轻朋友的生命与著作里所占的分量,远比我曾经猜想的要大。(6月1日)关于我身上这些蓄意的心境,见“三天之后”的3月27日(原书日期如此,与5月24日不合)的“格格不入者”一注(n\textsuperscript{o}~67$'$)。}——在《收获与播种》里并不少见。在每一个个案里,我最终检视了那份迟疑,并懂得它并无根据;其来源不是体贴,而是混乱,甚至怯懦。在我点名陈述他人的行为或态度的所有情形里(依我看),这些行为或态度绝非“保密”性质。它们关涉当事人的职业生活,连同其必然的种种后果——牵连其他同行(包括我自己)的职业生活(并由此,牵连他们的全部生活)。我点名的每一个人,对其行为与态度、及其含义的全部谱系(无论他乐于忽视与否),都与我对我自己的负有同样的责任。倘若其行为的某些后果以这样或那样的形态回到他身上——比如经由我的这个人在此充当媒介的一份公开的“问责”——他毫无根据感到被冒犯。倘若我的语言有时形象而辛辣,我的意图绝非论战,绝非冒犯或羞辱任何人,而宁可是描述事实、描述我感受它们的方式——作为一种敦促,促请每一个人(首先是我点名的那几位)从其自身一侧检视它们,而不是以这样或那样的方式把它们打发掉(一如《收获与播种》的反思之前我常常做的那样)。倘若被如此唤起的人选择感到被冒犯,那是关涉他自己的选择。这份选择可能使我难过——出自我所敬重、甚至怀有感情的人;但它不压垮我。我谈过的那份迟疑——我的事事观感中某种混乱的征兆——一经理解、并由此超越,便无踪无影。

在葬礼的反思里,没有任何时刻,我曾感到有一场针对我的著作、针对那些胆敢从中汲取灵感(而非仅仅借用工具、却讳言锻造并交到他们手里的工匠之名)的人而精心策划的宏大“阴谋”。没有阴谋;但有一个\textbf{共识}——在我所称的“上流”数学界里,它至今在我眼前毫无破绽。这个共识,除至多极罕见的例外,绝非由一份对我这个人、对我的著作的有意识的恶意喂养\fnote{24}{(*)}{我所认为的毫不含糊的恶意行为,我只在Deligne与Verdier的个案里得知。}\dmark{}\origpage{415}。仅在若干例外的个案里,它才表达为对前面各注所谈的四位“共同被埋葬者”中这一位或那一位的、毫不含糊的恶意。但肯定,这样的恶意之所以能在我昔日的某位学生身上滋蔓,之所以能不受阻拦地表达,全赖普遍共识的怂恿。

这份共识,在我的几乎全部旧友或旧学生身上,显现的不是“恶意”的态度,而是(我相信)全然无意识的机制:一种令人错愕的整齐划一、一种毫无破绽的效率,像扫麦秸一样扫走健全意识与数学家的健全本能,让位于纯然自动的\textbf{拒斥态度}\fnote{25}{(**)}{这些拒斥态度,当然从不以这样的面目出现——即便在我朋友Deligne或Verdier那样极端的个案里。在对我有意识的层面,它们几乎不可见;那个层面(如我已有机会说过的)在我昔日的朋友与学生身上,几乎总是(也许甚至总是)同情与敬重的心境(其中某位,还在尽力设法自我开脱)。这样的同情与敬重的心境,不只在有意识的“意见”的表浅层面在场,甚至在实际的好恶、以及对他人实际认识的更深层面上也在场(独立于人们力图把他人关进其中的那些形象)。我们在这里处在一个典型的\textbf{两价性}局面里(我几乎想说:集体性的两价性),粗看之下,什么也“看不见”!(可与“敌父(1)、(2)”(第29、30节)里的反思对读;在《收获与播种》里,我头一回在那里触及这个标记过我生命中许多关系的两价性侧面——不只在数学界。)然而,在具体显现的层面上(葬礼里已被充分检视),这些两价力量的“合力”,在我看来,不再有任何两价性;它分明带着“令人错愕的整齐划一与毫无破绽的效率”,以我正要进一步检视的“自动拒斥态度”的面目呈现。}。这样的自动态度,我疑心,不只由我这个人、由那些其数学“气味”多少让人想起我的人所触发——也对着任何不像已经持有某个“建制”的\textbf{心照不宣的担保}的数学家:无论他已身在其中的行列,还是俨然其中某一位的“门生”(借用Verdier笔下的说法)。我觉得\dmark{}\origpage{416}:在几乎全体数学家身上,最低限度的“数学开放”的心境——它是健全意识与数学健全本能得以启动的前提——\textbf{只对着已经持有这样一份担保的人才启动}。

这类机制,想必几乎是普遍的:不只在数学界,也在社会的所有部门,无一例外。它远远超出一切个案。倘若(依我看来)在我的个案里、以及在建制眼中俨然“我的门生”者的个案里,有某种例外的局面,那是因为:过去,我曾持有“他们自己人”之一的身份,连同对我、对“我的人”的那份惯常的“最低限度的开放”。这份身份,因我1970年的离去而被收回。或者更确切地说:经由我自己的选择——它在我离去后的年月里不止一次清楚地表达,并经由我直至今日的生存方式——我确实不再是一个“他们自己人”。事实上,我自己也不再感到自己是“他们自己人”;我离开了一个曾经共同的世界,绝无回归之念。今天,我的“回归数学”绝不是回归“他们中间”、回归建制,而是回归数学本身;更确切地说,回归一种持续的数学投资,回归对我数学反思的发表活动。

我才开始明白:我的离去,曾被我的旧友们、我的学生们感受为一种“开小差”,甚至一种“冒犯”\fnote{26}{(*)}{这样一种看待与感受事情的方式,在我的朋友Zoghman~Mebkhout身上表达得格外雄辩。正是由于这次开小差,我才要为他与上流数学界的种种挫折负责:唯有他,失去了那份曾环绕着今日乐于把他当拖鞋踢的人的“庇护”与支持。}。想必那是最简单的办法,用以打发我离去的意味、它可能在人们身上唤起的追问:借这样一份弥散的\textbf{受害}之感、借怨恨的自动反应,表达为一场\textbf{报复}之举(在意识的层面,它想必很少被当作报复、甚至当作行为来知觉):既然他与我们切割,我们就与他切割——我们不再给他与“他的人”以保留给“自己人”的“注意的自动机制”的收益——他与他们,将如初来乍到者一般,领受自动拒斥的严苛!

\dmark{}\origpage{417}局面(对我的旧友们与学生们)还更复杂:因为我不但曾是建制的一员,而且他们中任何一个人,要做一个数学家,不可能不随时使用我是作者的观念、想法、工具与结果。我不知道在我们的科学或任何别的科学的历史里,可曾有过如此难堪的悖论的例子!在这一光照下,那些电锯效应(绝不限于我的朋友Deligne)——为斩断任何发展带有我的印记的观念的意图(尽管这样的发展只会加深那份困惑)——如今在我眼前,显出被一种无情的内在逻辑所推动:自某个已然做出的选择——拒斥的选择——起的一份\textbf{必然}。同样,我看到人们四处做出的努力——把那些已进入公共遗产、人们无论愿意与否都离不开的观念、想法、工具与结果的来源,置于彻底的沉默之下——也属同一逻辑。我在Deligne那些非常扎眼的“战役”面前——他装作一项一项地把我的若干主要数学贡献的父权据为己有(或把残渣慷慨地记在某个形影不离的好哥们名下)——曾以为察觉的那份“漠然”——那绝不是漠然,而是一份\textbf{心照不宣的首肯}。Deligne不过是做着建制的集体无意识期待他做的事:\textbf{抹去}那个与所有人切割的人的名字,从而化解那个无法容忍的悖论——\textbf{以一份可容忍的假父权,替换一份真实却不可接受的父权}。

在这一光照下,主祭Deligne,显出的不再是一个按其生命的深层力量与行为塑造一场时尚的人,而分明是\textbf{工具}——凭借其“合法继承人”的角色天造地设的工具:服务于一份毫无破绽的\textbf{集体意志};那意志,执意于一项不可能的任务:把我的名字、我的个人做派,从当代数学中抹去。

我几乎不怀疑:这幅图景,就其实质而言,表达了事情的实在——至少在集体的层面上。肯定,我的“归来”——它出人意料地为一场曾进行得让所有人称心的下葬画上句点,或(即便不画句点)至少不合时宜、不可容忍地扰乱一场仿佛早已排定的仪式的进行——这份归来,将不只使主祭者中的某一位不适与不满,还将使为这场可悲场合齐聚的\origpage{418}全体教众为难!当然,我毫无概念:那鼎鼎大名的集体无意识,将发明怎样的“解数”,来打发这位可敬的亡者不期归来所造成的烂摊子——他从为其量身准备的舒适的棺材里突然爬出(不可容忍的丑闻),装模作样按他自己的方式主持自己的葬礼。但我信得过全体教众:他们会找出办法,打发掉数学大厦里这桩多出来的小小矛盾——它已经万事俱备,只差这个!

我如今,在每个人的意象与态度的层面上,已能相当清楚地知觉集体共识的反映及其总形态,知觉那份抹去、埋葬的集体意志。那是普遍使用的“两幅图景”体系——两幅相互矛盾的图景,人们同时依凭它们行事;我在《收获与播种》里,就我自己的个案,头一回谈到它。(见“功劳与轻蔑”一节,第12节。)我怀疑可有人会直言不讳、明明白白地说:“Grothendieck做的尽是冒牌数学,不提也罢,说正事吧。”照直说,那会太过分明地违背建制的公理——至少眼下如此。按事情预定的演变,二三十年后,这问题反正无从提出:届时连这个名字都无从提起——它早已被所有人遗忘。共同的战术——无论个人还是集体——是沉默:人们不想亡者,至少不作为数学家想;人们不谈他,不提他(除非不得已,借天赐的缩写SGA或EGA;且待这些引证被另一些没有亡者任何踪迹的引证替换……)。

然而,也有——想必例外——的场合,完全的沉默变得行不通。其中一个场合,我料想,曾是我的CNRS任职申请;它想必使不止一人为难\fnote{27}{(*)}{(5月26日)今天我恰好从Zoghman~Mebkhout的电话里得知:我在CNRS全国委员会的同事们为我出了一把力,为我安排了一个两年期的“接待职位”。我不知道他们是否怀着热情做的;总之,委员会里我的朋友中,没有一位把这份出力做到底——打来电话、捎来一句好话,把5月15日就已成定局的好消息告诉我。(9月)我到底经由CNRS一封8月16日的信得到通知:那是一项attaché de recherches职位,为期一年(而非两年)。}。另一场,将是《收获与播种》的预分发\fnotevariant{28}{(*)}{}{说的是由我的大学经手的一份限量印制(150份),在最近的同事与朋友中间分发。},且待其在《数学反思》第一卷里出版(倘若我的出版商不至于胆怯退缩,肯冒得罪整个科学建制的风险)。这些都是亡者本人不可容忍的出格所造成的场合:他不知好歹地走出了派给他的角色。另一个场合(对理解这场尚未被一个不守纪律的亡者搅乱的葬礼,也许更有教益),是去年“隆重”庆祝的IHES二十五周年庆典。作为“IHES的四枚菲尔兹奖章之首”,在这个庄严的场合,要把我全然略过,是困难的——尽管我在IHES存在的四个英雄年月里、为使它取得真实存在所扮演的角色,被略过了。为向我致敬而精心炮制的《葬礼悼词》——载于为此庆典印行的小册子(我已有机会两次指涉它)——依我看是这一类文体的典范:一种优雅而含蓄的方式,使人人满意地化解当代数学里的这份“小小矛盾”……

瞧,我忽然精神焕发——像开始嗅到马厩的马!将近两个星期前,我开始就这一富有教益的情节写一场反思,在一条当即得名“葬礼悼词——或恭维”的注里。经过几番该把它放在哪里的踌躇(它出自为葬礼所写的第一条注的一条迟到的脚注),看来最自然的安放之处,不是“编年”的位置,而是那场应为葬礼压轴的“葬礼仪式”。于是,不经我寻求,我三周来沿着最后三个行列——“研讨会”“学生”,以及刚刚入列的“灵车”(随葬礼的最终部分,即葬礼仪式而来)——所追的那根“线”,接上了头;这场仪式的首要印记,恰恰是那份我5月12日开始检视的《葬礼悼词》杰作\dmark{}\origpage{420},它如今是自然而然紧随本注之后的那条注\fnote{29}{(*)}{(1984年11月)由于一场意外的疾病情节,该注(n\textsuperscript{o}~104)与“本注”被一个新的行列隔开——“亡故者——仍未死去”(n\textsuperscript{o}~98--103)。}。

我终于(又一次?)触到目标了!同时,这场对一篇葬礼悼词的反思的开端,忽然取得一个新的维度。它不再只是一个强力的头脑为一桩执念效劳的精巧发明——在一场官方“重大场合”的贵宾们面前、对着他们的漠然或奉命的关注挥霍——它尤其是一份在所有场合中最微妙的场合、对着一份集体期待所做的完美而体贴周到的回答:那份期待,问的是对我的这个人应当采取的态度。倘若他那一代人中有人真正配得上全体教众毫无保留的感戴,那正是我的朋友Pierre~Deligne——以他那份绝无瑕疵的完美,履行了人们期待于他的角色。


\end{document}
